% Options for packages loaded elsewhere
\PassOptionsToPackage{unicode}{hyperref}
\PassOptionsToPackage{hyphens}{url}
\documentclass[
  a4paper,
  oneside]{scrbook}
\usepackage{xcolor}
\usepackage{amsmath,amssymb}
\setcounter{secnumdepth}{5}
\usepackage{iftex}
\ifPDFTeX
  \usepackage[T1]{fontenc}
  \usepackage[utf8]{inputenc}
  \usepackage{textcomp} % provide euro and other symbols
\else % if luatex or xetex
  \usepackage{unicode-math} % this also loads fontspec
  \defaultfontfeatures{Scale=MatchLowercase}
  \defaultfontfeatures[\rmfamily]{Ligatures=TeX,Scale=1}
\fi
\usepackage{lmodern}
\ifPDFTeX\else
  % xetex/luatex font selection
\fi
% Use upquote if available, for straight quotes in verbatim environments
\IfFileExists{upquote.sty}{\usepackage{upquote}}{}
\IfFileExists{microtype.sty}{% use microtype if available
  \usepackage[]{microtype}
  \UseMicrotypeSet[protrusion]{basicmath} % disable protrusion for tt fonts
}{}
\makeatletter
\@ifundefined{KOMAClassName}{% if non-KOMA class
  \IfFileExists{parskip.sty}{%
    \usepackage{parskip}
  }{% else
    \setlength{\parindent}{0pt}
    \setlength{\parskip}{6pt plus 2pt minus 1pt}}
}{% if KOMA class
  \KOMAoptions{parskip=half}}
\makeatother
\usepackage{color}
\usepackage{fancyvrb}
\newcommand{\VerbBar}{|}
\newcommand{\VERB}{\Verb[commandchars=\\\{\}]}
\DefineVerbatimEnvironment{Highlighting}{Verbatim}{commandchars=\\\{\}}
% Add ',fontsize=\small' for more characters per line
\usepackage{framed}
\definecolor{shadecolor}{RGB}{248,248,248}
\newenvironment{Shaded}{\begin{snugshade}}{\end{snugshade}}
\newcommand{\AlertTok}[1]{\textcolor[rgb]{0.94,0.16,0.16}{#1}}
\newcommand{\AnnotationTok}[1]{\textcolor[rgb]{0.56,0.35,0.01}{\textbf{\textit{#1}}}}
\newcommand{\AttributeTok}[1]{\textcolor[rgb]{0.13,0.29,0.53}{#1}}
\newcommand{\BaseNTok}[1]{\textcolor[rgb]{0.00,0.00,0.81}{#1}}
\newcommand{\BuiltInTok}[1]{#1}
\newcommand{\CharTok}[1]{\textcolor[rgb]{0.31,0.60,0.02}{#1}}
\newcommand{\CommentTok}[1]{\textcolor[rgb]{0.56,0.35,0.01}{\textit{#1}}}
\newcommand{\CommentVarTok}[1]{\textcolor[rgb]{0.56,0.35,0.01}{\textbf{\textit{#1}}}}
\newcommand{\ConstantTok}[1]{\textcolor[rgb]{0.56,0.35,0.01}{#1}}
\newcommand{\ControlFlowTok}[1]{\textcolor[rgb]{0.13,0.29,0.53}{\textbf{#1}}}
\newcommand{\DataTypeTok}[1]{\textcolor[rgb]{0.13,0.29,0.53}{#1}}
\newcommand{\DecValTok}[1]{\textcolor[rgb]{0.00,0.00,0.81}{#1}}
\newcommand{\DocumentationTok}[1]{\textcolor[rgb]{0.56,0.35,0.01}{\textbf{\textit{#1}}}}
\newcommand{\ErrorTok}[1]{\textcolor[rgb]{0.64,0.00,0.00}{\textbf{#1}}}
\newcommand{\ExtensionTok}[1]{#1}
\newcommand{\FloatTok}[1]{\textcolor[rgb]{0.00,0.00,0.81}{#1}}
\newcommand{\FunctionTok}[1]{\textcolor[rgb]{0.13,0.29,0.53}{\textbf{#1}}}
\newcommand{\ImportTok}[1]{#1}
\newcommand{\InformationTok}[1]{\textcolor[rgb]{0.56,0.35,0.01}{\textbf{\textit{#1}}}}
\newcommand{\KeywordTok}[1]{\textcolor[rgb]{0.13,0.29,0.53}{\textbf{#1}}}
\newcommand{\NormalTok}[1]{#1}
\newcommand{\OperatorTok}[1]{\textcolor[rgb]{0.81,0.36,0.00}{\textbf{#1}}}
\newcommand{\OtherTok}[1]{\textcolor[rgb]{0.56,0.35,0.01}{#1}}
\newcommand{\PreprocessorTok}[1]{\textcolor[rgb]{0.56,0.35,0.01}{\textit{#1}}}
\newcommand{\RegionMarkerTok}[1]{#1}
\newcommand{\SpecialCharTok}[1]{\textcolor[rgb]{0.81,0.36,0.00}{\textbf{#1}}}
\newcommand{\SpecialStringTok}[1]{\textcolor[rgb]{0.31,0.60,0.02}{#1}}
\newcommand{\StringTok}[1]{\textcolor[rgb]{0.31,0.60,0.02}{#1}}
\newcommand{\VariableTok}[1]{\textcolor[rgb]{0.00,0.00,0.00}{#1}}
\newcommand{\VerbatimStringTok}[1]{\textcolor[rgb]{0.31,0.60,0.02}{#1}}
\newcommand{\WarningTok}[1]{\textcolor[rgb]{0.56,0.35,0.01}{\textbf{\textit{#1}}}}
\usepackage{longtable,booktabs,array}
\newcounter{none} % for unnumbered tables
\usepackage{calc} % for calculating minipage widths
% Correct order of tables after \paragraph or \subparagraph
\usepackage{etoolbox}
\makeatletter
\patchcmd\longtable{\par}{\if@noskipsec\mbox{}\fi\par}{}{}
\makeatother
% Allow footnotes in longtable head/foot
\IfFileExists{footnotehyper.sty}{\usepackage{footnotehyper}}{\usepackage{footnote}}
\makesavenoteenv{longtable}
\usepackage{graphicx}
\makeatletter
\newsavebox\pandoc@box
\newcommand*\pandocbounded[1]{% scales image to fit in text height/width
  \sbox\pandoc@box{#1}%
  \Gscale@div\@tempa{\textheight}{\dimexpr\ht\pandoc@box+\dp\pandoc@box\relax}%
  \Gscale@div\@tempb{\linewidth}{\wd\pandoc@box}%
  \ifdim\@tempb\p@<\@tempa\p@\let\@tempa\@tempb\fi% select the smaller of both
  \ifdim\@tempa\p@<\p@\scalebox{\@tempa}{\usebox\pandoc@box}%
  \else\usebox{\pandoc@box}%
  \fi%
}
% Set default figure placement to htbp
\def\fps@figure{htbp}
\makeatother
\setlength{\emergencystretch}{3em} % prevent overfull lines
\providecommand{\tightlist}{%
  \setlength{\itemsep}{0pt}\setlength{\parskip}{0pt}}
\usepackage[]{natbib}
\bibliographystyle{apalike}
\usepackage{fancyhdr}
\usepackage{fancyvrb}
\usepackage{bm}
\usepackage[spanish]{babel}
\usepackage{fontawesome5}
%\usepackage[top=30mm, bottom=40mm, left=25mm, right=25mm, showframe,
%            foot=10mm, marginparsep=0mm]{geometry}
\usepackage[top=30mm, bottom=40mm, left=25mm, right=25mm, 
            foot=10mm, marginparsep=0mm]{geometry}


%%\usepackage{Alegreya,AlegreyaSans}
%%\usepackage{lipsum}
%%\usepackage{mathpazo}
%\usepackage{fourier}
%\usepackage[scale=.7]{sourcecodepro}
\usepackage[sc]{mathpazo}

\usepackage{booktabs}
\usepackage{amsthm}


%\usepackage{titlepic}
%\titlepic{\includegraphics[width=0.25\textwidth]{images/cover.jpg}}
\usepackage{eso-pic}
\newcommand\BackgroundPic{%
\put(0,0){%
\parbox[b][\paperheight]{\paperwidth}{%
\vfill
\centering
\includegraphics[width=\paperwidth,height=\paperheight,%
keepaspectratio]{images/coverAM.jpg}%
\vfill
}}}

% ***********************************************
% Márgenes
% ***********************************************
%%\topmargin=0in
%\evensidemargin=0cm
%\oddsidemargin=0cm
%\textwidth=16cm
%%\textheight=9.0in
%%\headsep=0.5in

% ***********************************************
% Encabezado y pie de página
% ***********************************************
\pagestyle{fancy}
\fancyhead[LE,LO]{Análisis Multivariante}
\fancyhead[RE,RO]{\nouppercase{\leftmark}}
\fancyfoot[CE,CO]{\thepage}

\renewcommand{\chaptermark}[1]{%
\markboth{#1}{}}

%\definecolor{excolor}{RGB}{60, 100, 120}
%\definecolor{excolor}{RGB}{136, 78, 160}
\definecolor{excolor}{RGB}{19, 141, 117}

\usepackage[nottoc,notlot,notlof]{tocbibind} % Bibliography in Table of Contents


\usepackage{framed,color}
\definecolor{shadecolor}{RGB}{248,248,248}


\renewcommand{\textfraction}{0.05}
\renewcommand{\topfraction}{0.8}
\renewcommand{\bottomfraction}{0.8}
\renewcommand{\floatpagefraction}{0.75}


\makeatletter
\newenvironment{kframe}{%
\medskip{}
\setlength{\fboxsep}{.8em}
 \def\at@end@of@kframe{}%
 \ifinner\ifhmode%
  \def\at@end@of@kframe{\end{minipage}}%
  \begin{minipage}{\columnwidth}%
 \fi\fi%
 \def\FrameCommand##1{\hskip\@totalleftmargin \hskip-\fboxsep
 \colorbox{shadecolor}{##1}\hskip-\fboxsep
     % There is no \\@totalrightmargin, so:
     \hskip-\linewidth \hskip-\@totalleftmargin \hskip\columnwidth}%
 \MakeFramed {\advance\hsize-\width
   \@totalleftmargin\z@ \linewidth\hsize
   \@setminipage}}%
 {\par\unskip\endMakeFramed%
 \at@end@of@kframe}
\makeatother

\makeatletter
\@ifundefined{Shaded}{
}{\renewenvironment{Shaded}{\begin{kframe}}{\end{kframe}}}
\makeatother
\newenvironment{rmdejemplo}
  {
  \begin{itemize}
  \renewcommand{\labelitemi}{
    \raisebox{-.7\height}[0pt][0pt]{
      {\setkeys{Gin}{width=3em,keepaspectratio}}
    }
  }
  \setlength{\fboxsep}{1em}
  \begin{kframe}
  \item
  }
  {
  \end{kframe}
  \end{itemize}
  }



\newenvironment{rmdblock}[1]
  {
  \begin{itemize}
  \renewcommand{\labelitemi}{
    \raisebox{-.7\height}[0pt][0pt]{
      {\setkeys{Gin}{width=2em,keepaspectratio}\includegraphics{images/#1}}
    }
  }
  \setlength{\fboxsep}{1em}
  \begin{kframe}
  \item
  }
  {
  \end{kframe}
  \end{itemize}
  }
\newenvironment{rmdnote}
  {\begin{rmdblock}{note}}
  {\end{rmdblock}}
\newenvironment{rmdcaution}
  {\begin{rmdblock}{caution}}
  {\end{rmdblock}}
\newenvironment{rmdimportant}
  {\begin{rmdblock}{important}}
  {\end{rmdblock}}
\newenvironment{rmdtip}
  {\begin{rmdblock}{tip}}
  {\end{rmdblock}}
\newenvironment{rmdwarning}
  {\begin{rmdblock}{warning}}
  {\end{rmdblock}}

\usepackage{makeidx}
\makeindex

%\urlstyle{tt}

\makeatletter
\def\thm@space@setup{%
  \thm@preskip=8pt plus 2pt minus 4pt
  \thm@postskip=\thm@preskip
}
\makeatother






\usepackage{bookmark}
\IfFileExists{xurl.sty}{\usepackage{xurl}}{} % add URL line breaks if available
\urlstyle{same}
% fallback for those not using the hyperref driver hyperxmp:
\makeatletter
\@ifundefined{xmpquote}{\newcommand{\xmpquote}[1]{#1}}{}
\makeatother
\hypersetup{
  hidelinks,
  pdfcreator={LaTeX via pandoc}}

\author{}
\date{\vspace{-2.5em}Last compiled on 24 septiembre, 2026}

\usepackage{amsthm}
\newtheorem{theorem}{Teorema}[chapter]
\newtheorem{lemma}{Lema}[chapter]
\newtheorem{corollary}{Corolario}[chapter]
\newtheorem{proposition}{Proposición}[chapter]
\newtheorem{conjecture}{Conjecture}[chapter]
\theoremstyle{definition}
\newtheorem{definition}{Definición}[chapter]
\theoremstyle{definition}
\newtheorem{example}{Ejemplo}[chapter]
\theoremstyle{definition}
\newtheorem{exercise}{Ejercicio}[chapter]
\theoremstyle{definition}
\newtheorem{hypothesis}{Hypothesis}[chapter]
\theoremstyle{remark}
\newtheorem*{remark}{Remark}
\newtheorem*{solution}{Solución }
\begin{document}

%\cleardoublepage\newpage\thispagestyle{empty}\null
%\cleardoublepage\newpage\thispagestyle{empty}\null
%\cleardoublepage\newpage

\thispagestyle{empty}
\AddToShipoutPicture*{\BackgroundPic}


\clearpage
\newpage
\mbox{}
\clearpage
\newpage

% \newpage
% \thispagestyle{empty}
% \begin{center}
% Material para el curso de Análisis Multivariante\\
% Máster en Técnicas Estadísticas
% %\includegraphics{images/dedication.pdf}
% \end{center}
%
% \setlength{\abovedisplayskip}{-5pt}
% \setlength{\abovedisplayshortskip}{-5pt}

{
\setcounter{tocdepth}{1}
\tableofcontents
}
\chapter*{Prefacio}\label{prefacio}
\addcontentsline{toc}{chapter}{Prefacio}

Este documento es el resultado de unificar el material de Análisis Multivariante, materia del Máster en Técnicas Estadísticas organizado por la Universidad de Santiago de Compostela (coordinadora), Universidad de A Coruña y Universidad de Vigo.

En esta materia se pretende que el alumno aprenda las técnicas más importantes en el análisis de datos multidimensionales. Se abordan tanto los aspectos metodológicos como del uso de software que proporcionan los paquetes estadísticos en el campo del análisis multivariante.

En la Bibliografía se incluye una lista de referencias con libros y artículos consultados para la elaboración de este material.

\begin{quote}
Agradecemos que nos ayudes a mejorar este documento. Si encuentras erratas, errores, o tienes alguna sugerencia para explicar mejor alguna sección, por favor escribe a \href{mailto:beatriz.pateiro@usc.es}{\nolinkurl{beatriz.pateiro@usc.es}}.

\begin{flushright}
Última actualización: 24 septiembre, 2026

\end{flushright}
\end{quote}

\section*{Autores}\label{autores}
\addcontentsline{toc}{section}{Autores}

\begin{itemize}
\tightlist
\item
  Beatriz Pateiro López. Departamento de Estadística, Análisis Matemático y Optimización. Universidad de Santiago de Compostela.
\item
  César Sánchez Sellero. Departamento de Estadística, Análisis Matemático y Optimización. Universidad de Santiago de Compostela
\end{itemize}

Con nuestro agradecimiento a Wenceslao González Manteiga, Catedrático de Estadística de la Universidad de Santiago de Compostela (USC) y profesor de la materia de Análisis Multivariante hasta el curso 2017/18.

\hfill\break

\pandocbounded{
\includegraphics[keepaspectratio,alt={Creative Commons License}]{images/by-nc-nd.png}}\\
Esta obra está bajo una \href{http://creativecommons.org/licenses/by-nc-nd/4.0}{Licencia Creative Commons Atribución-NoComercial-SinDerivadas 4.0 Internacional}

\section*{Prerrequisitos}\label{prerrequisitos}
\addcontentsline{toc}{section}{Prerrequisitos}

A lo largo del documento se incluyen ejemplos de aplicación de los distintos métodos del análisis multivariante presentados. El software estadístico utilizado es R \citep{R-base}. Puedes copiar y ejecutar los fragmentos de código de R para tu práctica. Simplemente debes asegurarte de tener instalados algunos paquetes que necesitaremos y que no están incluídos entre los paquetes que instala R por defecto.

\footnotesize

\begin{Shaded}
\begin{Highlighting}[]
\SpecialCharTok{\textgreater{}} \CommentTok{\# Paquetes utilizados}
\ErrorTok{\textgreater{}}\NormalTok{ packages }\OtherTok{\textless{}{-}} \FunctionTok{c}\NormalTok{(}\StringTok{"aplpack"}\NormalTok{, }\StringTok{"car"}\NormalTok{, }\StringTok{"moments"}\NormalTok{, }\StringTok{"MVN"}\NormalTok{, }\StringTok{"mvtnorm"}\NormalTok{, }\StringTok{"rgl"}\NormalTok{, }\StringTok{"scatterplot3d"}\NormalTok{)}
\SpecialCharTok{\textgreater{}} \FunctionTok{install.packages}\NormalTok{(packages)}
\SpecialCharTok{\textgreater{}} 
\end{Highlighting}
\end{Shaded}

\normalsize

Los conjuntos de datos utilizados en los ejemplos se encuentran disponibles en el siguiente enlace: \url{https://beatrizpateiro.github.io/materials/AM/data.zip}.

\chapter{Introducción al análisis multivariante}\label{introduccion-al-analisis-multivariante}

Son muchas las situaciones reales en las que necesitamos tener en cuenta varias variables de forma simultánea. Podemos pensar en problemas sencillos en todas las disciplinas en los que necesitemos analizar datos de distintas variables o características medidas sobre los individuos de una muestra.
En algunas ocasiones puede resultar adecuado estudiar cada una de las variables de interés de forma individual. Sin embargo, en general las variables están relacionadas entre sí de tal manera que los análisis individuales proporcionan poca información sobre la estructura del conjunto de datos.

\textbf{\includegraphics[width=1em,height=1em]{bookdown_AM_files/figure-latex/fa-icon-00ef1bb1d064b635425d49f06b331298.pdf} ~\textcolor{excolor}{Ejemplo.}} Un estudio pretende analizar diferentes medidas corporales del gorrión pantanero carileonado (\emph{Ammodramus caudacutus}), una especie de ave que habita en humedales de Norteamérica. Se seleccionan como variables del estudio el sexo, el tamaño de las alas, la longitud del tarso, el tamaño de la cabeza y el peso.

\textbf{\includegraphics[width=1em,height=1em]{bookdown_AM_files/figure-latex/fa-icon-00ef1bb1d064b635425d49f06b331298.pdf} ~\textcolor{excolor}{Ejemplo.}} Para llevar a cabo un estudio sobre lateralidad y aprendizaje se realizó un cuestionario a 237 estudiantes de Estadística de la Universidad de Adelaida. Entre las variables recogidas hay variables cualitativas (como el sexo, mano con la que escribe, etc.) y variables cuantitativas (anchura de la mano con la que escribe, anchura de la mano con la que no escribe, etc.).

\textbf{\includegraphics[width=1em,height=1em]{bookdown_AM_files/figure-latex/fa-icon-00ef1bb1d064b635425d49f06b331298.pdf} ~\textcolor{excolor}{Ejemplo.}} El conjunto de datos de iris de Fisher's es un conjunto de datos clásico que recoge la medida en cm. de las variables longitud y
anchura de sépalo y longitud y anchura de pétalo para flores de tres especies diferentes de iris (iris setosa, versicolor y virginica).

Se puede definir el \textbf{análisis multivariante} como el conjunto de métodos estadísticos que permiten analizar datos que surgen cuando se miden distintas variables en cada uno de los individuos de una o varias muestras.
Los avances en informática y en el procesamiento de grandes bancos de
datos ha favorecido el desarrollo del análisis multivariante en muchas disciplinas. Si bien tradicionalmente el análisis multivariante se aplicaba fundamentalmente en Biología, en la actualidad son muchos los campos de aplicación.
Las técnicas de análisis multivariante incluyen tanto métodos puramente descriptivos que tienen por objetivo extraer información de los datos disponibles, como métodos de inferencia que, a través de la construcción de modelos, pretenden obtener conclusiones sobre la población que ha generado los datos. A continuación describimos algunos de los métodos de análisis multivariante que podemos llevar a cabo cuando disponemos de medidas de un conjunto de variables.

\begin{itemize}
\tightlist
\item
  Contrastar la hipótesis de que las medias de las variables analizadas tienen un valor específico (inferencia sobre la media en poblaciones multivariantes).
\item
  Contrastar la hipótesis de que las variables son incorreladas y tienen varianza común (inferencia sobre la matriz de covarianzas).
\item
  Representar la información mediante un número menor de variables construidas como combinaciones lineales de las originales y que expliquen la mayor parte de la variabilidad original (análisis de componentes principales).
\item
  Encontrar un modelo que nos permita predecir un grupo de variables del conjunto original a partir de otro grupo de variables (modelos de regresión multivariante).
\item
  Comparar las medias de las variables en dos poblaciones (test de Hotelling).
\item
  Comparar las medias de las variables en más de dos poblaciones (MANOVA).
\item
  Si los individuos pertenecen a diferentes grupos, encontrar la combinación lineal de las variables que mejor separe dichos grupos (análisis discriminante).
\item
  Clasificar a un conjunto de observaciones en grupos homogéneos (técnicas de formación de grupos).
\end{itemize}

\section{Notación de vectores y matrices}\label{notacion-de-vectores-y-matrices}

Como hemos comentado, la información de partida en análisis multivariante es un conjunto de datos correspondiente a distintas variables medidas en los elementos de una muestra. Las observaciones correspondientes a cada individuo vendrán dadas por un vector y el conjunto de observaciones de la muestra por una matriz. Por lo tanto, la manipulación de datos multivariantes se simplificará mucho utilizando conceptos de matrices y sus propiedades.

Representamos el vector\footnote{Como es habitual en la literatura de análisis multivariante y para evitar confusión entre escalares, vectores y matrices, denotaremos los vectores con letras minúsculas en negrita y las matrices con letras mayúsculas en negrita. Según esta notación, tanto variables aleatorias como sus valores observados se representarán mediante letras minúsculas.} \(\bm{x}\in\mathbb{R}^d\) como el vector columna \(d\times 1\).

\[\bm{x}=\left(\begin{array}{c}x_1\\ x_2\\ \vdots \\ x_d \end{array}\right).\]

El vector \(d\times 1\) con todas sus componentes iguales a uno se
denotará como \(\bm{1_d}\). Representamos la matriz \(\bm{A}=(a_{ij})\) de
tamaño \(m\times n\) como

\[\bm{A}=\left(\begin{array}{cccc}a_{11}&a_{12}&\ldots&a_{1n}\\
a_{21}&a_{22}&\ldots&a_{2n}\\
\vdots&\vdots&\ddots&\vdots\\
a_{m1}&a_{m2}&\ldots&a_{mn}\end{array}\right).\]

La traspuesta de
\(\bm{A}\) se denota \(\bm{A}^\prime{}\). Si \(\bm{A}\) es una matriz
cuadrada, su determinante se denota \(\left|\bm{A}\right|\) y su traza
\(\textnormal{tr}(\bm{A})\). Si \(\bm{A}\) es no singular
(\(\left|\bm{A}\right|\neq 0\)), su inversa se denota como \(\bm{A^{-1}}\).
La matriz \(n\times n\) con elementos \(d_1,\ldots, d_n\) en la diagonal y
ceros fuera de ella se denota \(\textnormal{diag}(d_1,\ldots,d_n).\) La
matriz identidad \(n\times n\) se denota \(\bm{I_n}\). En el Anexo \ref{anexoAlg}
encontrarás algunos de los resultados fundamentales de álgebra lineal,
que nos servirán de ayuda en los próximos capítulos.

\section{Vectores aleatorios}\label{vectores-aleatorios}

Un vector aleatorio \(d\)-dimensional \(\bm{x}\in\mathbb{R}^d\) es una
colección de variables aleatorias \(\bm{x}=(x_1,\ldots,x_d)^\prime{}\)
medidas simultáneamente sobre el mismo individuo.

\textbf{\includegraphics[width=1em,height=1em]{bookdown_AM_files/figure-latex/fa-icon-00ef1bb1d064b635425d49f06b331298.pdf} ~\textcolor{excolor}{Ejemplo.}} En el estudio del gorrión pantanero carileonado
(\emph{Ammodramus caudacutus}), podemos considerar el vector aleatorio
\(5\)-dimensional \(\bm{x}=(x_1,\ldots,x_5)^\prime{}\) , siendo:

\begin{itemize}
\item
  \(x_1\)= Sexo.
\item
  \(x_2\)= Tamaño del ala.
\item
  \(x_3\)= Longitud del tarso.
\item
  \(x_4\)= Tamaño de cabeza.
\item
  \(x_5\)= Peso.
\end{itemize}

Cada una de las componentes de un vector aleatorio es una variable
aleatoria y, por lo tanto, se puede calcular su media, su varianza y su
distribución. Sin embargo, hay algunas propiedades conjuntas, como son
la covarianza o la distribución conjunta que reflejan la relación entre
las variables que forman el vector. Las medidas más importantes de las
distribuciones de probabilidad de vectores aleatorios son el vector de
medias \(\bm\mu\) y la matriz de covarianzas \(\bm\Sigma\).

Dado un vector aleatorio \(d\)-dimensional \(\bm{x}\in\mathbb{R}^d\), se
define el vector de medias como

\[{\bm{\mu}}=E({\bm{x}})=\left(\begin{array}{c}E(x_1)\\ \vdots \\ E(x_d) \end{array}\right).\]

Es decir, el vector de medias \(\bm\mu\) es el vector cuyas componentes
son las esperanzas, o medias, de las componentes del vector aleatorio.

La generalización de la varianza unidimensional al contexto
multivariante nos lleva a la definición de matriz de covarianzas. Dado
un vector aleatorio \(\bm{x}\in\mathbb{R}^d\) con vector de medias
\(\bm\mu\), se define la matriz de covarianzas \(\bm\Sigma\) como

\[\label{Sigma}
\bm\Sigma=E((\bm{x}-\bm\mu)(\bm{x}-\bm\mu)^\prime{}).\]

Así definida, la matriz de covarianzas \(\bm\Sigma\) es una matriz de
dimensión \(d\times d\), simétrica y semidefinida positiva. Podemos
comprobar que además los elementos de la diagonal son las varianzas de
cada una de las componentes del vector aleatorio y en cada entrada
\((j,k)\) se tiene la covarianza entre las variables \(x_j\) y \(x_k\),

\[\textnormal{Cov}(x_j,x_k)=\sigma_{jk}=E((x_j-E(x_j))(x_k-E(x_k))).\]

Fíjate que

\[\textnormal{Cov}(x_j,x_j)=\sigma_{jj}=E((x_j-E(x_j))^2)\]
representa la varianza de la variable \(x_j\), que denotaremos también
como \(\sigma_j^2\).En resumen,

\[\bm\Sigma=\left(\begin{array}{cccc}\textnormal{Var}(x_1)&\textnormal{Cov}(x_1,x_2)&\ldots&\textnormal{Cov}(x_1,x_d)\\
\textnormal{Cov}(x_2,x_1)&\textnormal{Var}(x_2)&\ldots&\textnormal{Cov}(x_2,x_d)\\
\vdots&\vdots&\ddots&\vdots\\
\textnormal{Cov}(x_d,x_1)&\textnormal{Cov}(x_d,x_2)&\ldots&\textnormal{Var}(x_d)\end{array}\right)=\left(\begin{array}{cccc}\sigma_1^2&\sigma_{12}&\ldots&\sigma_{1d}\\
\sigma_{21}&\sigma_{2}^2&\ldots&\sigma_{2d}\\
\vdots&\vdots&\ddots&\vdots\\
\sigma_{d1}&\sigma_{d2}&\ldots&\sigma_{d}^2\end{array}\right).\]

Puedes también demostrar como ejercicio que \(\bm\Sigma\) puede ser
expresada como

\[\bm\Sigma=E(\bm{x}\bm{x}^\prime)-\bm\mu\bm\mu^\prime.\]

Es difícil comparar covarianzas entre pares de variables ya que estas
dependen de la escala de medida. Para obtener una medida de la relación
lineal entre pares de variables que sea invariante ante cambios de
escala podemos estandarizar la covarianza dividiendo por el producto de
desviaciones típicas de las variables involucradas. Esta covarianza
estandarizada es lo que se denomina correlación. Así, la correlación
entre las variables \(x_i\) y \(x_j\) se calcula como

\[\rho_{ij}=\frac{\sigma_{ij}}{\sigma_i\sigma_j}.\]

La matriz de correlaciones será por lo tanto

\[\bm\rho=\left(\begin{array}{cccc}1&\rho_{12}&\ldots&\rho_{1d}\\
\rho_{21}&1&\ldots&\rho_{2d}\\
\vdots&\vdots&\ddots&\vdots\\
\rho_{d1}&\rho_{d2}&\ldots&1\end{array}\right).\]

Se verifica

\[\bm\rho=\bm{D_\Sigma}^{-1}\bm{\Sigma}\bm{D_\Sigma}^{-1}\]
siendo
\(\bm{D_\Sigma}\) la matriz diagonal de orden \(d\) construida colocando en
la diagonal principal las desviaciones típicas de las variables, es
decir, \(\bm{D_\Sigma}=\textnormal{diag}(\sigma_1,\ldots,\sigma_d)\).
Equivalentemente,

\[\bm{\Sigma}=\bm{D_\Sigma}\bm\rho\bm{D_\Sigma}.\]

\section{Representación de datos multivariantes}\label{representacion-de-datos-multivariantes}

En el contexto multivariante es conveniente representar las variables
aleatorias (o sus valores) en forma matricial. Supongamos que hemos
observado \(d\) variables en un conjunto de \(n\) elementos. Cada una de
estas \(d\) variables es una variable univariante y el conjunto de las \(d\)
variables constituye un vector aleatorio. Los valores del vector
aleatorio en cada uno de los \(n\) elementos pueden representarse mediante
una matriz de datos \(\bm{X}\) de tamaño \(n\times d\) de manera que
\(\bm{X}=(x_{ij})\), donde el índice \(i=1\ldots,n\) representa el individuo
y el índice \(j =1,\ldots,d\) representa la variable. Es decir, cada fila
de la matriz representa el valor del vector aleatorio (formado por las \(d\) variables)
sobre el individuo \(i\).

\[\bm{X}=\left(\begin{array}{cccc}x_{11}&x_{12}&\ldots&x_{1d}\\
x_{21}&x_{22}&\ldots&x_{2d}\\
\vdots&\vdots&\ddots&\vdots\\
x_{n1}&x_{n2}&\ldots&x_{nd}\end{array}\right)=\left(\begin{array}{c}\bm{x_1}^\prime{}\\ \bm{x_2}^\prime{}\\ \vdots \\ \bm{x_n}^\prime{} \end{array}\right).\]

\textbf{\includegraphics[width=1em,height=1em]{bookdown_AM_files/figure-latex/fa-icon-00ef1bb1d064b635425d49f06b331298.pdf} ~\textcolor{excolor}{Ejemplo.}} Se analizaron las variables del estudio del gorrión
pantanero carileonado en una muestra de 979 gorriones de dicha especie.

Los datos se encuentran recogidos en el fichero \texttt{gorrion.txt} y se
pueden leer en R como se muestra a continuación.

\footnotesize

\begin{Shaded}
\begin{Highlighting}[]
\SpecialCharTok{\textgreater{}}\NormalTok{ gorrion }\OtherTok{\textless{}{-}} \FunctionTok{read.table}\NormalTok{(}\StringTok{"data/gorrion.txt"}\NormalTok{, }\AttributeTok{header =} \ConstantTok{TRUE}\NormalTok{)}
\SpecialCharTok{\textgreater{}}\NormalTok{ gorrion}
\NormalTok{       Sex Wingcrd Tarsus Head   Wt}
\DecValTok{1}\NormalTok{     Male    }\FloatTok{58.0}   \FloatTok{21.7} \FloatTok{32.7} \FloatTok{20.3}
\DecValTok{2}\NormalTok{   Female    }\FloatTok{56.5}   \FloatTok{21.1} \FloatTok{31.4} \FloatTok{17.4}
\DecValTok{3}\NormalTok{     Male    }\FloatTok{59.0}   \FloatTok{21.0} \FloatTok{33.3} \FloatTok{21.0}
\DecValTok{4}\NormalTok{     Male    }\FloatTok{59.0}   \FloatTok{21.3} \FloatTok{32.5} \FloatTok{21.0}
\NormalTok{...}
\end{Highlighting}
\end{Shaded}

\normalsize

El conjunto de datos se puede representar mediante una matriz de 979
filas (número de individuos en la muestra) y 5 columnas (número de
variables).

\footnotesize

\begin{Shaded}
\begin{Highlighting}[]
\SpecialCharTok{\textgreater{}} \FunctionTok{class}\NormalTok{(gorrion)}
\NormalTok{[}\DecValTok{1}\NormalTok{] }\StringTok{"data.frame"}
\SpecialCharTok{\textgreater{}} \FunctionTok{dim}\NormalTok{(gorrion)}
\NormalTok{[}\DecValTok{1}\NormalTok{] }\DecValTok{979}   \DecValTok{5}
\end{Highlighting}
\end{Shaded}

\normalsize

El objeto utilizado en R para almacenar datos de este tipo es un
data.frame. Podremos editar los datos fácilmente utilizando el comando \texttt{fix} o
\texttt{edit}.

A continuación se muestra parte de uno de los muchos conjuntos de datos
multivariantes que puedes encontrar en R. Se trata del conjunto de datos
\texttt{survey} (librería MASS).

\textbf{\includegraphics[width=1em,height=1em]{bookdown_AM_files/figure-latex/fa-icon-00ef1bb1d064b635425d49f06b331298.pdf} ~\textcolor{excolor}{Ejemplo.}} El conjunto de datos \texttt{survey} (librería MASS) corresponde a
un cuestionario que se realizó a 237 estudiantes de Estadística de la
Universidad de Adelaida. Entre las variables recogidas hay variables
cualitativas (como el sexo, mano con la que escribe, etc.) y variables
cuantitativas (anchura de la mano con la que escribe, anchura de la mano
con la que no escribe, etc.) Puedes consultar más información sobre este
conjunto de datos escribiendo \texttt{help(survey)} en la consola de comandos.

\footnotesize

\begin{Shaded}
\begin{Highlighting}[]
\SpecialCharTok{\textgreater{}} \FunctionTok{library}\NormalTok{(MASS)}
\SpecialCharTok{\textgreater{}} \FunctionTok{data}\NormalTok{(survey)}
\SpecialCharTok{\textgreater{}}\NormalTok{ survey}
\NormalTok{       Sex Wr.Hnd NW.Hnd W.Hnd    Fold Pulse    Clap Exer Smoke Height      M.I}
\DecValTok{1}\NormalTok{   Female   }\FloatTok{18.5}   \FloatTok{18.0}\NormalTok{ Right  R on L    }\DecValTok{92}\NormalTok{    Left Some Never }\FloatTok{173.00}\NormalTok{   Metric}
\DecValTok{2}\NormalTok{     Male   }\FloatTok{19.5}   \FloatTok{20.5}\NormalTok{  Left  R on L   }\DecValTok{104}\NormalTok{    Left None Regul }\FloatTok{177.80}\NormalTok{ Imperial}
\DecValTok{3}\NormalTok{     Male   }\FloatTok{18.0}   \FloatTok{13.3}\NormalTok{ Right  L on R    }\DecValTok{87}\NormalTok{ Neither None Occas     }\ConstantTok{NA}     \SpecialCharTok{\textless{}}\ConstantTok{NA}\SpecialCharTok{\textgreater{}}
\DecValTok{4}\NormalTok{     Male   }\FloatTok{18.8}   \FloatTok{18.9}\NormalTok{ Right  R on L    }\ConstantTok{NA}\NormalTok{ Neither None Never }\FloatTok{160.00}\NormalTok{   Metric}
\DecValTok{5}\NormalTok{     Male   }\FloatTok{20.0}   \FloatTok{20.0}\NormalTok{ Right Neither    }\DecValTok{35}\NormalTok{   Right Some Never }\FloatTok{165.00}\NormalTok{   Metric}
\DecValTok{6}\NormalTok{   Female   }\FloatTok{18.0}   \FloatTok{17.7}\NormalTok{ Right  L on R    }\DecValTok{64}\NormalTok{   Right Some Never }\FloatTok{172.72}\NormalTok{ Imperial}
\DecValTok{7}\NormalTok{     Male   }\FloatTok{17.7}   \FloatTok{17.7}\NormalTok{ Right  L on R    }\DecValTok{83}\NormalTok{   Right Freq Never }\FloatTok{182.88}\NormalTok{ Imperial}
\DecValTok{8}\NormalTok{   Female   }\FloatTok{17.0}   \FloatTok{17.3}\NormalTok{ Right  R on L    }\DecValTok{74}\NormalTok{   Right Freq Never }\FloatTok{157.00}\NormalTok{   Metric}
\DecValTok{9}\NormalTok{     Male   }\FloatTok{20.0}   \FloatTok{19.5}\NormalTok{ Right  R on L    }\DecValTok{72}\NormalTok{   Right Some Never }\FloatTok{175.00}\NormalTok{   Metric}
\NormalTok{...}
\end{Highlighting}
\end{Shaded}

\normalsize

\subsection{Visualización de datos multivariantes}\label{visualizacion-de-datos-multivariantes}

Uno de los problemas con que nos encontramos a la hora de establecer
métodos gráficos para describir datos multivariantes es nuestro propio
sistema de percepción visual. Los gráficos de dispersión en dos
dimensiones son fáciles de entender. Somos también capaces de
interpretar gráficos de dispersión en tres dimensiones. Sin embargo,
cuando la dimensión es mayor necesitamos otros procedimientos para
visualizar los datos.

\textbf{\includegraphics[width=1em,height=1em]{bookdown_AM_files/figure-latex/fa-icon-00ef1bb1d064b635425d49f06b331298.pdf} ~\textcolor{excolor}{Ejemplo.}} Diagrama de dispersión para los datos de tamaño del ala
(Wingcrd), longitud de tarso (Tarsus) y tamaño de la cabeza (Head) de
gorrión. La función \texttt{scatterplot3d} de la librería \texttt{scatterplot3d} R
permite hacer un gráfico como el que se muestra en la Figura \ref{fig:scat3d}

\footnotesize

\begin{Shaded}
\begin{Highlighting}[]
\SpecialCharTok{\textgreater{}} \FunctionTok{library}\NormalTok{(scatterplot3d)}
\SpecialCharTok{\textgreater{}} \FunctionTok{scatterplot3d}\NormalTok{(gorrion[, }\DecValTok{2}\SpecialCharTok{:}\DecValTok{4}\NormalTok{], }\AttributeTok{color =} \DecValTok{4}\NormalTok{, }\AttributeTok{main =} \StringTok{"Diagrama de dispersión 3D."}\NormalTok{)}
\end{Highlighting}
\end{Shaded}

\begin{figure}

{\centering \resizebox{.4\linewidth}{!}{\includegraphics[width=0.5\linewidth]{bookdown_AM_files/figure-latex/scat3d-1} } 

}

\caption{Diagrama de dispersión 3D.}\label{fig:scat3d}
\end{figure}
\normalsize

También se pueden crear diagramas de dispersión en 3D interactivos
usando al función \texttt{plot3d} de la librería \texttt{rgl}, como se indica a
continuación:

\footnotesize

\begin{Shaded}
\begin{Highlighting}[]
\SpecialCharTok{\textgreater{}} \FunctionTok{library}\NormalTok{(rgl)}
\SpecialCharTok{\textgreater{}} \FunctionTok{plot3d}\NormalTok{(gorrion[, }\DecValTok{2}\SpecialCharTok{:}\DecValTok{4}\NormalTok{], }\AttributeTok{col =} \DecValTok{4}\NormalTok{)}
\end{Highlighting}
\end{Shaded}

\normalsize

Una forma de representar datos multivariantes es mediante \textbf{matrices de
diagramas de dispersión}. Se construye una cuadrícula con tantas filas
y columnas como variables. En cada casilla se representa el diagrama de
dispersión de las variables de la fila y columna correspondientes. Este
gráfico informa por lo tanto de cómo son las relaciones entre variables,
pero sólo dos a dos. En la Figura \ref{fig:scat} se muestran matrices de diagramas de dispersión
para los datos de gorrión.

\textbf{\includegraphics[width=1em,height=1em]{bookdown_AM_files/figure-latex/fa-icon-00ef1bb1d064b635425d49f06b331298.pdf} ~\textcolor{excolor}{Ejemplo.}} Matrices de diagramas de dispersión para los datos de
gorrión. La función \texttt{pairs} de R permite hacer un gráfico como el que se
muestra en la izquierda.

La función \texttt{scatterplotMatrix} de la librería \texttt{car} permite además
añadir en la diagonal los histogramas, boxplots, estimadores de la
densidad, etc. de las variables del conjunto de datos. También permite
añadir en cada gráfico el ajuste de regresión lineal o no paramétrico
para cada par de variables. Consulta la ayuda de la función para obtener
más información.

\footnotesize

\begin{Shaded}
\begin{Highlighting}[]
\SpecialCharTok{\textgreater{}} \FunctionTok{pairs}\NormalTok{(gorrion[, }\DecValTok{2}\SpecialCharTok{:}\DecValTok{5}\NormalTok{])}
\SpecialCharTok{\textgreater{}} \FunctionTok{library}\NormalTok{(car)}
\SpecialCharTok{\textgreater{}} \FunctionTok{scatterplotMatrix}\NormalTok{(gorrion[, }\DecValTok{2}\SpecialCharTok{:}\DecValTok{5}\NormalTok{], }\AttributeTok{diagonal =} \FunctionTok{list}\NormalTok{(}\AttributeTok{method =} \StringTok{"histogram"}\NormalTok{, }\AttributeTok{breaks =} \StringTok{"FD"}\NormalTok{),}
\SpecialCharTok{+}                   \AttributeTok{smooth =} \ConstantTok{FALSE}\NormalTok{, }\AttributeTok{regLine =} \ConstantTok{TRUE}\NormalTok{)}
\end{Highlighting}
\end{Shaded}

\begin{figure}

{\centering \resizebox{.4\linewidth}{!}{\includegraphics[width=0.5\linewidth]{bookdown_AM_files/figure-latex/scat-1} } \resizebox{.4\linewidth}{!}{\includegraphics[width=0.5\linewidth]{bookdown_AM_files/figure-latex/scat-2} } 

}

\caption{Matrices de diagramas de dispersión para los datos de gorrión.}\label{fig:scat}
\end{figure}
\normalsize

Cuando en un conjunto de datos tenemos diferentes variables categóricas,
puede ser interesante visualizar las relaciones entre pares de variables
teniendo en cuenta los distintos factores de las variables categóricas.
Podemos en este caso representar \textbf{gráficos de dispersión
condicionales} como el que se muestra en la Figura \ref{fig:con} para
los datos \texttt{survey}.

\textbf{\includegraphics[width=1em,height=1em]{bookdown_AM_files/figure-latex/fa-icon-00ef1bb1d064b635425d49f06b331298.pdf} ~\textcolor{excolor}{Ejemplo.}} Gráfico de dispersión condicional para los datos \texttt{survey}.
La función \texttt{coplot} nos permite visualizar la relación entre el ancho de
la mano con la que se escribe (variable \texttt{Wr.Hnd}) y el ancho de la mano
con la que no se escribe (variable \texttt{NW.Hnd}), en función del sexo
(variable \texttt{Sex}) y de la lateralidad (variable \texttt{W.Hnd}) del individuo.

\footnotesize

\begin{Shaded}
\begin{Highlighting}[]
\SpecialCharTok{\textgreater{}} \FunctionTok{coplot}\NormalTok{(Wr.Hnd }\SpecialCharTok{\textasciitilde{}}\NormalTok{ NW.Hnd }\SpecialCharTok{|}\NormalTok{ Sex }\SpecialCharTok{*}\NormalTok{ W.Hnd, }\AttributeTok{data =} \FunctionTok{na.omit}\NormalTok{(survey))}
\end{Highlighting}
\end{Shaded}

\begin{figure}

{\centering \resizebox{.4\linewidth}{!}{\includegraphics[width=0.5\linewidth]{bookdown_AM_files/figure-latex/con-1} } 

}

\caption{Gráfico de dispersión condicional para los datos de lateralidad.}\label{fig:con}
\end{figure}
\normalsize

Otra herramienta utilizada para representar datos multivariantes son los
\textbf{gráficos de estrellas}. Cada individuo se representa en una estrella,
con tantos rayos o ejes como variables queramos representar. Cada eje
representa el valor de la variable re-escalada de manera independiente
entre variables. Se muestra en la Figura \ref{fig:star} el
gráfico de estrellas para los datos de iris. A la vista de los gráficos
de estrella podemos identificar perfectamente las diferentes especies de
iris, representadas a su vez con distintos colores. Como conclusión,
parece que las medidas consideradas (longitud y anchura de pétalo y
sépalo) nos servirían para distinguir unas especies de otras.

\textbf{\includegraphics[width=1em,height=1em]{bookdown_AM_files/figure-latex/fa-icon-00ef1bb1d064b635425d49f06b331298.pdf} ~\textcolor{excolor}{Ejemplo.}} Gráfico de estrellas para los datos \texttt{iris} obtenido con la
función \texttt{stars} de R. Cada estrella representa una flor.

\footnotesize

\begin{Shaded}
\begin{Highlighting}[]
\SpecialCharTok{\textgreater{}} \FunctionTok{data}\NormalTok{(iris)}
\SpecialCharTok{\textgreater{}} \FunctionTok{stars}\NormalTok{(iris[, }\DecValTok{1}\SpecialCharTok{:}\DecValTok{4}\NormalTok{], }\AttributeTok{col.stars =}\NormalTok{ iris}\SpecialCharTok{$}\NormalTok{Species, }\AttributeTok{key.loc =} \FunctionTok{c}\NormalTok{(}\DecValTok{20}\NormalTok{, }\DecValTok{1}\NormalTok{))}
\end{Highlighting}
\end{Shaded}

\begin{figure}

{\centering \resizebox{.4\linewidth}{!}{\includegraphics[width=0.5\linewidth]{bookdown_AM_files/figure-latex/star-1} } 

}

\caption{Gráfico de estrellas para los datos de iris.}\label{fig:star}
\end{figure}
\normalsize

También es común representar datos multivariantes mediante \textbf{gráficos de
caras de Chernoff}. El gráfico de caras de Chernoff es como un gráfico
de estrellas, pero cada individuo ahora se representa en una cara y las
variables en los rasgos físicos (altura y anchura de la cara, forma de
la cara, altura y anchura de la boca, curva de la sonrisa, altura y
anchura de ojos, ...). En la Figura
\ref{fig:cara} se
representa el gráfico de caras de Chernoff para un subconjunto de los
datos de iris.

\textbf{\includegraphics[width=1em,height=1em]{bookdown_AM_files/figure-latex/fa-icon-00ef1bb1d064b635425d49f06b331298.pdf} ~\textcolor{excolor}{Ejemplo.}} Gráficos de caras de Chernoff para un subconjunto de los
datos \texttt{iris} obtenido con la función \texttt{faces} de la librería \texttt{aplpack} de
R. Cada cara representa una flor. Elegimos 4 flores de cada una de las
especies (setosa, versicolor, virginica).

\footnotesize

\begin{Shaded}
\begin{Highlighting}[]
\SpecialCharTok{\textgreater{}} \FunctionTok{library}\NormalTok{(aplpack)}
\SpecialCharTok{\textgreater{}} \FunctionTok{faces}\NormalTok{(iris[}\FunctionTok{c}\NormalTok{(}\DecValTok{1}\SpecialCharTok{:}\DecValTok{4}\NormalTok{, }\DecValTok{51}\SpecialCharTok{:}\DecValTok{54}\NormalTok{, }\DecValTok{101}\SpecialCharTok{:}\DecValTok{104}\NormalTok{), }\DecValTok{1}\SpecialCharTok{:}\DecValTok{4}\NormalTok{], }\AttributeTok{face.type =} \DecValTok{0}\NormalTok{)}
\NormalTok{effect of variables}\SpecialCharTok{:}
\NormalTok{ modified item       Var           }
 \StringTok{"height of face   "} \StringTok{"Sepal.Length"}
 \StringTok{"width of face    "} \StringTok{"Sepal.Width"} 
 \StringTok{"structure of face"} \StringTok{"Petal.Length"}
 \StringTok{"height of mouth  "} \StringTok{"Petal.Width"} 
 \StringTok{"width of mouth   "} \StringTok{"Sepal.Length"}
 \StringTok{"smiling          "} \StringTok{"Sepal.Width"} 
 \StringTok{"height of eyes   "} \StringTok{"Petal.Length"}
 \StringTok{"width of eyes    "} \StringTok{"Petal.Width"} 
 \StringTok{"height of hair   "} \StringTok{"Sepal.Length"}
 \StringTok{"width of hair   "}  \StringTok{"Sepal.Width"} 
 \StringTok{"style of hair   "}  \StringTok{"Petal.Length"}
 \StringTok{"height of nose  "}  \StringTok{"Petal.Width"} 
 \StringTok{"width of nose   "}  \StringTok{"Sepal.Length"}
 \StringTok{"width of ear    "}  \StringTok{"Sepal.Width"} 
 \StringTok{"height of ear   "}  \StringTok{"Petal.Length"}
\end{Highlighting}
\end{Shaded}

\begin{figure}

{\centering \resizebox{.4\linewidth}{!}{\includegraphics[width=0.5\linewidth]{bookdown_AM_files/figure-latex/cara-1} } 

}

\caption{Gráfico de caras de Chernoff para un subconjunto de los datos de iris.}\label{fig:cara}
\end{figure}
\normalsize

\subsection{Vector de medias y matriz de covarianzas muestrales}\label{vector-de-medias-y-matriz-de-covarianzas-muestrales}

Supongamos que disponemos de una muestra aleatoria
\(\bm{x_1},\ldots, \bm{x_n}\) de un vector aleatorio en \(\mathbb{R}^d\) con
media \(\bm\mu\) y matriz de covarianzas \(\bm\Sigma\). De forma análoga al
caso univariante, calcularemos el vector de medias muestral como
\[\overline{\bm{x}}=\frac{1}{n}\sum_{i=1}^n{\bm{x_i}}=\frac{1}{n}\bm{X}^\prime{}\bm{1_n}.\]
Recuerda que, según la notación establecida,
\(\bm{x_i}=(x_{i1},\ldots,x_{id})^\prime{}\) y \(\bm{X}\) es la matriz de
datos.

\textbf{\includegraphics[width=1em,height=1em]{bookdown_AM_files/figure-latex/fa-icon-00ef1bb1d064b635425d49f06b331298.pdf} ~\textcolor{excolor}{Ejemplo.}} Trabajaremos ahora con variables numéricas del estudio del
gorrión pantanero carileonado.

\footnotesize

\begin{Shaded}
\begin{Highlighting}[]
\SpecialCharTok{\textgreater{}}\NormalTok{ gmed }\OtherTok{\textless{}{-}}\NormalTok{ gorrion[, }\DecValTok{2}\SpecialCharTok{:}\DecValTok{5}\NormalTok{]}
\SpecialCharTok{\textgreater{}}\NormalTok{ gmed}
\NormalTok{    Wingcrd Tarsus Head   Wt}
\DecValTok{1}      \FloatTok{58.0}   \FloatTok{21.7} \FloatTok{32.7} \FloatTok{20.3}
\DecValTok{2}      \FloatTok{56.5}   \FloatTok{21.1} \FloatTok{31.4} \FloatTok{17.4}
\DecValTok{3}      \FloatTok{59.0}   \FloatTok{21.0} \FloatTok{33.3} \FloatTok{21.0}
\DecValTok{4}      \FloatTok{59.0}   \FloatTok{21.3} \FloatTok{32.5} \FloatTok{21.0}
\NormalTok{...}
\end{Highlighting}
\end{Shaded}

\normalsize

Calculamos el vector de medias muestral a partir de la matriz de
observaciones. Podemos hacerlo de varias maneras. A través de la función
\texttt{summary}.

\footnotesize

\begin{Shaded}
\begin{Highlighting}[]
\SpecialCharTok{\textgreater{}} \FunctionTok{summary}\NormalTok{(gmed)}
\NormalTok{    Wingcrd          Tarsus           Head             Wt       }
\NormalTok{ Min.   }\SpecialCharTok{:}\FloatTok{52.00}\NormalTok{   Min.   }\SpecialCharTok{:}\FloatTok{18.90}\NormalTok{   Min.   }\SpecialCharTok{:}\FloatTok{29.20}\NormalTok{   Min.   }\SpecialCharTok{:}\FloatTok{15.60}  
 \DecValTok{1}\NormalTok{st Qu.}\SpecialCharTok{:}\FloatTok{56.00}   \DecValTok{1}\NormalTok{st Qu.}\SpecialCharTok{:}\FloatTok{20.90}   \DecValTok{1}\NormalTok{st Qu.}\SpecialCharTok{:}\FloatTok{31.50}   \DecValTok{1}\NormalTok{st Qu.}\SpecialCharTok{:}\FloatTok{19.20}  
\NormalTok{ Median }\SpecialCharTok{:}\FloatTok{58.00}\NormalTok{   Median }\SpecialCharTok{:}\FloatTok{21.40}\NormalTok{   Median }\SpecialCharTok{:}\FloatTok{31.90}\NormalTok{   Median }\SpecialCharTok{:}\FloatTok{20.20}  
\NormalTok{ Mean   }\SpecialCharTok{:}\FloatTok{57.87}\NormalTok{   Mean   }\SpecialCharTok{:}\FloatTok{21.48}\NormalTok{   Mean   }\SpecialCharTok{:}\FloatTok{32.04}\NormalTok{   Mean   }\SpecialCharTok{:}\FloatTok{20.23}  
 \DecValTok{3}\NormalTok{rd Qu.}\SpecialCharTok{:}\FloatTok{59.00}   \DecValTok{3}\NormalTok{rd Qu.}\SpecialCharTok{:}\FloatTok{21.90}   \DecValTok{3}\NormalTok{rd Qu.}\SpecialCharTok{:}\FloatTok{32.40}   \DecValTok{3}\NormalTok{rd Qu.}\SpecialCharTok{:}\FloatTok{21.10}  
\NormalTok{ Max.   }\SpecialCharTok{:}\FloatTok{68.00}\NormalTok{   Max.   }\SpecialCharTok{:}\FloatTok{25.40}\NormalTok{   Max.   }\SpecialCharTok{:}\FloatTok{35.70}\NormalTok{   Max.   }\SpecialCharTok{:}\FloatTok{26.50}  
\end{Highlighting}
\end{Shaded}

\normalsize

A través de la función \texttt{apply}.

\footnotesize

\begin{Shaded}
\begin{Highlighting}[]
\SpecialCharTok{\textgreater{}} \FunctionTok{apply}\NormalTok{(gmed, }\DecValTok{2}\NormalTok{, mean)}
\NormalTok{ Wingcrd   Tarsus     Head       Wt }
\FloatTok{57.86599} \FloatTok{21.47620} \FloatTok{32.03790} \FloatTok{20.22646} 
\end{Highlighting}
\end{Shaded}

\normalsize

A través de la matriz de observaciones.

\footnotesize

\begin{Shaded}
\begin{Highlighting}[]
\SpecialCharTok{\textgreater{}}\NormalTok{ n }\OtherTok{\textless{}{-}} \FunctionTok{dim}\NormalTok{(gmed)[}\DecValTok{1}\NormalTok{]}
\SpecialCharTok{\textgreater{}}\NormalTok{ un }\OtherTok{\textless{}{-}} \FunctionTok{matrix}\NormalTok{(}\DecValTok{1}\NormalTok{, }\AttributeTok{nr =}\NormalTok{ n, }\AttributeTok{nc =} \DecValTok{1}\NormalTok{)}
\SpecialCharTok{\textgreater{}} \FunctionTok{t}\NormalTok{(gmed) }\SpecialCharTok{\%*\%}\NormalTok{ un}\SpecialCharTok{/}\NormalTok{n}
\NormalTok{            [,}\DecValTok{1}\NormalTok{]}
\NormalTok{Wingcrd }\FloatTok{57.86599}
\NormalTok{Tarsus  }\FloatTok{21.47620}
\NormalTok{Head    }\FloatTok{32.03790}
\NormalTok{Wt      }\FloatTok{20.22646}
\end{Highlighting}
\end{Shaded}

\normalsize

La matriz de covarianzas \(\bm\Sigma\) se estimará mediante

\[\bm{S}=\frac{1}{n}\sum_{i=1}^n\left(\bm{x_i}-\overline{\bm{x}}\right)\left(\bm{x_i}-\overline{\bm{x}}\right)^\prime{}.\]
La matriz de covarianzas muestral \(\bm{S}\) es una matriz cuadrada y
simétrica que contiene en la diagonal las varianzas muestrales y fuera
de la diagonal las covarianzas muestrales entre los distintos pares de
variables. Es decir

\[\bm{S}=\left(\begin{array}{cccc}s_1^2&s_{12}&\ldots&s_{1d}\\
s_{21}&s_{2}^2&\ldots&s_{2d}\\
\vdots&\vdots&\ddots&\vdots\\
s_{d1}&s_{d2}&\ldots&s_{d}^2\end{array}\right),\]
siendo

\[s_{jk}=\frac{1}{n}\sum_{i=1}^n(x_{ij}-\overline{x}_j)(x_{ik}-\overline{x}_k)\]
Las varianzas muestrales \(s_{jj}\) se denotan \(s_j^2\). La matriz de
correlaciones muestral \(\bm R\) se obtendrá dividiendo las covarianzas
muestrales por las desviaciones típicas de cada par de variables. Es
decir,

\[\bm{R}=\left(\begin{array}{cccc}1&r_{12}&\ldots&r_{1d}\\
r_{21}&1&\ldots&r_{2d}\\
\vdots&\vdots&\ddots&\vdots\\
r_{d1}&r_{d2}&\ldots&1\end{array}\right),\]
siendo

\[r_{jk}=\frac{s_{jk}}{s_js_k}.\]

De forma análoga al caso poblacional,

\[\bm{R}=\bm{D_S}^{-1}\bm{S}\bm{D_S}^{-1}\]
siendo ahora \(\bm{D_S}=\textnormal{diag}(s_1,\ldots,s_d)\), es decir, la
matriz diagonal de orden \(d\) construida colocando en la diagonal
principal las desviaciones típicas muestrales de las variables. Equivalentemente,

\[\bm{S}=\bm{D_S}\bm{R}\bm{D_S}.\]

\textbf{\includegraphics[width=1em,height=1em]{bookdown_AM_files/figure-latex/fa-icon-00ef1bb1d064b635425d49f06b331298.pdf} ~\textcolor{excolor}{Ejemplo.}} Volvemos al conjunto de datos de iris. Calculamos, para la
especie setosa, el vector de medias muestral, la matriz de covarianzas
muestral y la matriz de correlaciones muestral. Debemos tener en cuenta
que R calcula la matriz de covarianzas muestral corregida
\[\frac{1}{n-1}\sum_{i=1}^n\left(\bm{x_i}-\overline{\bm{x}}\right)\left(\bm{x_i}-\overline{\bm{x}}\right)^\prime{}.\]

\footnotesize

\begin{Shaded}
\begin{Highlighting}[]
\SpecialCharTok{\textgreater{}}\NormalTok{ setosa }\OtherTok{\textless{}{-}}\NormalTok{ iris[iris}\SpecialCharTok{$}\NormalTok{Species }\SpecialCharTok{==} \StringTok{"setosa"}\NormalTok{, }\DecValTok{1}\SpecialCharTok{:}\DecValTok{4}\NormalTok{]}
\SpecialCharTok{\textgreater{}} \FunctionTok{colMeans}\NormalTok{(setosa)}
\NormalTok{Sepal.Length  Sepal.Width Petal.Length  Petal.Width }
       \FloatTok{5.006}        \FloatTok{3.428}        \FloatTok{1.462}        \FloatTok{0.246} 
\SpecialCharTok{\textgreater{}} \FunctionTok{cov}\NormalTok{(setosa)  }\CommentTok{\# Matriz de covarianzas muestral corregida}
\NormalTok{             Sepal.Length Sepal.Width Petal.Length Petal.Width}
\NormalTok{Sepal.Length   }\FloatTok{0.12424898} \FloatTok{0.099216327}  \FloatTok{0.016355102} \FloatTok{0.010330612}
\NormalTok{Sepal.Width    }\FloatTok{0.09921633} \FloatTok{0.143689796}  \FloatTok{0.011697959} \FloatTok{0.009297959}
\NormalTok{Petal.Length   }\FloatTok{0.01635510} \FloatTok{0.011697959}  \FloatTok{0.030159184} \FloatTok{0.006069388}
\NormalTok{Petal.Width    }\FloatTok{0.01033061} \FloatTok{0.009297959}  \FloatTok{0.006069388} \FloatTok{0.011106122}
\SpecialCharTok{\textgreater{}} \FunctionTok{cor}\NormalTok{(setosa)}
\NormalTok{             Sepal.Length Sepal.Width Petal.Length Petal.Width}
\NormalTok{Sepal.Length    }\FloatTok{1.0000000}   \FloatTok{0.7425467}    \FloatTok{0.2671758}   \FloatTok{0.2780984}
\NormalTok{Sepal.Width     }\FloatTok{0.7425467}   \FloatTok{1.0000000}    \FloatTok{0.1777000}   \FloatTok{0.2327520}
\NormalTok{Petal.Length    }\FloatTok{0.2671758}   \FloatTok{0.1777000}    \FloatTok{1.0000000}   \FloatTok{0.3316300}
\NormalTok{Petal.Width     }\FloatTok{0.2780984}   \FloatTok{0.2327520}    \FloatTok{0.3316300}   \FloatTok{1.0000000}
\end{Highlighting}
\end{Shaded}

\normalsize

Comprobamos que la matriz de correlaciones muestral también se puede
calcular como \(\bm{R}=\bm{D_S}^{-1}\bm{S}\bm{D_S}^{-1}\).

\footnotesize

\begin{Shaded}
\begin{Highlighting}[]
\SpecialCharTok{\textgreater{}}\NormalTok{ desviaciones }\OtherTok{\textless{}{-}} \FunctionTok{apply}\NormalTok{(setosa, }\DecValTok{2}\NormalTok{, sd)}
\SpecialCharTok{\textgreater{}}\NormalTok{ Dinv }\OtherTok{\textless{}{-}} \FunctionTok{solve}\NormalTok{(}\FunctionTok{diag}\NormalTok{(desviaciones))}
\SpecialCharTok{\textgreater{}}\NormalTok{ Dinv }\SpecialCharTok{\%*\%} \FunctionTok{cov}\NormalTok{(setosa) }\SpecialCharTok{\%*\%}\NormalTok{ Dinv}
\NormalTok{          [,}\DecValTok{1}\NormalTok{]      [,}\DecValTok{2}\NormalTok{]      [,}\DecValTok{3}\NormalTok{]      [,}\DecValTok{4}\NormalTok{]}
\NormalTok{[}\DecValTok{1}\NormalTok{,] }\FloatTok{1.0000000} \FloatTok{0.7425467} \FloatTok{0.2671758} \FloatTok{0.2780984}
\NormalTok{[}\DecValTok{2}\NormalTok{,] }\FloatTok{0.7425467} \FloatTok{1.0000000} \FloatTok{0.1777000} \FloatTok{0.2327520}
\NormalTok{[}\DecValTok{3}\NormalTok{,] }\FloatTok{0.2671758} \FloatTok{0.1777000} \FloatTok{1.0000000} \FloatTok{0.3316300}
\NormalTok{[}\DecValTok{4}\NormalTok{,] }\FloatTok{0.2780984} \FloatTok{0.2327520} \FloatTok{0.3316300} \FloatTok{1.0000000}
\end{Highlighting}
\end{Shaded}

\normalsize

Algunos de los resultados que veremos irán encaminados a la obtención de
la distribución de \(\overline{\bm{x}}\) y \(\bm{S}\) en este contexto
multidimensional y bajo la hipótesis de normalidad.

\section{Vector de medias y matrices de covarianzas para subconjuntos de variables}\label{vector-de-medias-y-matrices-de-covarianzas-para-subconjuntos-de-variables}

En algunas ocasiones podemos estar interesados en dos o más subconjuntos
de variables medidas sobre las mismas unidades muestrales.

\textbf{\includegraphics[width=1em,height=1em]{bookdown_AM_files/figure-latex/fa-icon-00ef1bb1d064b635425d49f06b331298.pdf} ~\textcolor{excolor}{Ejemplo.}} Un estudio sobre diabetes analiza diferentes variables de
interés sobre pacientes que no presentan la enfermedad. Las variables de
mayor relevancia para el estudio son:

\begin{itemize}
\item
  \(x_1\)= Intolerancia a la glucosa
\item
  \(x_2\)= Respuesta de la insulina a la glucosa oral
\item
  \(x_3\)= Resistencia de la insulina
\end{itemize}

Se recogen igualmente otras dos variables de menor relevancia:

\begin{itemize}
\item
  \(y_1\)= Peso relativo
\item
  \(y_2\)= Glucemia basal
\end{itemize}

Supongamos, como en el ejemplo anterior, que tenemos dividido el vector
aleatorio original en dos subvectores \(\bm{y}\) y \(\bm{x}\) con \(p\) y \(q\)
variables, respectivamente. Tendremos entonces:

\[E\left(\begin{array}{c}\bm{y}\\ \bm{x}\end{array}\right)=\left(\begin{array}{c}E(\bm{y})\\ E(\bm{x})\end{array}\right)=\left(\begin{array}{c}\bm{\mu_y}\\ \bm{\mu_x}\end{array}\right)\]
y

\[\textnormal{Cov}\left(\begin{array}{c}\bm{y}\\ \bm{x}\end{array}\right)=\bm{\Sigma}=\left(\begin{array}{cc}\bm{\Sigma_{yy}}&\bm{\Sigma_{yx}}\\ \bm{\Sigma_{xy}}&\bm{\Sigma_{xx}}\end{array}\right).\]

En las expresiones anteriores
\(\left(\begin{array}{c}\bm{y}\\ \bm{x}\end{array}\right)\) denota el
vector original de \(p+q\) variables. La submatriz \(\bm{\Sigma_{yy}}\) es
la matriz de covarianzas \(p\times p\) del subvector \(\bm{y}\). Del mismo
modo, la submatriz \(\bm{\Sigma_{xx}}\) es la matriz de covarianzas
\(q\times q\) del subvector \(\bm{x}\). La submatriz \(\bm{\Sigma_{yx}}\) es
una matriz \(p\times q\) con las covarianzas de cada variable \(y_j\) con
\(x_k\). Nótese que \(\bm{\Sigma_{xy}}=\bm{\Sigma_{yx}}^\prime\).

Estas propiedades a nivel poblacional se trasladan directamente al
vector de medias y matriz de covarianzas muestrales. Cada observación en
la muestra se podrá particionar de la siguiente manera:

\[\left(\begin{array}{c}\bm{y_i}\\ \bm{x_i}\end{array}\right)=\left(\begin{array}{c}
y_{i1}\\
\vdots\\
y_{ip}\\
x_{i1}\\
\vdots\\
x_{iq}\\
\end{array}\right),\ \ i=1,\ldots,n.\]

Para la muestra de \(n\) observaciones, el vector de medias muestral y la
matriz de covarianzas muestrales tendrán el siguiente aspecto:
\[\left(\begin{array}{c}\overline{\bm{y}}\\ \overline{\bm{x}}\end{array}\right)=\left(\begin{array}{c}\overline{y}_1\\ \vdots\\\overline{y}_p \\\overline{x}_1\\ \vdots\\\overline{x}_q\end{array}\right),\ \ \ \ \ \ \ \ \bm{S}=\left(\begin{array}{cc}\bm{S_{yy}}&\bm{S_{yx}}\\ \bm{S_{xy}}&\bm{S_{xx}}\end{array}\right).\]

Lo visto hasta ahora se puede generalizar de forma natural al caso en
que tengamos el vector original dividido en más de dos subvectores.

\textbf{\includegraphics[width=1em,height=1em]{bookdown_AM_files/figure-latex/fa-icon-00ef1bb1d064b635425d49f06b331298.pdf} ~\textcolor{excolor}{Ejemplo.}} Los datos del estudio sobre diabetes se encuentran en el
fichero \texttt{diabetes.txt}.

\footnotesize

\begin{Shaded}
\begin{Highlighting}[]
\SpecialCharTok{\textgreater{}}\NormalTok{ diab }\OtherTok{\textless{}{-}} \FunctionTok{read.table}\NormalTok{(}\StringTok{"data/diabetes.txt"}\NormalTok{, }\AttributeTok{header =} \ConstantTok{TRUE}\NormalTok{)}
\SpecialCharTok{\textgreater{}} \FunctionTok{head}\NormalTok{(diab)}
\NormalTok{    y1  y2  x1  x2  x3}
\DecValTok{1} \FloatTok{0.81}  \DecValTok{80} \DecValTok{356} \DecValTok{124}  \DecValTok{55}
\DecValTok{2} \FloatTok{0.95}  \DecValTok{97} \DecValTok{289} \DecValTok{117}  \DecValTok{76}
\DecValTok{3} \FloatTok{0.94} \DecValTok{105} \DecValTok{319} \DecValTok{143} \DecValTok{105}
\DecValTok{4} \FloatTok{1.04}  \DecValTok{90} \DecValTok{356} \DecValTok{199} \DecValTok{108}
\DecValTok{5} \FloatTok{1.00}  \DecValTok{90} \DecValTok{323} \DecValTok{240} \DecValTok{143}
\DecValTok{6} \FloatTok{0.76}  \DecValTok{86} \DecValTok{381} \DecValTok{157} \DecValTok{165}
\SpecialCharTok{\textgreater{}} \FunctionTok{colMeans}\NormalTok{(diab)}
\NormalTok{         y1          y2          x1          x2          x3 }
  \FloatTok{0.9178261}  \FloatTok{90.4130435} \FloatTok{340.8260870} \FloatTok{171.3695652}  \FloatTok{97.7826087} 
\SpecialCharTok{\textgreater{}} \FunctionTok{cov}\NormalTok{(diab)}
\NormalTok{            y1         y2           x1           x2          x3}
\NormalTok{y1  }\FloatTok{0.01618184}   \FloatTok{0.216029}    \FloatTok{0.7871691}   \SpecialCharTok{{-}}\FloatTok{0.2138454}    \FloatTok{2.189072}
\NormalTok{y2  }\FloatTok{0.21602899}  \FloatTok{70.558937}   \FloatTok{26.2289855}  \SpecialCharTok{{-}}\FloatTok{23.9560386}  \SpecialCharTok{{-}}\FloatTok{20.841546}
\NormalTok{x1  }\FloatTok{0.78716908}  \FloatTok{26.228986} \FloatTok{1106.4135266}  \FloatTok{396.7323671}  \FloatTok{108.383575}
\NormalTok{x2 }\SpecialCharTok{{-}}\FloatTok{0.21384541} \SpecialCharTok{{-}}\FloatTok{23.956039}  \FloatTok{396.7323671} \FloatTok{2381.8826087} \FloatTok{1142.637681}
\NormalTok{x3  }\FloatTok{2.18907246} \SpecialCharTok{{-}}\FloatTok{20.841546}  \FloatTok{108.3835749} \FloatTok{1142.6376812} \FloatTok{2136.396135}
\end{Highlighting}
\end{Shaded}

\normalsize

\subsection{Combinaciones lineales de variables}\label{combinaciones-lineales-de-variables}

Con frecuencia estamos interesados en combinaciones lineales de las
variables de estudio \(x_1,\ldots, x_d\). Sean \(a_1,\ldots,a_d\) constantes
y consideremos la combinación lineal

\[z=a_1x_1+\ldots+a_dx_d=\bm{a}^\prime\bm{x},\]
donde
\(\bm{a}=(a_1,\ldots,a_d)^\prime\). Para la combinación
lineal\(z=\bm{a}^\prime\bm{x}\) se tiene que
\(E(z)=\bm{a}^\prime E(\bm{x})\) y
\(\textnormal{Var}(z)=\bm{a}^\prime\bm\Sigma\bm{a}\).

De forma más general, si consideramos \(k\) combinaciones lineales

\[\bm{z}=\left(\begin{array}{c}z_1\\z_2\\ \vdots\\ z_k\end{array}\right)=\left(\begin{array}{c}a_{11}x_1+\ldots+a_{1d}x_d\\a_{21}x_1+\ldots+a_{2d}x_d\\ \vdots\\ a_{k1}x_1+\ldots+a_{kd}x_d\end{array}\right)=\left(\begin{array}{c}\bm{a_1}^\prime\bm{x}\\\bm{a_2}^\prime\bm{x}\\ \vdots\\ \bm{a_k}^\prime\bm{x}\end{array}\right)=\bm{A}\bm{x},\]
entonces \(E(\bm{z})=\bm{A} E(\bm{x})\) y
\(\bm{\Sigma_z}=\bm{A}\bm\Sigma\bm{A}^\prime\). Para la transformación más
general \(\bm{z}=\bm{A}\bm{x}+\bm{b}\), se tiene
\(E(\bm{z})=\bm{A} E(\bm{x})+\bm{b}\) y
\(\bm{\Sigma_z}=\bm{A}\bm\Sigma\bm{A}^\prime\).

Consideremos ahora el análogo muestral. Si aplicamos la transformación a
todos los elementos de la muestra, es decir, si hacemos
\(\bm{z_i}=\bm{A}\bm{x_i}+\bm{b}\) para \(i=1,\ldots, n\), entonces
\(\overline{\bm{z}}=\bm{A} \overline{\bm{x}}+\bm{b}\) y
\(\bm{S_z}=\bm{A}\bm{S}\bm{A}^\prime\).

\textbf{\includegraphics[width=1em,height=1em]{bookdown_AM_files/figure-latex/fa-icon-00ef1bb1d064b635425d49f06b331298.pdf} ~\textcolor{excolor}{Ejemplo.}} Para el ejemplo del estudio sobre diabetes consideramos las
siguientes transformaciones lineales:

\begin{itemize}
\item
  \(z_1=x_1+2x_2-x3\)
\item
  \(z_2=y_1+4y_2\)
\end{itemize}

En primer lugar calculamos \(\overline{z}_1\), \(\overline{z}_2\),
\(s^2_{z_1}\), \(s^2_{z_2}\) y \(s_{z_1z_2}\) a partir de los datos
transformados

\footnotesize

\begin{Shaded}
\begin{Highlighting}[]
\SpecialCharTok{\textgreater{}} \FunctionTok{attach}\NormalTok{(diab)}
\SpecialCharTok{\textgreater{}}\NormalTok{ z1 }\OtherTok{\textless{}{-}}\NormalTok{ x1 }\SpecialCharTok{+} \DecValTok{2} \SpecialCharTok{*}\NormalTok{ x2 }\SpecialCharTok{{-}}\NormalTok{ x3}
\SpecialCharTok{\textgreater{}}\NormalTok{ z2 }\OtherTok{\textless{}{-}}\NormalTok{ y1 }\SpecialCharTok{+} \DecValTok{4} \SpecialCharTok{*}\NormalTok{ y2}
\SpecialCharTok{\textgreater{}} \FunctionTok{mean}\NormalTok{(z1)}
\NormalTok{[}\DecValTok{1}\NormalTok{] }\FloatTok{585.7826}
\SpecialCharTok{\textgreater{}} \FunctionTok{mean}\NormalTok{(z2)}
\NormalTok{[}\DecValTok{1}\NormalTok{] }\FloatTok{362.57}
\SpecialCharTok{\textgreater{}} \FunctionTok{var}\NormalTok{(z1)}
\NormalTok{[}\DecValTok{1}\NormalTok{] }\FloatTok{9569.952}
\SpecialCharTok{\textgreater{}} \FunctionTok{var}\NormalTok{(z2)}
\NormalTok{[}\DecValTok{1}\NormalTok{] }\FloatTok{1130.687}
\SpecialCharTok{\textgreater{}} \FunctionTok{cov}\NormalTok{(z1, z2)}
\NormalTok{[}\DecValTok{1}\NormalTok{] }\SpecialCharTok{{-}}\FloatTok{5.195778}
\end{Highlighting}
\end{Shaded}

\normalsize

De otro modo, podemos escribir

\[\bm{z}=\left(\begin{array}{c}z_1\\z_2\end{array}\right)=\left(\begin{array}{ccccc}0&0&1&2&-3\\1&4&0&0&0\end{array}\right)\left(\begin{array}{c}y_1\\y_2\\x_1\\x_2\\x_3\end{array}\right)=\bm{A}\left(\begin{array}{c}y_1\\y_2\\x_1\\x_2\\x_3\end{array}\right).\]

Teniendo en cuenta las propiedades de las transformaciones lineales
discutidas anteriormente, se tiene también:

\footnotesize

\begin{Shaded}
\begin{Highlighting}[]
\SpecialCharTok{\textgreater{}}\NormalTok{ A }\OtherTok{\textless{}{-}} \FunctionTok{rbind}\NormalTok{(}\FunctionTok{c}\NormalTok{(}\DecValTok{0}\NormalTok{, }\DecValTok{0}\NormalTok{, }\DecValTok{1}\NormalTok{, }\DecValTok{2}\NormalTok{, }\SpecialCharTok{{-}}\DecValTok{1}\NormalTok{), }\FunctionTok{c}\NormalTok{(}\DecValTok{1}\NormalTok{, }\DecValTok{4}\NormalTok{, }\DecValTok{0}\NormalTok{, }\DecValTok{0}\NormalTok{, }\DecValTok{0}\NormalTok{))}
\SpecialCharTok{\textgreater{}}\NormalTok{ A }\SpecialCharTok{\%*\%} \FunctionTok{colMeans}\NormalTok{(diab)}
\NormalTok{         [,}\DecValTok{1}\NormalTok{]}
\NormalTok{[}\DecValTok{1}\NormalTok{,] }\FloatTok{585.7826}
\NormalTok{[}\DecValTok{2}\NormalTok{,] }\FloatTok{362.5700}
\SpecialCharTok{\textgreater{}}\NormalTok{ A }\SpecialCharTok{\%*\%} \FunctionTok{cov}\NormalTok{(diab) }\SpecialCharTok{\%*\%} \FunctionTok{t}\NormalTok{(A)}
\NormalTok{            [,}\DecValTok{1}\NormalTok{]        [,}\DecValTok{2}\NormalTok{]}
\NormalTok{[}\DecValTok{1}\NormalTok{,] }\FloatTok{9569.951691}   \SpecialCharTok{{-}}\FloatTok{5.195778}
\NormalTok{[}\DecValTok{2}\NormalTok{,]   }\SpecialCharTok{{-}}\FloatTok{5.195778} \FloatTok{1130.687409}
\end{Highlighting}
\end{Shaded}

\normalsize

\section{Distancias estadísticas}\label{distancias-estadisticas}

El concepto de distancia entre objetos o individuos observados permite
interpretar geométricamente muchas técnicas de análisis multivariante.
En el caso unidimensional, la distancia entre dos puntos \(x\) e \(y\) se
mide de manera natural mediante la distancia euclídea
\(\left|x-y\right|\).

Cuando disponemos de una variable vectorial, cada dato es un punto en
\(\mathbb{R}^d\) y podríamos entonces pensar en utilizar de nuevo la
distancia euclídea en el contexto multidimensional. A partir de ahora
sean \(\bm{x_i}=(x_{i1},\ldots,x_{id})^\prime{}\) y
\(\bm{x_k}=(x_{k1},\ldots,x_{kd})^\prime{}\) las observaciones de dos
individuos \(i\), \(k\) obtenidas al medir el vector \(d\)-dimensional
\((x_1,\ldots, x_d)^\prime{}\).

\subsection{Distancia euclídea}\label{distancia-euclidea}

Se define la distancia euclídea entre \(\bm{x_i}\) y \(\bm{x_k}\) como

\[d_E(\bm{x_i},\bm{x_k})=\sqrt{\sum_{j=1}^d{(x_{ij}-x_{kj})^2}}.\]

La distancia euclídea es la más utilizada pero tiene como inconvenientes
que depende de las unidades de medida de las variables (no es invariante ante cambios de escala) y presupone que las variables son incorrelacionadas y de varianza unidad.

\textbf{\includegraphics[width=1em,height=1em]{bookdown_AM_files/figure-latex/fa-icon-00ef1bb1d064b635425d49f06b331298.pdf} ~\textcolor{excolor}{Ejemplo.}} Supongamos que disponemos de los datos de peso y estatura
de 300 mujeres con edades comprendidas entre 30 y 40 años. Queremos
determinar la posición con respecto a la media de tres nuevas mujeres a
partir de sus respectivos pesos y estaturas. Supongamos que la mujer \(A\)
pesa \(63\) kg. y mide \(180\) cm. La mujer \(B\) pesa \(60\) kg. y mide \(177\)
cm. La mujer \(C\) pesa \(70\) kg. y mide \(177\) cm. Los datos se muestran en
la Figura \ref{fig:dist}

Si bien parece evidente que la posición de las tres mujeres con respecto
a la media es diferente (la mujer \(C\) parece más ``cercana'' a la media
que la \(A\) y la \(B\)), la distancia euclídea de todos los puntos a la
media es la misma.

\footnotesize
\begin{figure}

{\centering \resizebox{.4\linewidth}{!}{\includegraphics[width=0.5\linewidth]{bookdown_AM_files/figure-latex/dist-1} } 

}

\caption{Peso (en kg.) y estatura (en cm.) de 300 mujeres con edades comprendidas entre 30 y 40 años. En rojo se representa el peso medio y la estatura media. En azul, datos de las mujeres A, B y C.}\label{fig:dist}
\end{figure}
\normalsize

\subsection{Distancia euclídea estandarizada por columnas}\label{distancia-euclidea-estandarizada-por-columnas}

Una manera de evitar el problema de que la distancia dependa de las
unidades de medida de las variables, consiste en dividir cada variable
por un término que elimine el efecto de la escala.

Se define la distancia euclídea estandarizada por columnas entre
\(\bm{x_i}\) y \(\bm{x_k}\) como

\[d_{Est}(\bm{x_i},\bm{x_k})=\sqrt{\sum_{j=1}^d{\left(\frac{x_{ij}-x_{kj}}{\sigma_j}\right)^2}}.\]
Su utilidad radica en que es una forma de determinar la similitud entre
dos variables aleatorias multidimensionales invariante ante cambios de
escala. Puede verse como una distancia euclídea donde cada coordenada se
pondera inversamente proporcional a la varianza.

Tiene como inconvenientes que presupone que las variables son incorrelacionadas.

\textbf{\includegraphics[width=1em,height=1em]{bookdown_AM_files/figure-latex/fa-icon-00ef1bb1d064b635425d49f06b331298.pdf} ~\textcolor{excolor}{Ejemplo.}} En el ejemplo anterior de los datos de peso y estatura de
300 mujeres, la distancia euclídea estandarizada por columnas de la
mujer \(A\) a la media es mayor que la distancia euclídea estandarizada
por columnas de las mujeres \(B\) y \(C\) a la media. Aun así, la distancia
euclídea estandarizada por columnas de las mujeres \(B\) y \(C\) a la media
es la misma.

\subsection{Distancia de Mahalanobis}\label{distancia-de-mahalanobis}

Se define la distancia de Mahalanobis entre \(\bm{x_i}\) y \(\bm{x_k}\) como

\[d_M(\bm{x_i},\bm{x_k})=\sqrt{(\bm{x_i}-\bm{x_k})^\prime{}\bm\Sigma^{-1}(\bm{x_i}-\bm{x_k})}.\]

Geométricamente equivale a girar la nube de puntos hasta eliminar las
correlaciones y luego calcular la distancia para los datos
estandarizados. Es adecuada como medida de discrepancia entre datos,
porque es invariante ante cambios de escala y tiene en cuenta las correlaciones entre las variables.

Además,

\begin{itemize}
\item
  \(d_M\equiv d_E\) cuando \(\bm\Sigma=\bm{I_d}\).
\item
  \(d_M\equiv d_{Est}\) cuando
  \(\bm\Sigma=\textnormal{diag}(\sigma_1^2,\ldots,\sigma_d^2)\).
\end{itemize}

En caso de que la matriz de covarianzas sea desconocida se calcula

\[d_M(\bm{x_i},\bm{x_k})=\sqrt{(\bm{x_i}-\bm{x_k})^\prime{}\bm{S}^{-1}(\bm{x_i}-\bm{x_k})}.\]

\textbf{\includegraphics[width=1em,height=1em]{bookdown_AM_files/figure-latex/fa-icon-00ef1bb1d064b635425d49f06b331298.pdf} ~\textcolor{excolor}{Ejemplo.}} En el ejemplo anterior de los datos de peso y estatura de
300 mujeres, la distancia de Mahalanobis de la mujer \(A\) a la media es
mayor que la distancia de Mahalanobis de la mujer \(B\) a la media, que a
su vez es mayor que la distancia de Mahalanobis de la mujer \(C\) a la
media.

\subsection{Estandarización}\label{estandarizacion}

Una transformación lineal que tiene un interés especial es la
\textbf{estandarización}. La estandarización (o tipificación) de una variable
aleatoria se consigue restando la media y dividiendo por la desviación
típica (raíz cuadrada de la varianza). De esta manera se consigue que la
nueva variable transformada tenga media cero y varianza uno.

En el caso de un vector aleatorio de dimensión \(d\), con \(d\) mayor que
uno, podemos pensar en estandarizar cada una de sus variables. Como
resultado de esta \textbf{estandarización univariante} se obtiene un vector
cuyas componentes tienen todas media cero y desviación típica uno. Sin
embargo, puede haber correlación entre las variables que componen el
vector.

\textbf{\includegraphics[width=1em,height=1em]{bookdown_AM_files/figure-latex/fa-icon-00ef1bb1d064b635425d49f06b331298.pdf} ~\textcolor{excolor}{Ejemplo.}} En la Figura \ref{fig:estandarizacion} se muestran dos diagramas de dispersión para
los datos de peso y estatura de mujeres. El primero de ellos representa
los datos tras haber estandarizado de manera univariante las dos
variables. Esto hace que se hayan contraído los ejes. Se mantiene la
correlación positiva, aunque la orientación de la nube de puntos ahora
es más próxima a la diagonal, pues ambas variables tienen la misma
dispersión.

El segundo gráfico muestra el diagrama de dispersión de los datos tras
haberles aplicado una estandarización bivariante.

\footnotesize
\begin{figure}

{\centering \resizebox{.4\linewidth}{!}{\includegraphics[width=0.5\linewidth]{bookdown_AM_files/figure-latex/estandarizacion-1} } \resizebox{.4\linewidth}{!}{\includegraphics[width=0.5\linewidth]{bookdown_AM_files/figure-latex/estandarizacion-2} } 

}

\caption{Estandarización univariante y multivariante}\label{fig:estandarizacion}
\end{figure}
\normalsize

La estandarización multivariante se obtiene así:

\[\bm{y}=\bm{\Sigma^{-1/2}}(\bm{x}-\bm\mu)\] donde \(\bm\mu\) denota el
vector de medias, \(\bm\Sigma\) es la matriz de covarianzas y
\(\bm{\Sigma^{-1/2}}\) es una matriz \(d\times d\) que definiremos más
adelante. Por las propiedades vistas en la sección anterior, se puede
comprobar que:

\begin{itemize}
\item
  El vector estandarizado tiene media cero:

  \[E(\bm{y})=E(\bm{\Sigma^{-1/2}}(\bm{x}-\bm\mu))=\bm{\Sigma^{-1/2}}E(\bm{x}-\bm\mu)=0.\]
\item
  La matriz de covarianzas de \(\bm{y}\) es la matriz identidad \(d\times d\), \(\bm{I_d}\):

  \[\mbox{Cov}(\bm{y},\bm{y})=\mbox{Cov}(\bm{\Sigma^{-1/2}}(\bm{x}-\bm\mu),\bm{\Sigma^{-1/2}}(\bm{x}-\bm\mu))=\bm{\Sigma^{-1/2}}\mbox{Cov}(\bm{x},\bm{x})(\bm{\Sigma^{-1/2}})^\prime{}
  \stackrel{(a)}{=}\bm{I_d}.\]
\end{itemize}

Para obtener la matriz \(\bm{\Sigma^{-1/2}}\), de manera que se cumpla la
igualdad (a) en el cálculo anterior de la matriz de covarianzas,
tendremos presente que toda matriz de covarianzas es una matriz
simétrica y semidefinida positiva. Así, se le puede aplicar la
descomposición en autovalores y autovectores, que vemos a continuación.

Si \(\bm{A}\) es una matriz simétrica, entonces

\[\bm{A}=\left(\bm{v_1},\ldots,\bm{v_d}\right)\left(\begin{array}{ccc}
\lambda_1 & & 0 \\
 & \ddots & \\
 0 & & \lambda_d
 \end{array}\right)\left(\begin{array}{c}
 \bm{v_1}^{\prime} \\
 \vdots \\
 \bm{v_d}^{\prime}
 \end{array}\right)\]
siendo \(\bm{v_1},\ldots,\bm{v_d}\) una base
ortonormal de autovectores de \(\bm{A}\) y \(\lambda_1,\ldots,\lambda_d\)
sus autovalores asociados. \(\bm{A}\) se dice definida positiva si todos
los autovalores de \(\bm{A}\) son positivos. En ese caso se puede emplear
para definir una norma (y una distancia):
\(\|\bm{x}\|^2=\bm{x}'\bm{A}\bm{x}\).

Entonces, se puede obtener \(\bm{\Sigma^{-1/2}}\) a partir de los
autovalores y autovectores de \(\bm\Sigma\), de la siguiente manera:

\[{\bm{\Sigma^{-1/2}}}=\left(\bm{v_1},\ldots,\bm{v_d}\right)\left(\begin{array}{ccc}
1/\sqrt{\lambda_1} & & 0 \\
 & \ddots & \\
 0 & & 1/\sqrt{\lambda_d}
 \end{array}\right)\left(\begin{array}{c}
 \bm{v_1}^{\prime} \\
 \vdots \\
 \bm{v_d}^{\prime}
 \end{array}\right).\]

\chapter{Inferencia en poblaciones normales}\label{inferencia-en-poblaciones-normales}

Muchos de los métodos de contrastes de hipótesis y construcción de
intervalos de confianza en el contexto unidimensional asumen que la
distribución de los datos es normal univariante. De la misma manera, la
mayor parte de los procedimientos multivariantes están basados en la
distribución normal multivariante. Una de las principales ventajas de la
distribución normal multivariante es que puede ser descrita totalmente
utilizando únicamente su vector de medias y matriz de covarianzas.
Además, las combinaciones lineales de normales siguen siendo normales e,
incluso cuando los datos no provienen exactamente de una distribución
normal, ésta puede servir como una buena aproximación (especialmente al
hacer inferencia que involucre el vector de medias muestrales, que será
aproximadamente normal multivariante por el Teorema Central del Límite).

Revisaremos en la primera parte de este capítulo la definición de
distribución normal multivariante junto con algunas de sus propiedades.
En el resto del tema veremos tareas de inferencia sobre el vector de
medias y la matriz de covarianzas de una población normal multivariante,
en base a una muestra aleatoria simple extraída de ella. También se
tratarán problemas que involucren a varias poblaciones. Muchos
procedimientos resultarán ser extensiones naturales de los métodos ya
conocidos para poblaciones normales univariantes, mientras que en algún
caso surgirán problemas nuevos, por ejemplo, comparación entre
componentes del vector de medias o cuestiones de inferencia simultánea;
en definitiva, situaciones provocadas por la dimensión múltiple.

Aunque el vector de medias muestral y la matriz de covarianzas muestral
son estimadores naturales de sus análogos poblacionales, también vamos a
ver que son los estimadores de máxima verosimilitud, y de paso,
introducimos la función de verosimilitud y en general la idea de
verosimilitud, que será empleada en varias ocasiones a lo largo de este
tema.

Por otro lado, al estudiar los procedimientos naturales de inferencia
sobre el vector de medias y la matriz de covarianzas, veremos que la
suposición de normalidad en la distribución del vector de variables
observadas es esencial en muchos de ellos. Por ejemplo, la forma
elíptica de las regiones de confianza es consustancial a la forma de la
distribución normal multivariante. Si la forma en que se presentaran las
observaciones en el espacio no fuera elíptica, tampoco lo debería ser
una región de confianza para el vector de medias. La última parte del
tema está dedicada a métodos de contraste de normalidad en observaciones
multivariantes. Existen varios procedimientos para evaluar si un
conjunto de datos proviene de una población normal multivariante. Una
primera posibilidad sería evaluar si cada una de las variables es normal
univariante. Como hemos comentado, la normalidad multivariante implica
la normalidad de distribuciones marginales unidimensionales. Por lo
tanto, si alguna de las variables no es normal univariante, el vector no
será normal multivariante. Tiene sentido entonces hacer un estudio
inicial sobre las variables de forma individual. En la Sección
\ref{secnorm} repasaremos en primer lugar los métodos de contraste de la normalidad en el caso
univariante, prestando especial atención a los métodos susceptibles de
ser extendidos al caso multivariante. De todos modos, los contrastes de
normalidad univariante no serán en general suficientes, ya que la
existencia de normalidad univariante de todas las variables no garantiza
la normalidad multivariante de los datos. A continuación se
abordan algunos métodos (los más representativos) para el contraste de
normalidad multivariante. Están basados en la extensión de las medidas
de asimetría, kurtosis, así como del test de Shapiro-Wilk al caso
multivariante.

\section{Estimación en la distribución normal multivariante}\label{estimacion-en-la-distribucion-normal-multivariante}

\subsection{La distribución normal multivariante}\label{la-distribucion-normal-multivariante}

Recordamos que si una variable aleatoria \(x\) unidimensional, con media
\(\mu\) y varianza \(\sigma^2\) tiene distribución normal, entonces su
función de densidad viene dada por

\[f(x)=\frac{1}{\sqrt{2\pi\sigma^2}}e^\frac{-({x}-\mu)^2}{2\sigma^2},\ \ \ -\infty<x<\infty.\]

Decimos en este caso que \(x\) tiene distribución \(N(\mu,\sigma^2)\).
Puedes ver la representación de la función de densidad de distintas
variables normales en la Figura
\ref{fig:norm}.

\begin{figure}

{\centering \resizebox{.4\linewidth}{!}{\includegraphics[width=0.5\linewidth]{bookdown_AM_files/figure-latex/norm-1} } 

}

\caption{Representación de la densidad de variables normales.}\label{fig:norm}
\end{figure}

Si el vector aleatorio \(d\)-dimensional \(\bm{x}\) tiene distribución
normal multivariante con vector de medias \(\bm{\mu}\) y matriz de
covarianzas \(\bm\Sigma\), entonces su función de densidad viene dada por

\[f(\bm{x})=\frac{1}{(\sqrt{2\pi})^d\left|\bm\Sigma\right|^{1/2}}e^\frac{{-(\bm{x}-\bm{\mu})^\prime\bm{\Sigma}^{-1}(\bm{x}-\bm{\mu})}}{2}.\]

Decimos en este caso que \(\bm{x}\) tiene distribución
\(N_d(\bm{\mu},\bm\Sigma)\). En la Figura \ref{fig:norm2} se
representan las funciones de densidad de distintas variables normales
bivariantes. En el Anexo \ref{anexoDist} encontrarás propiedades importantes
de esta distribución.

\begin{figure}

{\centering \resizebox{.4\linewidth}{!}{\includegraphics[width=0.5\linewidth]{bookdown_AM_files/figure-latex/norm2-1} } \resizebox{.4\linewidth}{!}{\includegraphics[width=0.5\linewidth]{bookdown_AM_files/figure-latex/norm2-2} } 

}

\caption{Densidad de vectores aleatorios con distribución normal bivariante. A la izquierda $N(\bm{0},\bm{I_2})$. A la derecha $N(\bm{0},\bm{\Sigma})$ con $\sigma_1^2=1$, $\sigma_2^2=3$ y $\sigma_{12}=1.5$.}\label{fig:norm2}
\end{figure}

Fíjate que en la densidad multivariante, \(\left|\bm\Sigma\right|^{1/2}\)
juega el papel de \(\sqrt{\sigma^2}\) en el caso unidimensional. Decimos
que \(\left|\bm\Sigma\right|\) es la \emph{varianza poblacional generalizada}.
En efecto, el determinante \(\left|\bm\Sigma\right|\) nos sirve como una
medida de variabilidad total. Dado un vector aleatorio \(\bm{x}\) con
media \(\bm{\mu}\) y matriz de covarianzas \(\bm{\Sigma}\), un valor pequeño
de \(\left|\bm\Sigma\right|\) indica que \(\bm{x}\) está concentrado cerca
de \(\bm{\mu}\) o que existe multicolinealidad entre las variables. Si
recordamos que el determinante de una matriz es el producto de sus
autovalores, podemos concebir el determinante como una medida de si una
matriz de covarianzas es grande o pequeña. Por ejemplo, si existe
multicolinealidad, uno o más autovalores de \(\bm\Sigma\) será próximo a
cero y en consecuencia, \(\left|\bm\Sigma\right|\) será pequeño.

Podemos representar una matriz de covarianzas \(2\times 2\) mediante una
elipse cuyos ejes principales siguen la dirección de sus autovectores.
Cada radio de la elipse es la raíz cuadrada de un autovalor de la matriz
de covarianzas, ver Figura \ref{fig:detsigma} (izquierda). Entonces, si denotamos por
\(\lambda_1\) y \(\lambda_2\) a los autovalores de \(\bm\Sigma\),

\[\textnormal{Area del rectángulo circunscrito en la elipse}=4\cdot \sqrt{\lambda_1}\sqrt{\lambda_2}=4\sqrt{\left|\bm\Sigma\right|}.\]
En la última igualdad hemos usado que el determinante es el producto de
los autovalores. En la Figura \ref{fig:detsigma} (derecha) se representan las elipses
correspondientes a distintas matrices de covarianzas. Para esta
representación hemos utilizado la función \texttt{ellipse} de la librería
\texttt{car}. En general, si la dimensión el espacio es \(d\), se tendrá que

\[\textnormal{Volumen del rectángulo circunscrito en el elipsoide}=2^d\sqrt{\left|\bm\Sigma\right|}.\]

Empleamos los términos volumen y elipsoide, pues son los que
corresponden en el espacio de dimensión tres. En espacios de dimensión
superior, que no son representables, se suele emplear igualmente esos
términos. Por tanto, el determinante es una forma de medir el volumen o
capacidad de un rectángulo (\(d=2\)), caja (\(d=3\)), ortoedro en general,
que contenga al grueso del vector.

\begin{figure}

{\centering \resizebox{.4\linewidth}{!}{\includegraphics[width=0.5\linewidth]{bookdown_AM_files/figure-latex/detsigma-1} } \resizebox{.4\linewidth}{!}{\includegraphics[width=0.5\linewidth]{bookdown_AM_files/figure-latex/detsigma-2} } 

}

\caption{Izquierda, representación de la elipse correspondiente a una matriz de covarianzas $2\times 2$. con autovectores $\bm{v_1}$, $\bm{v_2}$ y autovalores $\lambda_1$, $\lambda_2$. Cada radio de la elipse es la raíz cuadrada de un autovalor de la matriz de covarianzas. Derecha, representación de las elipses correspondientes a matrices de covarianzas $\bm{\Sigma_1}=\bm{I_2}$ (rojo), $\bm{\Sigma_2}=\mbox{diag}(3,1)$ (verde) y $\bm{\Sigma_3}=(\sigma_{ij})$ con $\sigma_{11}=3$, $\sigma_{22}=1$, $\sigma_{12}=1.4$ (azul).}\label{fig:detsigma}
\end{figure}

\subsection{Estimadores de máxima verosimilitud}\label{estimadores-de-maxima-verosimilitud}

Consideremos disponible una muestra aleatoria simple
\(\bm{x_1},\ldots,\bm{x_n}\in N_d(\bm\mu,\bm\Sigma)\) de vectores
aleatorios independientes y con la misma distribución normal
multivariante. Vamos a obtener los estimadores de máxima verosimilitud
para el vector de medias \(\bm\mu\) y para la matriz de covarianzas
\(\bm\Sigma\). La función de verosimilitud sería:

\[L(\bm{x},\bm\mu,\bm\Sigma)=(2\pi)^{-nd/2} |\bm\Sigma |^{-n/2}
\exp\left\{-\frac{1}{2}\sum_{i=1}^n
\left(\bm{x_i}-\bm\mu\right)^\prime{}\bm\Sigma^{-1} \left(\bm{x_i}-\bm\mu\right)\right\}.\]

Observamos que

\[\begin{aligned}
\sum_{i=1}^n \left(\bm{x_i}-\bm{\mu}\right)^\prime{}\bm\Sigma^{-1} \left(\bm{x_i}-\bm{\mu}\right)
&= \sum_{i=1}^n\left[ \left(\bm{x_i}-\overline{\bm{x}}\right)^\prime{}\bm\Sigma^{-1} \left(\bm{x_i}-\overline{\bm{x}}\right)\right.\\
&\left.\qquad + \left(\overline{\bm{x}}-\bm{\mu}\right)^\prime{}\bm\Sigma^{-1} \left(\overline{\bm{x}}-\bm{\mu}\right)+ 2 \left(\overline{\bm{x}}-\bm{\mu}\right)^\prime{}\bm\Sigma^{-1} \left(\bm{x_i}-\overline{\bm{x}}\right) \right] \\
&= \sum_{i=1}^n \left(\bm{x_i}-\overline{\bm{x}}\right)^\prime{}\bm\Sigma^{-1} \left(\bm{x_i}-\overline{\bm{x}}\right)+ n \left(\overline{\bm{x}}-\bm{\mu}\right)^\prime{}\bm\Sigma^{-1} \left(\overline{\bm{x}}-\bm{\mu}\right)
 \end{aligned},\]
ya que la suma de los dobles productos vale cero. Entonces la
log--verosimilitud se puede expresar así:

\[\begin{aligned}
\log L(\bm{x},\bm{\mu},\bm\Sigma) &= c - \frac{n}{2} \log |\bm\Sigma|
 - \frac{1}{2}\sum_{i=1}^n \left(\bm{x_i}-\bm{\mu}\right)^\prime{}\bm\Sigma^{-1} \left(\bm{x_i}-\bm{\mu}\right) \\
&= c - \frac{n}{2} \log |\bm\Sigma|
-\frac{1}{2} \sum_{i=1}^n \left(\bm{x_i}-\overline{\bm{x}}\right)^\prime{}\bm\Sigma^{-1} \left(\bm{x_i}-\overline{\bm{x}}\right)\\
&\qquad -\frac{n}{2}  \left(\overline{\bm{x}}-\bm{\mu}\right)^\prime{}\bm\Sigma^{-1} \left(\overline{\bm{x}}-\bm{\mu}\right)\end{aligned}\]
siendo \(c=-\frac{nd}{2}\log(2\pi)\).

Observamos que, por ser \(\bm\Sigma\) definida positiva, (y en
consecuencia, también lo será su inversa \(\bm\Sigma^{-1}\)), se tiene
\((\overline{\bm{x}}-\bm{\mu})^\prime{}\bm\Sigma^{-1}(\overline{\bm{x}}-\bm{\mu})>0\),
salvo que \(\bm{\mu}=\overline{\bm{x}}\), en cuyo caso vale cero. Por
tanto, la función de log--verosimilitud alcanza su máximo en
\(\hat{\bm{\mu}}=\overline{\bm{x}}\), que de este modo se convierte en el
estimador de máxima verosimilitud del vector de medias. Además,

\begin{equation}
\sup_{\bm\mu}\log L(\bm{x},\bm\mu,\bm\Sigma) = c - \frac{n}{2} \log |\bm\Sigma|
-\frac{1}{2} \sum_{i=1}^n \left(\bm{x_i}-\overline{\bm{x}}\right)^\prime{}\bm\Sigma^{-1} \left(\bm{x_i}-\overline{\bm{x}}\right)
\label{eq:supmu}
\end{equation}

para cualquier matriz de covarianzas \(\bm\Sigma\).

A continuación calcularemos el máximo de aquella función respecto de
\(\bm\Sigma\). Podemos expresar

\[\begin{aligned}
\sup_{\bm{\mu}}\log L(\bm{x},\bm{\mu},\bm\Sigma) &= c - \frac{n}{2} \log |\bm\Sigma|
-\frac{1}{2} \mbox{traza} \left[\sum_{i=1}^n
\left(\bm{x_i}-\overline{\bm{x}}\right)^\prime{}\bm\Sigma^{-1} \left(\bm{x_i}-\overline{\bm{x}}\right)\right] \nonumber \\
&= c - \frac{n}{2} \log |\bm\Sigma|
-\frac{1}{2} \sum_{i=1}^n \mbox{traza}
\left[\left(\bm{x_i}-\overline{\bm{x}}\right)^\prime{} \bm\Sigma^{-1} \left(\bm{x_i}-\overline{\bm{x}}\right)\right] \nonumber \\
&= c - \frac{n}{2} \log |\bm\Sigma|
-\frac{1}{2} \sum_{i=1}^n \mbox{traza}
\left[ \bm\Sigma^{-1} \left(\bm{x_i}-\overline{\bm{x}}\right) \left(\bm{x_i}-\overline{\bm{x}}\right)^\prime{} \right] \nonumber \\
&= c - \frac{n}{2}\left( \log |\bm\Sigma| + \mbox{traza} \left( \bm\Sigma^{-1} \bm{S} \right)\right)
\end{aligned}\]
donde hemos aplicado que
traza(\(\bm{A}\)+\(\bm{B}\))=traza(\(\bm{A}\))+traza(\(\bm{B}\)) y que
traza(\(\bm{AB}\))=traza(\(\bm{BA}\)). Ahora debemos obtener el máximo de
esta función respecto del argumento \(\bm\Sigma\). Para ello, apelamos al
resultado siguiente, del cual omitimos su demostración.

\begin{lemma}
\protect\hypertarget{lem:unnamed-chunk-15}{}\label{lem:unnamed-chunk-15}Supongamos una matriz \(\bm{A}\) definida positiva. La función

\[f(\bm\Sigma)=\log|\bm\Sigma|+{\textnormal{traza}}\left(\bm\Sigma^{-1}\bm{A}\right),\]
restringida a las matrices \(\bm\Sigma\) definidas positivas, alcanza su
mínimo en \(\bm\Sigma=\bm{A}\).
\end{lemma}

Entonces, aplicando este lema llegamos a la conclusión de que los
estimadores de máxima verosimilitud del vector de medias y la matriz de
covarianzas (sin restricciones) son \(\overline{\bm{x}}\) y \(\bm{S}\),
respectivamente. Asimismo, la función de verosimilitud tiene como
máximo:

\[\begin{aligned}
\sup_{\bm\Sigma}\sup_{\bm\mu}\log L(\bm{x},\bm{\mu},\bm\Sigma)
&= c - \frac{n}{2}
\left( \log |\bm{S}| + \mbox{traza} \left( \bm{S}^{-1} \bm{S} \right)\right)\\
&= c - \frac{n}{2} \left( \log |\bm{S}| + d\right).
\end{aligned}\]

\subsection{\texorpdfstring{Distribución de \(\overline{\bm{x}}\) y \(\bm{S}\)}{Distribución de \textbackslash overline\{\textbackslash bm\{x\}\} y \textbackslash bm\{S\}}}\label{distribucion-de-overline-bm-x-y-bm-s}

Sea entonces \(\bm{x_1},\ldots,\bm{x_n}\) una muestra aleatoria simple de
\(N_d(\bm\mu,\bm\Sigma)\). En el apartado anterior hemos visto que los
estimadores de máxima verosimilitud de \(\bm\mu\) y \(\bm\Sigma\) son
\(\overline{\bm{x}}\) y \(\bm{S}\), respectivamente. Se verifica que:

\[\overline{\bm{x}}=\frac{1}{n}\sum_{i=1}^n \bm{x_i} \in N_d\left(\bm{\mu},\frac{1}{n}\bm\Sigma\right).\]

Se puede ver este resultado como una generalización del caso
unidimensional. Recuerda que si \({x_1},\ldots,{x_n}\) es una muestra
aleatoria simple de \(N(\mu,\sigma^2)\), entonces
\(\overline{{x}}\in N\left({\mu},\sigma^2/n\right)\). Además en el caso
unidimensional se demostraba también que
\(ns^2/\sigma^2=\sum_{i=1}^n(x_i-\bar{x})^2/\sigma^2\in \chi_{n-1}^2\), es
decir, tiene distribución ji-cuadrado con \(n-1\) grados de libertad\footnote{La distribución ji-cuadrado se define formalmente como la suma de los cuadrados de variables normales estándar, es decir si \({x_1},\ldots,{x_n}\) son \(N(\mu,\sigma^2)\) independientes, entonces \(\sum_{i=1}^n(x_i-\mu)^2/\sigma^2\in \chi_{n}^2\).}.
Una forma de generalizar la suma de cuadrados anterior al caso
multivariante sería considerar

\[n\bm{S}=\sum_{i=1}^n(\bm{x_i}-\overline{\bm{x}})(\bm{x_i}-\overline{\bm{x}})^\prime{}.\]

Fíjate que \(n\bm{S}\) es una matriz aleatoria \(d\times d\) y por tanto su
distribución ya no podrá ser ji-cuadrado como en el caso
unidimensional. Decimos que \(n\bm{S}\) sigue una distribución de
Wishart\footnote{La distribución de Wishart es el análogo multivariante de la distribución ji-cuadrado y se define formalmente a partir de normales multivariantes estándar. Si \(\bm{x_1},\ldots,\bm{x_n}\) son \(N_d(\bm\mu,\bm\Sigma)\) independientes, entonces
  \[\sum_{i=1}^n(\bm{x_i}-\bm{\mu})(\bm{x_i}-\bm{\mu})^\prime{}\in W_d(\bm\Sigma,n).\]
  Cuando sustituimos \(\bm\mu\) por \(\overline{\bm{x}}\), la distribución de Wishart pierde un grado de libertad.} con parámetros \(\bm\Sigma\) y \(n-1\) y denotamos:

\[n\bm{S}=\sum_{i=1}^n(\bm{x_i}-\bm{\bar{x}})(\bm{x_i}-\bm{\bar{x}})^\prime{}\in W_d(\bm\Sigma,n-1).\]
En el Anexo \ref{anexoDist} puedes encontrar un resumen de las principales
propiedades de la distribución de Wishart.

\section{Contrastes multivariantes versus contrastes univariantes}\label{contrastes-multivariantes-versus-contrastes-univariantes}

Antes de presentar resultados sobre contrastes multivariantes deberíamos
discutir brevemente cual es la motivación para plantear contrastes
multivariantes en lugar de (o junto con) contrastes univariantes.
Consideremos como ejemplo un contraste acerca del vector de medias
\(\bm\mu=(\mu_1,\ldots,\mu_d)^\prime{}\).

\textbf{\includegraphics[width=1em,height=1em]{bookdown_AM_files/figure-latex/fa-icon-00ef1bb1d064b635425d49f06b331298.pdf} ~\textcolor{excolor}{Ejemplo.}} Un estudio sobre Psicología educativa pretende analizar si
las notas medias obtenidas por estudiantes universitarios en dos tests
de inteligencia diferentes fueron iguales a 100. El primer test está
orientado fundamentalmente a la valoración de habilidades aritméticas y
el segundo a la valoración de habilidades de memoria.

Denotaremos por \(\bm{x}=(x_1,x_2)^\prime{}\) el vector aleatorio
bidimensional correspondiente a las notas de los dos tests. Parece claro
que las medidas de habilidad analizadas estarán interrelacionadas y, por
lo tanto, es más natural contrastar si la media
\(\bm\mu=(\mu_1,\mu_2)^\prime\) del vector \(\bm{x}\) es igual a
\(\bm{\mu_0}=(100,100)^\prime{}\) que contrastar de forma individual si la
nota media de cada uno de los tests \(\mu_i\) es igual a 100. Es decir,
nos gustaría contrastar la hipótesis nula \(H_0:\bm{\mu}=\bm{\mu_0}\),
siendo \(\bm{\mu_0}=(100,100)^\prime{}\).

Además de que los contrastes univariantes ignoran la correlación entre
variables, lo cual en muchas ocasiones no es deseable (como se comenta
en el ejemplo anterior), existen otros motivos por los que puede ser más
conveniente plantear contrastes multivariantes. El uso de contrastes
univariantes tiende a \emph{inflar} el error de tipo I, que denotamos por
\(\alpha\). Volvemos al ejemplo sobre los tests de inteligencia. Si para
contrastar la hipótesis nula \(H_0:\bm{\mu}=\bm{\mu_0}\) planteamos dos
contrastes individuales de nivel \(\alpha=0.05\), entonces está claro que
rechazaremos \(H_0\) si rechazamos al menos uno de los dos contrastes
individuales. Pero bajo \(H_0\) (suponiendo que las variables fuesen
independientes),

\[\begin{aligned}
P(\textnormal{al menos un rechazo individual})&=1-P(\textnormal{ningún rechazo individual})\\
&=1-(1-0.05)^2\\
&=0.0975.\end{aligned}\]

Es decir, al plantear el problema a través de contrastes individuales
aumentamos la probabilidad del error de tipo I. Este inconveniente se
hace más patente al aumentar la dimensión \(d\). Además los test
multivariantes resultan en muchos casos más potentes (capaces de
rechazar \(H_0\) cuando es falsa).

\section{Contrastes sobre el vector de medias}\label{contrastes-sobre-el-vector-de-medias}

A continuación veremos cómo se puede usar el test de razón de
verosimilitudes para hacer inferencia en poblaciones normales
multivariantes. En esta sección ilustraremos el caso del problema de
inferencia sobre el vector de medias cuando la matriz de covarianzas es
conocida, y también cuando es desconocida.

\subsection{Contraste sobre el vector de medias con matriz de covarianzas conocida}\label{contraste-sobre-el-vector-de-medias-con-matriz-de-covarianzas-conocida}

En este apartado veremos como llevar a cabo un contraste sobre el vector
de medias de una población normal asumiendo que la matriz de covarianzas
\(\bm\Sigma\) es conocida. Antes, repasaremos la teoría en el caso
unidimensional.

\subsubsection{\texorpdfstring{Revisión del caso unidimensional \(H_0:\mu=\mu_0\) con \(\sigma^2\) conocida}{Revisión del caso unidimensional H\_0:\textbackslash mu=\textbackslash mu\_0 con \textbackslash sigma\^{}2 conocida}}\label{revision-del-caso-unidimensional-h-0-mu-mu-0-con-sigma-2-conocida}

La hipótesis de interés es que la media \(\mu\) de una variable aleatoria
\(x\) es igual a un valor dado \(\mu_0\) frente a la alternativa de que no
es igual a \(\mu_0\), es decir
\[H_0:\mu=\mu_0\textnormal{ frente a }H_1:\mu\neq\mu_0.\] Dada una
muestra \(x_1,\ldots,x_n\) de normales \(N(\mu,\sigma^2)\) con \(\sigma^2\)
conocida, se calcula la media muestral \(\overline{x}\) y se compara con
\(\mu_0\) a través del estadístico

\[z=\frac{\overline{x}-\mu_0}{\sigma/\sqrt{n}}\]

que, bajo la hipótesis
nula, se distribuye como una normal estándar. Por lo tanto, para un
nivel de significación \(\alpha=0.05\) rechazamos la hipótesis nula si
\(\left|z\right|\geq 1.96\). Equivalentemente

\[z^2=n(\overline{x}-\mu_0)\frac{1}{\sigma^2}(\overline{x}-\mu_0)\]
y podríamos utilizar que \(z^2\) se distribuye, bajo la hipótesis nula, como
una ji-cuadrado con un grado de libertad. Entonces rechazaríamos \(H_0\)
si \(z^2\geq 3.84\).

\subsubsection{\texorpdfstring{Contraste multivariante \(H_0:\bm\mu=\bm{\mu_0}\) con \(\bm\Sigma\) conocida}{Contraste multivariante H\_0:\textbackslash bm\textbackslash mu=\textbackslash bm\{\textbackslash mu\_0\} con \textbackslash bm\textbackslash Sigma conocida}}\label{contraste-multivariante-h-0-bm-mu-bm-mu-0-con-bm-sigma-conocida}

Partimos ahora de una muestra aleatoria simple \(\bm{x_1},\ldots,\bm{x_n}\in N_d(\bm{\mu},\bm\Sigma)\)
de vectores
aleatorios independientes y con la misma distribución normal
multivariante.

Suponiendo que la matriz de covarianzas \(\bm\Sigma\) es conocida,
deseamos llevar a cabo tareas de inferencia relativas al vector de
medias \(\bm{\mu}\). En concreto, podemos estar interesados en una región
de confianza para \(\bm{\mu}\), o podemos querer contrastar una hipótesis
nula del tipo \(H_0:\bm{\mu}=\bm{\mu_0}\).

Centrándonos en el contraste de la hipótesis nula
\(H_0:\bm{\mu}=\bm{\mu_0}\), vamos a abordar este problema mediante el
procedimiento de razón de verosimilitudes. En esta situación, el
estadístico de contraste sería:

\[-2\log \lambda(\bm{x})=-2\log \frac{L(\bm{x},\bm{\mu_0},\bm\Sigma)}{\sup_{\bm{\mu}} L(\bm{x},\bm{\mu},\bm\Sigma)}\]
donde la función de verosimilitud es la que se ha tratado en la sección
anterior.

De lo allí expuesto extraemos que, bajo la hipótesis nula,
\(H_0: \bm{\mu}=\bm{\mu_0}\), la función de log-verosimilitud adopta la
forma:

\[\log L(\bm{x},\bm{\mu_0},\bm\Sigma) = c - \frac{n}{2} \log |\bm\Sigma |
-\frac{1}{2} \sum_{i=1}^n \left(\bm{x_i}-\overline{\bm{x}}\right)^\prime{}\bm\Sigma^{-1} \left(\bm{x_i}-\overline{\bm{x}}\right)
-\frac{n}{2}  \left(\overline{\bm{x}}-\bm{\mu_0}\right)^\prime{}\bm\Sigma^{-1} \left(\overline{\bm{x}}-\bm{\mu_0}\right)\]
mientras que bajo la alternativa, aplicando la expresión \eqref{eq:supmu}, se tiene

\[\sup_{\bm{\mu}}\log L(\bm{x},\bm{\mu},\bm\Sigma) = c - \frac{n}{2} \log |\bm\Sigma|
-\frac{1}{2} \sum_{i=1}^n \left(\bm{x_i}-\overline{\bm{x}}\right)^\prime{}\bm\Sigma^{-1} \left(\bm{x_i}-\overline{\bm{x}}\right).\]

En definitiva, el estadístico de contraste resulta:

\[-2\log \lambda(\bm{x})=-2\log \frac{L(\bm{x},\bm{\mu_0},\bm\Sigma)}{\sup_{\bm{\mu}} L(\bm{x},\bm{\mu},\bm\Sigma)}
=n\left(\overline{\bm{x}}-\bm{\mu_0}\right)^\prime{}\bm\Sigma^{-1} \left(\overline{\bm{x}}-\bm{\mu_0}\right).\]

Observamos que si \(H_0:\bm{\mu}=\bm{\mu_0}\) es cierta,

\[\overline{\bm{x}}\in N_d\left(\bm{\mu_0},\bm\Sigma/n\right)\] y, en
consecuencia,

\[n\left(\overline{\bm{x}}-\bm{\mu_0}\right)^\prime{}\bm\Sigma^{-1} \left(\overline{\bm{x}}-\bm{\mu_0}\right)\in\chi_d^2.\]

Así, rechazaremos la hipótesis nula \(H_0:\bm{\mu}=\bm{\mu_0}\) cuando

\[\label{testchi}
n\left(\overline{\bm{x}}-\bm{\mu_0}\right)^\prime{}\bm\Sigma^{-1} \left(\overline{\bm{x}}-\bm{\mu_0}\right)>\chi_{d,\alpha}^2\]
siendo \(\chi_{d,\alpha}^2\) el cuantil \(1-\alpha\) de la distribución
\(\chi_d^2\). La región de rechazo se corresponde con el exterior de un
elipsoide en \(\mathbb{R}^d\).

\textbf{\includegraphics[width=1em,height=1em]{bookdown_AM_files/figure-latex/fa-icon-00ef1bb1d064b635425d49f06b331298.pdf} ~\textcolor{excolor}{Ejemplo.}} Volvemos al estudio sobre Psicología educativa. Se le han
realizado los dos tests de inteligencia a 20 alumnos de un centro
universitario. La nota media obtenida por los 20 alumnos en el test de
valoración de habilidades aritméticas fue de \(102\) puntos. La nota media
en el test de habilidades de memoria fue de \(98\) puntos. Asumamos además
que las notas obtenidas provienen de una distribución normal bivariante
\(N_2(\bm{\mu}, \bm\Sigma)\) siendo

\[\bm\Sigma=\left(\begin{array}{cc}
49 & 35 \\
35 & 52
\end{array}\right).\]
Queremos contrastar si las notas medias de
estudiantes universitarios en los dos tests de inteligencia son iguales
a 100. Es decir, queremos contrastar la hipótesis nula
\(H_0:\bm{\mu}=\bm{\mu_0}\), siendo \(\bm{\mu_0}=(100,100)^\prime{}\).

A partir de la muestra obtenemos que
\(\overline{\bm{x}}=(\overline{x_1},\overline{x_2})^\prime{}=(102,98)^\prime{}\).
Calculamos el valor del estadístico para el contraste multivariante.

\[n\left(\overline{\bm{x}}-\bm{\mu_0}\right)^\prime{}\bm\Sigma^{-1} \left(\overline{\bm{x}}-\bm{\mu_0}\right)=20\left(2,-2\right)\left(\begin{array}{cc}
49 & 35 \\
35 & 52
\end{array}\right)^{-1} \left(\begin{array}{r}
2 \\
-2
\end{array}\right)=10.34014.\]

\footnotesize

\begin{Shaded}
\begin{Highlighting}[]
\SpecialCharTok{\textgreater{}}\NormalTok{ Sigma }\OtherTok{\textless{}{-}} \FunctionTok{matrix}\NormalTok{(}\FunctionTok{c}\NormalTok{(}\DecValTok{49}\NormalTok{, }\DecValTok{35}\NormalTok{, }\DecValTok{35}\NormalTok{, }\DecValTok{52}\NormalTok{), }\AttributeTok{nrow =} \DecValTok{2}\NormalTok{)}
\SpecialCharTok{\textgreater{}}\NormalTok{ xbar }\OtherTok{\textless{}{-}} \FunctionTok{c}\NormalTok{(}\DecValTok{102}\NormalTok{, }\DecValTok{98}\NormalTok{)}
\SpecialCharTok{\textgreater{}}\NormalTok{ mu0 }\OtherTok{\textless{}{-}} \FunctionTok{c}\NormalTok{(}\DecValTok{100}\NormalTok{, }\DecValTok{100}\NormalTok{)}
\SpecialCharTok{\textgreater{}}\NormalTok{ estad }\OtherTok{\textless{}{-}} \DecValTok{20} \SpecialCharTok{*}\NormalTok{ (xbar }\SpecialCharTok{{-}}\NormalTok{ mu0) }\SpecialCharTok{\%*\%} \FunctionTok{solve}\NormalTok{(Sigma) }\SpecialCharTok{\%*\%}\NormalTok{ (xbar }\SpecialCharTok{{-}}\NormalTok{ mu0)}
\SpecialCharTok{\textgreater{}}\NormalTok{ estad}
\NormalTok{         [,}\DecValTok{1}\NormalTok{]}
\NormalTok{[}\DecValTok{1}\NormalTok{,] }\FloatTok{10.34014}
\end{Highlighting}
\end{Shaded}

\normalsize

Rechazamos la hipótesis nula (para un nivel de significación
\(\alpha=0.05\)) ya que el valor del estadístico es mayor que
\(\chi_{d,\alpha}^2=5.991465\). También podemos calcular el \(p\)-valor del
estadístico y razonar a partir de él (como el \(p\)-valor es menor que
\(\alpha=0.05\) se rechaza \(H_0\) para dicho nivel de significación).

\footnotesize

\begin{Shaded}
\begin{Highlighting}[]
\SpecialCharTok{\textgreater{}}\NormalTok{ alpha }\OtherTok{\textless{}{-}} \FloatTok{0.05}
\SpecialCharTok{\textgreater{}} \FunctionTok{qchisq}\NormalTok{(}\DecValTok{1} \SpecialCharTok{{-}}\NormalTok{ alpha, }\DecValTok{2}\NormalTok{)}
\NormalTok{[}\DecValTok{1}\NormalTok{] }\FloatTok{5.991465}
\SpecialCharTok{\textgreater{}}\NormalTok{ pvalue }\OtherTok{\textless{}{-}} \DecValTok{1} \SpecialCharTok{{-}} \FunctionTok{pchisq}\NormalTok{(estad, }\DecValTok{2}\NormalTok{)}
\SpecialCharTok{\textgreater{}}\NormalTok{ pvalue}
\NormalTok{            [,}\DecValTok{1}\NormalTok{]}
\NormalTok{[}\DecValTok{1}\NormalTok{,] }\FloatTok{0.005684182}
\end{Highlighting}
\end{Shaded}

\normalsize

Sin embargo, si planteamos los contrastes individuales \(H_0:\mu_1=100\) y
\(H_0:\mu_2=100\) se tiene que

\[\frac{\overline{x_1}-\mu_0}{\sigma_1/\sqrt{n}}=\frac{102-100}{\sqrt{49/20}}=1.2777\]
y

\[\frac{\overline{x_2}-\mu_0}{\sigma_2/\sqrt{n}}=\frac{98-100}{\sqrt{52/20}}= -1.24034.\]
En ambos casos, el valor absoluto del estadístico es menor que 1.96 y
por lo tanto no se rechazaría la hipótesis nula para ninguno de los dos
contrastes individuales. En la Figura
\ref{fig:test} se representan
las zonas de rechazo y no rechazo para los contrastes univariantes y
para el contraste bivariante correspondiente a los datos de este ejemplo
(\(n=20\)). Si \(\overline{\bm{x}}\) pertenece a la zona verde, no se
rechaza la hipótesis nula \(H_0:\bm\mu=\bm{\mu_0}\) siendo
\(\bm{\mu_0}=(100,100)^\prime\). Si \(\overline{\bm{x}}\) pertenece a la
zona azul, no se rechaza ni \(H_0:\mu_1=100\) ni \(H_0:\mu_2=100\). Fíjate
que el punto
\(\overline{\bm{x}}=(\overline{x_1},\overline{x_2})^\prime{}=(102,98)^\prime{}\)
está fuera de la elipse que representa la zona de no rechazo del test
multivariante y sin embargo, cae en la zona de no rechazo de los dos
contrastes univariantes.

\begin{figure}

{\centering \resizebox{.4\linewidth}{!}{\includegraphics[width=0.5\linewidth]{bookdown_AM_files/figure-latex/test-1} } 

}

\caption{Zonas de rechazo y no rechazo para los contrastes univariantes y para el contraste bivariante correspondiente a los datos de tests de inteligencia ($n=20$). Si $\overline{\bm{x}}$ pertenece a la zona verde, no se rechaza la hipótesis nula $H_0:\bm\mu=\bm{\mu_0}$ siendo $\bm{\mu_0}=(100,100)^\prime$. Si $\overline{\bm{x}}$ pertenece a la zona azul, no se rechaza ni $H_0:\mu_1=100$ ni $H_0:\mu_2=100$.}\label{fig:test}
\end{figure}

\subsection{Contraste sobre el vector de medias con matriz de covarianzas desconocida}\label{contraste-sobre-el-vector-de-medias-con-matriz-de-covarianzas-desconocida}

Veremos ahora como llevar a cabo un contraste sobre el vector de medias
de una población normal asumiendo que la matriz de covarianzas
\(\bm\Sigma\) es desconocida. Repasaremos de nuevo en primer lugar la
teoría en el caso unidimensional.

\subsubsection{\texorpdfstring{Revisión del caso unidimensional \(H_0:\mu=\mu_0\) con \(\sigma^2\) desconocida}{Revisión del caso unidimensional H\_0:\textbackslash mu=\textbackslash mu\_0 con \textbackslash sigma\^{}2 desconocida}}\label{revision-del-caso-unidimensional-h-0-mu-mu-0-con-sigma-2-desconocida}

La hipótesis de interés es que la media \(\mu\) de una variable aleatoria
\(x\) es igual a un valor dado \(\mu_0\) frente a la alternativa de que no
es igual a \(\mu_0\), es decir

\[H_0:\mu=\mu_0\textnormal{ frente a }H_1:\mu\neq\mu_0.\]
Dada una
muestra \(x_1,\ldots,x_n\) de normales \(N(\mu,\sigma^2)\) con \(\sigma^2\)
desconocida, se calcula la media muestral \(\overline{x}\) y se compara
con \(\mu_0\) a través del estadístico

\[t=\frac{\overline{x}-\mu_0}{s_c/\sqrt{n}}\]
siendo
\(s_c=\sum_{i=1}^n(x_i-\overline{x})^2/(n-1)\) la cuasivarianza muestral.
Bajo la hipótesis nula, el estadístico se distribuye como una \(t\) de
Student con \(n-1\) grados de libertad. Por lo tanto, para un nivel de
significación \(\alpha\) rechazamos la hipótesis nula si
\(\left|t\right|\geq t_{\alpha/2}\), siendo \(t_{\alpha/2}\) el cuantil
\(1-\alpha/2\) de la distribución \(t\) de Student con \(n-1\) grados de
libertad. Equivalentemente

\begin{equation}
t^2=n(\overline{x}-\mu_0)\frac{1}{s_c^2}(\overline{x}-\mu_0).
\label{eq:t2}
\end{equation}

Veremos
que esta forma de escribir el estadístico estará relacionada con la
versión multivariante del contraste.

\subsubsection{\texorpdfstring{Contraste multivariante \(H_0:\bm\mu=\bm{\mu_0}\) con \(\bm\Sigma\) desconocida}{Contraste multivariante H\_0:\textbackslash bm\textbackslash mu=\textbackslash bm\{\textbackslash mu\_0\} con \textbackslash bm\textbackslash Sigma desconocida}}\label{contraste-multivariante-h-0-bm-mu-bm-mu-0-con-bm-sigma-desconocida}

Disponemos ahora de una muestra aleatoria simple \(\bm{x_1},\ldots,\bm{x_n}\in N_d(\bm{\bm{\mu}},\bm\Sigma)\) y deseamos
realizar tareas de inferencia relativas al vector de medias \(\bm{\mu}\).
Suponemos que la matriz de covarianzas \(\bm\Sigma\) es desconocida.

El estadístico de razón de verosimilitudes para el contraste de la
hipótesis nula \(H_0:\bm{\mu}=\bm{\mu_0}\) sería:

\[-2\log \lambda(\bm{x})
=-2\log \frac{\sup_{\bm\Sigma} L(\bm{x},\bm{\mu_0},\bm\Sigma)}{\sup_{\bm{\mu},\bm\Sigma} L(\bm{x},\bm{\mu},\bm\Sigma)}.\]

Nótese que ahora, al ser \(\bm\Sigma\) desconocida, se convierte en un
parámetro tanto bajo la hipótesis nula como bajo la alternativa,
parámetro que será estimado por máxima verosimilitud.

Bajo la alternativa, hemos visto en la sección anterior que los
estimadores de máxima verosimilitud del vector de medias y la matriz de
covarianzas (sin restricciones) son \(\overline{\bm{x}}\) y \(\bm{S}\),
respectivamente. Asimismo, la función de verosimilitud tiene como
máximo:

\[\sup_{\bm\Sigma}\sup_{\bm{\mu}}\log L(\bm{x},\bm{\mu},\bm\Sigma)
= c - \frac{n}{2}
\left( \log |\bm{S}| + \mbox{traza} \left( \bm{S}^{-1} \bm{S} \right)\right)
= c - \frac{n}{2} \left( \log |\bm{S}| + d\right).\]

A continuación maximizamos la verosimilitud bajo la hipótesis nula. Para
ello basta con expresar la verosimilitud en una forma similar a la
anterior:

\[\begin{aligned}
\log L(\bm{x},\bm{\mu_0},\bm\Sigma) &= c - \frac{n}{2} \log |\bm\Sigma|
-\frac{1}{2} \mbox{traza} \left[\sum_{i=1}^n
\left(\bm{x_i}-\bm{\mu_0}\right)^\prime{}\bm\Sigma^{-1} \left(\bm{x_i}-\bm{\mu_0}\right)\right] \\
&= c - \frac{n}{2} \log |\bm\Sigma|
-\frac{1}{2} \sum_{i=1}^n \mbox{traza}
\left[ \bm\Sigma^{-1} \left(\bm{x_i}-\bm{\mu_0}\right)  \left(\bm{x_i}-\bm{\mu_0}\right)^\prime{} \right] \\
&= c - \frac{n}{2} \left( \log |\bm\Sigma|
+ \mbox{traza} \left( \bm\Sigma^{-1} \hat{\bm\Sigma}_{\bm{\mu_0}} \right)\right)\end{aligned}\]

siendo \(\hat{\bm\Sigma}_{\bm{\mu_0}}=\frac{1}{n}\sum_{i=1}^n
\left(\bm{x_i}-\bm{\mu_0}\right)  \left(\bm{x_i}-\bm{\mu_0}\right)^\prime{}\),
el cual resulta ser un estimador razonable de la matriz de covarianzas
bajo la hipótesis de que la media vale \(\bm{\mu_0}\). Por lo demás los
pasos son idénticos al caso anterior, salvo que se ha puesto
\(\bm{\mu_0}\) allí donde se hallaba \(\overline{\bm{x}}\). Aplicando de
nuevo el lema, concluimos que \(\hat{\bm\Sigma}_{\bm{\mu_0}}\) es el
estimador de máxima verosimilitud de la matriz de covarianzas bajo la
hipótesis nula, y que la función de verosimilitud bajo dicha hipótesis
alcanza el valor máximo:

\[\sup_{\bm\Sigma}\log L(\bm{x},\bm{\mu_0},\bm\Sigma)
= c - \frac{n}{2} \left( \log |\hat{\bm\Sigma}_{\bm{\mu_0}}| + d \right).\]

Entonces el estadístico de contraste mediante la razón de
verosimilitudes resulta:

\[-2\log \lambda(\bm{x})
=-2\log \frac{\sup_{\bm\Sigma} L(\bm{x},\bm{\mu_0},\bm\Sigma)}{\sup_{\bm{\mu},\bm\Sigma} L(\bm{x},\bm{\mu},\bm\Sigma)}
=n\left( \log |\hat{\bm\Sigma}_{\bm{\mu_0}}| - \log |\bm{S}| \right).\]

Descomponemos

\[\begin{aligned}
\hat{\bm\Sigma}_{\bm{\mu_0}} &=
\frac{1}{n} \sum_{i=1}^n \left(\bm{x_i}-\bm{\mu_0}\right) \left(\bm{x_i}-\bm{\mu_0}\right)^\prime{} \\
&= \frac{1}{n} \sum_{i=1}^n \left[ \left(\bm{x_i}-\overline{\bm{x}}\right) \left(\bm{x_i}-\overline{\bm{x}}\right)^\prime{}
 + \left(\overline{\bm{x}}-\bm{\mu_0}\right) \left(\overline{\bm{x}}-\bm{\mu_0}\right)^\prime{}
+ 2 \left(\overline{\bm{x}}-\bm{\mu_0}\right) \left(\bm{x_i}-\overline{\bm{x}}\right)^\prime{} \right] \\
&= \bm{S}+\bm{r}\bm{r}^\prime{}\end{aligned}\]
siendo
\(\bm{r}=\overline{\bm{x}}-\bm{\mu_0}\). Sustituyendo en el estadístico de
contraste obtenemos

\[\begin{aligned}
-2\log \lambda(\bm{x})
&= n\left( \log |\bm{S}+\bm{r}\bm{r}^\prime{}| - \log |\bm{S}| \right) \\
&= n\left( \log \left(|\bm{S}|\cdot \left|\bm{I_d}+\bm{S}^{-1}\bm{r}\bm{r}^\prime{}\right|\right) - \log |\bm{S}| \right)
\\
&=n \log \left|\bm{I_d}+\bm{S}^{-1}\bm{r}\bm{r}^\prime{}\right|.\end{aligned}\]

Estudiemos, pues, el determinante que aparece en el último término. En
(a) denotamos mediante \(\lambda_1,\ldots,\lambda_d\) a los autovalores de
\(\bm{S}^{-1}\bm{r}\bm{r}^\prime{}\), y observamos que
\(1+\lambda_1,\ldots,1+\lambda_d\) son los autovalores de
\(\bm{I_d}+\bm{S}^{-1}\bm{r}\bm{r}^\prime{}\). En (b) y (c) usamos que la
matriz \(\bm{S}^{-1}\bm{r}\bm{r}^\prime{}\) es de rango uno.

\[
\begin{aligned}
\left|\bm{I_d}+\bm{S}^{-1}\bm{r}\bm{r}^\prime{}\right|&\stackrel{(a)}{=}
\prod_{j=1}^d \left(1+\lambda_j\right)\\
&\stackrel{(b)}{=}
1+\lambda_1\\
&\stackrel{(c)}{=}
1+{\textnormal{traza}}\left(\bm{S}^{-1}\bm{r}\bm{r}^\prime{}\right)\\
&=1+{\textnormal{traza}}\left(\bm{r}^\prime{}\bm{S}^{-1}\bm{r}\right)\\
&=1+\bm{r}^\prime{}\bm{S}^{-1}\bm{r}.
\end{aligned}\]

Finalmente,

\[-2\log \lambda(\bm{x})=n\log\left(1+\bm{r}^\prime{}\bm{S}^{-1}\bm{r}\right)\]

será el estadístico de contraste y rechazaremos la hipótesis nula si
este estadístico toma un valor demasiado grande.

Será equivalente si consideramos el estadístico

\[\bm{r}^\prime{}\bm{S}^{-1}\bm{r}=\left(\overline{\bm{x}}-\bm{\mu_0}\right)^\prime{} \bm{S}^{-1} \left(\overline{\bm{x}}-\bm{\mu_0}\right)\]

y rechazamos la hipótesis nula cuando este nuevo estadístico toma un
valor demasiado grande. Nótese que el estadístico anterior se obtiene
tras aplicar una transformación creciente a este último.

El estadístico recuerda al comentado en \eqref{eq:t2} para el caso
univariante si tenemos en cuenta que:

\[T^2=(n-1)\left(\overline{\bm{x}}-\bm{\mu_0}\right)^\prime{} \bm{S}^{-1} \left(\overline{\bm{x}}-\bm{\mu_0}\right)=n\left(\overline{\bm{x}}-\bm{\mu_0}\right)^\prime{} \bm{S_c}^{-1} \left(\overline{\bm{x}}-\bm{\mu_0}\right)\]

donde \(\bm{S_c}=n/(n-1)\bm{S}\) denota la matriz de covarianzas
corregida. La distribución de \(T^2\) bajo la hipótesis nula fue obtenida
por \citet{hot36}. Se tiene que, si \(H_0\) es cierta, el estadístico
\(T^2\) sigue una distribución \(\Gamma^2\) de Hotelling\footnote{Formalmente, dadas \(\bm{Z}\in N_d(\bm{0},\bm{I_d})\) y
  \(\bm{R}\in W_d(\bm{I_d},m)\) independientes, con \(m\geq d\) se dice
  que la variable aleatoria unidimensional
  \(m\bm{Z}^{\prime} \bm{R}^{-1} \bm{Z} \in \Gamma^2(d,m)\) tiene una
  distribución \(\Gamma^2\) de Hotelling de parámetros \(d\) y \(m\).} de parámetros
\(d\) y \(n-1\). La distribución \(\Gamma^2\) de Hotelling juega un papel
semejante a la distribución \(t\) de Student en el caso multivariante.
Puedes encontrar propiedades y resultados sobre esta distribución en el
Anexo \ref{anexoDist}. Los valores críticos de la distribución \(\Gamma^2\) de
Hotelling están tabulados. Aun así, se puede convertir una variable con
distribución \(\Gamma^2\) de Hotelling en otra con distribución \(F\) de
Snédecor. De este modo se evita la necesidad de considerar nuevas tablas
de cuantiles. Simplemente se efectuaría la transformación y se apelaría
a las tablas y operaciones conocidas para la \(F\) de Snédecor. En nuestro
caso se tiene que, bajo \(H_0\),

\[\frac{n-d}{d} \left(\overline{\bm{x}}-\bm{\mu_0}\right)^\prime{} \bm{S}^{-1} \left(\overline{\bm{x}}-\bm{\mu_0}\right)
\in F_{d,n-d}.\]

En definitiva, rechazaremos la hipótesis nula \(H_0: \bm{\mu}=\bm{\mu_0}\)
si

\[\frac{n-d}{d} \left(\overline{\bm{x}}-\bm{\mu_0}\right)^\prime{} \bm{S}^{-1} \left(\overline{\bm{x}}-\bm{\mu_0}\right)
> f_{d,n-d,\alpha}\]
siendo \(f_{d,n-d,\alpha}\) el cuantil \(1-\alpha\) de
la distribución \(F_{d,n-d}\).

\textbf{\includegraphics[width=1em,height=1em]{bookdown_AM_files/figure-latex/fa-icon-00ef1bb1d064b635425d49f06b331298.pdf} ~\textcolor{excolor}{Ejemplo.}} El fichero \texttt{notas.txt} contiene las puntuaciones obtenidas
por los 20 alumnos que han realizado los tests de inteligencia citados
en los ejemplos previos. Asumamos ahora que las notas obtenidas
provienen de una distribución normal bivariante
\(N_2(\bm{\mu}, \bm\Sigma)\) pero que \(\bm\Sigma\) es desconocida. Queremos
contrastar la hipótesis nula \(H_0:\bm{\mu}=\bm{\mu_0}\), siendo
\(\bm{\mu_0}=(100,100)^\prime{}\).

\footnotesize

\begin{Shaded}
\begin{Highlighting}[]
\SpecialCharTok{\textgreater{}}\NormalTok{ tests }\OtherTok{\textless{}{-}} \FunctionTok{read.table}\NormalTok{(}\StringTok{"data/notas.txt"}\NormalTok{, }\AttributeTok{header =} \ConstantTok{TRUE}\NormalTok{)}
\SpecialCharTok{\textgreater{}}\NormalTok{ n }\OtherTok{\textless{}{-}} \FunctionTok{dim}\NormalTok{(tests)[}\DecValTok{1}\NormalTok{]}
\SpecialCharTok{\textgreater{}}\NormalTok{ d }\OtherTok{\textless{}{-}} \FunctionTok{dim}\NormalTok{(tests)[}\DecValTok{2}\NormalTok{]}
\SpecialCharTok{\textgreater{}}\NormalTok{ xbar }\OtherTok{\textless{}{-}} \FunctionTok{colMeans}\NormalTok{(tests)}
\SpecialCharTok{\textgreater{}}\NormalTok{ xbar}
\NormalTok{Test1 Test2 }
  \DecValTok{102}    \DecValTok{98} 
\SpecialCharTok{\textgreater{}}\NormalTok{ S }\OtherTok{\textless{}{-}}\NormalTok{ (n }\SpecialCharTok{{-}} \DecValTok{1}\NormalTok{)}\SpecialCharTok{/}\NormalTok{n }\SpecialCharTok{*} \FunctionTok{cov}\NormalTok{(tests)}
\SpecialCharTok{\textgreater{}}\NormalTok{ S}
\NormalTok{      Test1 Test2}
\NormalTok{Test1  }\FloatTok{84.1}  \FloatTok{38.3}
\NormalTok{Test2  }\FloatTok{38.3}  \FloatTok{32.6}
\SpecialCharTok{\textgreater{}}\NormalTok{ mu0 }\OtherTok{\textless{}{-}} \FunctionTok{c}\NormalTok{(}\DecValTok{100}\NormalTok{, }\DecValTok{100}\NormalTok{)}
\SpecialCharTok{\textgreater{}}\NormalTok{ estad }\OtherTok{\textless{}{-}}\NormalTok{ (n }\SpecialCharTok{{-}}\NormalTok{ d)}\SpecialCharTok{/}\NormalTok{d }\SpecialCharTok{*}\NormalTok{ (xbar }\SpecialCharTok{{-}}\NormalTok{ mu0) }\SpecialCharTok{\%*\%} \FunctionTok{solve}\NormalTok{(S) }\SpecialCharTok{\%*\%}\NormalTok{ (xbar }\SpecialCharTok{{-}}\NormalTok{ mu0)}
\SpecialCharTok{\textgreater{}}\NormalTok{ estad}
\NormalTok{         [,}\DecValTok{1}\NormalTok{]}
\NormalTok{[}\DecValTok{1}\NormalTok{,] }\FloatTok{5.458867}
\end{Highlighting}
\end{Shaded}

\normalsize

Rechazamos la hipótesis nula (para un nivel de significación
\(\alpha=0.05\)) ya que el valor del estadístico es mayor que
\(f_{d,n-d,\alpha}=3.554557\). También podemos calcular el \(p\)-valor del
estadístico y razonar a partir de él (como el \(p\)-valor es menor que
\(\alpha=0.05\) se rechaza \(H_0\) para dicho nivel de significación).

\footnotesize

\begin{Shaded}
\begin{Highlighting}[]
\SpecialCharTok{\textgreater{}}\NormalTok{ alpha }\OtherTok{\textless{}{-}} \FloatTok{0.05}
\SpecialCharTok{\textgreater{}} \FunctionTok{qf}\NormalTok{(}\DecValTok{1} \SpecialCharTok{{-}}\NormalTok{ alpha, d, n }\SpecialCharTok{{-}}\NormalTok{ d)}
\NormalTok{[}\DecValTok{1}\NormalTok{] }\FloatTok{3.554557}
\SpecialCharTok{\textgreater{}}\NormalTok{ pvalue }\OtherTok{\textless{}{-}} \DecValTok{1} \SpecialCharTok{{-}} \FunctionTok{pf}\NormalTok{(estad, d, n }\SpecialCharTok{{-}}\NormalTok{ d)}
\SpecialCharTok{\textgreater{}}\NormalTok{ pvalue}
\NormalTok{          [,}\DecValTok{1}\NormalTok{]}
\NormalTok{[}\DecValTok{1}\NormalTok{,] }\FloatTok{0.0140273}
\end{Highlighting}
\end{Shaded}

\normalsize

En la Figura \ref{fig:test3a} se representa en rojo la zona de no rechazo del
contraste bivariante correspondiente a los datos de este ejemplo
(\(n=20\)) en comparación con la zona de no rechazo cuando se suponía que
la matriz de covarianzas era conocida (elipse en verde).

\begin{figure}

{\centering \resizebox{.4\linewidth}{!}{\includegraphics[width=0.5\linewidth]{bookdown_AM_files/figure-latex/test3a-1} } 

}

\caption{Zonas de rechazo y no rechazo para el contraste bivariante correspondiente a los datos de tests de inteligencia ($n=20$). Para el caso en que se conoce la matriz de covarianzas, rechazamos la hipótesis nula si $\overline{\bm{x}}$ no pertenece a la elipse en verde. Para el caso en que se desconoce la matriz de covarianzas, rechazamos la hipótesis nula si $\overline{\bm{x}}$ no pertenece a la elipse en rojo.}\label{fig:test3a}
\end{figure}

\section{Regiones de confianza y comparaciones simultáneas}\label{regiones-de-confianza-y-comparaciones-simultaneas}

\subsection{Regiones de confianza}\label{regiones-de-confianza}

A partir del estadístico \(\Gamma^2\) de Hotelling, podemos obtener una
región de confianza para el vector de medias, de la forma:

\[\left\{\bm{\mu}\in{\Bbb R}^d:
\frac{n-d}{d}\left(\overline{\bm{x}}-\bm{\mu}\right)^\prime{}\bm{S}^{-1}\left(\overline{\bm{x}}-\bm{\mu}\right)<f_{d,n-d,\alpha}\right\}.\]

Esta región constituye un elipsoide en \({\Bbb R}^d\), centrado en
\(\overline{\bm{x}}\), cuyos ejes van en la dirección de los autovectores
de \(\bm{S}\) y la longitud de los radios (semilongitud de los ejes) viene
dada por

\[\sqrt{\lambda_j}\sqrt{\frac{d}{n-d}f_{d,n-d,\alpha}}\qquad j\in\{1,\ldots,d\}\]

siendo \(\lambda_1,\ldots,\lambda_d\) los autovalores de \(\bm{S}\). En la
Figura \ref{fig:region} se
muestra la región de confianza (\(\alpha=0.05\)) para el vector de notas
medias en el ejemplo de tests de inteligencia. El elipsoide se calcula
en base al estadístico \(\Gamma^2\) de Hotelling.

\begin{figure}

{\centering \resizebox{.4\linewidth}{!}{\includegraphics[width=0.5\linewidth]{bookdown_AM_files/figure-latex/region-1} } 

}

\caption{Región de confianza ($\alpha=0.05$) para el vector de notas medias en el ejemplo de tests de inteligencia.}\label{fig:region}
\end{figure}

\subsection{Intervalos de confianza individuales}\label{intervalos-de-confianza-individuales}

A continuación planteamos el problema de conseguir intervalos de
confianza para las componentes del vector de medias, o más en general,
para combinaciones lineales del tipo

\[\bm{a}^\prime{}\bm{\mu}=a_1\mu_1+\cdots+a_d\mu_d.\]

Observando que
\(\bm{a}^\prime{}\bm{x_1},\ldots,\bm{a}^\prime{}\bm{x_n}\in N(\bm{a}^\prime{}\bm{\mu},\bm{a}^\prime{}\bm\Sigma\bm{a})\)
y además son independientes, podemos abordar este problema, que ya es
univariante, mediante el procedimiento de la \(t\) de Student. Así, como
la media y la cuasivarianza muestrales calculadas sobre las
observaciones \(\bm{a}^\prime{}\bm{x_1},\ldots,\bm{a}^\prime{}\bm{x_n}\)
resultan ser \(\bm{a}^\prime{}\overline{\bm{x}}\) y
\(\bm{a}^\prime{}\bm{S_c}\bm{a}\), respectivamente, el intervalo de
confianza adopta la forma

\[\left(\bm{a}^\prime{}\overline{\bm{x}}-t_{n-1,\alpha/2}\frac{\sqrt{\bm{a}^\prime{}\bm{S_c}\bm{a}}}{\sqrt{n}},
\bm{a}^\prime{}\overline{\bm{x}}+t_{n-1,\alpha/2}\frac{\sqrt{\bm{a}^\prime{}\bm{S_c}\bm{a}}}{\sqrt{n}}\right)\]
siendo \(t_{n-1,\alpha/2}\) el cuantil \(1-\alpha/2\) de la distribución \(t\)
de Student con \(n-1\) grados de libertad. De este modo, para un \(\bm{a}\)
fijo, el intervalo anterior contiene a \(\bm{a}^\prime{}\bm{\mu}\) con una
probabilidad \(1-\alpha\). En particular, podemos pensar en un vector de
la forma \(\bm{a}=(1,0,\ldots,0)^\prime{}\) que serviría para extraer la
primera componente del vector aleatorio. Igual se haría con las demás
componentes mediante los vectores canónicos correspondientes. Así
obtendríamos \(p\) intervalos de confianza, uno para cada componente del
vector de medias. Sin embargo, el nivel de confianza se refiere a la
probabilidad individual de cada intervalo, de modo que la probabilidad
de que todos los intervalos simultáneamente contengan a la componente
correspondiente del vector de medias será en general inferior al nivel
de confianza fijado.

\subsection{Intervalos de confianza simultáneos por el método de Scheffé}\label{intervalos-de-confianza-simultaneos-por-el-metodo-de-scheffe}

Para satisfacer un nivel de confianza simultáneo, debemos modificar la
construcción de los intervalos haciéndolos más amplios. Vamos a plantear
este objetivo de manera simultánea en todos los vectores \(\bm{a}\). Si
seguimos partiendo como pivote de la media estudentizada, la idea podría
ser cambiar el valor \(t_{n-1,\alpha/2}\) por otra constante adecuada,
previsiblemente más grande. Así, si

\[P\left\{ \left|\frac{\sqrt{n}\left(\bm{a}^\prime{}\overline{\bm{x}}-\bm{a}^\prime{}\bm{\mu}\right)}{\sqrt{\bm{a}^\prime{}\bm{S_c}\bm{a}}}\right| < c
\qquad \forall \bm{a}\in{\Bbb R}^p\right\}=1-\alpha\] los intervalos de
confianza obtenidos al sustituir \(t_{n-1,\alpha/2}\) por \(c\) cumplirán el
nivel de confianza de manera simultánea. Enunciamos el lema siguiente.

\begin{lemma}
\protect\hypertarget{lem:unnamed-chunk-20}{}\label{lem:unnamed-chunk-20}Sea \(\bm{B}\) una matriz \(d\times d\), simétrica y definida positiva, y
\(\bm{r}\in {\Bbb R}^d\). Entonces

\[\max_{\bm{x}\in{\Bbb{R}}^d\backslash\{\bm{0}\}}\frac{\left(\bm{x}^\prime{}\bm{r}\right)^2}{\bm{x}^\prime{}\bm{B}\bm{x}}=\bm{r}^\prime{}\bm{B}^{-1}\bm{r}\]
y este máximo se alcanza cuando \(\bm{x}=c\bm{B}^{-1}\bm{r}\) para
cualquier \(c\in {\Bbb R}\backslash \{0\}\).
\end{lemma}

Aplicando este lema obtenemos

\[\max_{\bm{a}\in{\Bbb R}^d}\frac{n\left(\bm{a}^\prime{}\left(\overline{\bm{x}}-\bm{\mu}\right)\right)^2}{\bm{a}^\prime{}\bm{S_c}\bm{a}}
=n\left(\overline{\bm{x}}-\bm{\mu}\right)^\prime{}\bm {S_c}^{-1}\left(\overline{\bm{x}}-\bm{\mu}\right)\in \Gamma^2(d,n-1).\]
De este resultado se puede extraer el valor de \(c\) y finalmente se
obtienen los que se conocen como intervalos de confianza simultáneos por
el método de Scheffé:

\[\left(\ \bm{a}^\prime{}\overline{\bm{x}} - \sqrt{\frac{d(n-1)}{n-d}f_{d,n-d,\alpha}}\frac{\sqrt{\bm{a}^\prime{}\bm{S_c}\bm{a}}}{\sqrt{n}}\ ,
 \bm{a}^\prime{}\overline{\bm{x}} + \sqrt{\frac{d(n-1)}{n-d}f_{d,n-d,\alpha}}\frac{\sqrt{\bm{a}^\prime{}\bm {S_c}\bm{a}}}{\sqrt{n}}\ \right)\]

o equivalentemente

\[\left(\ \bm{a}^\prime{}\overline{\bm{x}} - \sqrt{\frac{d}{n-d}f_{d,n-d,\alpha}\bm{a}^\prime{}\bm{S}\bm{a}}\ ,
 \bm{a}^\prime{}\overline{\bm{x}} + \sqrt{\frac{d}{n-d}f_{d,n-d,\alpha}\bm{a}^\prime{}\bm{S}\bm{a}}\ \right)\]

\textbf{Observación.} Se puede demostrar que los intervalos de confianza
simultáneos por el método de Scheffé coinciden con la proyección de la
elipse de confianza sobre cada una de las direcciones determinadas por
los vectores \(\bm{a}\). En concreto, si en \(\bm{a}\) se consideran las
direcciones de los ejes, se obtendrían intervalos de confianza para cada
componente.

La tabla siguiente permite comparar los valores de \(c\) para el cálculo
de los intervalos de confianza, extraídos de la \(t\) de Student frente a
los que se obtienen mediante el método de Scheffé.

{\def\LTcaptype{none} % do not increment counter
\begin{longtable}[]{@{}llcl@{}}
\toprule\noalign{}
\endhead
\bottomrule\noalign{}
\endlastfoot
& & \(\sqrt{\frac{d(n-1)}{n-d}f_{d,n-d,0.05}}\) & \\
\(n\) & \(t_{n-1,0.025}\) & \(d=4\) & \(d=10\) \\
15 & 2.145 & 4.14 & 11.52 \\
25 & 2.064 & 3.60 & 6.39 \\
50 & 2.010 & 3.31 & 5.05 \\
100 & 1.970 & 3.19 & 4.61 \\
\(\infty\) & 1.960 & 3.08 & 4.28 \\
\end{longtable}
}

\subsection{Intervalos de confianza simultáneos por el método de Bonferroni}\label{intervalos-de-confianza-simultaneos-por-el-metodo-de-bonferroni}

Otro método para obtener intervalos de confianza simultáneos es el
método de Bonferroni. Es una alternativa válida en cualquier contexto en
el que se requiera una cantidad finita de intervalos simultáneos, ya que
no se basa en la naturaleza probabilística del problema en cuestión
(como sí lo hace el método de Scheffé), sino que su fundamento radica
simplemente en la subaditividad de la probabilidad.

Si \(C_1,\ldots,C_m\) consisten en los sucesos respectivos de que cada
intervalo de confianza contenga al parámetro correspondiente,

\[\begin{aligned}
P(C_i\ \mbox{cierto}\ \forall i) &= 1-P(C_i\ \mbox{falso para algún}\ i) \\
&\geq  1-\sum_{i=1}^m P(C_i\ \mbox{falso})=1-\left(\alpha_1+\cdots +\alpha_m\right)\end{aligned}\]

siendo \(1-\alpha_1,\ldots,1-\alpha_m\) los niveles de confianza
individuales de cada intervalo. Así, para alcanzar un nivel de confianza
simultáneo \(1-\alpha\) basta con tomar \(\alpha_1,\ldots,\alpha_m\) de modo
que \(\alpha_1+\cdots+\alpha_m=\alpha\), por ejemplo mediante
\(\alpha_1=\cdots=\alpha_m=\alpha/m\).

La tabla siguiente muestra el cociente de longitudes (para
\(1-\alpha=0.95\)), cuando se construyen intervalos para cada componente:

\[\frac{\mbox{Longitud del intervalo de Bonferroni}}
{\mbox{Longitud del intervalo de Scheffé}}
=\frac{t_{n-1,\alpha/(2m)}}{\sqrt{\frac{d(n-1)}{n-d}f_{d,n-d,\alpha}}}.\]

{\def\LTcaptype{none} % do not increment counter
\begin{longtable}[]{@{}llrl@{}}
\toprule\noalign{}
\endhead
\bottomrule\noalign{}
\endlastfoot
& & \(m=d\) & \\
\(n\) & 2 & 4 & 10 \\
15 & 0.88 & 0.69 & 0.29 \\
25 & 0.90 & 0.75 & 0.48 \\
50 & 0.91 & 0.78 & 0.58 \\
100 & 0.91 & 0.80 & 0.62 \\
\(\infty\) & 0.91 & 0.81 & 0.66 \\
\end{longtable}
}

\textbf{\includegraphics[width=1em,height=1em]{bookdown_AM_files/figure-latex/fa-icon-00ef1bb1d064b635425d49f06b331298.pdf} ~\textcolor{excolor}{Ejemplo.}} Sobre el ejemplo de las puntuaciones de tests de
inteligencia, vamos a calcular los intervalos de confianza para la
calificación media del test 1 y del test 2, al nivel de confianza del
95\%, obtenidos de manera individual, y simultáneos por el método de
Scheffé y por el método de Bonferroni. Así, en la Figura \ref{fig:tema2figura1} se
representan mediante una elipse en color rojo la región de confianza
para el vector de medias \(\bm{\mu}=(\bm{\mu_1}, \bm{\mu_2})^\prime{}\).
En el centro de la elipse se sitúa el vector con las medias muestrales
\(\overline{\bm{x}}=(102, 98)^\prime{}\).

Al proyectar la elipse sobre los ejes horizontal y vertical se obtienen
los intervalos representados en color verde, que serían intervalos de
confianza simultáneos por el método de Scheffé. Los intervalos de
confianza simultáneos por el método de Bonferroni vienen representados
en color negro. Como cabía esperar, son más pequeños que los intervalos
simultáneos por el método de Scheffé. Por último, los intervalos de
confianza individuales han sido representados con trazo azul. Son los
más pequeños, también conforme a lo previsible.

\begin{figure}

{\centering \resizebox{.4\linewidth}{!}{\includegraphics[width=0.5\linewidth]{bookdown_AM_files/figure-latex/tema2figura1-1} } 

}

\caption{Región de confianza para el vector de medias (elipse), junto con intervalos de confianza individuales (trazo azul), simultáneos por el método de Bonferroni (trazo negro) y simultáneos por el método de Scheffé (trazo verde).}\label{fig:tema2figura1}
\end{figure}

\section{Generalización del contraste sobre el vector de medias}\label{generalizacion-del-contraste-sobre-el-vector-de-medias}

En esta sección veremos cómo se puede generalizar el contraste sobre el
vector de medias, al caso de restricciones más genéricas sobre
\(\bm{\mu}\), más generales que la hipótesis nula,
\(H_0:\bm{\mu}=\bm{\mu_0}\). El resultado básico lo enunciamos como un
teorema. Después, como aplicación más común de este resultado, veremos
el contraste de restricciones lineales sobre \(\bm{\mu}\), entre las
cuales tiene un interés especial el contraste de igualdad de las
componentes del vector de medias.

\subsection{Resultado general}\label{resultado-general}

\begin{theorem}
\protect\hypertarget{thm:unnamed-chunk-21}{}\label{thm:unnamed-chunk-21}Sea \(\bm{x_1},\ldots,\bm{x_n}\) una muestra aleatoria simple de
\(N_d(\bm{\mu},\bm\Sigma)\). Si las hipótesis \(H_0\) y \(H_a\) conducen a los
estimadores de máxima verosimilitud \(\hat{\bm{\mu}}\) y
\(\overline{\bm{x}}\), respectivamente, y bajo ninguna de las dos
hipótesis hay restricciones para \(\bm\Sigma\), entonces los estimadores
de máxima verosimilitud de \(\bm\Sigma\) son
\(\bm{S}+\bm{r}\bm{r}^\prime{}\) y \(\bm{S}\), bajo \(H_0\) y \(H_a\)
respectivamente, siendo \(\bm{r}=\overline{\bm{x}}-\hat{\bm{\mu}}\).

Además, el test de razón de verosimilitudes para contrastar \(H_0\) frente
a \(H_a\) viene dado por

\begin{equation}
-2\log\lambda(\bm{x})=n\bm{r}^\prime{}\bm\Sigma^{-1}\bm{r} \qquad \mbox{si}\ \bm\Sigma\ \mbox{es conocida}
  \label{eq:conocido}
\end{equation}

y

\begin{equation}
-2\log\lambda(\bm{x})=n\log\left(1+\bm{r}^\prime{}\bm{S}^{-1}\bm{r}\right)
 \qquad \mbox{si}\ \bm\Sigma\ \mbox{es desconocida.}
  \label{eq:desconocido}
\end{equation}
\end{theorem}

\begin{proof}
La demostración seguiría los mismos pasos que en los casos
anteriores, donde contrastábamos una hipótesis nula simple sobre el
vector de medias con matriz de covarianzas conocida o desconocida.
\end{proof}

\subsection{Contraste de restricciones lineales}\label{contraste-de-restricciones-lineales}

Supongamos que \(\bm\Sigma\) es conocida y deseamos contrastar la
hipótesis nula

\[H_0: \bm{A}\bm{\mu}=\bm{b}\] siendo \(\bm{A}\) una matriz
conocida de orden \(q\times d\) y rango máximo \(q\), y \(\bm{b}\) un vector
conocido.

A este problema de contraste le podemos aplicar el teorema anterior.
Para ello, tenemos que obtener el estimador de máxima verosimilitud bajo
\(H_0\), que denotaremos mediante \(\hat{\bm{\mu}}\). La función de log--verosimilitud se puede escribir así:

\[\begin{aligned}
l(\bm{x},\bm{\mu},\bm\Sigma)&=\log L(\bm{x},\bm{\mu},\bm\Sigma)\\
&= c-\frac{n}{2} \log |\bm\Sigma|
-\frac{1}{2} \sum_{i=1}^n \left(\bm{x_i}-\overline{\bm{x}}\right)^\prime{}\bm\Sigma^{-1}\left(\bm{x_i}-\overline{\bm{x}}\right)\\
&\qquad -\frac{n}{2} \left(\overline{\bm{x}}-\bm{\mu}\right)^\prime{}\bm\Sigma^{-1}\left(\overline{\bm{x}}-\bm{\mu}\right).
\end{aligned}\]

En tal caso, el problema consiste en:

\[
\begin{aligned}
  \textnormal{Maximizar  }&\ \ l(\bm{x}, \bm{\mu},\bm\Sigma)\\
  \textnormal{sujeto a  }&\ \ \bm{A}\bm{\mu}=\bm{b}
\end{aligned}
\]

Consideramos la función

\[\bm{a}^+=\bm{a}-n\bm{\lambda}^\prime{}(\bm{A}\bm{\mu}-\bm{b})\] siendo
\(\bm{\lambda}\) un vector de multiplicadores de Lagrange. Derivando

\[\frac{\partial \bm{a}^+}{\partial \bm{\mu}}=n\left(\overline{\bm{x}}-\bm{\mu}\right)^\prime{}\bm\Sigma^{-1}-n\bm{\lambda}^\prime{}\bm{A}=0.\]
De donde

\begin{equation}
\overline{\bm{x}}-\bm{\mu}=\bm\Sigma \bm{A}^\prime{} \bm{\lambda}
\label{eq:verosirestricta}
\end{equation}

ecuación que debemos añadir a la restricción \(\bm{A}\bm{\mu}=\bm{b}\),
para obtener las soluciones para \(\bm{\lambda}\) y \(\bm{\mu}\).
Multiplicando por \(\bm{A}\),
\[\bm{A}\overline{\bm{x}}-\bm{A}\bm{\mu}=\bm{A}\overline{\bm{x}}-\bm{b}=\left(\bm{A}\bm\Sigma \bm{A}^\prime{}\right)\bm{\lambda}\]
lo cual nos permite despejar
\(\bm{\lambda}=\left(\bm{A}\bm\Sigma \bm{A}^\prime{}\right)^{-1}\left(\bm{A}\overline{\bm{x}}-\bm{b}\right)\)
que, sustituido en la ecuación \eqref{eq:verosirestricta}, da lugar al estimador de máxima
verosimilitud

\[\hat{\bm{\mu}}=\overline{\bm{x}}-\bm\Sigma \bm{A}^\prime{} \left(\bm{A}\bm\Sigma \bm{A}^\prime{}\right)^{-1}\left(\bm{A}\overline{\bm{x}}-\bm{b}\right).\]

El test de razón de verosimilitudes viene dado por \eqref{eq:conocido}, donde

\[\bm{r}=\overline{\bm{x}}-\hat{\bm{\mu}}=\bm\Sigma \bm{A}^\prime{} \left(\bm{A}\bm\Sigma \bm{A}^\prime{}\right)^{-1}\left(\bm{A}\overline{\bm{x}}-\bm{b}\right)\]
de modo que finalmente adopta la forma

\[-2\log \lambda(\bm{x})=n\left(\bm{A}\overline{\bm{x}}-\bm{b}\right)^\prime{}
\left(\bm{A}\bm\Sigma \bm{A}^\prime{}\right)^{-1}\left(\bm{A}\overline{\bm{x}}-\bm{b}\right).\]

Bajo la hipótesis nula, \(H_0\),
\(\bm{A}\bm{x_1},\ldots,\bm{A}\bm{x_n}\in N_q(\bm{b},\bm{A}\bm\Sigma \bm{A}^\prime{})\)
y son independientes, siendo \(q\) la dimensión de \(\bm{b}\), y por tanto
la distribución del estadístico de contraste es \(\chi_q^2\). En resumen,
rechazaremos la hipótesis nula \(H_0: \bm{A}\bm{\mu}=\bm{b}\) si

\[n\left(\bm{A}\overline{\bm{x}}-\bm{b}\right)^\prime{}
\left(\bm{A}\bm\Sigma \bm{A}^\prime{}\right)^{-1}\left(\bm{A}\overline{\bm{x}}-\bm{b}\right)>\chi_{q,\alpha}^2\]
siendo \(\chi_{q,\alpha}^2\) el cuantil \(1-\alpha\) de la distribución
\(\chi_q^2\).

Si la matriz de covarianzas es desconocida, los desarrollos que conducen
al estimador de máxima verosimilitud de \(\bm{\mu}\) bajo \(H_0\) son algo
más complejos. Por ello, omitimos los detalles. El resultado es el
siguiente:

\[\hat{\bm{\mu}}=\overline{\bm{x}}- \bm{S} \bm{A}^\prime{} \left(\bm{A}\bm{S}\bm{A}^\prime{}\right)^{-1}\left(\bm{A}\overline{\bm{x}}-\bm{b}\right).\]

El test de razón de verosimilitudes viene dado por \eqref{eq:desconocido}, donde

\[\bm{r}=\overline{\bm{x}}-\hat{\bm{\mu}}=\bm{S}\bm{A}^\prime{} \left(\bm{A}\bm{S}\bm{A}^\prime{}\right)^{-1}\left(\bm{A}\overline{\bm{x}}-\bm{b}\right)\]
de modo que finalmente tomamos como estadístico de contraste

\[(n-1)\bm{r}^\prime{}\bm{S}^{-1}\bm{r}=(n-1)
\left(\bm{A}\overline{\bm{x}}-\bm{b}\right)^\prime{}
\left(\bm{A}\bm{S} \bm{A}^\prime{}\right)^{-1}\left(\bm{A}\overline{\bm{x}}-\bm{b}\right)\]
cuya distribución es \(\Gamma^2(q,n-1)\). Empleando que la distribución
\(\Gamma^2\) de Hotelling se puede expresar como una \(F\) de Snédecor,
salvo constantes, se llega a

\[\frac{n-q}{q}\left(\bm{A}\overline{\bm{x}}-\bm{b}\right)^\prime{}
\left(\bm{A}\bm{S}\bm{A}^\prime{}\right)^{-1}\left(\bm{A}\overline{\bm{x}}-\bm{b}\right)\in F_{q,n-q}.\]

Por lo tanto, si \(\bm\Sigma\) es desconocida, rechazaremos la hipótesis
nula \(H_0: \bm{A}\bm{\mu}=\bm{b}\) si
\[\frac{n-q}{q}\left(\bm{A}\overline{\bm{x}}-\bm{b}\right)^\prime{}
\left(\bm{A}\bm{S}\bm{A}^\prime{}\right)^{-1}\left(\bm{A}\overline{\bm{x}}-\bm{b}\right)> f_{q,n-q, \alpha}\]
siendo \(f_{q,n-q,\alpha}\) el cuantil \(1-\alpha\) de la distribución
\(F_{q,n-q}\).

\subsubsection{Caso particular. Contraste de igualdad de las componentes del vector de medias}\label{caso-particular-contraste-de-igualdad-de-las-componentes-del-vector-de-medias}

El contraste de la hipótesis nula de que las \(d\) componentes del vector
de medias, \(\bm{\mu}=(\mu_1,\ldots,\mu_d)^\prime{}\), son iguales, se
puede ver como un caso particular del contraste de restricciones
lineales. Para ello, basta considerar la siguiente matriz

\[\bm{A}=\left(\begin{array}{ccccc}
1      &     -1 &      0 & \cdots &      0 \\
1      &      0 &     -1 & \ddots & \vdots \\
\vdots & \vdots & \ddots & \ddots &      0 \\
1      &      0 & \cdots &      0 &     -1
\end{array}\right)\]
de modo que \(H_0:\bm{A}\bm{\mu}=\bm{0}\) equivale a
la igualdad de las \(d\) medias. Nótese que hay otras matrices que también
servirían para efectuar este contraste. En concreto, la matriz \(\bm{A}\)
que acabamos de proponer, efectúa las diferencias entre la media de la
primera componente y cada una de las demás medias. En este sentido,
además de servir para el contraste, permite estimar la discrepancia
entre las medias por comparación con la primera de ellas. Si se emplea
otro tipo de matriz, se obtendrían las posibles discrepancias entre las
medias en una presentación diferente.

\textbf{\includegraphics[width=1em,height=1em]{bookdown_AM_files/figure-latex/fa-icon-00ef1bb1d064b635425d49f06b331298.pdf} ~\textcolor{excolor}{Ejemplo.}} En \citet{mkb79}, página 12, se pueden
encontrar los datos de depósitos de corcho obtenidos en 28 árboles y
extraídos en las cuatro direcciones: Norte, Sur, Este y Oeste. Se está
estudiando si la cantidad media de corcho que se llega a recoger, es
similar en las cuatro direcciones. Vamos a efectuar el contraste de esta
hipótesis usando el test propuesto en esta sección.

\footnotesize

\begin{Shaded}
\begin{Highlighting}[]
\SpecialCharTok{\textgreater{}}\NormalTok{ corcho }\OtherTok{\textless{}{-}} \FunctionTok{read.table}\NormalTok{(}\StringTok{"data/corcho.txt"}\NormalTok{, }\AttributeTok{header =} \ConstantTok{TRUE}\NormalTok{)}
\SpecialCharTok{\textgreater{}}\NormalTok{ n }\OtherTok{\textless{}{-}} \FunctionTok{dim}\NormalTok{(corcho)[}\DecValTok{1}\NormalTok{]}
\SpecialCharTok{\textgreater{}}\NormalTok{ d }\OtherTok{\textless{}{-}} \FunctionTok{dim}\NormalTok{(corcho)[}\DecValTok{2}\NormalTok{]}
\SpecialCharTok{\textgreater{}}\NormalTok{ xbar }\OtherTok{\textless{}{-}} \FunctionTok{colMeans}\NormalTok{(corcho)}
\SpecialCharTok{\textgreater{}}\NormalTok{ xbar}
\NormalTok{       N        E        S        W }
\FloatTok{50.53571} \FloatTok{46.17857} \FloatTok{49.67857} \FloatTok{45.17857} 
\SpecialCharTok{\textgreater{}}\NormalTok{ Sc }\OtherTok{\textless{}{-}} \FunctionTok{cov}\NormalTok{(corcho)}
\SpecialCharTok{\textgreater{}}\NormalTok{ S }\OtherTok{\textless{}{-}}\NormalTok{ (n }\SpecialCharTok{{-}} \DecValTok{1}\NormalTok{)}\SpecialCharTok{/}\NormalTok{n }\SpecialCharTok{*}\NormalTok{ Sc}
\SpecialCharTok{\textgreater{}}\NormalTok{ S}
\NormalTok{         N        E        S        W}
\NormalTok{N }\FloatTok{280.0344} \FloatTok{215.7615} \FloatTok{278.1365} \FloatTok{218.1901}
\NormalTok{E }\FloatTok{215.7615} \FloatTok{212.0753} \FloatTok{220.8788} \FloatTok{165.2538}
\NormalTok{S }\FloatTok{278.1365} \FloatTok{220.8788} \FloatTok{337.5038} \FloatTok{250.2717}
\NormalTok{W }\FloatTok{218.1901} \FloatTok{165.2538} \FloatTok{250.2717} \FloatTok{217.9324}
\end{Highlighting}
\end{Shaded}

\normalsize

Se consideró una matriz \(\bm{A}\) que calcula las diferencias respecto de
la dirección norte. A continuación se muestran los valores de esta
matriz.

\footnotesize

\begin{Shaded}
\begin{Highlighting}[]
\SpecialCharTok{\textgreater{}}\NormalTok{ A }\OtherTok{\textless{}{-}} \FunctionTok{rbind}\NormalTok{(}\FunctionTok{c}\NormalTok{(}\DecValTok{1}\NormalTok{, }\SpecialCharTok{{-}}\DecValTok{1}\NormalTok{, }\DecValTok{0}\NormalTok{, }\DecValTok{0}\NormalTok{), }\FunctionTok{c}\NormalTok{(}\DecValTok{1}\NormalTok{, }\DecValTok{0}\NormalTok{, }\SpecialCharTok{{-}}\DecValTok{1}\NormalTok{, }\DecValTok{0}\NormalTok{), }\FunctionTok{c}\NormalTok{(}\DecValTok{1}\NormalTok{, }\DecValTok{0}\NormalTok{, }\DecValTok{0}\NormalTok{, }\SpecialCharTok{{-}}\DecValTok{1}\NormalTok{))}
\SpecialCharTok{\textgreater{}}\NormalTok{ A}
\NormalTok{     [,}\DecValTok{1}\NormalTok{] [,}\DecValTok{2}\NormalTok{] [,}\DecValTok{3}\NormalTok{] [,}\DecValTok{4}\NormalTok{]}
\NormalTok{[}\DecValTok{1}\NormalTok{,]    }\DecValTok{1}   \SpecialCharTok{{-}}\DecValTok{1}    \DecValTok{0}    \DecValTok{0}
\NormalTok{[}\DecValTok{2}\NormalTok{,]    }\DecValTok{1}    \DecValTok{0}   \SpecialCharTok{{-}}\DecValTok{1}    \DecValTok{0}
\NormalTok{[}\DecValTok{3}\NormalTok{,]    }\DecValTok{1}    \DecValTok{0}    \DecValTok{0}   \SpecialCharTok{{-}}\DecValTok{1}
\end{Highlighting}
\end{Shaded}

\normalsize

El estadístico \(F\) y su nivel crítico se muestran a continuación. Como
el nivel crítico es pequeño se rechaza la igualdad de las cuatro
direcciones en cuanto al depósito medio de corcho.

\footnotesize

\begin{Shaded}
\begin{Highlighting}[]
\SpecialCharTok{\textgreater{}}\NormalTok{ q }\OtherTok{\textless{}{-}} \FunctionTok{nrow}\NormalTok{(A)}
\SpecialCharTok{\textgreater{}}\NormalTok{ difmed }\OtherTok{\textless{}{-}}\NormalTok{ A }\SpecialCharTok{\%*\%}\NormalTok{ xbar}
\SpecialCharTok{\textgreater{}}\NormalTok{ covdif }\OtherTok{\textless{}{-}}\NormalTok{ A }\SpecialCharTok{\%*\%}\NormalTok{ S }\SpecialCharTok{\%*\%} \FunctionTok{t}\NormalTok{(A)}
\SpecialCharTok{\textgreater{}}\NormalTok{ estadF }\OtherTok{\textless{}{-}}\NormalTok{ ((n }\SpecialCharTok{{-}}\NormalTok{ q)}\SpecialCharTok{/}\NormalTok{q) }\SpecialCharTok{*} \FunctionTok{t}\NormalTok{(difmed) }\SpecialCharTok{\%*\%} \FunctionTok{solve}\NormalTok{(covdif) }\SpecialCharTok{\%*\%}\NormalTok{ difmed}
\SpecialCharTok{\textgreater{}}\NormalTok{ estadF}
\NormalTok{         [,}\DecValTok{1}\NormalTok{]}
\NormalTok{[}\DecValTok{1}\NormalTok{,] }\FloatTok{6.401857}
\SpecialCharTok{\textgreater{}}\NormalTok{ alpha }\OtherTok{\textless{}{-}} \FloatTok{0.05}
\SpecialCharTok{\textgreater{}} \FunctionTok{qf}\NormalTok{(}\DecValTok{1} \SpecialCharTok{{-}}\NormalTok{ alpha, q, n }\SpecialCharTok{{-}}\NormalTok{ q)}
\NormalTok{[}\DecValTok{1}\NormalTok{] }\FloatTok{2.991241}
\SpecialCharTok{\textgreater{}}\NormalTok{ pvalor }\OtherTok{\textless{}{-}} \DecValTok{1} \SpecialCharTok{{-}} \FunctionTok{pf}\NormalTok{(estadF, q, n }\SpecialCharTok{{-}}\NormalTok{ q)}
\SpecialCharTok{\textgreater{}}\NormalTok{ pvalor}
\NormalTok{            [,}\DecValTok{1}\NormalTok{]}
\NormalTok{[}\DecValTok{1}\NormalTok{,] }\FloatTok{0.002280399}
\end{Highlighting}
\end{Shaded}

\normalsize

Como ejercicio adicional, se podría modificar la matriz \(\bm{A}\), por
ejemplo de modo que se calculen las diferencias respecto de otra
dirección, y efectuar todo el procedimiento del mismo modo con la nueva
matriz \(\bm{A}\). El vector \texttt{difmed} y la matriz \texttt{covdif} tomarán valores
diferentes. Sin embargo, el estadístico \(F\) (\texttt{estadF}) y su nivel
crítico (\texttt{pvalor}) tendrán los mismos valores que hemos mostrado, pues,
como hemos dicho, el test no depende de la matriz \(\bm{A}\) escogida,
siempre que se mantenga la naturaleza de la hipótesis que se va a
contrastar. En este caso, la hipótesis es la igualdad de las cuatro
medias, independientemente de cómo se exprese.

\section{Inferencia sobre la matriz de covarianzas}\label{inferencia-sobre-la-matriz-de-covarianzas}

Consideramos ahora contrastes de hipótesis que involucran la estructura
de covarianzas de una población normal. En general, este tipo de
contrastes se llevan a cabo para comprobar hipótesis de otros modelos.
Los procedimientos que estudiaremos nos permitirán contrastar por
ejemplo si la matriz de covarianzas tiene una estructura determinada o
si algunos elementos de la matriz de covarianzas son cero (implicando
así la independencia de los subvectores correspondientes).

\subsection{\texorpdfstring{Contraste \(H_0: \bm\Sigma=\bm{\Sigma_0}\)}{Contraste H\_0: \textbackslash bm\textbackslash Sigma=\textbackslash bm\{\textbackslash Sigma\_0\}}}\label{contraste-h-0-bm-sigma-bm-sigma-0}

Disponemos de una muestra aleatoria simple
\(\bm{x_1},\ldots,\bm{x_n}\in N_d(\bm{\mu},\bm\Sigma)\) de vectores
aleatorios independientes y con la misma distribución normal
multivariante. Suponemos que el vector de medias \(\bm{\mu}\) es
desconocido. Deseamos llevar a cabo tareas de inferencia relativas a la
matriz de covarianzas. En particular, deseamos contrastar

\[H_0: \bm\Sigma=\bm{\Sigma_0} \textnormal{ frente a } H_1: \bm\Sigma\neq\bm{\Sigma_0}.\]

En este contraste la matriz de covarianzas en la hipótesis alternativa
no está sujeta a restricciones. Aplicando el procedimiento de razón de
verosimilitudes, resulta el estadístico de contraste:

\[-2\log \lambda(\bm{x})
=-2\log \frac{\sup_{\bm{\mu}} L\left(\bm{x}, \bm{\mu},\bm{\Sigma_0}\right)}
{\sup_{\bm{\mu},\bm\Sigma} L\left(\bm{x}, \bm{\mu},\bm\Sigma\right)}.\]

Recordamos que la función de verosimilitud alcanza su máximo en
\(\bm{\mu}=\overline{\bm{x}}\) para cualquier matriz de covarianzas
\(\bm\Sigma\). Se tiene por lo tanto,

\[\sup_{\bm{\mu}} \log L\left(\bm{x}, \bm{\mu},\bm{\Sigma_0}\right)
= c - \frac{n}{2} \left( \log|\bm{\Sigma_0}|
+ \mbox{traza}\left(\bm{\Sigma_0}^{-1}\bm{S}\right) \right).\]

También
habíamos visto que,

\[\sup_{\bm\Sigma}\sup_{\bm{\mu}}\log L(\bm{x}, \bm{\mu},\bm\Sigma)
= c - \frac{n}{2} \left( \log |\bm{S}| + d\right)\]
de modo que el
estadístico de contraste adopta la forma:

\[\begin{aligned}
-2\log \lambda(\bm{x}) &= n\left(\log \left|\bm{\Sigma_0}\right|
  + \mbox{traza}\left(\bm{\Sigma_0}^{-1}\bm{S}\right) - \log |\bm{S}| - d \right)\\
  &= n\left( \mbox{traza}\left(\bm{\Sigma_0}^{-1}\bm{S}\right)
    - \log\left|\bm{\Sigma_0}^{-1}\bm{S}\right| - d \right) \\
&= n\left(\sum_{j=1}^d \lambda_j - \log\left(\prod_{j=1}^d \lambda_j\right) - d\right)\\
&= n\left(da-\log\left(g^d\right)-d\right) \\
&= nd \left(a-\log g -1\right)\end{aligned}\]
siendo \(\lambda_1,\ldots,\lambda_d\) los autovalores de la matriz
\(\bm{\Sigma_0}^{-1}\bm{S}\), \(a\) la media aritmética de tales autovalores
y \(g\) su media geométrica. La distribución exacta de este estadístico
bajo la hipótesis nula no se encuentra disponible. En su lugar, usaremos
la distribución asintótica que presenta por ser un estadístico de razón
de verosimilitudes:

\[-2\log \lambda(\bm{x}) = nd \left(a-\log g -1\right)\sim \chi_m^2\]
siendo el número de grados de libertad, la diferencia entre el número de
parámetros independientes bajo la hipótesis alternativa y bajo la
hipótesis nula, que en este caso resulta, \(m=\frac{1}{2}d(d+1)\), pues es
el número de parámetros independientes en una matriz de covarianzas.

Por haberse construido como cociente de verosimilitudes bajo la
hipótesis nula y bajo la alternativa, rechazaremos la hipótesis nula
cuando este estadístico sea grande o, mejor dicho, cuando supere el
cuantil \((1-\alpha)\) de la distribución \(\chi_m^2\), denotado por
\(\chi_{m,\alpha}^2\), siendo \(\alpha\) el nivel de significación fijado de
antemano.

\textbf{\includegraphics[width=1em,height=1em]{bookdown_AM_files/figure-latex/fa-icon-00ef1bb1d064b635425d49f06b331298.pdf} ~\textcolor{excolor}{Ejemplo.}} Recordemos el ejemplo sobre Psicología educativa en el que
se analizaba si las notas medias obtenidas por estudiantes
universitarios en dos tests de inteligencia fueron iguales a 100. Cuando
planteamos por primera vez el contraste sobre el vector de medias,
suponíamos que la matriz de covarianzas era conocida e igual a

\[\bm{\Sigma}=\left(\begin{array}{cc}
49 & 35 \\
35 & 52
\end{array}\right).\]

Utilizando los datos del fichero \texttt{notas.txt} vamos a contrastar si esa
es efectivamente la estructura de la matriz de covarianzas del vector
\(\bm{x}=(x_1,x_2)^\prime{}\) correspondiente a las notas de los dos
tests. En primer lugar leemos los datos, calculamos la matriz de
covarianzas muestral \(S\) y definimos la matriz de covarianzas bajo la
hipótesis nula \(\bm{\Sigma_0}\).

\footnotesize

\begin{Shaded}
\begin{Highlighting}[]
\SpecialCharTok{\textgreater{}}\NormalTok{ tests }\OtherTok{\textless{}{-}} \FunctionTok{read.table}\NormalTok{(}\StringTok{"data/notas.txt"}\NormalTok{, }\AttributeTok{header =} \ConstantTok{TRUE}\NormalTok{)}
\SpecialCharTok{\textgreater{}}\NormalTok{ n }\OtherTok{\textless{}{-}} \FunctionTok{dim}\NormalTok{(tests)[}\DecValTok{1}\NormalTok{]}
\SpecialCharTok{\textgreater{}}\NormalTok{ d }\OtherTok{\textless{}{-}} \FunctionTok{dim}\NormalTok{(tests)[}\DecValTok{2}\NormalTok{]}
\SpecialCharTok{\textgreater{}}\NormalTok{ S }\OtherTok{\textless{}{-}}\NormalTok{ (n }\SpecialCharTok{{-}} \DecValTok{1}\NormalTok{)}\SpecialCharTok{/}\NormalTok{n }\SpecialCharTok{*} \FunctionTok{cov}\NormalTok{(tests)}
\SpecialCharTok{\textgreater{}}\NormalTok{ Sigma\_0 }\OtherTok{\textless{}{-}} \FunctionTok{matrix}\NormalTok{(}\FunctionTok{c}\NormalTok{(}\DecValTok{49}\NormalTok{, }\DecValTok{35}\NormalTok{, }\DecValTok{35}\NormalTok{, }\DecValTok{52}\NormalTok{), }\AttributeTok{nrow =} \DecValTok{2}\NormalTok{)}
\end{Highlighting}
\end{Shaded}

\normalsize

Calculamos ahora los autovalores de la matriz
\(\bm{\Sigma_0}^{-1}\bm{S}\). El estadístico de contraste se define a
partir de la media aritmética \(a\) y la media geométrica \(g\) de los
autovalores.

\footnotesize

\begin{Shaded}
\begin{Highlighting}[]
\SpecialCharTok{\textgreater{}}\NormalTok{ lambda }\OtherTok{\textless{}{-}} \FunctionTok{eigen}\NormalTok{(}\FunctionTok{solve}\NormalTok{(Sigma\_0) }\SpecialCharTok{\%*\%}\NormalTok{ S)}\SpecialCharTok{$}\NormalTok{values}
\SpecialCharTok{\textgreater{}}\NormalTok{ a }\OtherTok{\textless{}{-}} \FunctionTok{mean}\NormalTok{(lambda)}
\SpecialCharTok{\textgreater{}}\NormalTok{ g }\OtherTok{\textless{}{-}} \FunctionTok{prod}\NormalTok{(lambda)}\SpecialCharTok{\^{}}\NormalTok{(}\DecValTok{1}\SpecialCharTok{/}\NormalTok{d)}
\SpecialCharTok{\textgreater{}}\NormalTok{ estad }\OtherTok{\textless{}{-}}\NormalTok{ n }\SpecialCharTok{*}\NormalTok{ d }\SpecialCharTok{*}\NormalTok{ (a }\SpecialCharTok{{-}} \FunctionTok{log}\NormalTok{(g) }\SpecialCharTok{{-}} \DecValTok{1}\NormalTok{)}
\SpecialCharTok{\textgreater{}}\NormalTok{ estad}
\NormalTok{[}\DecValTok{1}\NormalTok{] }\FloatTok{10.47213}
\end{Highlighting}
\end{Shaded}

\normalsize

Bajo \(H_0\), la distribución asintótica del estadístico es \(\chi_m^2\),
siendo \(m=d(d+1)/2\).

\footnotesize

\begin{Shaded}
\begin{Highlighting}[]
\SpecialCharTok{\textgreater{}}\NormalTok{ m }\OtherTok{\textless{}{-}} \FloatTok{0.5} \SpecialCharTok{*}\NormalTok{ d }\SpecialCharTok{*}\NormalTok{ (d }\SpecialCharTok{+} \DecValTok{1}\NormalTok{)}
\SpecialCharTok{\textgreater{}}\NormalTok{ pval }\OtherTok{\textless{}{-}} \DecValTok{1} \SpecialCharTok{{-}} \FunctionTok{pchisq}\NormalTok{(estad, m)}
\SpecialCharTok{\textgreater{}}\NormalTok{ pval}
\NormalTok{[}\DecValTok{1}\NormalTok{] }\FloatTok{0.01495119}
\end{Highlighting}
\end{Shaded}

\normalsize

Para un nivel de significación inferior al del nivel crítico (\(0.0149\))
podremos aceptar que la matriz de covarianzas sea la que se planteó como
hipótesis.

\textbf{Nota.} Debemos observar que si se hubiera supuesto que el vector de
medias es conocido, siguiendo los mismos pasos habríamos llegado al
estadístico de contraste

\[-2\log \lambda(\bm{x}) = nd \left(a-\log g -1\right)\] siendo \(a\) y
\(g\) las medias aritmética y geométrica, respectivamente, de los
autovalores de la matriz \(\bm{\Sigma_0}^{-1}\hat{\bm\Sigma}_{\bm{\mu}}\).
La única diferencia radica en la sustitución de \(\bm{S}\) por el
estimador

\[\hat{\bm\Sigma}_{\bm{\mu}}=\frac{1}{n}\sum_{i=1}^n \left(\bm{x_i}-\bm{\mu}\right)\left(\bm{x_i}-\bm{\mu}\right)^\prime{}.\]

Nuevamente tenemos los mismos problemas con la distribución del
estadístico de contraste y apelamos a la distribución asintótica, que es
\(\chi_m^2\) con el mismo número de grados de libertad,
\(m=\frac{1}{2}d(d+1)\).

\section{Generalización del contraste sobre la matriz de covarianzas}\label{generalizacion-del-contraste-sobre-la-matriz-de-covarianzas}

En esta sección se va a generalizar el test obtenido para el
contraste de la matriz de covarianzas. En primer lugar ofrecemos el
enunciado de un teorema de generalización, cuya demostración es
innecesaria, pues consiste en la constatación de los mismos argumentos
de máxima verosimilitud ya empleados, y de los desarrollos
subsiguientes.

\begin{theorem}
\protect\hypertarget{thm:unnamed-chunk-29}{}\label{thm:unnamed-chunk-29}Sea \(\bm{x_1},\ldots,\bm{x_n}\) una muestra aleatoria simple de
\(N_d(\bm{\mu},\bm\Sigma)\). Si las hipótesis \(H_0\) y \(H_a\) conducen a los
estimadores de máxima verosimilitud \(\hat{\bm\Sigma}\) y \(\bm{S}\),
respectivamente, y si \(\overline{\bm{x}}\) es el estimador de máxima
verosimilitud para \(\bm{\mu}\) bajo cualquiera de las dos hipótesis,
entonces el estadístico de razón de verosimilitudes para contrastar
\(H_0\) frente a \(H_a\) viene dado por

\[-2\log\lambda(\bm{x})=nd \left(a-\log g -1\right)\] siendo \(a\) y \(g\)
las medias aritmética y geométrica, respectivamente, de los autovalores
de \(\hat{\bm\Sigma}^{-1}\bm{S}\).
\end{theorem}

Una vez establecido este resultado general, trataremos diversas
situaciones en las cuales se puede aplicar este teorema. En todos los
casos supondremos que disponemos de una muestra aleatoria simple
\(\bm{x_1},\ldots,\bm{x_n}\in N_d(\bm{\mu},\bm\Sigma)\) y que el vector de
medias \(\bm{\mu}\) es desconocido.

\subsection{\texorpdfstring{Contraste de la hipótesis nula \(H_0: \bm\Sigma=k\bm{\Sigma_0},\ \ k\in (0,+\infty)\)}{Contraste de la hipótesis nula H\_0: \textbackslash bm\textbackslash Sigma=k\textbackslash bm\{\textbackslash Sigma\_0\},\textbackslash{} \textbackslash{} k\textbackslash in (0,+\textbackslash infty)}}\label{contraste-de-la-hipotesis-nula-h-0-bm-sigma-k-bm-sigma-0-k-in-0-infty}

Supongamos que queremos contrastar una hipótesis nula compuesta sobre la
matriz de covarianzas

\[H_0: \bm\Sigma= k \bm{\Sigma_0}\qquad k\in (0,+\infty)\]
siendo
\(\bm{\Sigma_0}\) una matriz de covarianzas fijada, frente a una
alternativa en la que la matriz de covarianzas no está sujeta a
restricciones. El vector de medias carece de restricciones tanto bajo la
hipótesis nula como bajo la alternativa. Estamos en las condiciones del
teorema anterior, por lo que sólo nos falta calcular el estimador de la
matriz de covarianzas bajo la hipótesis nula,
\(\hat{\bm\Sigma}=\hat{k}\bm{\Sigma_0}\). Se puede comprobar que
\(\hat{k}=a_0\) es el estimador de máxima verosimilitud de \(k\), siendo
\(a_0\) la media aritmética de los autovalores de
\(\bm{\Sigma_0}^{-1}\bm{S}\).

Entonces, aplicando la expresión que figura en el teorema anterior, el
estadístico de razón de verosimilitudes adopta la forma:

\[-2\log\lambda(\bm{x})=nd \left(a-\log g -1\right)\] siendo \(a\) y \(g\)
las medias aritmética y geométrica, respectivamente, de los autovalores
de la matriz
\(\hat{\bm\Sigma}^{-1}\bm{S}=\frac{1}{a_0}\bm{\Sigma_0}^{-1}\bm{S}\). En
consecuencia, \(a=1\) y \(g=\frac{1}{a_0}g_0\), siendo \(g_0\) la media
geométrica de los autovalores de \(\bm{\Sigma_0}^{-1}\bm{S}\).
Sustituyendo los valores de \(a\) y \(g\) obtenemos

\[-2\log\lambda(\bm{x})=nd \left(1-\log\left(\frac{1}{a_0}g_0\right)-1\right)
=nd\log\frac{a_0}{g_0}.\]

Por último, no estando disponible la distribución exacta de este
estadístico, la aproximamos por una \(\chi_m^2\), siendo el número de
grados de libertad \(m=\frac{1}{2}d(d+1)-1=\frac{1}{2}(d-1)(d+2)\).

\subsubsection{\texorpdfstring{Test de esfericidad \(H_0: \bm\Sigma=\sigma^2\bm{I_d}\)}{Test de esfericidad H\_0: \textbackslash bm\textbackslash Sigma=\textbackslash sigma\^{}2\textbackslash bm\{I\_d\}}}\label{test-de-esfericidad-h-0-bm-sigma-sigma-2-bm-i-d}

Hay un caso particular de este tipo de contraste que tiene un interés
especial. Es el test de esfericidad, que consiste en contrastar la
hipótesis nula de que las variables son incorreladas y tienen la misma
varianza. La hipótesis nula es de la forma

\[H_0: \bm\Sigma=\sigma^2\bm{I_d},\] siendo \(\sigma^2\) la varianza común
desconocida. Nótese que la incorrelación equivale a independencia cuando
se trata de variables normales. Es inmediato que estamos ante un caso
particular del test anterior. Para verlo basta con tomar
\(\bm{\Sigma_0}=\bm{I_d}\). Por tanto, el estadístico de contraste sería:

\[-2\log\lambda(\bm{x})=nd\log\frac{a_0}{g_0}\sim \chi_m^2\] siendo
\(a_0\) y \(g_0\) las medias aritmética y geométrica, respectivamente, de
los autovalores de la matriz \(\bm{\Sigma_0}^{-1}\bm{S}=\bm{S}\), y
\(m=\frac{1}{2}(d-1)(d+2)\).

\textbf{\includegraphics[width=1em,height=1em]{bookdown_AM_files/figure-latex/fa-icon-00ef1bb1d064b635425d49f06b331298.pdf} ~\textcolor{excolor}{Ejemplo.}} Vamos a contrastar la esfericidad de datos generados de distintas
normales multivariantes. Para ello utilizaremos la función \texttt{rmvnorm} de
la librería \texttt{mvtnorm}. En primer lugar generamos una muestra de tamaño
\(n=100\) de una normal \(N_3(\bm{0},\bm{I_3})\).

\footnotesize

\begin{Shaded}
\begin{Highlighting}[]
\SpecialCharTok{\textgreater{}} \FunctionTok{set.seed}\NormalTok{(}\DecValTok{123456}\NormalTok{)}
\SpecialCharTok{\textgreater{}} \FunctionTok{library}\NormalTok{(mvtnorm)}
\SpecialCharTok{\textgreater{}}\NormalTok{ d }\OtherTok{\textless{}{-}} \DecValTok{3}
\SpecialCharTok{\textgreater{}}\NormalTok{ mu }\OtherTok{\textless{}{-}} \FunctionTok{c}\NormalTok{(}\DecValTok{0}\NormalTok{, }\DecValTok{0}\NormalTok{, }\DecValTok{0}\NormalTok{)}
\SpecialCharTok{\textgreater{}}\NormalTok{ n }\OtherTok{\textless{}{-}} \DecValTok{100}
\SpecialCharTok{\textgreater{}}\NormalTok{ Sigma }\OtherTok{\textless{}{-}} \FunctionTok{diag}\NormalTok{(d)}
\SpecialCharTok{\textgreater{}}\NormalTok{ x }\OtherTok{\textless{}{-}} \FunctionTok{rmvnorm}\NormalTok{(n, mu, Sigma)}
\end{Highlighting}
\end{Shaded}

\normalsize

Una vez generados los datos, vamos a realizar el contraste de
esfericidad. Para ello calculamos \(a_0\) y
\(g_0\) las medias aritmética y geométrica, respectivamente, de los
autovalores de la matriz \(\bm{S}\).

\footnotesize

\begin{Shaded}
\begin{Highlighting}[]
\SpecialCharTok{\textgreater{}}\NormalTok{ S }\OtherTok{\textless{}{-}} \FunctionTok{cov}\NormalTok{(x) }\SpecialCharTok{*}\NormalTok{ (n }\SpecialCharTok{{-}} \DecValTok{1}\NormalTok{)}\SpecialCharTok{/}\NormalTok{n}
\SpecialCharTok{\textgreater{}}\NormalTok{ lambda }\OtherTok{\textless{}{-}} \FunctionTok{eigen}\NormalTok{(S)}\SpecialCharTok{$}\NormalTok{values}
\SpecialCharTok{\textgreater{}}\NormalTok{ a0 }\OtherTok{\textless{}{-}} \FunctionTok{mean}\NormalTok{(lambda)}
\SpecialCharTok{\textgreater{}}\NormalTok{ g0 }\OtherTok{\textless{}{-}} \FunctionTok{prod}\NormalTok{(lambda)}\SpecialCharTok{\^{}}\NormalTok{(}\DecValTok{1}\SpecialCharTok{/}\NormalTok{d)}
\SpecialCharTok{\textgreater{}}\NormalTok{ estad }\OtherTok{\textless{}{-}}\NormalTok{ n }\SpecialCharTok{*}\NormalTok{ d }\SpecialCharTok{*} \FunctionTok{log}\NormalTok{(a0}\SpecialCharTok{/}\NormalTok{g0)}
\SpecialCharTok{\textgreater{}}\NormalTok{ estad}
\NormalTok{[}\DecValTok{1}\NormalTok{] }\FloatTok{4.483253}
\end{Highlighting}
\end{Shaded}

\normalsize

Bajo \(H_0\), la distribución asintótica del estadístico es \(\chi_m^2\),
siendo \(m=(d-1)(d+2)/2\). En este caso, para un nivel de significación
\(\alpha=0.05\), no rechazamos la hipótesis nula de esfericidad, ya que el
valor del estadístico es menor que \(\chi_{m,\alpha}^2\). De hecho el
valor crítico del estadístico es considerablemente grande.

\footnotesize

\begin{Shaded}
\begin{Highlighting}[]
\SpecialCharTok{\textgreater{}}\NormalTok{ m }\OtherTok{\textless{}{-}} \FloatTok{0.5} \SpecialCharTok{*}\NormalTok{ (d }\SpecialCharTok{{-}} \DecValTok{1}\NormalTok{) }\SpecialCharTok{*}\NormalTok{ (d }\SpecialCharTok{+} \DecValTok{2}\NormalTok{)}
\SpecialCharTok{\textgreater{}} \FunctionTok{qchisq}\NormalTok{(}\FloatTok{0.95}\NormalTok{, m)}
\NormalTok{[}\DecValTok{1}\NormalTok{] }\FloatTok{11.0705}
\SpecialCharTok{\textgreater{}}\NormalTok{ pval }\OtherTok{\textless{}{-}} \DecValTok{1} \SpecialCharTok{{-}} \FunctionTok{pchisq}\NormalTok{(estad, m)}
\SpecialCharTok{\textgreater{}}\NormalTok{ pval}
\NormalTok{[}\DecValTok{1}\NormalTok{] }\FloatTok{0.4821273}
\end{Highlighting}
\end{Shaded}

\normalsize

Trabajaremos ahora con datos generados de una normal multivariante
\(N_3(\bm{0},4\bm{I_3})\).

\footnotesize

\begin{Shaded}
\begin{Highlighting}[]
\SpecialCharTok{\textgreater{}}\NormalTok{ Sigma }\OtherTok{\textless{}{-}} \DecValTok{4} \SpecialCharTok{*} \FunctionTok{diag}\NormalTok{(d)}
\SpecialCharTok{\textgreater{}}\NormalTok{ x }\OtherTok{\textless{}{-}} \FunctionTok{rmvnorm}\NormalTok{(n, mu, Sigma)}
\end{Highlighting}
\end{Shaded}

\normalsize

Planteamos el contraste de esfericidad y, de nuevo, no rechazamos la
hipótesis nula.

\footnotesize

\begin{Shaded}
\begin{Highlighting}[]
\SpecialCharTok{\textgreater{}}\NormalTok{ S }\OtherTok{\textless{}{-}} \FunctionTok{cov}\NormalTok{(x) }\SpecialCharTok{*}\NormalTok{ (n }\SpecialCharTok{{-}} \DecValTok{1}\NormalTok{)}\SpecialCharTok{/}\NormalTok{n}
\SpecialCharTok{\textgreater{}}\NormalTok{ lambda }\OtherTok{\textless{}{-}} \FunctionTok{eigen}\NormalTok{(S)}\SpecialCharTok{$}\NormalTok{values}
\SpecialCharTok{\textgreater{}}\NormalTok{ a0 }\OtherTok{\textless{}{-}} \FunctionTok{mean}\NormalTok{(lambda)}
\SpecialCharTok{\textgreater{}}\NormalTok{ g0 }\OtherTok{\textless{}{-}} \FunctionTok{prod}\NormalTok{(lambda)}\SpecialCharTok{\^{}}\NormalTok{(}\DecValTok{1}\SpecialCharTok{/}\NormalTok{d)}
\SpecialCharTok{\textgreater{}}\NormalTok{ estad }\OtherTok{\textless{}{-}}\NormalTok{ n }\SpecialCharTok{*}\NormalTok{ d }\SpecialCharTok{*} \FunctionTok{log}\NormalTok{(a0}\SpecialCharTok{/}\NormalTok{g0)}
\SpecialCharTok{\textgreater{}}\NormalTok{ estad}
\NormalTok{[}\DecValTok{1}\NormalTok{] }\FloatTok{3.251188}
\SpecialCharTok{\textgreater{}} \FunctionTok{qchisq}\NormalTok{(}\FloatTok{0.95}\NormalTok{, m)}
\NormalTok{[}\DecValTok{1}\NormalTok{] }\FloatTok{11.0705}
\SpecialCharTok{\textgreater{}}\NormalTok{ pval }\OtherTok{\textless{}{-}} \DecValTok{1} \SpecialCharTok{{-}} \FunctionTok{pchisq}\NormalTok{(estad, m)}
\SpecialCharTok{\textgreater{}}\NormalTok{ pval}
\NormalTok{[}\DecValTok{1}\NormalTok{] }\FloatTok{0.6613227}
\end{Highlighting}
\end{Shaded}

\normalsize

Como último ejemplo, vamos a generar datos de un vector normal que no
verifique la hipótesis nula de esfericidad. Para ello elegiremos una
matriz de covarianzas diagonal pero en la que no todos las variables
tengan la misma varianza.

\footnotesize

\begin{Shaded}
\begin{Highlighting}[]
\SpecialCharTok{\textgreater{}}\NormalTok{ Sigma }\OtherTok{\textless{}{-}} \FunctionTok{diag}\NormalTok{(}\FunctionTok{c}\NormalTok{(}\DecValTok{3}\NormalTok{, }\DecValTok{5}\NormalTok{, }\DecValTok{1}\NormalTok{))}
\SpecialCharTok{\textgreater{}}\NormalTok{ x }\OtherTok{\textless{}{-}} \FunctionTok{rmvnorm}\NormalTok{(n, mu, Sigma)}
\end{Highlighting}
\end{Shaded}

\normalsize

Calculamos el estadístico de contraste y obtenemos los siguientes
resultados.

\footnotesize

\begin{Shaded}
\begin{Highlighting}[]
\SpecialCharTok{\textgreater{}}\NormalTok{ S }\OtherTok{\textless{}{-}} \FunctionTok{cov}\NormalTok{(x) }\SpecialCharTok{*}\NormalTok{ (n }\SpecialCharTok{{-}} \DecValTok{1}\NormalTok{)}\SpecialCharTok{/}\NormalTok{n}
\SpecialCharTok{\textgreater{}}\NormalTok{ lambda }\OtherTok{\textless{}{-}} \FunctionTok{eigen}\NormalTok{(S)}\SpecialCharTok{$}\NormalTok{values}
\SpecialCharTok{\textgreater{}}\NormalTok{ a0 }\OtherTok{\textless{}{-}} \FunctionTok{mean}\NormalTok{(lambda)}
\SpecialCharTok{\textgreater{}}\NormalTok{ g0 }\OtherTok{\textless{}{-}} \FunctionTok{prod}\NormalTok{(lambda)}\SpecialCharTok{\^{}}\NormalTok{(}\DecValTok{1}\SpecialCharTok{/}\NormalTok{d)}
\SpecialCharTok{\textgreater{}}\NormalTok{ estad }\OtherTok{\textless{}{-}}\NormalTok{ n }\SpecialCharTok{*}\NormalTok{ d }\SpecialCharTok{*} \FunctionTok{log}\NormalTok{(a0}\SpecialCharTok{/}\NormalTok{g0)}
\SpecialCharTok{\textgreater{}}\NormalTok{ estad}
\NormalTok{[}\DecValTok{1}\NormalTok{] }\FloatTok{69.92047}
\SpecialCharTok{\textgreater{}} \FunctionTok{qchisq}\NormalTok{(}\FloatTok{0.95}\NormalTok{, m)}
\NormalTok{[}\DecValTok{1}\NormalTok{] }\FloatTok{11.0705}
\SpecialCharTok{\textgreater{}}\NormalTok{ pval }\OtherTok{\textless{}{-}} \DecValTok{1} \SpecialCharTok{{-}} \FunctionTok{pchisq}\NormalTok{(estad, m)}
\SpecialCharTok{\textgreater{}}\NormalTok{ pval}
\NormalTok{[}\DecValTok{1}\NormalTok{] }\FloatTok{1.064704e{-}13}
\end{Highlighting}
\end{Shaded}

\normalsize

La esfericidad permite una mayor o menor varianza, pero esta varianza
tendría que ser la misma en las tres variables y no debe existir
correlación. Esto se pone en evidencia mediante los autovalores de la
matriz de covarianzas, que deben ser iguales. Como se puede comprobar,
los autovalores para este último ejemplo muestran una gran desigualdad.
Esto hace que las medias aritmética y geométrica sean muy diferentes. El
nivel crítico es muy pequeño, por lo que hay
pruebas muy significativas de ausencia de esfericidad.

\subsection{\texorpdfstring{Contraste de independencia de dos subvectores \(H_0: \bm{\Sigma_{12}}=\bm{0}\)}{Contraste de independencia de dos subvectores H\_0: \textbackslash bm\{\textbackslash Sigma\_\{12\}\}=\textbackslash bm\{0\}}}\label{contraste-de-independencia-de-dos-subvectores-h-0-bm-sigma-12-bm-0}

Supongamos que el vector de interés \(\bm{x}\) se divide en dos
subvectores con \(d_1\) y \(d_2\) variables (\(d_1+d_2=d\)). Por ejemplo, si
\(\bm{x_1}\) denota el subvector de las primeras \(d_1\) componentes de
\(\bm{x}\) y \(\bm{x_2}\) denota el subvector con las restantes \(d_2\)
componentes, entonces podemos escribir:

\[\bm{x}=\left(
\begin{array}{c}
\bm{x_1}\\ 
\bm{x_2}
\end{array}
\right),\ \ \ \bm{\mu}=\left(\begin{array}{c}\bm{\mu_1}\\ \bm{\mu_2}\end{array}\right),\ \ \ \bm{\Sigma}=\left(\begin{array}{cc}\bm{\Sigma_{11}}&\bm{\Sigma_{12}}\\ \bm{\Sigma_{21}}&\bm{\Sigma_{22}}\end{array}\right).\]

Queremos contrastar la hipótesis de que los dos conjuntos de variables
son independientes entre sí, lo cual, en esta situación donde se supone
normalidad, equivale a incorrelación, esto es, a que
\(\bm{\Sigma_{12}}=\bm{0}\). Bajo la hipótesis nula
\(H_0: \bm{\Sigma_{12}}=\bm{0}\), la verosimilitud se descompone en el
producto de dos factores correspondientes a las verosimilitudes que
provienen de cada conjunto de variables. De este modo, bajo la hipótesis
nula, los estimadores de \(\bm{\mu_1}\) y \(\bm{\Sigma_{11}}\) por un lado,
y de \(\bm{\mu_2}\) y \(\bm{\Sigma_{22}}\) por otro, se obtienen maximizando
la verosimilitud de cada conjunto de variables por separado. Así,
suponiendo que el vector de medias
\(\bm{\mu}=(\bm{\mu_1},\bm{\mu_2})^\prime{}\) es desconocido, su estimador
de máxima verosimilitud será
\(\hat{\bm{\mu}}=(\bm{\overline{x}_1},\bm{\overline{x}_2})^\prime{}=\overline{\bm{x}}\),
mientras que

\[\hat{\bm\Sigma}=\left(\begin{array}{cc} \bm{S_{11}} & \bm{0} \\
\bm{0} & \bm{S_{22}} \end{array}\right).\]

Bajo la alternativa, no hay restricciones ni para el vector de medias ni
para la matriz de covarianzas, de modo que los estimadores de máxima
verosimilitud serán \(\overline{\bm{x}}\) y \(\bm{S}\), respectivamente.

Entonces el estadístico de razón de verosimilitudes adopta la forma que
figura en la expresión del teorema anterior, donde

\[\hat{\bm\Sigma}^{-1}\bm{S}=\left(\begin{array}{cc} \bm{S_{11}} & \bm{0} \\
\bm{0} & \bm{S_{22}} \end{array}\right)^{-1} \left(\begin{array}{cc} \bm{S_{11}} & \bm{S_{12}} \\
\bm{S_{21}} & \bm{S_{22}} \end{array}\right)
=\left(\begin{array}{cc} \bm{I_{d_1}} & \bm{S_{11}}^{-1}\bm{S_{12}} \\
\bm{S_{22}}^{-1}\bm{S_{21}} & \bm{I_{d_2}} \end{array}\right).\]

En tal caso,
\(a=\frac{1}{d}\mbox{traza}\left(\hat{\bm\Sigma}^{-1}\bm{S}\right)=\frac{1}{d}d=1\).
Por otro lado,

\[
\begin{aligned}
g^d&=\left|\hat{\bm\Sigma}^{-1}\bm{S}\right|\\
&=\frac{|\bm{S}|}{|\bm{S_{11}}|\cdot |\bm{S_{22}}|}\\
&=\frac{|\bm{S_{22}}-\bm{S_{21}}\bm{S_{11}}^{-1}\bm{S_{12}}|}{|\bm{S_{22}}|}\\
&=\left|\bm{I_{d_2}}-\bm{S_{22}}^{-1}\bm{S_{21}}\bm{S_{11}}^{-1}\bm{S_{12}}\right|\\
&=\left|\bm{I_{d_2}}-\bm{R_{22}}^{-1}\bm{R_{21}}\bm{R_{11}}^{-1}\bm{R_{12}}\right|.
\end{aligned}\]

En el último paso hemos sustituido las matrices de covarianzas por
matrices de correlaciones, que se construyen a partir de las anteriores
así: \(\bm{R}=\bm{D^{-1/2}}\bm{S}\bm{D^{-1/2}}\), siendo \(\bm{D}\) una
matriz diagonal que contiene las varianzas. Entonces, el estadístico de
contraste será:

\[-2\log \lambda(\bm{x})=-nd\log g=-n\log\left|\bm{I_{d_2}}-\bm{S_{22}}^{-1}\bm{S_{21}}\bm{S_{11}}^{-1}\bm{S_{12}}\right|.\]

Ahora podríamos apelar a la distribución asintótica del estadístico
\(-2\log \lambda(\bm{x})\), como hemos hecho en el resto de contrastes.
Sin embargo, se puede obtener la distribución exacta
\(\lambda(\bm{x})^{2/n}\) bajo \(H_0\). Se tiene que, si
\(H_0: \bm{\Sigma_{12}}=\bm{0}\) es cierta, \(\lambda(\bm{x})^{2/n}\) sigue
una distribución \(\Lambda\) de Wilks\footnote{La distribución \(\Lambda\) de Wilks sirve para extender al caso
  multivariante la distribución \(F\) de Snédecor. Así, si la \(F\) de
  Snédecor se obtiene mediante el cociente de dos variables de tipo
  ji-cuadrado, la \(\Lambda\) de Wilks surge como el cociente de los
  determinantes de dos matrices de covarianzas.} de parámetros \(d_2\), \(d_1\),
\(n-1-d_1\). Es decir,

\[\lambda(\bm{x})^{2/n}=\frac{|\bm{S}|}{|\bm{S_{11}}|\cdot |\bm{S_{22}}|}=\left|\bm{I_{d_2}}-\bm{S_{22}}^{-1}\bm{S_{21}}\bm{S_{11}}^{-1}\bm{S_{12}}\right|
\in \Lambda\left(d_2,d_1,n-1-d_1\right).\]

Como corresponde a un test de razón de verosimilitudes, se rechazará la
hipótesis nula \(H_0: \bm{\Sigma_{12}}=\bm{0}\) cuando \(\lambda(\bm{x})\)
sea pequeño, o equivalentemente, cuando
\(|\bm{I_{d_2}}-\bm{S_{22}}^{-1}\bm{S_{21}}\bm{S_{11}}^{-1}\bm{S_{12}}|\)
sea pequeño. La distribución \(\Lambda\) de Wilks se puede transformar (ver Anexo \ref{anexoDist}) en
una \(F\) de Snédecor\footnote{En una \(\Lambda(d,m_H,m_E)\), si \(m_H=2\) (\(d\geq 2\)), \[\frac{1-\sqrt{\Lambda(d,2,m_E)}}{\sqrt{\Lambda(d,2,m_E)}}\stackrel{d}{=} \frac{d}{m_E-d+1} F_{2d,2(m_E-d+1)}.\]}. Puesto que las funciones para
transformar una \(\Lambda\) de Wilks en una \(F\) de Snédecor son
decrecientes, se rechazará la hipótesis nula
\(H_0: \bm{\Sigma_{12}}=\bm{0}\) cuando su transformación en una \(F\) de
Snédecor sea grande.

\textbf{\includegraphics[width=1em,height=1em]{bookdown_AM_files/figure-latex/fa-icon-00ef1bb1d064b635425d49f06b331298.pdf} ~\textcolor{excolor}{Ejemplo.}} En el primer capítulo habíamos presentado un estudio sobre diabetes en el
que se analizaban diferentes variables de interés sobre pacientes que no
presentan la enfermedad. Las variables de mayor relevancia para el
estudio eran:

\begin{itemize}
\item
  \(x_1\)= Intolerancia a la glucosa
\item
  \(x_2\)= Respuesta de la insulina a la glucosa oral
\item
  \(x_3\)= Resistencia de la insulina
\end{itemize}

Se recogían igualmente otras dos variables de menor relevancia:

\begin{itemize}
\item
  \(y_1\)= Peso relativo
\item
  \(y_2\)= Glucemia basal
\end{itemize}

Suponiendo que el vector \(\bm{v}=(y_1,y_2,x_1,x_2,x_3)^\prime{}\) sigue
una distribución normal multivariante, nos interesa contrastar la
independencia de los subvectores \(\bm{y}=(y_1,y_2)^\prime{}\) y
\(\bm{x}=(x_1,x_2,x_3)^\prime\). Es decir, planteamos la hipótesis nula
\(H_0: \bm{\Sigma_{yx}}=\bm{0}\). Como habíamos visto al comentar los
datos, podemos particionar la matriz de covarianzas muestral \(\bm{S}\) y
definir a partir de ella las submatrices \(\bm{S_{11}}\), \(\bm{S_{22}}\),
\(\bm{S_{12}}\) y \(\bm{S_{21}}\).

\footnotesize

\begin{Shaded}
\begin{Highlighting}[]
\SpecialCharTok{\textgreater{}}\NormalTok{ diab }\OtherTok{\textless{}{-}} \FunctionTok{read.table}\NormalTok{(}\StringTok{"data/diabetes.txt"}\NormalTok{, }\AttributeTok{header =} \ConstantTok{TRUE}\NormalTok{)}
\SpecialCharTok{\textgreater{}}\NormalTok{ n }\OtherTok{\textless{}{-}} \FunctionTok{dim}\NormalTok{(diab)[}\DecValTok{1}\NormalTok{]}
\SpecialCharTok{\textgreater{}}\NormalTok{ S }\OtherTok{\textless{}{-}} \FunctionTok{cov}\NormalTok{(diab) }\SpecialCharTok{*}\NormalTok{ (n }\SpecialCharTok{{-}} \DecValTok{1}\NormalTok{)}\SpecialCharTok{/}\NormalTok{n}
\SpecialCharTok{\textgreater{}}\NormalTok{ S}
\NormalTok{            y1          y2           x1           x2          x3}
\NormalTok{y1  }\FloatTok{0.01583006}   \FloatTok{0.2113327}    \FloatTok{0.7700567}   \SpecialCharTok{{-}}\FloatTok{0.2091966}    \FloatTok{2.141484}
\NormalTok{y2  }\FloatTok{0.21133270}  \FloatTok{69.0250473}   \FloatTok{25.6587902}  \SpecialCharTok{{-}}\FloatTok{23.4352552}  \SpecialCharTok{{-}}\FloatTok{20.388469}
\NormalTok{x1  }\FloatTok{0.77005671}  \FloatTok{25.6587902} \FloatTok{1082.3610586}  \FloatTok{388.1077505}  \FloatTok{106.027410}
\NormalTok{x2 }\SpecialCharTok{{-}}\FloatTok{0.20919660} \SpecialCharTok{{-}}\FloatTok{23.4352552}  \FloatTok{388.1077505} \FloatTok{2330.1025520} \FloatTok{1117.797732}
\NormalTok{x3  }\FloatTok{2.14148393} \SpecialCharTok{{-}}\FloatTok{20.3884688}  \FloatTok{106.0274102} \FloatTok{1117.7977316} \FloatTok{2089.952741}
\SpecialCharTok{\textgreater{}}\NormalTok{ S11 }\OtherTok{\textless{}{-}}\NormalTok{ S[}\DecValTok{1}\SpecialCharTok{:}\DecValTok{2}\NormalTok{, }\DecValTok{1}\SpecialCharTok{:}\DecValTok{2}\NormalTok{]}
\SpecialCharTok{\textgreater{}}\NormalTok{ S12 }\OtherTok{\textless{}{-}}\NormalTok{ S[}\DecValTok{1}\SpecialCharTok{:}\DecValTok{2}\NormalTok{, }\DecValTok{3}\SpecialCharTok{:}\DecValTok{5}\NormalTok{]}
\SpecialCharTok{\textgreater{}}\NormalTok{ S21 }\OtherTok{\textless{}{-}}\NormalTok{ S[}\DecValTok{3}\SpecialCharTok{:}\DecValTok{5}\NormalTok{, }\DecValTok{1}\SpecialCharTok{:}\DecValTok{2}\NormalTok{]}
\SpecialCharTok{\textgreater{}}\NormalTok{ S22 }\OtherTok{\textless{}{-}}\NormalTok{ S[}\DecValTok{3}\SpecialCharTok{:}\DecValTok{5}\NormalTok{, }\DecValTok{3}\SpecialCharTok{:}\DecValTok{5}\NormalTok{]}
\end{Highlighting}
\end{Shaded}

\normalsize

Calculamos el estadístico del contraste \(\lambda(\bm{v})^{2/n}\). Lo
podemos hacer de dos formas:

\footnotesize

\begin{Shaded}
\begin{Highlighting}[]
\SpecialCharTok{\textgreater{}}\NormalTok{ estad }\OtherTok{\textless{}{-}} \FunctionTok{det}\NormalTok{(S)}\SpecialCharTok{/}\NormalTok{(}\FunctionTok{det}\NormalTok{(S11) }\SpecialCharTok{*} \FunctionTok{det}\NormalTok{(S22))}
\SpecialCharTok{\textgreater{}}\NormalTok{ estad}
\NormalTok{[}\DecValTok{1}\NormalTok{] }\FloatTok{0.7240096}
\SpecialCharTok{\textgreater{}}\NormalTok{ estad }\OtherTok{\textless{}{-}} \FunctionTok{det}\NormalTok{(}\FunctionTok{diag}\NormalTok{(}\DecValTok{3}\NormalTok{) }\SpecialCharTok{{-}} \FunctionTok{solve}\NormalTok{(S22) }\SpecialCharTok{\%*\%}\NormalTok{ S21 }\SpecialCharTok{\%*\%} \FunctionTok{solve}\NormalTok{(S11) }\SpecialCharTok{\%*\%}\NormalTok{ S12)}
\SpecialCharTok{\textgreater{}}\NormalTok{ estad}
\NormalTok{[}\DecValTok{1}\NormalTok{] }\FloatTok{0.7240096}
\end{Highlighting}
\end{Shaded}

\normalsize

Recordamos que si \(H_0\) es cierta, \(\lambda(\bm{v})^{2/n}\) sigue una
distribución \(\Lambda\) de Wilks de parámetros \(d_2\), \(d_1\), \(n-1-d_1\).
En nuestro ejemplo, \(d_2=3\), \(d_1=2\) y \(n=46\), por lo tanto,
\(\lambda(\bm{v})^{2/n}\) sigue una distribución \(\Lambda(3,2,43)\). Se
verifica además (ver Anexo \ref{anexoDist}) que

\[\frac{41(1-\sqrt{\Lambda(3,2,43)})}{3\sqrt{\Lambda(3,2,43)}}\stackrel{d}{=} F_{6,82}\]

Teniendo en cuenta que rechazamos \(H_0\) cuando \(\lambda(\bm{x})\) es
pequeño y que la transformación anterior es una función decreciente,
rechararemos la hipótesis de independencia (para una significación
\(\alpha\)) si

\[\frac{41(1-\sqrt{\Lambda(3,2,43)})}{3\sqrt{\Lambda(3,2,43)}}>f_{6,82,\alpha}\]
siendo \(f_{6,82,\alpha}\) el cuantil \(1-\alpha\) de la distribución
\(F_{6,82}\).

\footnotesize

\begin{Shaded}
\begin{Highlighting}[]
\SpecialCharTok{\textgreater{}}\NormalTok{ d1 }\OtherTok{\textless{}{-}} \DecValTok{2}
\SpecialCharTok{\textgreater{}}\NormalTok{ d2 }\OtherTok{\textless{}{-}} \DecValTok{3}
\SpecialCharTok{\textgreater{}} \CommentTok{\# Parámetros de la distribución Lambda de Wilks (d,mH,mE)}
\ErrorTok{\textgreater{}}\NormalTok{ d }\OtherTok{\textless{}{-}}\NormalTok{ d2}
\SpecialCharTok{\textgreater{}}\NormalTok{ mH }\OtherTok{\textless{}{-}}\NormalTok{ d1}
\SpecialCharTok{\textgreater{}}\NormalTok{ mE }\OtherTok{\textless{}{-}}\NormalTok{ n }\SpecialCharTok{{-}} \DecValTok{1} \SpecialCharTok{{-}}\NormalTok{ d1}
\SpecialCharTok{\textgreater{}} \CommentTok{\# Transformación del estadístico a F de Snedecor}
\ErrorTok{\textgreater{}}\NormalTok{ estadF }\OtherTok{\textless{}{-}}\NormalTok{ (mE }\SpecialCharTok{{-}}\NormalTok{ d }\SpecialCharTok{+} \DecValTok{1}\NormalTok{)}\SpecialCharTok{/}\NormalTok{d }\SpecialCharTok{*}\NormalTok{ (}\DecValTok{1} \SpecialCharTok{{-}} \FunctionTok{sqrt}\NormalTok{(estad))}\SpecialCharTok{/}\FunctionTok{sqrt}\NormalTok{(estad)}
\SpecialCharTok{\textgreater{}}\NormalTok{ estadF}
\NormalTok{[}\DecValTok{1}\NormalTok{] }\FloatTok{2.394993}
\SpecialCharTok{\textgreater{}} \CommentTok{\# Grados de libertad F de Snedecor}
\ErrorTok{\textgreater{}}\NormalTok{ glF1 }\OtherTok{\textless{}{-}} \DecValTok{2} \SpecialCharTok{*}\NormalTok{ d}
\SpecialCharTok{\textgreater{}}\NormalTok{ glF2 }\OtherTok{\textless{}{-}} \DecValTok{2} \SpecialCharTok{*}\NormalTok{ (mE }\SpecialCharTok{{-}}\NormalTok{ d }\SpecialCharTok{+} \DecValTok{1}\NormalTok{)}
\SpecialCharTok{\textgreater{}} \FunctionTok{qf}\NormalTok{(}\FloatTok{0.95}\NormalTok{, glF1, glF2)}
\NormalTok{[}\DecValTok{1}\NormalTok{] }\FloatTok{2.211303}
\SpecialCharTok{\textgreater{}}\NormalTok{ pval }\OtherTok{\textless{}{-}} \DecValTok{1} \SpecialCharTok{{-}} \FunctionTok{pf}\NormalTok{(estadF, glF1, glF2)}
\SpecialCharTok{\textgreater{}}\NormalTok{ pval}
\NormalTok{[}\DecValTok{1}\NormalTok{] }\FloatTok{0.03499556}
\end{Highlighting}
\end{Shaded}

\normalsize

Es decir, rechazamos la hipótesis nula de independencia.

\subsubsection{Caso particular. Coeficiente de correlación múltiple}\label{caso-particular-coeficiente-de-correlacion-multiple}

Consideremos que uno de los conjuntos de variables tenga un único
elemento, por ejemplo, \(d_1=1\) y \(d_2=d-1\). En esta situación,
\(\bm{R_{11}}=1\), y si denotamos \(\bm{\alpha}=\bm{R_{21}}\), éste será un
vector (\(d-1\))--dimensional. De este modo, el estadístico de contraste
resulta:

\[\lambda(\bm{x})^{2/n} = \left|{\bm{I_{d_2}}}-\bm{R_{22}}^{-1}\bm{\alpha}\bm{\alpha}^\prime{}\right|
\stackrel{(a)}{=} 1-\bm{\alpha}^\prime{}\bm{R_{22}}^{-1}\bm{\alpha}
= 1-\bm{R_{12}}\bm{R_{22}}^{-1}\bm{R_{21}} = 1-R^2\in \Lambda\left(d_2,1,n-2\right)\]
siendo \(R\) el coeficiente de correlación múltiple entre la primera
variable y las restantes. Usando que en general

\[\frac{1-\Lambda(d,1,m)}{\Lambda(d,1,m)}=\frac{d}{m-d+1}F_{d,m-d+1},\]
tenemos que

\[\frac{R^2}{1-R^2}\in \frac{d-1}{n-d}F_{d-1,n-d}\]

Finalmente, el coeficiente de correlación múltiple se considerará
significativo cuando el estadístico anterior \(R^2/(1-R^2)\) sea grande,
comparado con la distribución \(F\) de Snédecor.

\subsection{\texorpdfstring{Contraste de independencia de todas las variables \(H_0:\bm\Sigma\) es diagonal}{Contraste de independencia de todas las variables H\_0:\textbackslash bm\textbackslash Sigma es diagonal}}\label{contraste-de-independencia-de-todas-las-variables-h-0-bm-sigma-es-diagonal}

La hipótesis nula consiste en suponer que las variables son
incorrelacionadas, pero, a diferencia del test de esfericidad, no
exigimos que tengan la misma varianza. De nuevo, la incorrelación
equivale a independencia en un contexto de normalidad. Así, bajo la
hipótesis nula, se maximiza la verosimilitud separadamente para cada
variable, dando lugar a los estimadores de la media y la varianza de
dichas variables y en consecuencia a los estimadores del vector de
medias y matriz de covarianzas:

\[\overline{\bm{x}}\qquad \mbox{y} \qquad \hat{\bm\Sigma}=\textnormal{diag}(s_{1}^2, \ldots, s_d^2)\]

Bajo la alternativa, no hay ninguna clase de restricciones sobre los
parámetros \(\bm{\mu}\) y \(\bm\Sigma\) que, por tanto, admiten como
estimadores de máxima verosimilitud \(\overline{\bm{x}}\) y \(\bm{S}\).

Entonces, aplicando el teorema general, el estadístico de contraste
será:

\[-2\log \lambda(\bm{x})=nd\left(a-\log g-1\right)=-n\log |\bm{R}|\]
siendo \(a\) y \(g\) las medias aritmética y geométrica, respectivamente, de
los autovalores de la matriz \(\hat{\bm\Sigma}^{-1}\bm{S}\), y \(\bm{R}\) la
matriz de correlaciones.

El último paso de la expresión anterior se debe a que
\(\bm{R}=\hat{\bm{\Sigma}}^{\bm{-1/2}}\bm{S}\hat{\bm\Sigma}^{\bm{-1/2}}\)
y esta matriz, aún siendo distinta de \(\hat{\bm\Sigma}^{-1}\bm{S}\),
tiene los mismos autovalores que élla. Además, como la diagonal de una
matriz de correlaciones está formada por unos, la traza vale \(d\) y, en
consecuencia, la media de los autovalores vale uno, \(a=1\).

Por último, aproximamos la distribución del estadístico así:

\[-n\log|\bm{R}|\sim \chi_{\frac{1}{2}d(d-1)}^2\] donde el número de
grados de libertad resulta de la diferencia del número de parámetros
independientes bajo la hipótesis nula y bajo la alternativa:
\(d+\frac{1}{2}d(d+1)-(d+d)=\frac{1}{2}d(d-1)\).

\textbf{\includegraphics[width=1em,height=1em]{bookdown_AM_files/figure-latex/fa-icon-00ef1bb1d064b635425d49f06b331298.pdf} ~\textcolor{excolor}{Ejemplo.}} Para ilustrar los contrastes de esfericidad habíamos generado, entre
otros, datos de un vector normal multivariante con matriz de covarianzas
diagonal y tal que no todas las variables tenían la misma varianza. En
concreto, habíamos ejecutado el siguiente código para generar la muestra.

\footnotesize

\begin{Shaded}
\begin{Highlighting}[]
\SpecialCharTok{\textgreater{}} \FunctionTok{library}\NormalTok{(mvtnorm)}
\SpecialCharTok{\textgreater{}}\NormalTok{ d }\OtherTok{\textless{}{-}} \DecValTok{3}
\SpecialCharTok{\textgreater{}}\NormalTok{ mu }\OtherTok{\textless{}{-}} \FunctionTok{c}\NormalTok{(}\DecValTok{0}\NormalTok{, }\DecValTok{0}\NormalTok{, }\DecValTok{0}\NormalTok{)}
\SpecialCharTok{\textgreater{}}\NormalTok{ Sigma }\OtherTok{\textless{}{-}} \FunctionTok{diag}\NormalTok{(}\FunctionTok{c}\NormalTok{(}\DecValTok{3}\NormalTok{, }\DecValTok{5}\NormalTok{, }\DecValTok{1}\NormalTok{))}
\SpecialCharTok{\textgreater{}}\NormalTok{ n }\OtherTok{\textless{}{-}} \DecValTok{100}
\SpecialCharTok{\textgreater{}}\NormalTok{ x }\OtherTok{\textless{}{-}} \FunctionTok{rmvnorm}\NormalTok{(n, mu, Sigma)}
\end{Highlighting}
\end{Shaded}

\normalsize

Si queremos contrastar \(H_0:\bm\Sigma\) es diagonal, debemos calcular la
matriz de correlaciones muestral y su determinante. Como hemos visto,
bajo \(H_0\) la distribución asintótica del estadístico es \(\chi_m^2\),
siendo \(m=d(d-1)/2\). En este caso, para un nivel de significación
\(\alpha=0.05\), no rechazamos la hipótesis nula, ya que el valor del
estadístico es menor que \(\chi_{m,\alpha}^2\). El valor crítico del
estadístico nos lleva a aceptar sin lugar a dudas la hipotesis nula.

\footnotesize

\begin{Shaded}
\begin{Highlighting}[]
\SpecialCharTok{\textgreater{}}\NormalTok{ R }\OtherTok{\textless{}{-}} \FunctionTok{cor}\NormalTok{(x)}
\SpecialCharTok{\textgreater{}}\NormalTok{ estad }\OtherTok{\textless{}{-}} \SpecialCharTok{{-}}\NormalTok{n }\SpecialCharTok{*} \FunctionTok{log}\NormalTok{(}\FunctionTok{det}\NormalTok{(R))}
\SpecialCharTok{\textgreater{}}\NormalTok{ estad}
\NormalTok{[}\DecValTok{1}\NormalTok{] }\FloatTok{0.6970712}
\SpecialCharTok{\textgreater{}}\NormalTok{ m }\OtherTok{\textless{}{-}} \FloatTok{0.5} \SpecialCharTok{*}\NormalTok{ d }\SpecialCharTok{*}\NormalTok{ (d }\SpecialCharTok{{-}} \DecValTok{1}\NormalTok{)}
\SpecialCharTok{\textgreater{}} \FunctionTok{qchisq}\NormalTok{(}\FloatTok{0.95}\NormalTok{, m)}
\NormalTok{[}\DecValTok{1}\NormalTok{] }\FloatTok{7.814728}
\SpecialCharTok{\textgreater{}}\NormalTok{ pval }\OtherTok{\textless{}{-}} \DecValTok{1} \SpecialCharTok{{-}} \FunctionTok{pchisq}\NormalTok{(estad, m)}
\SpecialCharTok{\textgreater{}}\NormalTok{ pval}
\NormalTok{[}\DecValTok{1}\NormalTok{] }\FloatTok{0.8738926}
\end{Highlighting}
\end{Shaded}

\normalsize

\section{Comparación de poblaciones normales multivariantes}\label{comparacion-de-poblaciones-normales-multivariantes}

Vamos a considerar ahora varias poblaciones normales multivariantes, de
las cuales extraemos muestras de manera independiente. En base a estas
muestras contrastaremos hipótesis de igualdad o comparación entre los
parámetros de las diferentes poblaciones. Sean entonces

\[\begin{aligned}
\bm{x_{11}},\ldots,\bm{x_{1n_1}}&\in & N_d(\bm{\mu_1},\bm{\Sigma_1}) \qquad \mbox{independientes},\\
\ldots && \ldots \\
\bm{x_{k1}},\ldots,\bm{x_{kn_k}}&\in & N_d(\bm{\mu_k},\bm{\Sigma_k}) \qquad \mbox{independientes},\end{aligned}\]
donde a su vez las \(k\) muestras independientes entre sí. El abanico de
posibles problemas de comparación de estas \(k\) poblaciones es muy
amplio. En lo que sigue hemos seleccionado algunos casos que nos parecen
más interesantes.

\subsection{Contraste de igualdad de medias de dos poblaciones normales multivariantes con matrices de covarianzas iguales}\label{contraste-de-igualdad-de-medias-de-dos-poblaciones-normales-multivariantes-con-matrices-de-covarianzas-iguales}

En primer lugar trataremos el contraste de igualdad de medias de dos
poblaciones normales multivariantes, suponiendo que la matrices de
covarianzas son desconocidas, pero asumiendo que son iguales. El caso de
más de dos muestras será tratado en el capítulo dedicado al análisis
multivariante de la varianza (MANOVA). Respecto de cómo realizar el
contraste si las dos matrices de covarianzas son distintas, es un
problema que ya estaba presente en observaciones univariantes. Se conoce
como problema de Behrens-Fisher. La complejidad de este problema así
como de sus soluciones, no es esencialmente mayor en el caso
multivariante. No nos adentraremos en más detalles en este momento y nos
ceñiremos al caso de matrices de covarianzas iguales. Repasemos en
primer lugar el contraste de igualdad de medias de dos poblaciones
normales univariantes suponiendo que las varianzas son desconocidas pero
iguales.

\subsubsection{\texorpdfstring{Revisión del caso unidimensional \(H_0:\mu_1=\mu_2\) con \(\sigma_1^2=\sigma_2^2\) desconocidas}{Revisión del caso unidimensional H\_0:\textbackslash mu\_1=\textbackslash mu\_2 con \textbackslash sigma\_1\^{}2=\textbackslash sigma\_2\^{}2 desconocidas}}\label{revision-del-caso-unidimensional-h-0-mu-1-mu-2-con-sigma-1-2-sigma-2-2-desconocidas}

Disponemos de una muestra formada por \(n_1\) variables independientes y
con la misma distribución \(N(\mu_1, \sigma_1^2)\) y una muestra formada
por \(n_2\) variables independientes y con la misma distribución
\(N(\mu_2, \sigma_2^2)\). Es decir

\[\begin{aligned}
{x_{11}},\ldots,{x_{1n_1}}&\in  N({\mu_1},\sigma_1^2) \qquad \mbox{independientes},\\
{x_{21}},\ldots,{x_{2n_2}}&\in  N({\mu_2},\sigma_2^2) \qquad \mbox{independientes}.\end{aligned}\]
Suponemos que las varianzas \(\sigma_1^2\) y \(\sigma_2^2\) son desconocidas
pero iguales. Si suponemos que las varianzas de las dos poblaciones son
iguales, el mejor estimador de la varianza es:

\[{s}_{P}^2 = \frac{(n_1 - 1) {s}_{c1}^2 + (n_2 - 1) {s}_{c2}^2 }{n_1 + n_2 -2}.\]

En la ecuación anterior, \({s}_{c1}^2\) y \({s}_{c2}^2\) denotan la
cuasivarianza muestral de la primera y segunda población,
respectivamente. Si la hipótesis nula \(H_0:\mu_1=\mu_2\) es cierta,
entonces

\[t=\frac{\overline{x}_1 -\overline{x}_2}{\sqrt{\frac{s_{P}^2}{n_1} +\frac{s_P^2}{n_2} }} \in t_{n_1+n_2-2}.\]

Fijado un nivel de significación \(\alpha\), rechazamos la hipótesis nula
\(H_0:\mu_1=\mu_2\) frente a \(H_1:\mu_1\neq\mu_2\) si
\(\left|t\right|\geq t_{\alpha/2}\) siendo \(t_{\alpha/2}\) el cuantil
\(1-\alpha/2\) de la distribución \(t\) de Student con \(n_1+n_2-2\) grados de
libertad. Equivalentemente

\begin{equation}
t^2=\frac{n_1n_2}{n_1+n_2}(\overline{x}_1 -\overline{x}_2)\frac{1}{s_{P}^2}(\overline{x}_1 -\overline{x}_2).
\label{eq:t22}
\end{equation}

\subsubsection{\texorpdfstring{Caso multidimensional \(H_0:\bm{\mu_1}=\bm{\mu_2}\) con \(\bm{\Sigma_1}=\bm{\Sigma_2}\) desconocidas}{Caso multidimensional H\_0:\textbackslash bm\{\textbackslash mu\_1\}=\textbackslash bm\{\textbackslash mu\_2\} con \textbackslash bm\{\textbackslash Sigma\_1\}=\textbackslash bm\{\textbackslash Sigma\_2\} desconocidas}}\label{caso-multidimensional-h-0-bm-mu-1-bm-mu-2-con-bm-sigma-1-bm-sigma-2-desconocidas}

En el caso multivariante disponemos de \[\begin{aligned}
\bm{x_{11}},\ldots,\bm{x_{1n_1}}&\in & N_d(\bm{\mu_1},\bm{\Sigma_1}) \qquad \mbox{independientes}\\
\bm{x_{21}},\ldots,\bm{x_{2n_2}}&\in & N_d(\bm{\mu_2},\bm{\Sigma_2}) \qquad \mbox{independientes}\end{aligned}\]
siendo a su vez las dos muestras independientes entre sí. Además
supondremos que las dos matrices de covarianzas son iguales,
\(\bm{\Sigma_1}=\bm{\Sigma_2}\). En estas condiciones pretendemos
contrastar la hipótesis nula

\[H_0:\bm{\mu_1}=\bm{\mu_2}\]

Parece razonable basar el procedimiento de contraste en una medida de
discrepancia o distancia entre estimadores de \(\bm{\mu_1}\) y de
\(\bm{\bm{\mu_2}}\). De cada una de las dos poblaciones se pueden obtener
estimadores de los parámetros correspondientes a cada población,

\[\bm{\overline{x}_1}=\frac{1}{n_1}\sum_{j=1}^{n_1} \bm{x_{1j}}
\qquad \bm{S_{c1}}=\frac{1}{n_1-1}
\sum_{j=1}^{n_1} \left(\bm{x_{1j}}-\bm{\overline{x}_1}\right)\left(\bm{x_{1j}}-\bm{\overline{x}_1}\right)^\prime{}\]

\[\bm{\overline{x}_2}=\frac{1}{n_2}\sum_{j=1}^{n_2} \bm{x_{2j}}
\qquad \bm{S_{c2}}=\frac{1}{n_2-1}
\sum_{j=1}^{n_2} \left(\bm{x_{2j}}-\bm{\overline{x}_2}\right)\left(\bm{x_{2j}}-\bm{\overline{x}_2}\right)^\prime{}\]

Por ser las dos muestras independientes entre sí, también lo son los
estadísticos procedentes de cada una de ellas. Además, la extensión del
teorema de Fisher al caso multivariante garantiza que

\[\bm{\overline{x}_1}\in N_d\left(\bm{\mu_1},\bm\Sigma/n_1\right),
\ \left(n_1-1\right)\bm{S_{c1}} \in \mbox{Wishart}_d\left(\bm\Sigma,n_1-1\right)
\ \mbox{y son independientes}\]

\[\bm{\overline{x}_2}\in N_d\left(\bm{\mu_2},\bm\Sigma/n_2\right),\ \left(n_2-1\right)\bm{S_{c2}} \in \mbox{Wishart}_d\left(\bm\Sigma,n_2-1\right)\ 
\mbox{y son independientes}.\]

Además, como las dos muestras son
independientes entre sí, también lo son los estadísticos obtenidos en
base a cada una de ellas. Por tanto,

\[\bm{\overline{x}_1}-\bm{\overline{x}_2}
\in N_d\left(\bm{\mu_1}-\bm{\mu_2},\left(\frac{1}{n_1}+\frac{1}{n_2}\right)\bm\Sigma\right),\]

\[\left(n_1-1\right)\bm{S_{c1}}+\left(n_2-1\right)\bm{S_{c2}}
\in \mbox{Wishart}_d\left(\bm\Sigma,n_1+n_2-2\right)\] y además son
independientes. Denotemos mediante

\[\bm{S_{P}}=\frac{\left(n_1-1\right)\bm{S_{c1}}+\left(n_2-1\right)\bm{S_{c2}}}{n_1+n_2-2}\]
una media ponderada de \(\bm{S_{c1}}\) y \(\bm{S_{c2}}\), que sirve como
estimador de la matriz de covarianzas común \(\bm\Sigma\).

Entonces, bajo la hipótesis nula \(H_0:\bm{\mu_1}=\bm{\mu_2}\),

\[T^2=\frac{n_1n_2}{n_1+n_2}
\left(\bm{\overline{x}_1}-\bm{\overline{x}_2}\right)^\prime{}\bm {S_P}^{-1}\left(\bm{\overline{x}_1}-\bm{\overline{x}_2}\right)
\in \Gamma^2\left(d,n_1+n_2-2\right)\]
lo cual lo convierte en un
estadístico adecuado para el contraste de dicha hipótesis. Además, el
estadístico \(T^2\) se puede transformar fácilmente a una \(F\) de Snédecor
utilizando que

\[\frac{n_1+n_2-d-1}{(n_1+n_2-2)d}T^2=F_{d,n_1+n_2-d-1}\]

\textbf{\includegraphics[width=1em,height=1em]{bookdown_AM_files/figure-latex/fa-icon-00ef1bb1d064b635425d49f06b331298.pdf} ~\textcolor{excolor}{Ejemplo.}} Se ha medido la longitud, la anchura y la altura del caparazón de 48
tortugas, 24 hembras y 24 machos. Los vectores de medias y matrices de
covarianzas respectivos son:

\[\bm{\overline{x}_1}=\left(\begin{array}{c} 136.00 \\ 102.58 \\ 51.96 \end{array}\right)
\qquad\qquad \bm{S_1}=\left(\begin{array}{ccc} 432.58 & 259.87 & 161.67 \\ & 164.57 & 98.99 \\ && 63.87 \end{array}\right)\]

\[\bm{\overline{x}_2}=\left(\begin{array}{c} 113.38 \\ 88.29 \\ 40.71 \end{array}\right)
\qquad\qquad \bm{S_2}=\left(\begin{array}{ccc} 132.99 & 75.85 & 35.82 \\ & 47.96 & 20.75 \\ && 10.79 \end{array}\right)\]

Vamos a contrastar que los vectores de medias son iguales en ambos
sexos, suponiendo que en cada sexo las observaciones son normales y que
presentan la misma matriz de covarianzas. En primer lugar introducimos
los datos en R:

\footnotesize

\begin{Shaded}
\begin{Highlighting}[]
\SpecialCharTok{\textgreater{}}\NormalTok{ n1 }\OtherTok{\textless{}{-}} \DecValTok{24}
\SpecialCharTok{\textgreater{}}\NormalTok{ n2 }\OtherTok{\textless{}{-}} \DecValTok{24}
\SpecialCharTok{\textgreater{}}\NormalTok{ d }\OtherTok{\textless{}{-}} \DecValTok{3}
\SpecialCharTok{\textgreater{}}\NormalTok{ bx1 }\OtherTok{\textless{}{-}} \FunctionTok{c}\NormalTok{(}\DecValTok{136}\NormalTok{, }\FloatTok{102.58}\NormalTok{, }\FloatTok{51.96}\NormalTok{)}
\SpecialCharTok{\textgreater{}}\NormalTok{ bx2 }\OtherTok{\textless{}{-}} \FunctionTok{c}\NormalTok{(}\FloatTok{113.38}\NormalTok{, }\FloatTok{88.29}\NormalTok{, }\FloatTok{40.71}\NormalTok{)}
\SpecialCharTok{\textgreater{}}\NormalTok{ S1 }\OtherTok{\textless{}{-}} \FunctionTok{matrix}\NormalTok{(}\FunctionTok{c}\NormalTok{(}\FloatTok{432.58}\NormalTok{, }\FloatTok{259.87}\NormalTok{, }\FloatTok{161.67}\NormalTok{, }\FloatTok{259.87}\NormalTok{, }\FloatTok{164.57}\NormalTok{, }\FloatTok{98.99}\NormalTok{, }\FloatTok{161.67}\NormalTok{, }\FloatTok{98.99}\NormalTok{, }\FloatTok{63.87}\NormalTok{),}
\SpecialCharTok{+}     \AttributeTok{nrow =} \DecValTok{3}\NormalTok{, }\AttributeTok{ncol =} \DecValTok{3}\NormalTok{)}
\SpecialCharTok{\textgreater{}}\NormalTok{ S2 }\OtherTok{\textless{}{-}} \FunctionTok{matrix}\NormalTok{(}\FunctionTok{c}\NormalTok{(}\FloatTok{132.99}\NormalTok{, }\FloatTok{75.85}\NormalTok{, }\FloatTok{35.82}\NormalTok{, }\FloatTok{75.85}\NormalTok{, }\FloatTok{47.96}\NormalTok{, }\FloatTok{20.75}\NormalTok{, }\FloatTok{35.82}\NormalTok{, }\FloatTok{20.75}\NormalTok{, }\FloatTok{10.79}\NormalTok{), }\AttributeTok{nrow =} \DecValTok{3}\NormalTok{,}
\SpecialCharTok{+}     \AttributeTok{ncol =} \DecValTok{3}\NormalTok{)}
\end{Highlighting}
\end{Shaded}

\normalsize

Calculamos ahora el estimador de la matriz de covarianzas común
\(\bm{S_P}\).

\footnotesize

\begin{Shaded}
\begin{Highlighting}[]
\SpecialCharTok{\textgreater{}}\NormalTok{ Sp }\OtherTok{\textless{}{-}}\NormalTok{ (n1 }\SpecialCharTok{*}\NormalTok{ S1 }\SpecialCharTok{+}\NormalTok{ n2 }\SpecialCharTok{*}\NormalTok{ S2)}\SpecialCharTok{/}\NormalTok{(n1 }\SpecialCharTok{+}\NormalTok{ n2 }\SpecialCharTok{{-}} \DecValTok{2}\NormalTok{)}
\end{Highlighting}
\end{Shaded}

\normalsize

Por último calculamos el valor del estadístico y su transformación a una
\(F\) de Snédecor:

\footnotesize

\begin{Shaded}
\begin{Highlighting}[]
\SpecialCharTok{\textgreater{}}\NormalTok{ estad }\OtherTok{\textless{}{-}}\NormalTok{ n1 }\SpecialCharTok{*}\NormalTok{ n2}\SpecialCharTok{/}\NormalTok{(n1 }\SpecialCharTok{+}\NormalTok{ n2) }\SpecialCharTok{*} \FunctionTok{t}\NormalTok{(bx1 }\SpecialCharTok{{-}}\NormalTok{ bx2) }\SpecialCharTok{\%*\%} \FunctionTok{solve}\NormalTok{(Sp) }\SpecialCharTok{\%*\%}\NormalTok{ (bx1 }\SpecialCharTok{{-}}\NormalTok{ bx2)}
\SpecialCharTok{\textgreater{}}\NormalTok{ estad}
\NormalTok{         [,}\DecValTok{1}\NormalTok{]}
\NormalTok{[}\DecValTok{1}\NormalTok{,] }\FloatTok{66.85048}
\SpecialCharTok{\textgreater{}}\NormalTok{ f }\OtherTok{\textless{}{-}}\NormalTok{ (n1 }\SpecialCharTok{+}\NormalTok{ n2 }\SpecialCharTok{{-}}\NormalTok{ d }\SpecialCharTok{{-}} \DecValTok{1}\NormalTok{)}\SpecialCharTok{/}\NormalTok{((n1 }\SpecialCharTok{+}\NormalTok{ n2 }\SpecialCharTok{{-}} \DecValTok{2}\NormalTok{) }\SpecialCharTok{*}\NormalTok{ d) }\SpecialCharTok{*}\NormalTok{ estad}
\SpecialCharTok{\textgreater{}}\NormalTok{ f}
\NormalTok{         [,}\DecValTok{1}\NormalTok{]}
\NormalTok{[}\DecValTok{1}\NormalTok{,] }\FloatTok{21.31464}
\SpecialCharTok{\textgreater{}}\NormalTok{ pval }\OtherTok{\textless{}{-}} \DecValTok{1} \SpecialCharTok{{-}} \FunctionTok{pf}\NormalTok{(f, d, n1 }\SpecialCharTok{+}\NormalTok{ n2 }\SpecialCharTok{{-}}\NormalTok{ d }\SpecialCharTok{{-}} \DecValTok{1}\NormalTok{)}
\SpecialCharTok{\textgreater{}}\NormalTok{ pval}
\NormalTok{             [,}\DecValTok{1}\NormalTok{]}
\NormalTok{[}\DecValTok{1}\NormalTok{,] }\FloatTok{1.119998e{-}08}
\end{Highlighting}
\end{Shaded}

\normalsize

El nivel crítico obtenido es muy pequeño, \ensuremath{1.1199977\times 10^{-8}}, por lo
que se rechaza la igualdad de vectores de medias. Realmente ya podíamos
observar que las medias muestrales eran diferentes, y que las varianzas
no parecían suficientes para ocultar esas diferencias, pero aún así es
conveniente efectuar un contraste como éste, que aporta un criterio
objetivo.

\subsection{Contraste de igualdad de matrices de covarianzas de normales multivariantes}\label{contraste-de-igualdad-de-matrices-de-covarianzas-de-normales-multivariantes}

Recordemos el modelo inicial de esta sección

\[\begin{aligned}
\bm{x_{11}},\ldots,\bm{x_{1n_1}}&\in  N_d(\bm{\mu_1},\bm{\Sigma_1}) \qquad \mbox{independientes}\\
\ldots & \ldots \\
\bm{x_{k1}},\ldots,\bm{x_{kn_k}}&\in  N_d(\bm{\mu_k},\bm{\Sigma_k}) \qquad \mbox{independientes}\end{aligned}\]
siendo a su vez las \(k\) muestras independientes entre sí. Nos planteamos
el contraste de la hipótesis nula
\(H_0:\bm{\Sigma_1}=\cdots=\bm{\Sigma_k}\). Los vectores de medias están
exentos de restricciones tanto bajo la hipótesis nula como bajo la
alternativa.

Vamos a obtener el test de razón de verosimilitudes. Lo primero que
observamos es que, como las muestras son independientes entre sí, la
función de verosimilitud se puede expresar como producto de las
funciones de verosimilitud de cada una de las \(k\) poblaciones. Bajo la
alternativa,

\[\begin{aligned}
\sup_{\bm{\mu_1},\bm\Sigma_1,\ldots,\bm{\mu_k},\bm {\Sigma_k}}
\log \prod_{i=1}^k L\left(\bm{\mu_i},\bm{\Sigma_i}\right)
&=\sum_{i=1}^k \sup_{\bm{\mu_i},\bm {\Sigma_i}} \log L\left(\bm{\mu_i},\bm {\Sigma_i}\right)\\
&=\sum_{i=1}^k \log L\left(\bm{\overline{x}_i},\bm{S_i}\right)  \\
&=\sum_{i=1}^k \left[-\frac{n_id}{2}\log (2\pi)
-\frac{1}{2} n_i\log \left|\bm{S_i}\right|-\frac{n_id}{2}\right]\\
&=-\frac{nd}{2}\log (2\pi)
-\frac{1}{2}\sum_{i=1}^k n_i\log \left|\bm{S_i}\right|-\frac{nd}{2}\end{aligned}\]

Bajo la hipótesis nula, denotemos
\(\bm\Sigma=\bm{\Sigma_1}=\cdots=\bm {\Sigma_k}\) a la matriz de
covarianzas común,

\[\begin{aligned}
\sup_{\bm{\mu_1},\ldots,\bm{\mu_k},\bm\Sigma}
\log \prod_{i=1}^k L\left(\bm{\mu_i},\bm\Sigma\right)
&=\sup_{\bm\Sigma}\sum_{i=1}^k \sup_{\bm{\mu_i}} \log L\left(\bm{\mu_i},\bm\Sigma\right)\\
&=\sup_{\bm\Sigma}\sum_{i=1}^k \log L\left(\bm{\overline{x}_i},\bm\Sigma\right) \\
&=\sup_{\bm\Sigma}\sum_{i=1}^k \left[-\frac{n_id}{2}\log (2\pi)
-\frac{n_i}{2}
\left(\log |\bm\Sigma|+\mbox{traza}\left(\bm\Sigma^{-1}\bm{S_i}\right)\right)\right] \\
&=\sup_{\bm\Sigma}\left[-\frac{nd}{2}\log (2\pi)
-\frac{n}{2}
\left(\log |\bm\Sigma|+\mbox{traza}\left(\bm\Sigma^{-1}\frac{\bm{Q}}{n}\right)\right)\right]\\
&=-\frac{nd}{2}\log (2\pi)
-\frac{n}{2} \log \left|\frac{\bm{Q}}{n}\right|-\frac{nd}{2}\end{aligned}\]
siendo \(\bm{Q}=\sum_{i=1}^k n_i\bm{S_i}\). Nótese que \(\bm{Q}/n\) es el
estimador de máxima verosimilitud de la matriz de covarianzas común,
\(\bm\Sigma\), y resulta ser una media ponderada de las matrices de
covarianzas muestrales provenientes de cada población.

Finalmente, el estadístico de razón de verosimilitudes adopta la forma:

\[\begin{aligned}
-2\log \frac{\sup_{\bm{\mu_1},\ldots,\bm{\mu_k},\bm\Sigma} \prod_{i=1}^k L\left(\bm{\mu_i},\bm\Sigma\right)}
{\sup_{\bm{\mu_1},\bm{\Sigma_1},\ldots,\bm{\mu_k},\bm {\Sigma_k}} \prod_{i=1}^k L\left(\bm{\mu_i},\bm {\Sigma_i}\right)}
&= n\log \left|\frac{\bm{Q}}{n}\right|-\sum_{i=1}^k n_i\log \left|\bm{S_i}\right|  \\
&= \sum_{i=1}^k n_i\log \left|\frac{\bm{Q}}{n}\bm{S_i}^{-1}\right|
\sim \chi_{\frac{1}{2}d(d+1)(k-1)}^2\end{aligned}\] cuya distribución
hemos aproximado por una ji-cuadrado cuyos grados de libertad resultan
de la diferencia de parámetros independientes entre la hipótesis nula y
la alternativa.

\textbf{\includegraphics[width=1em,height=1em]{bookdown_AM_files/figure-latex/fa-icon-00ef1bb1d064b635425d49f06b331298.pdf} ~\textcolor{excolor}{Ejemplo.}} Para el ejemplo anterior de las tortugas, vamos a contrastar que las
matrices de covarianzas son iguales en ambos sexos, suponiendo que en
cada sexo las observaciones son normales.

Calculamos el estimador de la matriz de covarianzas común bajo \(H_0\)

\footnotesize

\begin{Shaded}
\begin{Highlighting}[]
\SpecialCharTok{\textgreater{}}\NormalTok{ Q }\OtherTok{\textless{}{-}}\NormalTok{ n1 }\SpecialCharTok{*}\NormalTok{ S1 }\SpecialCharTok{+}\NormalTok{ n2 }\SpecialCharTok{*}\NormalTok{ S2}
\SpecialCharTok{\textgreater{}}\NormalTok{ n }\OtherTok{\textless{}{-}}\NormalTok{ n1 }\SpecialCharTok{+}\NormalTok{ n2}
\SpecialCharTok{\textgreater{}}\NormalTok{ Qn }\OtherTok{\textless{}{-}}\NormalTok{ Q}\SpecialCharTok{/}\NormalTok{n}
\end{Highlighting}
\end{Shaded}

\normalsize

Calculamos ahora el valor del estadístico y el \(p\)-valor
correspondiente.

\footnotesize

\begin{Shaded}
\begin{Highlighting}[]
\SpecialCharTok{\textgreater{}}\NormalTok{ estad }\OtherTok{\textless{}{-}}\NormalTok{ n1 }\SpecialCharTok{*} \FunctionTok{log}\NormalTok{(}\FunctionTok{det}\NormalTok{(Qn }\SpecialCharTok{\%*\%} \FunctionTok{solve}\NormalTok{(S1))) }\SpecialCharTok{+}\NormalTok{ n2 }\SpecialCharTok{*} \FunctionTok{log}\NormalTok{(}\FunctionTok{det}\NormalTok{(Qn }\SpecialCharTok{\%*\%} \FunctionTok{solve}\NormalTok{(S2)))}
\SpecialCharTok{\textgreater{}}\NormalTok{ estad}
\NormalTok{[}\DecValTok{1}\NormalTok{] }\FloatTok{26.98611}
\SpecialCharTok{\textgreater{}}\NormalTok{ k }\OtherTok{\textless{}{-}} \DecValTok{2}
\SpecialCharTok{\textgreater{}}\NormalTok{ gl }\OtherTok{\textless{}{-}} \FloatTok{0.5} \SpecialCharTok{*}\NormalTok{ d }\SpecialCharTok{*}\NormalTok{ (d }\SpecialCharTok{+} \DecValTok{1}\NormalTok{) }\SpecialCharTok{*}\NormalTok{ (k }\SpecialCharTok{{-}} \DecValTok{1}\NormalTok{)}
\SpecialCharTok{\textgreater{}}\NormalTok{ pval }\OtherTok{\textless{}{-}} \DecValTok{1} \SpecialCharTok{{-}} \FunctionTok{pchisq}\NormalTok{(estad, gl)}
\SpecialCharTok{\textgreater{}}\NormalTok{ pval}
\NormalTok{[}\DecValTok{1}\NormalTok{] }\FloatTok{0.0001456775}
\end{Highlighting}
\end{Shaded}

\normalsize

El nivel crítico obtenido es bastante pequeño, \ensuremath{1.4567753\times 10^{-4}}, por lo que se
rechaza la igualdad de matrices de covarianzas. A la vista de este
resultado, y como en el ejemplo anterior se supuso igualdad de matrices
de covarianzas para el contraste de igualdad de los vectores de medias,
tendría sentido reconsiderar el procedimiento aplicado en aquella
ocasión. Para ello, en lugar de emplear la estimación de una matriz de
covarianzas común, habría que considerar las dos matrices de manera más
separada. No entraremos en más detalle en este momento.

\section{Contrastes de normalidad}\label{secnorm}

\subsection{Revisión del caso univariante}\label{secuni}

Se han desarrollado muchos procedimientos para el contraste de la
normalidad en el caso univariante. Un procedimiento gráfico sencillo
para evaluar la normalidad de un conjunto de observaciones de una
variable es el \(Q\)-\(Q\) plot, que compara los cuantiles de la muestra
con los cuantiles de la distribución normal univariante. En la Figura
\ref{fig:qq} se muestran los
\(Q\)-\(Q\) plots correspondientes a muestras generadas a partir de
distintas distribuciones. Cuando los puntos se acercan a la línea recta
representada (como en la primera gráfica), no hay evidencia de falta de
normalidad. La desviación de los puntos con respecto a la recta indica
falta de normalidad. Por ejemplo, un \(Q\)-\(Q\) plot como el de los datos
generados a partir de la \(t\) de Student, es típico en datos que
provienen de distribuciones con colas más pesadas que la Normal. Para
datos que provienen de distribuciones con colas menos pesadas que la
Normal se suelen obtener \(Q\)-\(Q\) plots como el de los datos uniformes
en la Figura \ref{fig:qq}. Un
\(Q\)-\(Q\) plot como el de los datos generados a partir de la \(\chi^2\) es
típico en datos que provienen de distribuciones asimétricas. Si bien el
\(Q\)-\(Q\) plot nos puede servir para valorar de forma visual la
normalidad de una muestra, es adecuado hacer uso de procedimientos
objetivos. Veremos a continuación algunos de dichos procedimientos,
aquellos que se extienden de manera más natural al caso multivariante.
Serán los métodos basados en la asimetría, la kurtosis, y el test de
Shapiro-Wilk.

\begin{figure}

{\centering \resizebox{.4\linewidth}{!}{\includegraphics[width=0.5\linewidth]{bookdown_AM_files/figure-latex/qq-1} } \resizebox{.4\linewidth}{!}{\includegraphics[width=0.5\linewidth]{bookdown_AM_files/figure-latex/qq-2} } 

}

\caption{Funciones de densidad correspondientes a una distribución normal de media cero y varianza $5/3$ (rojo), $t$ de Student con 5 grados de libertad (azul), uniforme en $\left[-\sqrt{5},\sqrt{5}\right]$ (verde) y $\chi^2$ con 5 grados de libertad (rosa) y $Q$-$Q$ plots correspondientes a muestras generadas a partir de dichas distribuciones.}\label{fig:qq}
\end{figure}

\subsubsection{Test basado en la asimetría}\label{test-basado-en-la-asimetria}

Es bien conocido que la distribución normal es simétrica en torno a su
media. Como además, la forma más común de desviarse respecto de la
normalidad es por falta de simetría, por ejemplo en variables positivas,
parece lógico construir un método de contraste en base a cierta medida
de asimetría.

Si \(x_1,\ldots,x_n\) es una muestra aleatoria simple (observaciones
independientes e idénticamente distribuidas), el coeficiente de
asimetría muestral se define como

\[A=\frac{\frac{1}{n}\sum_{i=1}^n\left(x_i-\overline{x}\right)^3}{s^3}
=\frac{1}{n}\sum_{i=1}^n\left(\frac{x_i-\overline{x}}{s}\right)^3\]
siendo \(\overline{x}\) la media muestral y \(s^2\) la varianza muestral. El
coeficiente de asimetría es la media de las potencias cúbicas (también
llamado momento de orden tres) de los valores estandarizados. Al estar
estandarizados, habrá observaciones a ambos lados de la media,
manteniendo un equilibrio (en orden uno). Sin embargo, la potencia tres
que se emplea en el coeficiente de asimetría, rompe el equilibrio en
caso de distribución asimétrica. En la Figura \ref{fig:asim2} se muestran
los coeficiente de asimetría (para calcularlos puedes utilizar la
función \texttt{skewness} de la librería \texttt{moments}) de datos generados de
distintas distribuciones.

\begin{figure}

{\centering \resizebox{.28\linewidth}{!}{\includegraphics[width=0.33\linewidth]{bookdown_AM_files/figure-latex/asim2-1} } \resizebox{.28\linewidth}{!}{\includegraphics[width=0.33\linewidth]{bookdown_AM_files/figure-latex/asim2-2} } \resizebox{.28\linewidth}{!}{\includegraphics[width=0.33\linewidth]{bookdown_AM_files/figure-latex/asim2-3} } 

}

\caption{Histogramas y coeficiente de asimetría correspondientes a datos generados de una distribución normal (izquierda), de una exponencial (centro) y de los mismos datos exponenciales tras cambiarles el signo (derecha). Los datos se representan en azul. En rojo, media muestral.}\label{fig:asim2}
\end{figure}

Si la distribución de los datos muestrales es normal, entonces el
coeficiente de asimetría tiene distribución asintótica normal de media
cero y varianza \(6/n\), por lo que se puede emplear como estadístico de
contraste el siguiente:

\[\sqrt{\frac{n}{6}}\ A\approx N(0,1)\]

Se rechazará la normalidad, en base al coeficiente de asimetría, cuando
el estadístico anterior sea muy grande (en positivo o negativo), en
comparación con los cuantiles de la \(N(0,1)\).

\textbf{\includegraphics[width=1em,height=1em]{bookdown_AM_files/figure-latex/fa-icon-00ef1bb1d064b635425d49f06b331298.pdf} ~\textcolor{excolor}{Ejemplo.}} Vamos a presentar los resultados del contraste de normalidad basado en
asimetría aplicado a datos simulados de distintas distribuciones.
Mostramos en primer lugar los valores del coeficiente de asimetría, el
estadístico de contraste y nivel crítico para datos simulados de la
distribución normal.

\footnotesize

\begin{Shaded}
\begin{Highlighting}[]
\SpecialCharTok{\textgreater{}} \FunctionTok{library}\NormalTok{(moments)}
\SpecialCharTok{\textgreater{}}\NormalTok{ n }\OtherTok{\textless{}{-}} \DecValTok{30}
\SpecialCharTok{\textgreater{}}\NormalTok{ x }\OtherTok{\textless{}{-}} \FunctionTok{rnorm}\NormalTok{(n)}
\SpecialCharTok{\textgreater{}}\NormalTok{ A }\OtherTok{\textless{}{-}} \FunctionTok{skewness}\NormalTok{(x)}
\SpecialCharTok{\textgreater{}}\NormalTok{ A}
\NormalTok{[}\DecValTok{1}\NormalTok{] }\SpecialCharTok{{-}}\FloatTok{0.09408815}
\SpecialCharTok{\textgreater{}}\NormalTok{ estad }\OtherTok{\textless{}{-}} \FunctionTok{sqrt}\NormalTok{(n}\SpecialCharTok{/}\DecValTok{6}\NormalTok{) }\SpecialCharTok{*}\NormalTok{ A}
\SpecialCharTok{\textgreater{}}\NormalTok{ estad}
\NormalTok{[}\DecValTok{1}\NormalTok{] }\SpecialCharTok{{-}}\FloatTok{0.2103875}
\SpecialCharTok{\textgreater{}}\NormalTok{ pval }\OtherTok{\textless{}{-}} \DecValTok{2} \SpecialCharTok{*}\NormalTok{ (}\DecValTok{1} \SpecialCharTok{{-}} \FunctionTok{pnorm}\NormalTok{(}\FunctionTok{abs}\NormalTok{(estad)))}
\SpecialCharTok{\textgreater{}}\NormalTok{ pval}
\NormalTok{[}\DecValTok{1}\NormalTok{] }\FloatTok{0.8333653}
\end{Highlighting}
\end{Shaded}

\normalsize

Como cabía esperar, la pequeña asimetría muestral de los datos normales
no resulta significativa. Veamos ahora lo que ocurre cuando generamos
datos de una distribución asimétrica como la exponencial.

\footnotesize

\begin{Shaded}
\begin{Highlighting}[]
\SpecialCharTok{\textgreater{}}\NormalTok{ x }\OtherTok{\textless{}{-}} \FunctionTok{rexp}\NormalTok{(n)}
\SpecialCharTok{\textgreater{}}\NormalTok{ A }\OtherTok{\textless{}{-}} \FunctionTok{skewness}\NormalTok{(x)}
\SpecialCharTok{\textgreater{}}\NormalTok{ A}
\NormalTok{[}\DecValTok{1}\NormalTok{] }\FloatTok{1.754585}
\SpecialCharTok{\textgreater{}}\NormalTok{ estad }\OtherTok{\textless{}{-}} \FunctionTok{sqrt}\NormalTok{(n}\SpecialCharTok{/}\DecValTok{6}\NormalTok{) }\SpecialCharTok{*}\NormalTok{ A}
\SpecialCharTok{\textgreater{}}\NormalTok{ estad}
\NormalTok{[}\DecValTok{1}\NormalTok{] }\FloatTok{3.923371}
\SpecialCharTok{\textgreater{}}\NormalTok{ pval }\OtherTok{\textless{}{-}} \DecValTok{2} \SpecialCharTok{*}\NormalTok{ (}\DecValTok{1} \SpecialCharTok{{-}} \FunctionTok{pnorm}\NormalTok{(}\FunctionTok{abs}\NormalTok{(estad)))}
\SpecialCharTok{\textgreater{}}\NormalTok{ pval}
\NormalTok{[}\DecValTok{1}\NormalTok{] }\FloatTok{8.73186e{-}05}
\end{Highlighting}
\end{Shaded}

\normalsize

\subsubsection{Test basado en la kurtosis}\label{test-basado-en-la-kurtosis}

Partimos de la muestra \(x_1,\ldots,x_n\). El momento de orden cuatro de
los valores estandarizados se calcula como:

\[\frac{1}{n}\sum_{i=1}^n\left(\frac{x_i-\overline{x}}{s}\right)^4,\]

siendo \(\overline{x}\) la media muestral y \(s^2\) la varianza muestral.
Recordemos que las características de posición se calculan mediante un
momento de orden uno (la media), las de dispersión a través de un
momento de orden dos (la varianza) y la asimetría con un momento de
orden tres (el coeficiente de asimetría). El momento de orden cuatro
permite medir características de forma, que son un refinamiento todavía
más detallado en el análisis de la distribución. Si bien en algunos
textos la kurtosis se define como el momento de orden 4, es más habitual
definir la kurtosis como:

\[K=\frac{1}{n}\sum_{i=1}^n\left(\frac{x_i-\overline{x}}{s}\right)^4-3.\]
En esta definición se resta tres para que el coeficiente valga cero en
la distribución normal. Así definida, \(K\) también se conoce como \emph{exceso
de kurtosis}.

En la Figura \ref{fig:qq} se
representaban tres distribuciones con la misma media, desviación típica
y asimetría (asimetría cero, pues son simétricas). Son la \(t\) de Student
con 5 grados de libertad, la normal y la uniforme, ambas con la misma
media y desviación típica que la \(t_5\). Para ello, se toma la normal con
media cero y desviación típica \(\sqrt{5/3}\), y la uniforme entre
\(-\sqrt{5}\) y \(\sqrt{5}\). Vemos que la distribución \(t\) de Student es la
que presenta una mayor concentración de probabilidad en las proximidades
de la media, así como en las colas (lejanías extremas)\footnote{El fenómeno de la cola es difícilmente perceptible en la Figura
  \ref{fig:qq}. Aún así, se
  puede observar que el trazo rojo, correspondiente a la \(t_5\), se
  encuentra ligeramente más elevado que el azul, correspondiente a la
  normal, en las colas. Es una pequeña diferencia a la vista, porque
  ambas densidades convergen a cero en las colas. Sin embargo, la
  consecuencia es que la \(t_5\) presenta una frecuencia no despreciable
  de observaciones muy grandes, cosa que no ocurre con la normal.
  Estas observaciones grandes (o extremas), al elevarlas a la potencia
  cuarta, aportan una contribución considerable a la kurtosis.}. Esto hace
que tenga un coeficiente de apuntamiento más grande. Por el contrario,
la distribución uniforme es la que presenta menor densidad en las
proximidades de la media, y carece de colas. En este caso, la densidad
es más grande en los valores con desviaciones intermedias \((\mu-\sigma)\)
y \((\mu+\sigma)\). Como resultado, el coeficiente de apuntamiento será
más pequeño para la distribución uniforme. La distribución normal
presenta un comportamiento intermedio entre la \(t_5\) y la uniforme.

Si una distribución tiene un apuntamiento similar al de la normal
decimos que es \textbf{mesocúrtica}, si el apuntamiento es mayor se dice
\textbf{leptocúrtica}, y si es menor se dice \textbf{platicúrtica}. La \(t\) de
Student es una distribución leptocúrtica y la uniforme platicúrtica.

Volviendo al propósito de contraste de la normalidad, si estamos ante
una distribución simétrica, podemos valorar una posible desviación del
modelo normal en función de su apuntamiento o kurtosis. Rechazaremos la
normalidad si la kurtosis muestral es mucho mayor o mucho menor que
cero, que es lo que corresponde a la distribución normal.

Se puede demostrar que si la distribución de los datos muestrales es
normal, entonces la kurtosis tiene distribución asintótica normal de
media cero y varianza \(24/n\), por lo que se puede emplear como
estadístico de contraste el siguiente:

\[\sqrt{\frac{n}{24}}\ K\approx N(0,1).\]

Se rechazará la normalidad, en
base a la kurtosis, cuando el estadístico anterior sea muy grande (en
positivo o negativo), en comparación con los cuantiles de la \(N(0,1)\).

\textbf{\includegraphics[width=1em,height=1em]{bookdown_AM_files/figure-latex/fa-icon-00ef1bb1d064b635425d49f06b331298.pdf} ~\textcolor{excolor}{Ejemplo.}} Vamos a extraer muestras simuladas de tamaño 100 de cada una de las tres
distribuciones que se han comentado en esta sección, y cuyas densidades
se representan en la Figura \hyperref[fig:qq]{2}: \(t_5\) de Student, normal y uniforme. Aplicamos a
continuación el contraste de normalidad basado en la kurtosis. En primer
lugar, generamos datos de una normal con varianza \(5/3\):

\footnotesize

\begin{Shaded}
\begin{Highlighting}[]
\SpecialCharTok{\textgreater{}} \FunctionTok{library}\NormalTok{(moments)}
\SpecialCharTok{\textgreater{}}\NormalTok{ n }\OtherTok{\textless{}{-}} \DecValTok{100}
\SpecialCharTok{\textgreater{}}\NormalTok{ x }\OtherTok{\textless{}{-}} \FunctionTok{rnorm}\NormalTok{(n, }\AttributeTok{sd =} \FunctionTok{sqrt}\NormalTok{(}\DecValTok{5}\SpecialCharTok{/}\DecValTok{3}\NormalTok{))}
\SpecialCharTok{\textgreater{}}\NormalTok{ K }\OtherTok{\textless{}{-}} \FunctionTok{kurtosis}\NormalTok{(x) }\SpecialCharTok{{-}} \DecValTok{3}
\SpecialCharTok{\textgreater{}}\NormalTok{ K}
\NormalTok{[}\DecValTok{1}\NormalTok{] }\FloatTok{0.6656975}
\SpecialCharTok{\textgreater{}}\NormalTok{ estad }\OtherTok{\textless{}{-}} \FunctionTok{sqrt}\NormalTok{(n}\SpecialCharTok{/}\DecValTok{24}\NormalTok{) }\SpecialCharTok{*}\NormalTok{ K}
\SpecialCharTok{\textgreater{}}\NormalTok{ estad}
\NormalTok{[}\DecValTok{1}\NormalTok{] }\FloatTok{1.358849}
\SpecialCharTok{\textgreater{}}\NormalTok{ pval }\OtherTok{\textless{}{-}} \DecValTok{2} \SpecialCharTok{*}\NormalTok{ (}\DecValTok{1} \SpecialCharTok{{-}} \FunctionTok{pnorm}\NormalTok{(}\FunctionTok{abs}\NormalTok{(estad)))}
\SpecialCharTok{\textgreater{}}\NormalTok{ pval}
\NormalTok{[}\DecValTok{1}\NormalTok{] }\FloatTok{0.1741943}
\end{Highlighting}
\end{Shaded}

\normalsize

Aplicamos ahora el contraste a una muestra generada a partir de una
\(t_5\):

\footnotesize

\begin{Shaded}
\begin{Highlighting}[]
\SpecialCharTok{\textgreater{}}\NormalTok{ x }\OtherTok{\textless{}{-}} \FunctionTok{rt}\NormalTok{(n, }\DecValTok{5}\NormalTok{)}
\SpecialCharTok{\textgreater{}}\NormalTok{ K }\OtherTok{\textless{}{-}} \FunctionTok{kurtosis}\NormalTok{(x) }\SpecialCharTok{{-}} \DecValTok{3}
\SpecialCharTok{\textgreater{}}\NormalTok{ K}
\NormalTok{[}\DecValTok{1}\NormalTok{] }\FloatTok{1.095394}
\SpecialCharTok{\textgreater{}}\NormalTok{ estad }\OtherTok{\textless{}{-}} \FunctionTok{sqrt}\NormalTok{(n}\SpecialCharTok{/}\DecValTok{24}\NormalTok{) }\SpecialCharTok{*}\NormalTok{ K}
\SpecialCharTok{\textgreater{}}\NormalTok{ estad}
\NormalTok{[}\DecValTok{1}\NormalTok{] }\FloatTok{2.235964}
\SpecialCharTok{\textgreater{}}\NormalTok{ pval }\OtherTok{\textless{}{-}} \DecValTok{2} \SpecialCharTok{*}\NormalTok{ (}\DecValTok{1} \SpecialCharTok{{-}} \FunctionTok{pnorm}\NormalTok{(}\FunctionTok{abs}\NormalTok{(estad)))}
\SpecialCharTok{\textgreater{}}\NormalTok{ pval}
\NormalTok{[}\DecValTok{1}\NormalTok{] }\FloatTok{0.02535411}
\end{Highlighting}
\end{Shaded}

\normalsize

Por último, aplicamos el contraste a una muestra generada a partir de la
distribución uniforme en el intervalo \(\left[-\sqrt{5},\sqrt{5}\right]\).

\footnotesize

\begin{Shaded}
\begin{Highlighting}[]
\SpecialCharTok{\textgreater{}}\NormalTok{ x }\OtherTok{\textless{}{-}} \FunctionTok{runif}\NormalTok{(n, }\SpecialCharTok{{-}}\FunctionTok{sqrt}\NormalTok{(}\DecValTok{5}\NormalTok{), }\FunctionTok{sqrt}\NormalTok{(}\DecValTok{5}\NormalTok{))}
\SpecialCharTok{\textgreater{}}\NormalTok{ K }\OtherTok{\textless{}{-}} \FunctionTok{kurtosis}\NormalTok{(x) }\SpecialCharTok{{-}} \DecValTok{3}
\SpecialCharTok{\textgreater{}}\NormalTok{ K}
\NormalTok{[}\DecValTok{1}\NormalTok{] }\SpecialCharTok{{-}}\FloatTok{1.031331}
\SpecialCharTok{\textgreater{}}\NormalTok{ estad }\OtherTok{\textless{}{-}} \FunctionTok{sqrt}\NormalTok{(n}\SpecialCharTok{/}\DecValTok{24}\NormalTok{) }\SpecialCharTok{*}\NormalTok{ K}
\SpecialCharTok{\textgreater{}}\NormalTok{ estad}
\NormalTok{[}\DecValTok{1}\NormalTok{] }\SpecialCharTok{{-}}\FloatTok{2.105195}
\SpecialCharTok{\textgreater{}}\NormalTok{ pval }\OtherTok{\textless{}{-}} \DecValTok{2} \SpecialCharTok{*}\NormalTok{ (}\DecValTok{1} \SpecialCharTok{{-}} \FunctionTok{pnorm}\NormalTok{(}\FunctionTok{abs}\NormalTok{(estad)))}
\SpecialCharTok{\textgreater{}}\NormalTok{ pval}
\NormalTok{[}\DecValTok{1}\NormalTok{] }\FloatTok{0.03527432}
\end{Highlighting}
\end{Shaded}

\normalsize

Los resultados son los que cabía esperar. Así, el coeficiente de
apuntamiento muestral es próximo a cero para la normal, negativo para la
uniforme y positivo para la \(t_5\). Al estandarizarlos, se obtienen
cantidades comparables con la \(N(0,1)\), que arrojan los niveles críticos
que se muestran en las salidas. Las discrepancias más significativas se
obtienen con la \(t_5\) de Student.

\subsubsection{Test de Shapiro-Wilk}\label{test-de-shapiro-wilk}

Veamos ahora uno de los contrastes más utilizados para contrastar
normalidad. Se trata del test de Shapiro~-Wilk, que plantea como
hipótesis nula que la muestra \(x_1,\ldots,x_n\) proviene de una población
normal. Este procedimiento es más general, pues permite detectar tanto
defectos de simetría como de kurtosis, o incluso de otro tipo no
considerado por los tests anteriores. Consideremos los datos
estandarizados

\[z_i=\frac{x_i-\overline{x}}{s}\qquad i\in\{1,\ldots,n\}.\]

El estadístico de Shapiro-Wilk se construye de la siguiente manera:

\[W=\sum_{i=1}^{[n/2]}a_{i,n} \left(z_{(n-i+1):n}-z_{i:n}\right)\]

siendo \(z_{1:n}<\ldots<z_{n:n}\) la muestra ordenada de los datos
estandarizados y \(a_{i,n}\) ciertas constantes. Consiste en calcular las
distancias entre los datos de la muestra ordenada, simétricos respecto
de la mediana, esto es, la distancia entre el primero y el último, el
segundo y el penúltimo, y así sucesivamente; en general el \(z_{i:n}\) y
el \(z_{(n-i+1):n}\). El propósito es comparar estas distancias con las
que habría en una muestra de observaciones normales. Podremos llevar a
cabo el test de Shapiro~-Wilk en R mediante la función \texttt{shapiro.test},
como se muestra en el siguiente ejemplo.

\textbf{\includegraphics[width=1em,height=1em]{bookdown_AM_files/figure-latex/fa-icon-00ef1bb1d064b635425d49f06b331298.pdf} ~\textcolor{excolor}{Ejemplo.}} Vamos a extraer muestras simuladas de tamaño 100 de cada una
de las distribuciones vistas en las secciones anteriores: normal,
\(\chi^2\), uniforme y \(t\) de Student, ver Figura \ref{fig:qq}. Los resultados
del test de Shapiro-Wilk se muestran a continuación:

\footnotesize

\begin{Shaded}
\begin{Highlighting}[]
\SpecialCharTok{\textgreater{}}\NormalTok{ n }\OtherTok{\textless{}{-}} \DecValTok{30}
\SpecialCharTok{\textgreater{}}\NormalTok{ x1 }\OtherTok{\textless{}{-}} \FunctionTok{rnorm}\NormalTok{(n, }\AttributeTok{sd =} \FunctionTok{sqrt}\NormalTok{(}\DecValTok{5}\SpecialCharTok{/}\DecValTok{3}\NormalTok{))}
\SpecialCharTok{\textgreater{}}\NormalTok{ x2 }\OtherTok{\textless{}{-}} \FunctionTok{rt}\NormalTok{(n, }\DecValTok{5}\NormalTok{)}
\SpecialCharTok{\textgreater{}}\NormalTok{ x3 }\OtherTok{\textless{}{-}} \FunctionTok{runif}\NormalTok{(n, }\SpecialCharTok{{-}}\FunctionTok{sqrt}\NormalTok{(}\DecValTok{5}\NormalTok{), }\FunctionTok{sqrt}\NormalTok{(}\DecValTok{5}\NormalTok{))}
\SpecialCharTok{\textgreater{}}\NormalTok{ x4 }\OtherTok{\textless{}{-}} \FunctionTok{rchisq}\NormalTok{(n, }\DecValTok{5}\NormalTok{)}
\SpecialCharTok{\textgreater{}} \FunctionTok{shapiro.test}\NormalTok{(x1)}

\NormalTok{    Shapiro}\SpecialCharTok{{-}}\NormalTok{Wilk normality test}

\NormalTok{data}\SpecialCharTok{:}\NormalTok{  x1}
\NormalTok{W }\OtherTok{=} \FloatTok{0.9515}\NormalTok{, p}\SpecialCharTok{{-}}\NormalTok{value }\OtherTok{=} \FloatTok{0.1854}
\SpecialCharTok{\textgreater{}} \FunctionTok{shapiro.test}\NormalTok{(x2)}

\NormalTok{    Shapiro}\SpecialCharTok{{-}}\NormalTok{Wilk normality test}

\NormalTok{data}\SpecialCharTok{:}\NormalTok{  x2}
\NormalTok{W }\OtherTok{=} \FloatTok{0.88543}\NormalTok{, p}\SpecialCharTok{{-}}\NormalTok{value }\OtherTok{=} \FloatTok{0.003765}
\SpecialCharTok{\textgreater{}} \FunctionTok{shapiro.test}\NormalTok{(x3)}

\NormalTok{    Shapiro}\SpecialCharTok{{-}}\NormalTok{Wilk normality test}

\NormalTok{data}\SpecialCharTok{:}\NormalTok{  x3}
\NormalTok{W }\OtherTok{=} \FloatTok{0.90155}\NormalTok{, p}\SpecialCharTok{{-}}\NormalTok{value }\OtherTok{=} \FloatTok{0.009167}
\SpecialCharTok{\textgreater{}} \FunctionTok{shapiro.test}\NormalTok{(x4)}

\NormalTok{    Shapiro}\SpecialCharTok{{-}}\NormalTok{Wilk normality test}

\NormalTok{data}\SpecialCharTok{:}\NormalTok{  x4}
\NormalTok{W }\OtherTok{=} \FloatTok{0.89688}\NormalTok{, p}\SpecialCharTok{{-}}\NormalTok{value }\OtherTok{=} \FloatTok{0.007053}
\end{Highlighting}
\end{Shaded}

\normalsize

En los tres casos donde los datos no proceden de una distribución normal
los niveles críticos conducen al rechazo de la normalidad.

\subsection{Contraste de la normalidad multivariante}\label{secmul}

En esta sección veremos cómo se pueden extender los conceptos de
asimetría y kurtosis al caso multivariante. Asimismo, también se verá la
manera de extender el test de Shapiro-Wilk con datos multivariantes.
Empezamos estudiando ciertos conceptos previos en relación con este tipo
de contrastes.

\subsubsection{Invariancia, estandarización, y distancias y ángulos de Mahalanobis}\label{invariancia-estandarizacion-y-distancias-y-angulos-de-mahalanobis}

Sabemos que la normalidad se conserva frente a traslaciones y
transformaciones lineales, esto es:

Para cualesquiera \(\bm{b}\) vector de dimensión \(d\) y \(\bm{A}\) matriz
\(d\times d\) no singular,

\[\bm{x}\ \ \mbox{tiene distribución normal} \Leftrightarrow\bm{A}\bm{x}+\bm{b}\ \ \mbox{tiene distribución normal}\]

Por tanto, si estamos contrastando que el vector \(\bm{x}\) tenga
distribución normal, tanta evidencia habrá a favor o en contra de esta
hipótesis, si empleamos los datos de \(\bm{x}\) como si empleamos los
datos transformados \(\bm{A}\bm{x}+\bm{b}\).

Diremos que un test de normalidad multivariante es \textbf{invariante} si
proporciona los mismos resultados con los datos originales que con datos
transformados por traslaciones o transformaciones lineales.

Nos parece que la invariancia es una propiedad deseable para un test de
normalidad. Ahora, si somos libres de transformar los datos, se nos
ocurre aplicar la estandarización multivariante, pues es una
transformación que los simplifica al suprimirles la media y la matriz de
covarianzas. Pasamos entonces de los datos originales \(\bm{x_1},\ldots,\bm{x_n}\) a los datos estandarizados:

\[\bm{z_i}=\bm{S^{-1/2}}\left(\bm{x_i}-\overline{\bm{x}}\right)\] siendo
\(\overline{\bm{x}}\) la media muestral y \(\bm{S}\) la matriz de covarianzas
muestral. La idea es que si los datos originales
\(\bm{x_1},\ldots,\bm{x_n}\) provienen de una distribución normal, los
datos estandarizados \(\bm{z_1},\ldots,\bm{z_n}\) deberían presentar una
distribución normal estándar, esto es, deberían mostrar un aspecto
esférico, sin especial predilección por unas direcciones frente a otras,
y con mayor concentración en la zona central.

Las distancias de Mahalanobis de las observaciones al vector de medias
se pueden expresar así:

\[r_{ii}=\left(\bm{x_i}-\overline{\bm{x}}\right)^\prime{}\bm{S}^{-1}\left(\bm{x_i}-\overline{\bm{x}}\right)=\bm{z_i}^\prime{}\bm{z_i}=\|\bm{z_i}\|^2.\]

Consideramos ahora los valores

\[r_{ij}=\left(\bm{x_i}-\overline{\bm{x}}\right)^\prime{}\bm{S}^{-1}\left(\bm{x_j}-\overline{\bm{x}}\right)=\bm{z_i}^\prime{}\bm{z_j}.\]
Además, como \(\bm{z_i}^\prime{}\bm{z_j}\) representa un producto escalar
de vectores, se tiene también

\[r_{ij}=\bm{z_i}^\prime{}\bm{z_j}=\|\bm{z_i}\|\cdot\|\bm{z_j}\|\cdot \cos(\theta_{ij})\]
siendo \(\theta_{ij}\) el ángulo que forman las observaciones
estandarizadas \(\bm{z_i}\) y \(\bm{z_j}\) en el espacio usual. El ángulo
\(\theta_{ij}\) se conoce como \textbf{ángulo de Mahalanobis} que forman las
observaciones \(\bm{x_i}\) y \(\bm{x_j}\) en el espacio original. En la
Figura \ref{fig:Tema3Figura5} se representa una muestra de observaciones,
que han sido estandarizadas. Por tanto, las distancias y los ángulos de
Mahalanobis coinciden con las distancias y ángulos ordinarias en esta
representación. El punto central, representado en color azul, es la
posición del vector de medias. En el gráfico se han destacado también
dos puntos, denotados por (a) y (b), respectivamente.

Las longitudes de las flechas, que son sus distancias al vector de
medias, serían las raíces cuadradas de sus valores \(r_{ii}\) y \(r_{jj}\).
El valor \(r_{ij}\) para este par de puntos sería el producto de las
longitudes, multiplicado por el coseno del ángulo que forman las dos
flechas. Así, si el ángulo es de 90 grados, \(r_{ij}\) vale cero, si es
menor de 90 grados \(r_{ij}\) es positivo, y si es mayor de 90 grados es
negativo. De este modo, \(r_{ij}\) refleja si los dos puntos (a) y (b)
están del mismo lado o si se encuentran en lados opuestos. Esto será
utilizado en la sección siguiente para definir una medida de asimetría.

\begin{figure}

{\centering \resizebox{.4\linewidth}{!}{\includegraphics[width=0.5\linewidth]{bookdown_AM_files/figure-latex/Tema3Figura5-1} } 

}

\caption{Representación de los valores $r_{ij}$.}\label{fig:Tema3Figura5}
\end{figure}

\subsubsection{Tests basados en medidas de asimetría y kurtosis}\label{tests-basados-en-medidas-de-asimetria-y-kurtosis}

Existen diversas formas de definir medidas de asimetría y kurtosis con
datos multivariantes. Consideraremos aquí las introducidas por \citet{mardia70}.

\paragraph{Coeficiente de asimetría multivariante}\label{coeficiente-de-asimetria-multivariante}

Así, calcularemos la asimetría de la siguiente manera:

\[A_m=\frac{1}{n^2}\sum_{i=1}^n\sum_{j=1}^n r_{ij}^3\]

En la Figura \ref{fig:Tema3Figura5} se ofrece una representación que resulta muy
útil para interpretar los valores \(r_{ij}\). Como vimos en la sección
anterior, \(r_{ij}\) es positivo si los dos puntos en cuestión, (a) y (b)
en el gráfico, forman un ángulo inferior a 90 grados. Por el contrario,
\(r_{ij}\) será negativo si el ángulo es superior a 90 grados. Siempre se
cumple que

\[\sum_{i=1}^n r_{ij}=0=\sum_{j=1}^n r_{ij}\] esto es, se
cancelan los valores positivos con los negativos, lo cual es muy
comprensible, pues el vector de media ocupa la posición central. Esto
también se cumplía en el caso univariante.

Sin embargo, al elevar al cubo los \(r_{ij}\), se desequilibra este
balance. Para mantener el equilibrio, por cada punto, como el (b),
tendría que haber otro punto, que se representa como -(b) en el gráfico,
situado en la posición contraria a (b) respecto del vector de medias.
Los puntos (b) y -(b) tendrían los mismos valores \(r_{ij}\) con el resto
de la muestra, pero con signos cambiados. Al efectuar cualquier potencia
impar, como el cubo, se cancelarían. Así, si cada punto de la muestra
tuviera otro punto simétrico en la propia muestra, el coeficiente de
asimetría valdría cero. En otro caso, si la simetría se cumple
aceptablemente, el coeficiente sería sólo algo mayor que cero, mientras
que si hay un comportamiento claramente distinto en lados opuestos, el
coeficiente de asimetría sería más grande.

En la Figura \ref{fig:asim} se presentan ejemplos de muestras bivariantes generadas a partir de
diferentes distribuciones. Los puntos en negro representan los datos
muestrales y en rojo se representa el vector de medias muestral y la
elipse asociada a la matriz de covarianzas muestral.

Los dos primeros gráficos representan datos simulados de distintas
distribuciones normales. Como se esperaba, en ambos casos el valor del
coeficiente de asimetría muestral \(A_m\) se aproxima a cero ya que, a
nivel teórico, la distribución normal multivariante es simétrica.

A continuación se representan datos bivariantes con marginales
exponenciales e independientes. La asimetría del modelo se refleja en el
valor del coeficiente de asimetría muestral. En el cuarto gráfico se han
simulados tres grupos de datos normales multivariantes. Esto hace que la
distribución de la muestra completa no sea normal. De hecho, la
presencia de varios grupos es una causa muy común de falta de
normalidad. También es común que produzca un efecto de asimetría, que
permite detectar esa falta de normalidad. Por último, se han simulado
observaciones uniformes en el cuadrado unidad. El modelo teórico es
simétrico y de hecho el valor muestral de \(A_m\) es próximo a cero.

Como test de normalidad, y dado que la distribución normal es simétrica,
rechazaremos la normalidad, por falta de simetría, si el coeficiente
\(A_m\) es demasiado grande, comparado con ciertos valores tabulados
correspondientes a muestras normales. Aproximaremos estos valores por
simulación.

\begin{figure}

{\centering \resizebox{.28\linewidth}{!}{\includegraphics[width=0.33\linewidth]{bookdown_AM_files/figure-latex/asim-1} } \resizebox{.28\linewidth}{!}{\includegraphics[width=0.33\linewidth]{bookdown_AM_files/figure-latex/asim-2} } \resizebox{.28\linewidth}{!}{\includegraphics[width=0.33\linewidth]{bookdown_AM_files/figure-latex/asim-3} } \resizebox{.28\linewidth}{!}{\includegraphics[width=0.33\linewidth]{bookdown_AM_files/figure-latex/asim-4} } \resizebox{.28\linewidth}{!}{\includegraphics[width=0.33\linewidth]{bookdown_AM_files/figure-latex/asim-5} } 

}

\caption{Ejemplos de muestras bivariantes para su contraste de normalidad.}\label{fig:asim}
\end{figure}

\paragraph{Coeficiente de kurtosis multivariante}\label{coeficiente-de-kurtosis-multivariante}

El coeficiente de kurtosis propuesto por \citet{mardia70} adopta la forma
siguiente: \[K_m=\frac{1}{n}\sum_{i=1}^n r_{ii}^2.\]

A la vista de esta definición, el coeficiente de kurtosis está basado
únicamente en los valores \(r_{ii}\). Como \(r_{ii}\) es el cuadrado de la
distancia del dato i-ésimo al vector de medias, los valores \(r_{ii}^2\)
que figuran en la definición de \(K_m\) son las potencias cuartas de estas
distancias.

Por tanto, el coeficiente \(K_m\) generaliza la medida de kurtosis al caso
multivariante como estas potencias cuartas de la distancia a la media.
De esta manera, para la kurtosis no importará la dirección o sentido de
la desviación respecto de la media, sino únicamente la magnitud de dicha
desviación. Se trata, pues, de detectar si los datos se agrupan en torno
a la Media\(\pm\)Desviación típica, que en el caso multivariante sería la
elipse en rojo de la Figura \hyperref[fig:asim]{11}, o si por el contrario se agrupan o muy cerca de
la media o muy lejos de ésta. Si están cerca de la elipse la kurtosis
será pequeña (distribución platicúrtica), y si están o muy cerca de la
media o muy lejos de ella, la kurtosis será grande y diremos que la
distribución es leptocúrtica.

Igual que en el caso univariante, la distribución normal multivariante
presenta un valor intermedio de kurtosis (diríamos que es mesocúrtica).
Se puede demostrar que un vector normal \(d\)-dimensional tiene kurtosis
multivariante igual a \(d(d+2)\). En la Figura
\ref{fig:asim} se observa
que el valor del coeficiente de kurtosis \(K_m\) es próximo a \(d(d+2)=8\)
para las muestras normales. Los datos generados de la mixtura de
normales y de la uniforme presentan valores menores, y serían
platicúrticos. Los datos generados de la exponencial son leptocúrticos.

Respecto del contraste de normalidad, el test basado en la kurtosis
rechazará la normalidad tanto para valores demasiado grandes como para
valores demasiado pequeños de \(K_m\). De nuevo, los valores críticos
están tabulados y los aproximaremos por simulación.

\subsubsection{Extensión multivariantes del test de Shapiro-Wilk}\label{extension-multivariantes-del-test-de-shapiro-wilk}

Existen diversas maneras de extender el test de Shapiro-Wilk. Nos
centraremos en una idea muy sencilla, que surge del concepto de
invariancia. Así, el vector aleatorio es normal multivariante si y sólo
si su estandarización lo es, esto es:

\[\bm{x}\in N_d(\bm{\mu},\bm{\Sigma}) \Longleftrightarrow \bm{z}=\bm{\Sigma^{-1/2}}(\bm{x}-\bm{\mu})\in N_d(\bm{0},\bm{I_d})\]

Lo interesante es que \(\bm{z}\) tiene sus \(d\) componentes (las \(d\)
variables) incorrelacionadas, y entonces \(\bm{z}\) es normal si y sólo si
sus componentes son normales. Ante tal resultado, basta con comprobar si
las \(d\) variables resultantes de la estandarización multivariante son
normales. Volvemos entonces a los datos estandarizados:

\[\bm{z_i}=\bm{S^{-1/2}}\left(\bm{x_i}-\overline{\bm{x}}\right)\] siendo
\(\overline{\bm{x}}\) la media muestral y \(\bm{S}\) la matriz de covarianzas
muestral.

Si \(W_1,\ldots,W_d\) son los estadísticos de Shapiro-Wilk de cada
componente de \(\bm{z_1},\ldots,\bm{z_n}\), entonces podemos considerar el
estadístico de Shapiro-Wilk multivariante:

\[W_{VG}=\frac{1}{d}\sum_{j=1}^d W_j\] Este estadístico ha sido
propuesto por Villasenor Alva y González Estrada (2009). Se rechazará la
normalidad cuando \(W_{VG}\) sea pequeño, en comparación con ciertos
valores tabulados, que de nuevo podemos aproximar por simulación.

\chapter{Análisis Multivariante de la Varianza}\label{analisis-multivariante-de-la-varianza}

En este capítulo vemos la extensión del análisis de la varianza univariante
(ANOVA) al caso de que la variable respuesta es multivariante (MANOVA).
En primer lugar repasaremos los aspectos fundamentales del caso
unidimensional. Para el caso multidimensional veremos que, en lugar del
test \(F\), surgen los test de la \(\Lambda\) de Wilks y sus alternativas.

\section{Revisión del modelo de análisis de la varianza univariante ANOVA}\label{revision-del-modelo-de-analisis-de-la-varianza-univariante-anova}

En el modelo de análisis de la varianza univariante, disponemos de \(k\)
muestras independientes que provienen de poblaciones normales con la
misma varianza (modelo homocedástico). Es decir,

\[\begin{array}{cccccc}
y_{11} & y_{12} & \cdots & y_{1\,n_1} & \mbox{de una población} & N\left(\mu_1,\sigma^2\right), \\
y_{21} & y_{22} & \cdots & y_{2\,n_2} & \mbox{de una población} & N\left(\mu_2,\sigma^2\right), \\
\cdots & \cdots & \cdots & \cdots & \cdots & \cdots \\
y_{k1} & y_{k2} & \cdots & y_{k\,n_k} &\mbox{de una población} & N\left(\mu_k,\sigma^2\right).
\end{array}\]

Cada una de las \(k\) muestras (también denominadas grupos) está formada
por variables independientes y con la misma distribución. Se supone
además que las \(k\) muestras son independientes entre sí. Los grupos
pueden corresponder, por ejemplo, a diferentes tratamientos aplicados
por el investigador en un experimento. El primer objetivo de un modelo
de análisis de la varianza es contrastar la hipótesis nula de que todas
las medias son iguales. Planteamos así el contraste:

\[H_0: \mu_1=\mu_2=\cdots=\mu_k.\]

\textbf{\includegraphics[width=1em,height=1em]{bookdown_AM_files/figure-latex/fa-icon-00ef1bb1d064b635425d49f06b331298.pdf} ~\textcolor{excolor}{Ejemplo.}} En el trabajo clásico de \citet{andrews1985collection}, se
comparan manzanos sobre seis tipos de portainjertos. En concreto, se
analizan ocho manzanos sobre cada tipo de portainjerto. Para cada
manzano se determina la circunferencia del tronco (mm \(\times\) 100) a
los 4 años. Los datos se encuentran en el fichero \texttt{tronco.txt}.

\footnotesize

\begin{Shaded}
\begin{Highlighting}[]
\SpecialCharTok{\textgreater{}}\NormalTok{ tronco }\OtherTok{\textless{}{-}} \FunctionTok{read.table}\NormalTok{(}\StringTok{"data/tronco.txt"}\NormalTok{, }\AttributeTok{header =} \ConstantTok{TRUE}\NormalTok{)}
\SpecialCharTok{\textgreater{}} \FunctionTok{head}\NormalTok{(tronco)}
\NormalTok{  portainjerto cir4}
\DecValTok{1}            \DecValTok{1} \FloatTok{1.11}
\DecValTok{2}            \DecValTok{1} \FloatTok{1.19}
\DecValTok{3}            \DecValTok{1} \FloatTok{1.09}
\DecValTok{4}            \DecValTok{1} \FloatTok{1.25}
\DecValTok{5}            \DecValTok{1} \FloatTok{1.11}
\DecValTok{6}            \DecValTok{1} \FloatTok{1.08}
\SpecialCharTok{\textgreater{}} \FunctionTok{table}\NormalTok{(tronco[, }\DecValTok{1}\NormalTok{])}

\DecValTok{1} \DecValTok{2} \DecValTok{3} \DecValTok{4} \DecValTok{5} \DecValTok{6} 
\DecValTok{8} \DecValTok{8} \DecValTok{8} \DecValTok{8} \DecValTok{8} \DecValTok{8} 
\end{Highlighting}
\end{Shaded}

\normalsize

En este caso \(k=6\) (tipos de portainjertos) y \(n_i=8\) para
\(i=1,\ldots,6\) (número de árboles para cada tipo de portainjerto). El
objetivo es determinar si existen diferencias significativas en el
crecimiento de los manzanos debidas al tipo de portainjerto, ver Figura
\ref{fig:boxm}

\begin{figure}

{\centering \resizebox{.4\linewidth}{!}{\includegraphics[width=0.5\linewidth]{bookdown_AM_files/figure-latex/boxm-1} } 

}

\caption{Boxplot para la circunferencia de tronco a los 4 años en cada tipo de portainjerto.}\label{fig:boxm}
\end{figure}

Denotamos mediante \(\overline{y}_{i\bullet}\) a la media muestral
procedente de la población \(i\)-ésima, es decir,

\[\overline{y}_{i\bullet}=\frac{1}{n_i}\sum_{j=1}^{n_i} y_{ij}\] para
\(i=1,\ldots,k\). Sea \(n=n_1+\ldots+n_k\) el número total de datos.
Entonces, la media global de la muestra se calcula como:

\[\overline{y}_{\bullet\bullet}
=\frac{1}{n}\sum_{i=1}^k \sum_{j=1}^{n_i} y_{ij}.\]

Se puede expresar el modelo de análisis multivariante de la varianza
como un caso particular del modelo de regresión lineal. El modelo para
cada observación es:

\[y_{ij}=\mu+\alpha_i+\epsilon_{ij}=\mu_i+\epsilon_{ij}\] con
\(i=1,\ldots,k\) y \(j=1,\ldots, n_i\). En el modelo, \(\mu_i=\mu+\alpha_i\)
es la media de la población \(i\)-ésima.

Queremos comparar las medias muestrales \(\overline{y}_{i\bullet}\) para
ver si son suficientemente distintas como para llevarnos a pensar que
las medias de cada población difieren. Para determinar si hay
diferencias significativas entre las respuestas medias a distintos
niveles del factor, el ANOVA descompone la variabilidad total. La
variabilidad total, medida a través de la desviación cuadrática de los
datos a la media global, admite la siguiente descomposición:

\[\sum_{i=1}^k \sum_{j=1}^{n_i} \left( y_{ij} - \overline y_{\bullet\bullet}\right)^2
=\sum_{i=1}^k \sum_{j=1}^{n_i} \left( y_{ij} - \overline{y}_{i\bullet} \right)^2
 + \sum_{i=1}^k \sum_{j=1}^{n_i} \left( \overline{y}_{i\bullet} - \overline y_{\bullet\bullet} \right)^2.\]

Al igual que en regresión, lo representamos mediante la \textbf{tabla de
análisis de la varianza}.

{\def\LTcaptype{none} % do not increment counter
\begin{longtable}[]{@{}
  >{\centering\arraybackslash}p{(\linewidth - 4\tabcolsep) * \real{0.2941}}
  >{\centering\arraybackslash}p{(\linewidth - 4\tabcolsep) * \real{0.4118}}
  >{\centering\arraybackslash}p{(\linewidth - 4\tabcolsep) * \real{0.2941}}@{}}
\toprule\noalign{}
\begin{minipage}[b]{\linewidth}\centering
Fuente de variación
\end{minipage} & \begin{minipage}[b]{\linewidth}\centering
Suma de cuadrados
\end{minipage} & \begin{minipage}[b]{\linewidth}\centering
Grados de libertad
\end{minipage} \\
\midrule\noalign{}
\endhead
\bottomrule\noalign{}
\endlastfoot
Entre poblaciones (VE) & \({\sum_{i=1}^k {n_i} \left(\bar y_{i\bullet}-\bar y_{\bullet\bullet}\right)^2}\) & \(k-1\) \\
Error (VNE) & \({\sum_{i=1}^k \sum_{j=1}^{n_i} \left( y_{ij}-\bar y_{i\bullet}\right)^2}\) & \(n-k\) \\
Total (VT) & \({\sum_{i=1}^k \sum_{j=1}^{n_i} \left( y_{ij}-\bar y_{\bullet\bullet}\right)^2}\) & \(n-1\) \\
\end{longtable}
}

Nos fijamos en que VE representa las desviaciones de las medias
muestrales de cada población respecto a la media global. Sirve por lo
tanto como medición de la variabilidad entre poblaciones (inter-grupos).
Además, VNE representa las desviaciones de cada dato respecto a la media
muestral de la población de la que procede. Es útil como medida de la
variabilidad interna (intra-grupos), presente entre los individuos de la
misma población. Para efectuar el contraste de igualdad de las medias
debemos considerar un estadístico que mida la discrepancia respecto a la
hipótesis nula de igualdad. Si las medias fueran iguales, las
desviaciones entre poblaciones no deberían ser muy grandes, comparadas
con las desviaciones dentro de cada población. Un estadístico razonable
para el contraste es

\[F=\frac{VE/(k-1)}{{VNE}/(n-k)}=\frac{\sum_{i=1}^k \sum_{j=1}^{n_i} \left(\bar y_{i\bullet}-\bar y_{\bullet\bullet}\right)^2/(k-1)}
{\sum_{i=1}^k \sum_{j=1}^{n_i} \left( y_{ij}-\bar y_{i\bullet}\right)^2/(n-k)}.\]
Si la hipótesis nula \(H_0: \mu_1=\mu_2=\cdots=\mu_k\) es cierta, \(F\)
presenta una distribución \(F\) de Snédecor \(F\in F_{(k-1),(n-k)}\).
Rechazaremos \(H_0\) cuando para valores grandes de \(F\).

\textbf{\includegraphics[width=1em,height=1em]{bookdown_AM_files/figure-latex/fa-icon-00ef1bb1d064b635425d49f06b331298.pdf} ~\textcolor{excolor}{Ejemplo.}} Volvemos sobre el ejemplo de manzanos sobre diferentes
tipos de portainjertos. Existen distintas formas de realizar un ANOVA en
R. Esta es una de ellas:

\footnotesize

\begin{Shaded}
\begin{Highlighting}[]
\SpecialCharTok{\textgreater{}}\NormalTok{ x }\OtherTok{\textless{}{-}} \FunctionTok{factor}\NormalTok{(tronco[, }\DecValTok{1}\NormalTok{])}
\SpecialCharTok{\textgreater{}}\NormalTok{ y }\OtherTok{\textless{}{-}}\NormalTok{ tronco[, }\DecValTok{2}\NormalTok{]}
\SpecialCharTok{\textgreater{}}\NormalTok{ z }\OtherTok{\textless{}{-}} \FunctionTok{lm}\NormalTok{(y }\SpecialCharTok{\textasciitilde{}}\NormalTok{ x)}
\SpecialCharTok{\textgreater{}} \FunctionTok{anova}\NormalTok{(z)}
\NormalTok{Analysis of Variance Table}

\NormalTok{Response}\SpecialCharTok{:}\NormalTok{ y}
\NormalTok{          Df  Sum Sq   Mean Sq F value }\FunctionTok{Pr}\NormalTok{(}\SpecialCharTok{\textgreater{}}\NormalTok{F)}
\NormalTok{x          }\DecValTok{5} \FloatTok{0.07356} \FloatTok{0.0147121}   \FloatTok{1.931} \FloatTok{0.1094}
\NormalTok{Residuals }\DecValTok{42} \FloatTok{0.31999} \FloatTok{0.0076187}               
\end{Highlighting}
\end{Shaded}

\normalsize

Vemos que el \(p\)-valor indica que no hay evidencias significativas para
rechazar la hipótesis nula de igualdad de medias.

\section{MANOVA con un factor de variación}\label{manova-con-un-factor-de-variacion}

En general, cuando disponemos de una muestra, no medimos una única
variable. Lo más habitual es realizar mediciones de diversas variables
dependientes sobre cada unidad muestral.

\textbf{\includegraphics[width=1em,height=1em]{bookdown_AM_files/figure-latex/fa-icon-00ef1bb1d064b635425d49f06b331298.pdf} ~\textcolor{excolor}{Ejemplo.}} En el trabajo clásico de \citet{andrews1985collection}, se
comparan manzanos sobre seis tipos de portainjertos. En concreto, se
analizan ocho manzanos sobre cada tipo de portainjerto. Para cada
manzano se determina \(\bm{y}=(y_1,\ldots,y_4)^\prime\), siendo \(y_1\) la
circunferencia del tronco (mm \(\times\) 100) a los 4 años, \(y_2\) el
crecimiento (m) a los 4 años, \(y_3\) la circunferencia del tronco (mm
\(\times\) 100) a los 15 años e \(y_4\) el peso del árbol (lb \(\times\) 1000)
a los 15 años. Los datos se encuentran en el fichero \texttt{manzanos.txt}.

\footnotesize

\begin{Shaded}
\begin{Highlighting}[]
\SpecialCharTok{\textgreater{}}\NormalTok{ manzanos }\OtherTok{\textless{}{-}} \FunctionTok{read.table}\NormalTok{(}\StringTok{"data/manzanos.txt"}\NormalTok{, }\AttributeTok{header =} \ConstantTok{TRUE}\NormalTok{)}
\SpecialCharTok{\textgreater{}} \FunctionTok{head}\NormalTok{(manzanos)}
\NormalTok{  portainjerto cir4 crec4 cir15 peso15}
\DecValTok{1}            \DecValTok{1} \FloatTok{1.11} \FloatTok{2.569}  \FloatTok{3.58}  \FloatTok{0.760}
\DecValTok{2}            \DecValTok{1} \FloatTok{1.19} \FloatTok{2.928}  \FloatTok{3.75}  \FloatTok{0.821}
\DecValTok{3}            \DecValTok{1} \FloatTok{1.09} \FloatTok{2.865}  \FloatTok{3.93}  \FloatTok{0.928}
\DecValTok{4}            \DecValTok{1} \FloatTok{1.25} \FloatTok{3.844}  \FloatTok{3.94}  \FloatTok{1.009}
\DecValTok{5}            \DecValTok{1} \FloatTok{1.11} \FloatTok{3.027}  \FloatTok{3.60}  \FloatTok{0.766}
\DecValTok{6}            \DecValTok{1} \FloatTok{1.08} \FloatTok{2.336}  \FloatTok{3.51}  \FloatTok{0.726}
\end{Highlighting}
\end{Shaded}

\normalsize

Se trata de comparar las medias de \(k\) poblaciones normales
multivariantes independientes y con matriz de dispersión común.
Consideremos entonces \(k\) muestras independientes

\[\begin{array}{cccccc}
\bm{y_{11}} & \bm{y_{12}} & \cdots & \bm{y_{1n_1}} & \mbox{de una población} & N_d(\bm{\mu_1},\bm\Sigma), \\
\bm{y_{21}} & \bm{y_{22}} & \cdots & \bm{y_{2n_2}} & \mbox{de una población} & N_d(\bm{\mu_2},\bm\Sigma), \\
\cdots & \cdots & \cdots & \cdots & \cdots & \cdots \\
\bm{y_{k1}} & \bm{y_{k2}} & \cdots & \bm{y_{kn_k}} &\mbox{de una población} & N_d(\bm{\mu_k},\bm\Sigma).
\end{array}\]

Cada una de las \(k\) muestras está formada por vectores independientes y
con la misma distribución. Se trata, por tanto, de \(k\) muestras
aleatorias simples. Además se supone que las \(k\) muestras son, entre sí,
independientes. La hipótesis de homocedasticidad consiste ahora en que
la matriz de covarianzas se supone que es la misma en todas las
poblaciones. Recuerda que la homocedasticidad se puede contrastar
mediante un test de igualdad de las matrices de covarianzas.

\subsection{Descomposición de la variabilidad. Contraste de igualdad de medias}\label{descomposicion-de-la-variabilidad-contraste-de-igualdad-de-medias}

Igual que ocurría en el caso univariante, se puede expresar el modelo de
análisis multivariante de la varianza como un caso particular del modelo
lineal múltiple (la variable respuesta es un vector \(d\)-dimensional).
Para cada observación se tiene,

\[\bm{y_{ij}}=\bm{\mu}+\bm{\alpha_i}+\bm{\epsilon_{ij}}=\bm{\mu_i}+\bm{\epsilon_{ij}}\]
con \(i=1,\ldots,k\) y \(j=1,\ldots, n_i\). En el modelo,
\(\bm{\mu_i}=\bm{\mu}+\bm{\alpha_i}\) es el vector de medias de la
población \(i\)-ésima.

El objetivo del modelo de análisis de la varianza multivariante es
contrastar la hipótesis nula de que los vectores de medias son iguales.
Planteamos así el contraste:

\[H_0: \bm{\mu_1}=\bm{\mu_2}=\cdots=\bm{\mu_k}.\]

Es decir,

\[H_0: 
\left(\begin{array}{c}\mu_{11}\\ \mu_{12}\\\vdots\\ \mu_{1d}\end{array}\right)=\left(\begin{array}{c}\mu_{21}\\ \mu_{22}\\\vdots\\ \mu_{2d}\end{array}\right)=\cdots=\left(\begin{array}{c}\mu_{k1}\\ \mu_{k2}\\\vdots\\ \mu_{kd}\end{array}\right).\]

Al igual que en el caso univariante, podremos calcular la media global
de la muestra y las medias muestrales por grupo. Así denotamos mediante
\(\bm{\overline{y}_{i\bullet}}\) a la media muestral procedente de la
población \(i\)-ésima, es decir,

\[\bm{\overline{y}_{i\bullet}}=\frac{1}{n_i}\sum_{j=1}^{n_i} \bm{y_{ij}}\]
para \(i=1,\ldots,k\). Sea \(n=n_1+\ldots+n_k\) el número total de datos.
Entonces, la media global de la muestra se calcula como:
\[\bm{\overline{y}_{\bullet\bullet}}
=\frac{1}{n}\sum_{i=1}^k \sum_{j=1}^{n_i} \bm{y_{ij}}\]

Para determinar si hay diferencias significativas entre las respuestas
medias a distintos niveles del factor, el MANOVA descompone la
variabilidad total de forma similar a como se hacía en el ANOVA. La
diferencia está en que ahora la variabilidad total se medirá mediante la
siguiente matriz \(d\times d\):

\[\sum_{i=1}^k \sum_{j=1}^{n_i}
\left( \bm{y_{ij}}-\bm{\overline{y}_{\bullet\bullet}}\right)
\left( \bm{y_{ij}}-\bm{\overline{y}_{\bullet\bullet}}\right)^\prime{}.\]
Por analogía con el caso univariante, en el que se comparaban la
variabilidad inter-grupos con la variabilidad intra-grupos, se calculan
ahora las matrices

\[\bm{H}=\sum_{i=1}^kn_i(\bm{\overline{y}_{i\bullet}}-\bm{\overline{y}_{\bullet\bullet}})(\bm{\overline{y}_{i\bullet}}-\bm{\overline{y}_{\bullet\bullet}})^\prime{}\]
y

\[\bm{E} = \sum_{i=1}^k \sum_{j=1}^{n_i}
\left( \bm{y_{ij}}- \bm{\overline{y}_{i\bullet}}\right)
\left( \bm{y_{ij}}- \bm{\overline{y}_{i\bullet}}\right)^\prime{}.\]

Se muestra a continuación la \textbf{tabla del análisis multivariante de la
varianza} o tabla MANOVA:

{\def\LTcaptype{none} % do not increment counter
\begin{longtable}[]{@{}
  >{\centering\arraybackslash}p{(\linewidth - 4\tabcolsep) * \real{0.2667}}
  >{\centering\arraybackslash}p{(\linewidth - 4\tabcolsep) * \real{0.4667}}
  >{\centering\arraybackslash}p{(\linewidth - 4\tabcolsep) * \real{0.2667}}@{}}
\toprule\noalign{}
\begin{minipage}[b]{\linewidth}\centering
Fuente de variación
\end{minipage} & \begin{minipage}[b]{\linewidth}\centering
Matriz de covarianzas
\end{minipage} & \begin{minipage}[b]{\linewidth}\centering
Grados de libertad
\end{minipage} \\
\midrule\noalign{}
\endhead
\bottomrule\noalign{}
\endlastfoot
Entre poblaciones & \({ \bm{H}=\sum_{i=1}^kn_i(\bm{\overline{y}_{i\bullet}}-\bm{\overline{y}_{\bullet\bullet}})(\bm{\overline{y}_{i\bullet}}-\bm{\overline{y}_{\bullet\bullet}})^\prime{}}\) & \(k-1\) \\
Error & \({\bm{E} = \sum_{i=1}^k \sum_{j=1}^{n_i}\left( \bm{y_{ij}}- \bm{\overline{y}_{i\bullet}}\right)\left( \bm{y_{ij}}- \bm{\overline{y}_{i\bullet}}\right)^\prime{}}\) & \({\sum_{i=1}^k \left(n_i-1\right)}\) \\
Total & \({\bm{E_H} = \sum_{i=1}^k \sum_{j=1}^{n_i}\left( \bm{y_{ij}}-\bm{\overline{y}_{\bullet\bullet}}\right)\left( \bm{y_{ij}}-\bm{\overline{y}_{\bullet\bullet}}\right)^\prime{}}\) & \({\sum_{i=1}^k n_i-1}\) \\
\end{longtable}
}

Observamos que esta tabla no es más que una extensión de la tabla ANOVA
al caso multivariante. El sentido común nos invita a rechazar la
hipótesis nula cuando la variabilidad proveniente de las diferencias
entre poblaciones (que medimos mediante la matriz \(\bm{H}\)) sea grande
comparada con la proveniente del error (medida por \(\bm{E}\)).

\subsection{\texorpdfstring{El estadístico \(\Lambda\) de Wilks}{El estadístico \textbackslash Lambda de Wilks}}\label{el-estadistico-lambda-de-wilks}

El contraste de razón de verosimiltudes para el contraste
\(H_0: \bm{\mu_1}=\bm{\mu_2}=\cdots=\bm{\mu_k}\) viene dado por

\[\Lambda=\frac{|\bm{E}|}{|\bm{E}+\bm{H}|}.\]

Este estadístico, conocido
como estadístico de Wilks, compara la matriz \(\bm{E}\), que mide la
variabilidad intra-grupos, con la matriz \(\bm{E_H}=\bm{E}+\bm{H}\), que
mide la variabilidad total\footnote{La razón para no comparar directamente \(\bm{H}\) con \(\bm{E}\), sino
  con \(\bm{E}+\bm{H}\) es que \(\bm{H}\) podría tener algún autovalor
  cero y eso haría que su determinante fuera cero.}. Se puede demostrar que, bajo la
hipótesis nula,

\[\begin{aligned}
\bm{H} &\in & \mbox{Wishart}_d(\bm\Sigma,k-1) \\
\bm{E} &\in & \mbox{Wishart}_d(\bm\Sigma,n-k)\end{aligned}\] y además
son independientes. Entonces tenemos que\footnote{Ver Anexo \ref{anexoDist} de distribuciones multidimensionales para la
  definición formal de la distribución \(\Lambda\) de Wilks}

\[\frac{|\bm{E}|}{|\bm{E}+\bm{H}|}\in \Lambda(d,k-1,n-k).\] Por tanto,
rechazaremos la hipótesis nula cuando el estadístico
\(|\bm{E}|/|\bm{E}+\bm{H}|\) tome un valor menor que el cuantil \(\alpha\)
de la distribución \(\Lambda(d,k-1,n-k)\), siendo \(\alpha\) el nivel de
significación fijado de antemano. El estadístico se puede aproximar por
una distribución \(F\) de Snédecor.

\textbf{\includegraphics[width=1em,height=1em]{bookdown_AM_files/figure-latex/fa-icon-00ef1bb1d064b635425d49f06b331298.pdf} ~\textcolor{excolor}{Ejemplo.}} Volvemos al ejemplo sobre manzanos. Se trata de contrastar
\(H_0: \bm{\mu_1}=\bm{\mu_2}=\cdots=\bm{\mu_6}\), donde para cada \(i\) el
vector \(\bm{\mu_i}=(\mu_{i1},\ldots,\mu{i4})^\prime\) representa el
vector de medias de \(\bm{y}=(y_1,\ldots,y_4)^\prime{}\) para el tipo \(i\)
de portainjerto. Los datos son:

\footnotesize

\begin{Shaded}
\begin{Highlighting}[]
\SpecialCharTok{\textgreater{}}\NormalTok{ x }\OtherTok{\textless{}{-}} \FunctionTok{factor}\NormalTok{(manzanos[, }\DecValTok{1}\NormalTok{])}
\SpecialCharTok{\textgreater{}}\NormalTok{ y }\OtherTok{\textless{}{-}} \FunctionTok{as.matrix}\NormalTok{(manzanos[, }\DecValTok{2}\SpecialCharTok{:}\DecValTok{5}\NormalTok{])}
\SpecialCharTok{\textgreater{}}\NormalTok{ n }\OtherTok{\textless{}{-}} \FunctionTok{dim}\NormalTok{(y)[}\DecValTok{1}\NormalTok{]}
\SpecialCharTok{\textgreater{}}\NormalTok{ d }\OtherTok{\textless{}{-}} \FunctionTok{dim}\NormalTok{(y)[}\DecValTok{2}\NormalTok{]}
\SpecialCharTok{\textgreater{}}\NormalTok{ n}
\NormalTok{[}\DecValTok{1}\NormalTok{] }\DecValTok{48}
\SpecialCharTok{\textgreater{}}\NormalTok{ d}
\NormalTok{[}\DecValTok{1}\NormalTok{] }\DecValTok{4}
\end{Highlighting}
\end{Shaded}

\normalsize

Calculamos a continuación la media global de la muestra
\(\bm{\overline{y}_{\bullet\bullet}}\) y las medias
\(\bm{\overline{y}_{i\bullet}}\) correspondientes a cada grupo:

\footnotesize

\begin{Shaded}
\begin{Highlighting}[]
\SpecialCharTok{\textgreater{}}\NormalTok{ medy }\OtherTok{\textless{}{-}} \FunctionTok{apply}\NormalTok{(y, }\DecValTok{2}\NormalTok{, mean)}
\SpecialCharTok{\textgreater{}}\NormalTok{ medy}
\NormalTok{    cir4    crec4    cir15   peso15 }
\FloatTok{1.102708} \FloatTok{2.758854} \FloatTok{4.087292} \FloatTok{1.083000} 
\SpecialCharTok{\textgreater{}}\NormalTok{ medyi }\OtherTok{\textless{}{-}} \FunctionTok{t}\NormalTok{(}\FunctionTok{apply}\NormalTok{(y, }\DecValTok{2}\NormalTok{, tapply, x, mean))}
\SpecialCharTok{\textgreater{}}\NormalTok{ medyi}
              \DecValTok{1}        \DecValTok{2}        \DecValTok{3}       \DecValTok{4}       \DecValTok{5}        \DecValTok{6}
\NormalTok{cir4   }\FloatTok{1.137500} \FloatTok{1.157500} \FloatTok{1.107500} \FloatTok{1.09750} \FloatTok{1.08000} \FloatTok{1.036250}
\NormalTok{crec4  }\FloatTok{2.977125} \FloatTok{3.109125} \FloatTok{2.815250} \FloatTok{2.87975} \FloatTok{2.55725} \FloatTok{2.214625}
\NormalTok{cir15  }\FloatTok{3.738750} \FloatTok{4.515000} \FloatTok{4.455000} \FloatTok{3.90625} \FloatTok{4.31250} \FloatTok{3.596250}
\NormalTok{peso15 }\FloatTok{0.871125} \FloatTok{1.280500} \FloatTok{1.391375} \FloatTok{1.03900} \FloatTok{1.18100} \FloatTok{0.735000}
\end{Highlighting}
\end{Shaded}

\normalsize

A continuación, realizamos el MANOVA:

\footnotesize

\begin{Shaded}
\begin{Highlighting}[]
\SpecialCharTok{\textgreater{}}\NormalTok{ z }\OtherTok{\textless{}{-}} \FunctionTok{lm}\NormalTok{(y }\SpecialCharTok{\textasciitilde{}}\NormalTok{ x)}
\SpecialCharTok{\textgreater{}} \FunctionTok{summary}\NormalTok{(}\FunctionTok{manova}\NormalTok{(z), }\AttributeTok{test =} \FunctionTok{c}\NormalTok{(}\StringTok{"Wilks"}\NormalTok{))}
\NormalTok{          Df   Wilks approx F num Df den Df    }\FunctionTok{Pr}\NormalTok{(}\SpecialCharTok{\textgreater{}}\NormalTok{F)    }
\NormalTok{x          }\DecValTok{5} \FloatTok{0.15401}   \FloatTok{4.9369}     \DecValTok{20}  \FloatTok{130.3} \FloatTok{7.714e{-}09} \SpecialCharTok{**}\ErrorTok{*}
\NormalTok{Residuals }\DecValTok{42}                                             
\SpecialCharTok{{-}{-}{-}}
\NormalTok{Signif. codes}\SpecialCharTok{:}  \DecValTok{0} \StringTok{\textquotesingle{}***\textquotesingle{}} \FloatTok{0.001} \StringTok{\textquotesingle{}**\textquotesingle{}} \FloatTok{0.01} \StringTok{\textquotesingle{}*\textquotesingle{}} \FloatTok{0.05} \StringTok{\textquotesingle{}.\textquotesingle{}} \FloatTok{0.1} \StringTok{\textquotesingle{} \textquotesingle{}} \DecValTok{1}
\end{Highlighting}
\end{Shaded}

\normalsize

Comprobamos que el valor del estadístico es \(\Lambda=0.15401\). Para
ello, calculamos en primer lugar la matriz \(\bm{E}\)

\footnotesize

\begin{Shaded}
\begin{Highlighting}[]
\SpecialCharTok{\textgreater{}}\NormalTok{ res }\OtherTok{\textless{}{-}}\NormalTok{ z}\SpecialCharTok{$}\NormalTok{residuals}
\SpecialCharTok{\textgreater{}}\NormalTok{ E }\OtherTok{\textless{}{-}} \FunctionTok{crossprod}\NormalTok{(res)  }\CommentTok{\# Equivalente a (n {-} 1) * cov(res)}
\SpecialCharTok{\textgreater{}}\NormalTok{ E}
\NormalTok{            cir4     crec4     cir15   peso15}
\NormalTok{cir4   }\FloatTok{0.3199875}  \FloatTok{1.696564} \FloatTok{0.5540875} \FloatTok{0.217140}
\NormalTok{crec4  }\FloatTok{1.6965637} \FloatTok{12.142790} \FloatTok{4.3636125} \FloatTok{2.110214}
\NormalTok{cir15  }\FloatTok{0.5540875}  \FloatTok{4.363612} \FloatTok{4.2908125} \FloatTok{2.481656}
\NormalTok{peso15 }\FloatTok{0.2171400}  \FloatTok{2.110214} \FloatTok{2.4816562} \FloatTok{1.722525}
\end{Highlighting}
\end{Shaded}

\normalsize

Calculamos ahora la matriz \(\bm{E_H}=\bm{E}+\bm{H}\) aplicando la
definición.

\footnotesize

\begin{Shaded}
\begin{Highlighting}[]
\SpecialCharTok{\textgreater{}}\NormalTok{ resh }\OtherTok{\textless{}{-}}\NormalTok{ y }\SpecialCharTok{{-}} \FunctionTok{matrix}\NormalTok{(medy, }\AttributeTok{nrow =}\NormalTok{ n, }\AttributeTok{ncol =}\NormalTok{ d, }\AttributeTok{byrow =} \ConstantTok{TRUE}\NormalTok{)}
\SpecialCharTok{\textgreater{}}\NormalTok{ EH }\OtherTok{\textless{}{-}} \FunctionTok{crossprod}\NormalTok{(resh)  }\CommentTok{\# Equivalente a (n {-} 1) * cov(resh)}
\SpecialCharTok{\textgreater{}}\NormalTok{ EH}
\NormalTok{            cir4     crec4      cir15   peso15}
\NormalTok{cir4   }\FloatTok{0.3935479}  \FloatTok{2.233949}  \FloatTok{0.8863521} \FloatTok{0.425610}
\NormalTok{crec4  }\FloatTok{2.2339490} \FloatTok{16.342452}  \FloatTok{6.7190010} \FloatTok{3.747322}
\NormalTok{cir15  }\FloatTok{0.8863521}  \FloatTok{6.719001} \FloatTok{10.4047479} \FloatTok{6.262700}
\NormalTok{peso15 }\FloatTok{0.4256100}  \FloatTok{3.747322}  \FloatTok{6.2627000} \FloatTok{4.215616}
\end{Highlighting}
\end{Shaded}

\normalsize

Efectivamente, se tiene

\footnotesize

\begin{Shaded}
\begin{Highlighting}[]
\SpecialCharTok{\textgreater{}} \FunctionTok{det}\NormalTok{(E)}\SpecialCharTok{/}\FunctionTok{det}\NormalTok{(EH)}
\NormalTok{[}\DecValTok{1}\NormalTok{] }\FloatTok{0.1540077}
\end{Highlighting}
\end{Shaded}

\normalsize

En base a la aproximación a la \(F\) de Snédecor proporcionada por R y al
\(p\)-valor del estadístico, se rechaza la hipótesis nula de igualdad de
los vectores de medias en los 6 tipos de portainjerto.

\subsection{Test de Roy}\label{test-de-roy}

Existen otros procedimientos para contrastar
\(H_0: \bm{\mu_1}=\bm{\mu_2}=\cdots=\bm{\mu_k}\) que dan lugar a
estadísticos de contraste diferentes al \(\Lambda\) de Wilks. Uno de esos
procedimientos alternativos es el de \emph{unión-intersección}.

Un test de unión-intersección para contrastar una hipótesis
\(H_0=\cap_a H_{0a}\) es un test cuya región de rechazo es de la forma
\(R=\cup_a R_a\), siendo \(R_a\) la región de rechazo para \(H_{0a}\).

Esta idea se aplica en problemas de inferencia multivariante, expresando
la hipótesis multivariante como intersección de hipótesis univariantes,
para las cuales ya hay tests disponibles. Siguiendo este procedimiento,
el test de unión-intersección de Roy tiene como estadístico de
contraste:

\[\lambda_1=\mbox{máximo autovalor de}\ \bm{H}\bm{E}^{-1}\] y
rechazaremos la hipótesis nula para valores grandes de \(\lambda_1\). De
nuevo, se utilizarán aproximaciones de este estadístico a la \(F\) de
Snédecor.

\textbf{\includegraphics[width=1em,height=1em]{bookdown_AM_files/figure-latex/fa-icon-00ef1bb1d064b635425d49f06b331298.pdf} ~\textcolor{excolor}{Ejemplo.}} Calculamos, para el ejemplo sobre manzanos, el estadístico
de Roy:

\footnotesize

\begin{Shaded}
\begin{Highlighting}[]
\SpecialCharTok{\textgreater{}} \FunctionTok{summary}\NormalTok{(}\FunctionTok{manova}\NormalTok{(z), }\AttributeTok{test =} \FunctionTok{c}\NormalTok{(}\StringTok{"Roy"}\NormalTok{))}
\NormalTok{          Df    Roy approx F num Df den Df    }\FunctionTok{Pr}\NormalTok{(}\SpecialCharTok{\textgreater{}}\NormalTok{F)    }
\NormalTok{x          }\DecValTok{5} \FloatTok{1.8757}   \FloatTok{15.756}      \DecValTok{5}     \DecValTok{42} \FloatTok{1.002e{-}08} \SpecialCharTok{**}\ErrorTok{*}
\NormalTok{Residuals }\DecValTok{42}                                            
\SpecialCharTok{{-}{-}{-}}
\NormalTok{Signif. codes}\SpecialCharTok{:}  \DecValTok{0} \StringTok{\textquotesingle{}***\textquotesingle{}} \FloatTok{0.001} \StringTok{\textquotesingle{}**\textquotesingle{}} \FloatTok{0.01} \StringTok{\textquotesingle{}*\textquotesingle{}} \FloatTok{0.05} \StringTok{\textquotesingle{}.\textquotesingle{}} \FloatTok{0.1} \StringTok{\textquotesingle{} \textquotesingle{}} \DecValTok{1}
\end{Highlighting}
\end{Shaded}

\normalsize

En base a este estadístico también se rechaza la hipótesis nula de
igualdad de medias. Comprobamos que el valor del estadístico es
\(\lambda_1=1.8757\). Para ello, calculamos los autovalores de
\(\bm{H}\bm{E}^{-1}.\)

\footnotesize

\begin{Shaded}
\begin{Highlighting}[]
\SpecialCharTok{\textgreater{}}\NormalTok{ H }\OtherTok{\textless{}{-}}\NormalTok{ EH }\SpecialCharTok{{-}}\NormalTok{ E}
\SpecialCharTok{\textgreater{}}\NormalTok{ lambda }\OtherTok{\textless{}{-}} \FunctionTok{eigen}\NormalTok{(H }\SpecialCharTok{\%*\%} \FunctionTok{solve}\NormalTok{(E))}\SpecialCharTok{$}\NormalTok{values}
\SpecialCharTok{\textgreater{}}\NormalTok{ lambda1 }\OtherTok{\textless{}{-}} \FunctionTok{max}\NormalTok{(lambda)}
\SpecialCharTok{\textgreater{}}\NormalTok{ lambda1}
\NormalTok{[}\DecValTok{1}\NormalTok{] }\FloatTok{1.875671}
\end{Highlighting}
\end{Shaded}

\normalsize

\subsection{Tests de Pillai y Lawley-Hotelling}\label{tests-de-pillai-y-lawley-hotelling}

Existen otros dos tests adicionales para contrastar
\(H_0: \bm{\mu_1}=\bm{\mu_2}=\cdots=\bm{\mu_k}\) basados en los
autovalores de la matriz \(\bm{H}\bm{E}^{-1}\). El estadístico de Pillai
viene dado por

\[V^{(^d)}=\textnormal{traza}\left[\bm{H}(\bm{E}+\bm{H})^{-1}\right]=\sum_{i=1}^d\theta_i=\sum_{i=1}^d\frac{\lambda_i}{1+\lambda_i}\]
siendo \(\theta_1,\ldots,\theta_d\) los autovalores de
\(\bm{H}(\bm{E}+\bm{H})^{-1}\) y \(\lambda_1,\ldots,\lambda_d\) los
autovalores de \(\bm{H}\bm{E}^{-1}\).

El estadístico de Lawley-Hotelling se define como

\[U^{(^d)}=\textnormal{traza}\left[\bm{H}\bm{E}^{-1}\right]=\sum_{i=1}^d\lambda_i\]
donde \(\lambda_1,\ldots,\lambda_d\) son los autovalores de
\(\bm{H}\bm{E}^{-1}\).

\textbf{\includegraphics[width=1em,height=1em]{bookdown_AM_files/figure-latex/fa-icon-00ef1bb1d064b635425d49f06b331298.pdf} ~\textcolor{excolor}{Ejemplo.}} Calculamos, para el ejemplo sobre manzanos, los
estadísticos de Pillai y Lawley-Hotelling y comprobamos los valores
obtenidos:

\footnotesize

\begin{Shaded}
\begin{Highlighting}[]
\SpecialCharTok{\textgreater{}} \FunctionTok{summary}\NormalTok{(}\FunctionTok{manova}\NormalTok{(z), }\AttributeTok{test =} \FunctionTok{c}\NormalTok{(}\StringTok{"Pillai"}\NormalTok{))}
\NormalTok{          Df Pillai approx F num Df den Df    }\FunctionTok{Pr}\NormalTok{(}\SpecialCharTok{\textgreater{}}\NormalTok{F)    }
\NormalTok{x          }\DecValTok{5} \FloatTok{1.3055}   \FloatTok{4.0697}     \DecValTok{20}    \DecValTok{168} \FloatTok{1.983e{-}07} \SpecialCharTok{**}\ErrorTok{*}
\NormalTok{Residuals }\DecValTok{42}                                            
\SpecialCharTok{{-}{-}{-}}
\NormalTok{Signif. codes}\SpecialCharTok{:}  \DecValTok{0} \StringTok{\textquotesingle{}***\textquotesingle{}} \FloatTok{0.001} \StringTok{\textquotesingle{}**\textquotesingle{}} \FloatTok{0.01} \StringTok{\textquotesingle{}*\textquotesingle{}} \FloatTok{0.05} \StringTok{\textquotesingle{}.\textquotesingle{}} \FloatTok{0.1} \StringTok{\textquotesingle{} \textquotesingle{}} \DecValTok{1}
\SpecialCharTok{\textgreater{}} \FunctionTok{summary}\NormalTok{(}\FunctionTok{manova}\NormalTok{(z), }\AttributeTok{test =} \FunctionTok{c}\NormalTok{(}\StringTok{"Hotelling{-}Lawley"}\NormalTok{))}
\NormalTok{          Df Hotelling}\SpecialCharTok{{-}}\NormalTok{Lawley approx F num Df den Df    }\FunctionTok{Pr}\NormalTok{(}\SpecialCharTok{\textgreater{}}\NormalTok{F)    }
\NormalTok{x          }\DecValTok{5}           \FloatTok{2.9214}   \FloatTok{5.4776}     \DecValTok{20}    \DecValTok{150} \FloatTok{2.568e{-}10} \SpecialCharTok{**}\ErrorTok{*}
\NormalTok{Residuals }\DecValTok{42}                                                      
\SpecialCharTok{{-}{-}{-}}
\NormalTok{Signif. codes}\SpecialCharTok{:}  \DecValTok{0} \StringTok{\textquotesingle{}***\textquotesingle{}} \FloatTok{0.001} \StringTok{\textquotesingle{}**\textquotesingle{}} \FloatTok{0.01} \StringTok{\textquotesingle{}*\textquotesingle{}} \FloatTok{0.05} \StringTok{\textquotesingle{}.\textquotesingle{}} \FloatTok{0.1} \StringTok{\textquotesingle{} \textquotesingle{}} \DecValTok{1}
\end{Highlighting}
\end{Shaded}

\normalsize

En base a ambos estadísticos también se rechaza la hipótesis nula de
igualdad de medias. Comprobamos el valor del estadístico de Pillai y de Lawley-Hotelling

\footnotesize

\begin{Shaded}
\begin{Highlighting}[]
\SpecialCharTok{\textgreater{}} \FunctionTok{sum}\NormalTok{(lambda}\SpecialCharTok{/}\NormalTok{(}\DecValTok{1} \SpecialCharTok{+}\NormalTok{ lambda))  }\CommentTok{\# estadístico de Pillai}
\NormalTok{[}\DecValTok{1}\NormalTok{] }\FloatTok{1.305472}
\SpecialCharTok{\textgreater{}} \FunctionTok{sum}\NormalTok{(lambda)  }\CommentTok{\# estadístico de Lawley{-}Hotelling}
\NormalTok{[}\DecValTok{1}\NormalTok{] }\FloatTok{2.921368}
\end{Highlighting}
\end{Shaded}

\normalsize

\section{MANOVA con dos factores de variación}\label{manova-con-dos-factores-de-variacion}

En esta sección suponemos que hay dos factores: A y B. Del factor A
podemos distinguir \(a\) niveles, mientras que en el factor B podemos
encontrar \(b\) niveles. En cada una de las \(a\cdot b\) posibilidades
realizamos \(n\) observaciones de un vector aleatorio
\(\bm{y}=(y_1,\ldots,y_d)^\prime{}\). El objetivo será estudiar la
influencia de los factores A y B, o de su interacción, en la media del
vector \(\bm{y}\). Así, planteamos el siguiente modelo:

\[\bm{y_{ijk}}=\bm{\mu}+\bm{\alpha_i}+\bm{\beta_j}+\bm{\gamma_{ij}}+{\bm{\epsilon_{ijk}}} \qquad i=1,\ldots,a, \ \ j=1,\ldots,b, \ \ k=1,\ldots,n.\]
Se asume que \({\bm{\epsilon_{ijk}}}\in N_d(\bm{0},\bm\Sigma)\). El
parámetro \(\bm{\mu}\) representa la media global, los parámetros
\(\bm{\alpha_i}\) representan el efecto principal del factor A, los
parámetros \(\bm{\beta_j}\) representan el efecto principal del factor B y
los parámetros \(\bm{\gamma_{ij}}\) representan la interacción de los
factores A y B. Además verifican

\[\sum_{i=1}^a \bm{\alpha_i}=\sum_{j=1}^b \bm{\beta_j}=\sum_{i=1}^a \bm{\gamma_{ij}}
=\sum_{j=1}^b\bm{\gamma_{ij}}=\bm{0}\]

Este modelo también se puede ver como un caso particular del modelo
lineal general multivariante. Las labores de inferencia sobre el modelo
se realizan en base a la siguiente descomposición de la variabilidad,
mediante lo que llamaremos tabla MANOVA II:

{\def\LTcaptype{none} % do not increment counter
\begin{longtable}[]{@{}
  >{\centering\arraybackslash}p{(\linewidth - 4\tabcolsep) * \real{0.1765}}
  >{\centering\arraybackslash}p{(\linewidth - 4\tabcolsep) * \real{0.4706}}
  >{\centering\arraybackslash}p{(\linewidth - 4\tabcolsep) * \real{0.3529}}@{}}
\toprule\noalign{}
\begin{minipage}[b]{\linewidth}\centering
Fuente de variación
\end{minipage} & \begin{minipage}[b]{\linewidth}\centering
Matriz de covarianzas
\end{minipage} & \begin{minipage}[b]{\linewidth}\centering
Grados de libertad
\end{minipage} \\
\midrule\noalign{}
\endhead
\bottomrule\noalign{}
\endlastfoot
Factor A & \({ \bm{H_A} = nb\sum_{i=1}^a\left(\bm{\overline y_{i\bullet\bullet}}-\bm{\overline y_{\bullet\bullet\bullet}}\right)  \left(\bm{\overline y_{i\bullet\bullet}}-\bm{\overline y_{\bullet\bullet\bullet}}\right)^\prime{}}\) & \(a-1\) \\
Factor B & \({ \bm{H_B} = na\sum_{j=1}^b\left(\bm{\overline y_{\bullet j \bullet}}-\bm{\overline y_{\bullet\bullet\bullet}}\right)\left(\bm{\overline y_{\bullet j \bullet}}-\bm{\overline y_{\bullet\bullet\bullet}}\right)^\prime{}}\) & \(b-1\) \\
Interacción & \({ \bm{H_{AB}} = n\sum_{i=1}^a\sum_{j=1}^b\left(\bm{\overline{y}_{ij\bullet}}-\bm{\overline y_{i\bullet\bullet}}-\bm{\overline y_{\bullet j\bullet}}+\bm{\overline y_{\bullet\bullet\bullet}}\right)}\) & \((a-1)(b-1)\) \\
& \({\times \left(\bm{\overline{y}_{ij\bullet}}-\bm{\overline y_{i\bullet\bullet}} -\bm{\overline y_{\bullet j\bullet}}+\bm{\overline y_{\bullet\bullet\bullet}}\right)^\prime{}}\) & \\
Error & \({ \bm{E} = \sum_{i=1}^a \sum_{j=1}^b \sum_{k=1}^n\left( \bm{y_{ijk}}-\bm{\overline y_{ij\bullet}}\right)\left( \bm{y_{ijk}}-\bm{\overline y_{ij\bullet}}\right)^\prime{}}\) & \(ab(n-1)\) \\
Total & \({\sum_{i=1}^a \sum_{j=1}^b \sum_{k=1}^n\left( \bm{y_{ijk}}-\bm{\overline y_{\bullet\bullet\bullet}}\right)\left(  \bm{y_{ijk}}-\bm{\overline y_{\bullet\bullet\bullet}}\right)^\prime{}}\) & \(abn-1\) \\
\end{longtable}
}

La hipótesis relativa a la interacción
\(H_{AB}:\bm{\gamma_{ij}}=\bm{0}\ \forall i,j\) se contrasta en base al
estadístico

\[\frac{|\bm{E}|}{|\bm{E}+\bm{H_{AB}}|}\in \Lambda\left(d,(a-1)(b-1),ab(n-1)\right).\]

Si se obtiene que la interacción no es significativa, se pasaría a un
modelo sin interacción, sobre el cual se pueden cuestionar los efectos
principales. Esto sería el proceder habitual dentro de un planteamiento
jerárquico, cuestión que ya surgía en los modelos univariantes.

\textbf{\includegraphics[width=1em,height=1em]{bookdown_AM_files/figure-latex/fa-icon-00ef1bb1d064b635425d49f06b331298.pdf} ~\textcolor{excolor}{Ejemplo.}} En una especie de césped, denominada Paspalum, se está
investigando el efecto que experimenta al ser infectada con un hongo. Al
mismo tiempo se tiene en cuenta la temperatura, dentro de un diseño con
cuatro valores diferentes: \(14, 18, 22, 26^oC\). En cada realización del
experimento se miden tres variables:

\[\begin{aligned}
y_1&=\mbox{El peso fresco de las raíces (medido en gramos)} \\
y_2&=\mbox{La longitud máxima de las raíces (medida en milímetros)} \\
y_3&=\mbox{El peso fresco de las hojas (medido en gramos)}\end{aligned}\]
Los datos se encuentran en el fichero \texttt{hongos.txt}. Vamos a contrastar
el efecto del tratamiento con hongos, el efecto de la temperatura y la
posible interacción entre ambos efectos.

\footnotesize

\begin{Shaded}
\begin{Highlighting}[]
\SpecialCharTok{\textgreater{}}\NormalTok{ hongos }\OtherTok{\textless{}{-}} \FunctionTok{read.table}\NormalTok{(}\AttributeTok{file =} \StringTok{"data/hongos.txt"}\NormalTok{, }\AttributeTok{header =} \ConstantTok{TRUE}\NormalTok{)}
\SpecialCharTok{\textgreater{}} \FunctionTok{head}\NormalTok{(hongos)}
\NormalTok{  tratamiento temperatura  y1   y2  y3}
\DecValTok{1}           \DecValTok{1}           \DecValTok{1} \FloatTok{2.2} \FloatTok{23.5} \FloatTok{1.7}
\DecValTok{2}           \DecValTok{1}           \DecValTok{1} \FloatTok{3.0} \FloatTok{27.0} \FloatTok{2.3}
\DecValTok{3}           \DecValTok{1}           \DecValTok{1} \FloatTok{3.3} \FloatTok{24.5} \FloatTok{3.2}
\DecValTok{4}           \DecValTok{1}           \DecValTok{1} \FloatTok{2.2} \FloatTok{20.5} \FloatTok{1.5}
\DecValTok{5}           \DecValTok{1}           \DecValTok{1} \FloatTok{2.0} \FloatTok{19.0} \FloatTok{2.0}
\DecValTok{6}           \DecValTok{1}           \DecValTok{1} \FloatTok{3.5} \FloatTok{23.5} \FloatTok{2.9}
\SpecialCharTok{\textgreater{}}\NormalTok{ y }\OtherTok{\textless{}{-}} \FunctionTok{cbind}\NormalTok{(hongos}\SpecialCharTok{$}\NormalTok{y1, hongos}\SpecialCharTok{$}\NormalTok{y2, hongos}\SpecialCharTok{$}\NormalTok{y3)}
\SpecialCharTok{\textgreater{}}\NormalTok{ tratamiento }\OtherTok{\textless{}{-}} \FunctionTok{factor}\NormalTok{(hongos}\SpecialCharTok{$}\NormalTok{tratamiento)}
\SpecialCharTok{\textgreater{}}\NormalTok{ temperatura }\OtherTok{\textless{}{-}} \FunctionTok{factor}\NormalTok{(hongos}\SpecialCharTok{$}\NormalTok{temperatura)}
\end{Highlighting}
\end{Shaded}

\normalsize

Empezando por el modelo más general, con ambos efectos principales y la
interacción, si se contrasta la interacción resulta:

\footnotesize

\begin{Shaded}
\begin{Highlighting}[]
\SpecialCharTok{\textgreater{}}\NormalTok{ m12i }\OtherTok{=} \FunctionTok{lm}\NormalTok{(y }\SpecialCharTok{\textasciitilde{}}\NormalTok{ tratamiento }\SpecialCharTok{*}\NormalTok{ temperatura)}
\SpecialCharTok{\textgreater{}}\NormalTok{ m12 }\OtherTok{=} \FunctionTok{lm}\NormalTok{(y }\SpecialCharTok{\textasciitilde{}}\NormalTok{ tratamiento }\SpecialCharTok{+}\NormalTok{ temperatura)}
\SpecialCharTok{\textgreater{}} \FunctionTok{anova}\NormalTok{(m12, m12i, }\AttributeTok{test =} \StringTok{"Wilks"}\NormalTok{)}
\NormalTok{Analysis of Variance Table}

\NormalTok{Model }\DecValTok{1}\SpecialCharTok{:}\NormalTok{ y }\SpecialCharTok{\textasciitilde{}}\NormalTok{ tratamiento }\SpecialCharTok{+}\NormalTok{ temperatura}
\NormalTok{Model }\DecValTok{2}\SpecialCharTok{:}\NormalTok{ y }\SpecialCharTok{\textasciitilde{}}\NormalTok{ tratamiento }\SpecialCharTok{*}\NormalTok{ temperatura}
\NormalTok{  Res.Df Df Gen.var.   Wilks approx F num Df den Df }\FunctionTok{Pr}\NormalTok{(}\SpecialCharTok{\textgreater{}}\NormalTok{F)}
\DecValTok{1}     \DecValTok{43}      \FloatTok{12.852}                                      
\DecValTok{2}     \DecValTok{40} \SpecialCharTok{{-}}\DecValTok{3}   \FloatTok{12.299} \FloatTok{0.70548}   \FloatTok{1.5864}      \DecValTok{9} \FloatTok{92.633} \FloatTok{0.1307}
\SpecialCharTok{\textgreater{}} \FunctionTok{anova}\NormalTok{(m12, m12i, }\AttributeTok{test =} \StringTok{"Roy"}\NormalTok{)}
\NormalTok{Analysis of Variance Table}

\NormalTok{Model }\DecValTok{1}\SpecialCharTok{:}\NormalTok{ y }\SpecialCharTok{\textasciitilde{}}\NormalTok{ tratamiento }\SpecialCharTok{+}\NormalTok{ temperatura}
\NormalTok{Model }\DecValTok{2}\SpecialCharTok{:}\NormalTok{ y }\SpecialCharTok{\textasciitilde{}}\NormalTok{ tratamiento }\SpecialCharTok{*}\NormalTok{ temperatura}
\NormalTok{  Res.Df Df Gen.var.     Roy approx F num Df den Df  }\FunctionTok{Pr}\NormalTok{(}\SpecialCharTok{\textgreater{}}\NormalTok{F)  }
\DecValTok{1}     \DecValTok{43}      \FloatTok{12.852}                                         
\DecValTok{2}     \DecValTok{40} \SpecialCharTok{{-}}\DecValTok{3}   \FloatTok{12.299} \FloatTok{0.25155}   \FloatTok{3.3541}      \DecValTok{3}     \DecValTok{40} \FloatTok{0.02819} \SpecialCharTok{*}
\SpecialCharTok{{-}{-}{-}}
\NormalTok{Signif. codes}\SpecialCharTok{:}  \DecValTok{0} \StringTok{\textquotesingle{}***\textquotesingle{}} \FloatTok{0.001} \StringTok{\textquotesingle{}**\textquotesingle{}} \FloatTok{0.01} \StringTok{\textquotesingle{}*\textquotesingle{}} \FloatTok{0.05} \StringTok{\textquotesingle{}.\textquotesingle{}} \FloatTok{0.1} \StringTok{\textquotesingle{} \textquotesingle{}} \DecValTok{1}
\SpecialCharTok{\textgreater{}} \FunctionTok{anova}\NormalTok{(m12, m12i, }\AttributeTok{test =} \StringTok{"Hotelling{-}Lawley"}\NormalTok{)}
\NormalTok{Analysis of Variance Table}

\NormalTok{Model }\DecValTok{1}\SpecialCharTok{:}\NormalTok{ y }\SpecialCharTok{\textasciitilde{}}\NormalTok{ tratamiento }\SpecialCharTok{+}\NormalTok{ temperatura}
\NormalTok{Model }\DecValTok{2}\SpecialCharTok{:}\NormalTok{ y }\SpecialCharTok{\textasciitilde{}}\NormalTok{ tratamiento }\SpecialCharTok{*}\NormalTok{ temperatura}
\NormalTok{  Res.Df Df Gen.var. Hotelling}\SpecialCharTok{{-}}\NormalTok{Lawley approx F num Df den Df }\FunctionTok{Pr}\NormalTok{(}\SpecialCharTok{\textgreater{}}\NormalTok{F)}
\DecValTok{1}     \DecValTok{43}      \FloatTok{12.852}                                               
\DecValTok{2}     \DecValTok{40} \SpecialCharTok{{-}}\DecValTok{3}   \FloatTok{12.299}          \FloatTok{0.38126}   \FloatTok{1.5533}      \DecValTok{9}    \DecValTok{110} \FloatTok{0.1384}
\SpecialCharTok{\textgreater{}} \FunctionTok{anova}\NormalTok{(m12, m12i, }\AttributeTok{test =} \StringTok{"Pillai"}\NormalTok{)}
\NormalTok{Analysis of Variance Table}

\NormalTok{Model }\DecValTok{1}\SpecialCharTok{:}\NormalTok{ y }\SpecialCharTok{\textasciitilde{}}\NormalTok{ tratamiento }\SpecialCharTok{+}\NormalTok{ temperatura}
\NormalTok{Model }\DecValTok{2}\SpecialCharTok{:}\NormalTok{ y }\SpecialCharTok{\textasciitilde{}}\NormalTok{ tratamiento }\SpecialCharTok{*}\NormalTok{ temperatura}
\NormalTok{  Res.Df Df Gen.var.  Pillai approx F num Df den Df }\FunctionTok{Pr}\NormalTok{(}\SpecialCharTok{\textgreater{}}\NormalTok{F)}
\DecValTok{1}     \DecValTok{43}      \FloatTok{12.852}                                      
\DecValTok{2}     \DecValTok{40} \SpecialCharTok{{-}}\DecValTok{3}   \FloatTok{12.299} \FloatTok{0.32058}   \FloatTok{1.5953}      \DecValTok{9}    \DecValTok{120} \FloatTok{0.1241}
\end{Highlighting}
\end{Shaded}

\normalsize

Únicamente el test de Roy produce un nivel crítico algo pequeño,
mientras que los otros tests no muestran significación. En consecuencia,
se suprime la interacción del modelo. En el modelo sin interacción y con
los dos efectos principales, se contrasta la significación de cada uno
de los efectos, con los resultados siguientes. Para el tratamiento:

\footnotesize

\begin{Shaded}
\begin{Highlighting}[]
\SpecialCharTok{\textgreater{}}\NormalTok{ m1 }\OtherTok{=} \FunctionTok{lm}\NormalTok{(y }\SpecialCharTok{\textasciitilde{}}\NormalTok{ tratamiento)}
\SpecialCharTok{\textgreater{}} \FunctionTok{anova}\NormalTok{(m1, m12, }\AttributeTok{test =} \StringTok{"Wilks"}\NormalTok{)}
\NormalTok{Analysis of Variance Table}

\NormalTok{Model }\DecValTok{1}\SpecialCharTok{:}\NormalTok{ y }\SpecialCharTok{\textasciitilde{}}\NormalTok{ tratamiento}
\NormalTok{Model }\DecValTok{2}\SpecialCharTok{:}\NormalTok{ y }\SpecialCharTok{\textasciitilde{}}\NormalTok{ tratamiento }\SpecialCharTok{+}\NormalTok{ temperatura}
\NormalTok{  Res.Df Df Gen.var.    Wilks approx F num Df den Df    }\FunctionTok{Pr}\NormalTok{(}\SpecialCharTok{\textgreater{}}\NormalTok{F)    }
\DecValTok{1}     \DecValTok{46}      \FloatTok{29.940}                                              
\DecValTok{2}     \DecValTok{43} \SpecialCharTok{{-}}\DecValTok{3}   \FloatTok{12.852} \FloatTok{0.064601}    \FloatTok{23.12}      \DecValTok{9} \FloatTok{99.934} \SpecialCharTok{\textless{}} \FloatTok{2.2e{-}16} \SpecialCharTok{**}\ErrorTok{*}
\SpecialCharTok{{-}{-}{-}}
\NormalTok{Signif. codes}\SpecialCharTok{:}  \DecValTok{0} \StringTok{\textquotesingle{}***\textquotesingle{}} \FloatTok{0.001} \StringTok{\textquotesingle{}**\textquotesingle{}} \FloatTok{0.01} \StringTok{\textquotesingle{}*\textquotesingle{}} \FloatTok{0.05} \StringTok{\textquotesingle{}.\textquotesingle{}} \FloatTok{0.1} \StringTok{\textquotesingle{} \textquotesingle{}} \DecValTok{1}
\SpecialCharTok{\textgreater{}} \FunctionTok{anova}\NormalTok{(m1, m12, }\AttributeTok{test =} \StringTok{"Roy"}\NormalTok{)}
\NormalTok{Analysis of Variance Table}

\NormalTok{Model }\DecValTok{1}\SpecialCharTok{:}\NormalTok{ y }\SpecialCharTok{\textasciitilde{}}\NormalTok{ tratamiento}
\NormalTok{Model }\DecValTok{2}\SpecialCharTok{:}\NormalTok{ y }\SpecialCharTok{\textasciitilde{}}\NormalTok{ tratamiento }\SpecialCharTok{+}\NormalTok{ temperatura}
\NormalTok{  Res.Df Df Gen.var.    Roy approx F num Df den Df    }\FunctionTok{Pr}\NormalTok{(}\SpecialCharTok{\textgreater{}}\NormalTok{F)    }
\DecValTok{1}     \DecValTok{46}      \FloatTok{29.940}                                            
\DecValTok{2}     \DecValTok{43} \SpecialCharTok{{-}}\DecValTok{3}   \FloatTok{12.852} \FloatTok{5.1828}   \FloatTok{74.287}      \DecValTok{3}     \DecValTok{43} \SpecialCharTok{\textless{}} \FloatTok{2.2e{-}16} \SpecialCharTok{**}\ErrorTok{*}
\SpecialCharTok{{-}{-}{-}}
\NormalTok{Signif. codes}\SpecialCharTok{:}  \DecValTok{0} \StringTok{\textquotesingle{}***\textquotesingle{}} \FloatTok{0.001} \StringTok{\textquotesingle{}**\textquotesingle{}} \FloatTok{0.01} \StringTok{\textquotesingle{}*\textquotesingle{}} \FloatTok{0.05} \StringTok{\textquotesingle{}.\textquotesingle{}} \FloatTok{0.1} \StringTok{\textquotesingle{} \textquotesingle{}} \DecValTok{1}
\SpecialCharTok{\textgreater{}} \FunctionTok{anova}\NormalTok{(m1, m12, }\AttributeTok{test =} \StringTok{"Hotelling{-}Lawley"}\NormalTok{)}
\NormalTok{Analysis of Variance Table}

\NormalTok{Model }\DecValTok{1}\SpecialCharTok{:}\NormalTok{ y }\SpecialCharTok{\textasciitilde{}}\NormalTok{ tratamiento}
\NormalTok{Model }\DecValTok{2}\SpecialCharTok{:}\NormalTok{ y }\SpecialCharTok{\textasciitilde{}}\NormalTok{ tratamiento }\SpecialCharTok{+}\NormalTok{ temperatura}
\NormalTok{  Res.Df Df Gen.var. Hotelling}\SpecialCharTok{{-}}\NormalTok{Lawley approx F num Df den Df    }\FunctionTok{Pr}\NormalTok{(}\SpecialCharTok{\textgreater{}}\NormalTok{F)    }
\DecValTok{1}     \DecValTok{46}      \FloatTok{29.940}                                                      
\DecValTok{2}     \DecValTok{43} \SpecialCharTok{{-}}\DecValTok{3}   \FloatTok{12.852}           \FloatTok{6.6858}   \FloatTok{29.467}      \DecValTok{9}    \DecValTok{119} \SpecialCharTok{\textless{}} \FloatTok{2.2e{-}16} \SpecialCharTok{**}\ErrorTok{*}
\SpecialCharTok{{-}{-}{-}}
\NormalTok{Signif. codes}\SpecialCharTok{:}  \DecValTok{0} \StringTok{\textquotesingle{}***\textquotesingle{}} \FloatTok{0.001} \StringTok{\textquotesingle{}**\textquotesingle{}} \FloatTok{0.01} \StringTok{\textquotesingle{}*\textquotesingle{}} \FloatTok{0.05} \StringTok{\textquotesingle{}.\textquotesingle{}} \FloatTok{0.1} \StringTok{\textquotesingle{} \textquotesingle{}} \DecValTok{1}
\SpecialCharTok{\textgreater{}} \FunctionTok{anova}\NormalTok{(m1, m12, }\AttributeTok{test =} \StringTok{"Pillai"}\NormalTok{)}
\NormalTok{Analysis of Variance Table}

\NormalTok{Model }\DecValTok{1}\SpecialCharTok{:}\NormalTok{ y }\SpecialCharTok{\textasciitilde{}}\NormalTok{ tratamiento}
\NormalTok{Model }\DecValTok{2}\SpecialCharTok{:}\NormalTok{ y }\SpecialCharTok{\textasciitilde{}}\NormalTok{ tratamiento }\SpecialCharTok{+}\NormalTok{ temperatura}
\NormalTok{  Res.Df Df Gen.var. Pillai approx F num Df den Df   }\FunctionTok{Pr}\NormalTok{(}\SpecialCharTok{\textgreater{}}\NormalTok{F)    }
\DecValTok{1}     \DecValTok{46}      \FloatTok{29.940}                                           
\DecValTok{2}     \DecValTok{43} \SpecialCharTok{{-}}\DecValTok{3}   \FloatTok{12.852} \FloatTok{1.4391}   \FloatTok{13.215}      \DecValTok{9}    \DecValTok{129} \FloatTok{8.45e{-}15} \SpecialCharTok{**}\ErrorTok{*}
\SpecialCharTok{{-}{-}{-}}
\NormalTok{Signif. codes}\SpecialCharTok{:}  \DecValTok{0} \StringTok{\textquotesingle{}***\textquotesingle{}} \FloatTok{0.001} \StringTok{\textquotesingle{}**\textquotesingle{}} \FloatTok{0.01} \StringTok{\textquotesingle{}*\textquotesingle{}} \FloatTok{0.05} \StringTok{\textquotesingle{}.\textquotesingle{}} \FloatTok{0.1} \StringTok{\textquotesingle{} \textquotesingle{}} \DecValTok{1}
\end{Highlighting}
\end{Shaded}

\normalsize

Para la temperatura:

\footnotesize

\begin{Shaded}
\begin{Highlighting}[]
\SpecialCharTok{\textgreater{}}\NormalTok{ m2 }\OtherTok{=} \FunctionTok{lm}\NormalTok{(y }\SpecialCharTok{\textasciitilde{}}\NormalTok{ temperatura)}
\SpecialCharTok{\textgreater{}} \FunctionTok{anova}\NormalTok{(m2, m12, }\AttributeTok{test =} \StringTok{"Wilks"}\NormalTok{)}
\NormalTok{Analysis of Variance Table}

\NormalTok{Model }\DecValTok{1}\SpecialCharTok{:}\NormalTok{ y }\SpecialCharTok{\textasciitilde{}}\NormalTok{ temperatura}
\NormalTok{Model }\DecValTok{2}\SpecialCharTok{:}\NormalTok{ y }\SpecialCharTok{\textasciitilde{}}\NormalTok{ tratamiento }\SpecialCharTok{+}\NormalTok{ temperatura}
\NormalTok{  Res.Df Df Gen.var.  Wilks approx F num Df den Df  }\FunctionTok{Pr}\NormalTok{(}\SpecialCharTok{\textgreater{}}\NormalTok{F)  }
\DecValTok{1}     \DecValTok{44}      \FloatTok{13.439}                                        
\DecValTok{2}     \DecValTok{43} \SpecialCharTok{{-}}\DecValTok{1}   \FloatTok{12.852} \FloatTok{0.8163}   \FloatTok{3.0756}      \DecValTok{3}     \DecValTok{41} \FloatTok{0.03807} \SpecialCharTok{*}
\SpecialCharTok{{-}{-}{-}}
\NormalTok{Signif. codes}\SpecialCharTok{:}  \DecValTok{0} \StringTok{\textquotesingle{}***\textquotesingle{}} \FloatTok{0.001} \StringTok{\textquotesingle{}**\textquotesingle{}} \FloatTok{0.01} \StringTok{\textquotesingle{}*\textquotesingle{}} \FloatTok{0.05} \StringTok{\textquotesingle{}.\textquotesingle{}} \FloatTok{0.1} \StringTok{\textquotesingle{} \textquotesingle{}} \DecValTok{1}
\SpecialCharTok{\textgreater{}} \FunctionTok{anova}\NormalTok{(m2, m12, }\AttributeTok{test =} \StringTok{"Roy"}\NormalTok{)}
\NormalTok{Analysis of Variance Table}

\NormalTok{Model }\DecValTok{1}\SpecialCharTok{:}\NormalTok{ y }\SpecialCharTok{\textasciitilde{}}\NormalTok{ temperatura}
\NormalTok{Model }\DecValTok{2}\SpecialCharTok{:}\NormalTok{ y }\SpecialCharTok{\textasciitilde{}}\NormalTok{ tratamiento }\SpecialCharTok{+}\NormalTok{ temperatura}
\NormalTok{  Res.Df Df Gen.var.     Roy approx F num Df den Df  }\FunctionTok{Pr}\NormalTok{(}\SpecialCharTok{\textgreater{}}\NormalTok{F)  }
\DecValTok{1}     \DecValTok{44}      \FloatTok{13.439}                                         
\DecValTok{2}     \DecValTok{43} \SpecialCharTok{{-}}\DecValTok{1}   \FloatTok{12.852} \FloatTok{0.22504}   \FloatTok{3.0756}      \DecValTok{3}     \DecValTok{41} \FloatTok{0.03807} \SpecialCharTok{*}
\SpecialCharTok{{-}{-}{-}}
\NormalTok{Signif. codes}\SpecialCharTok{:}  \DecValTok{0} \StringTok{\textquotesingle{}***\textquotesingle{}} \FloatTok{0.001} \StringTok{\textquotesingle{}**\textquotesingle{}} \FloatTok{0.01} \StringTok{\textquotesingle{}*\textquotesingle{}} \FloatTok{0.05} \StringTok{\textquotesingle{}.\textquotesingle{}} \FloatTok{0.1} \StringTok{\textquotesingle{} \textquotesingle{}} \DecValTok{1}
\SpecialCharTok{\textgreater{}} \FunctionTok{anova}\NormalTok{(m2, m12, }\AttributeTok{test =} \StringTok{"Hotelling{-}Lawley"}\NormalTok{)}
\NormalTok{Analysis of Variance Table}

\NormalTok{Model }\DecValTok{1}\SpecialCharTok{:}\NormalTok{ y }\SpecialCharTok{\textasciitilde{}}\NormalTok{ temperatura}
\NormalTok{Model }\DecValTok{2}\SpecialCharTok{:}\NormalTok{ y }\SpecialCharTok{\textasciitilde{}}\NormalTok{ tratamiento }\SpecialCharTok{+}\NormalTok{ temperatura}
\NormalTok{  Res.Df Df Gen.var. Hotelling}\SpecialCharTok{{-}}\NormalTok{Lawley approx F num Df den Df  }\FunctionTok{Pr}\NormalTok{(}\SpecialCharTok{\textgreater{}}\NormalTok{F)  }
\DecValTok{1}     \DecValTok{44}      \FloatTok{13.439}                                                  
\DecValTok{2}     \DecValTok{43} \SpecialCharTok{{-}}\DecValTok{1}   \FloatTok{12.852}          \FloatTok{0.22504}   \FloatTok{3.0756}      \DecValTok{3}     \DecValTok{41} \FloatTok{0.03807} \SpecialCharTok{*}
\SpecialCharTok{{-}{-}{-}}
\NormalTok{Signif. codes}\SpecialCharTok{:}  \DecValTok{0} \StringTok{\textquotesingle{}***\textquotesingle{}} \FloatTok{0.001} \StringTok{\textquotesingle{}**\textquotesingle{}} \FloatTok{0.01} \StringTok{\textquotesingle{}*\textquotesingle{}} \FloatTok{0.05} \StringTok{\textquotesingle{}.\textquotesingle{}} \FloatTok{0.1} \StringTok{\textquotesingle{} \textquotesingle{}} \DecValTok{1}
\SpecialCharTok{\textgreater{}} \FunctionTok{anova}\NormalTok{(m2, m12, }\AttributeTok{test =} \StringTok{"Pillai"}\NormalTok{)}
\NormalTok{Analysis of Variance Table}

\NormalTok{Model }\DecValTok{1}\SpecialCharTok{:}\NormalTok{ y }\SpecialCharTok{\textasciitilde{}}\NormalTok{ temperatura}
\NormalTok{Model }\DecValTok{2}\SpecialCharTok{:}\NormalTok{ y }\SpecialCharTok{\textasciitilde{}}\NormalTok{ tratamiento }\SpecialCharTok{+}\NormalTok{ temperatura}
\NormalTok{  Res.Df Df Gen.var. Pillai approx F num Df den Df  }\FunctionTok{Pr}\NormalTok{(}\SpecialCharTok{\textgreater{}}\NormalTok{F)  }
\DecValTok{1}     \DecValTok{44}      \FloatTok{13.439}                                        
\DecValTok{2}     \DecValTok{43} \SpecialCharTok{{-}}\DecValTok{1}   \FloatTok{12.852} \FloatTok{0.1837}   \FloatTok{3.0756}      \DecValTok{3}     \DecValTok{41} \FloatTok{0.03807} \SpecialCharTok{*}
\SpecialCharTok{{-}{-}{-}}
\NormalTok{Signif. codes}\SpecialCharTok{:}  \DecValTok{0} \StringTok{\textquotesingle{}***\textquotesingle{}} \FloatTok{0.001} \StringTok{\textquotesingle{}**\textquotesingle{}} \FloatTok{0.01} \StringTok{\textquotesingle{}*\textquotesingle{}} \FloatTok{0.05} \StringTok{\textquotesingle{}.\textquotesingle{}} \FloatTok{0.1} \StringTok{\textquotesingle{} \textquotesingle{}} \DecValTok{1}
\end{Highlighting}
\end{Shaded}

\normalsize

Por tanto, la temperatura tiene un efecto muy significativo sobre las
variables que se están midiendo. El tratamiento también resulta
significativo, aunque no es tan notorio. Los niveles críticos de los
cuatro test coinciden para el contraste del tratamiento. Este hecho se
debe a que hay únicamente dos tratamientos, por lo que la matriz \(H\)
asociada al contraste tiene rango uno.

\chapter{Análisis de componentes principales}\label{analisis-de-componentes-principales}

El análisis de componentes principales se concibe como una técnica de
reducción de la dimensión, pues permite pasar de una gran cantidad de
variables interrelacionadas a unas pocas componentes principales. El
método consiste en buscar combinaciones lineales de las variables
originales que representen lo mejor posible a la variabilidad presente
en los datos. De este modo, con unas pocas combinaciones lineales, que
serán las componentes principales, sería suficiente para entender la
información contenida en los datos. Al mismo tiempo, la forma en que se
construyen las componentes, y su relación con unas u otras variables
originales, sirven para entender la estructura de correlación inherente
a los datos. Por último, las componentes principales, que forman un
vector aleatorio de dimensión menor, pueden ser empleadas en análisis
estadísticos posteriores.

\section{Descomposición de un vector aleatorio en sus componentes principales}\label{descomposicion-de-un-vector-aleatorio-en-sus-componentes-principales}

En esta sección se va a definir el concepto de componentes principales,
y cómo se obtienen a partir de la matriz de covarianzas de un vector
aleatorio. Estas ideas y procedimientos se aplican de igual modo tanto a
un vector aleatorio como a un conjunto de datos. En este último caso, se
puede pensar que el análisis se está aplicando a una distribución de
probabilidad discreta equiprobable sobre los vectores observados. Por
este motivo, esta sección se desarrolla sobre un vector aleatorio
general, con la única condición de que exista la matriz de covarianzas.

\begin{definition}
\protect\hypertarget{def:unnamed-chunk-70}{}\label{def:unnamed-chunk-70}Sea \(\bm{x}=(x_1,\ldots,x_d)^\prime{}\) un vector aleatorio
\(d\)-dimensional con vector de medias \(\bm\mu=E({\bm{x}})\) y matriz de
covarianzas \(\bm\Sigma=E((\bm{x}-\bm\mu)(\bm{x}-\bm\mu)^\prime{})\). Se
define la primera componente principal de \(\bm{x}\) como una variable
aleatoria \(z_1\) tal que

\[z_1=\bm{v_1}^{\prime}\bm{x}=v_{11}x_1+\cdots+v_{d1}x_d\ \ \ \mbox{con}
\ \ \ \bm{v_1}=(v_{11},\ldots,v_{d1})^{\prime} \in {\Bbb R}^d,\]

\[\textnormal{Var}(z_1)=\max\{\textnormal{Var}(\bm{v}^{\prime}\bm{x}):\bm{v}\in {\Bbb R}^d,\ \bm{v}^{\prime}\bm{v}=1\}.\]
La primera componente principal es una combinación lineal normalizada de
las variables de \(\bm{x}\) y, de entre todas las combinaciones lineales
normalizadas, es la que tiene mayor varianza.
\end{definition}

\textbf{Nota.} Debe observarse que
\(\textnormal{Var}(\alpha \bm{v}^{\prime}\bm{x})=\alpha^2\textnormal{Var}(\bm{v}^{\prime}\bm{x})\)
con \(\alpha\in {\Bbb R}\). De modo que multiplicando por una constante
podemos hacer la varianza tan grande o tan pequeña como queramos. Por
ello es necesario normalizar, pidiendo que
\(\|\bm{v}\|=\sqrt{\bm{v}^{\prime}\bm{v}}=1\).

\begin{theorem}
\protect\hypertarget{thm:unnamed-chunk-71}{}\label{thm:unnamed-chunk-71}La primera componente principal de \(\bm{x}\) adopta la forma

\[z_1=\bm{v_1}^{\prime}\bm{x},\] siendo \(\lambda_1\) el mayor autovalor
de la matriz de covarianzas \(\bm\Sigma\) y \(\bm{v_1}\) un autovector
asociado a \(\lambda_1\) de norma uno (\(\bm{v_1}^{\prime}\bm{v_1}=1\)).
Además,

\[\textnormal{Var}(z_1)=\lambda_1.\]
\end{theorem}

\begin{proof}
Consideremos \(z=\bm{v}^{\prime}\bm{x}\). Entonces podemos
calcular su varianza

\[\textnormal{Var}(z)=\textnormal{Cov}(\bm{v}^{\prime}\bm{x},\bm{v}^{\prime}\bm{x})=\bm{v}^{\prime}\textnormal{Cov}(\bm{x},\bm{x})\bm{v}=\bm{v}^{\prime}\bm\Sigma \bm{\bm{v}}.\]
Nuestro problema consiste en calcular

\[\begin{array}{rl}
\max & \bm{v}^{\prime}\bm\Sigma \bm{v} \\
\mbox{ sujeto a} & \bm{v}^{\prime}\bm{v}=1.
\end{array}\] Resolvemos este problema por el método de los
multiplicadores de Lagrange.

\[L=\bm{v}^{\prime}\bm\Sigma \bm{v}-\lambda(\bm{v}^{\prime}\bm{v}-1)\]
Calculamos la matriz jacobiana de \(L\) como función de \(\bm{v}\) e
igualamos a cero

\[\frac{\partial L}{\partial \bm{v}}=2\bm{v}^{\prime}\bm\Sigma-2\lambda \bm{v}^{\prime}=0.\]
Equivalentemente \[\bm\Sigma \bm{v}=\lambda \bm{v}.\] Luego el vector
\(\bm{v_1}\) que maximice la función sujeto a la restricción ha de ser un
autovector de \(\bm\Sigma\), y su autovalor asociado es \(\lambda_1\).
Multiplicando los dos términos de la ecuación por la izquierda por \(\bm{v_1}^{\prime}\) resulta

\[\bm{v_1}^{\prime}\bm\Sigma \bm{v_1}=\bm{v_1}^{\prime}\lambda_1 \bm{v_1}=\lambda_1 \bm{v_1}^{\prime} \bm{v_1}=\lambda_1.\]
Luego el autovalor \(\lambda_1\) es la función objetivo en el máximo y por
tanto, \(\textnormal{Var}(z_1)=\lambda_1\).
\end{proof}

\begin{definition}
\protect\hypertarget{def:unnamed-chunk-73}{}\label{def:unnamed-chunk-73}Se define la segunda componente principal de \(\bm{x}\) como una variable
aleatoria \(z_2\) tal que

\[z_2=\bm{v_2}^{\prime}\bm{x}=v_{12}x_1+\cdots+v_{d2}x_d\ \ \ \mbox{con}\ \ \ \bm{v_2}=(v_{12},\ldots,v_{d2})^{\prime}
\in {\Bbb R}^d,\]

\[\textnormal{Var}(z_2)=\max\{\textnormal{Var}(\bm{v}^{\prime}\bm{x}):\bm{v}\in {\Bbb R}^d,\ \bm{v}^{\prime}\bm{v}=1,\ \bm{v}^{\prime}\bm{v_1}=0\}.\]
La segunda componente principal es otra combinación lineal de las
variables de \(\bm{x}\) y, de entre todas las combinaciones lineales
formadas por vectores unitarios ortogonales a \(\bm{v_1}\), es la que
tiene mayor varianza.
\end{definition}

Nótese que

\[\textnormal{Cov}(\bm{v}^{\prime}\bm{x},z_1)=\textnormal{Cov}(\bm{v}^{\prime}\bm{x}, \bm{v_1}^{\prime}\bm{x})=\bm{v}^{\prime}\textnormal{Cov}(\bm{x},\bm{x})\bm{v_1}=\bm{v}^{\prime}\bm\Sigma \bm{v_1}.\]

Enunciamos a continuación un resultado sencillo.

\textbf{Resultado}. Sea \(\bm{w}\in {\Bbb R}^d\) un autovector de \(\bm\Sigma\)
asociado a un autovalor no nulo \(\lambda\neq 0\), y sea
\(\bm{v}\in {\Bbb R}^d\). Entonces

\[\bm{v}^{\prime} \bm\Sigma \bm{w}=0 \Leftrightarrow \bm{v}^{\prime} \bm{w}=0.\]

\begin{proof}
En primer lugar,

\[\bm{v}^{\prime}\bm\Sigma \bm{w}=\bm{v}^{\prime}\lambda \bm{w}=\lambda \bm{v}^{\prime} \bm{w}=0.\]

En el otro sentido,

\[\bm{v}^{\prime} \bm{w}=\bm{v}^{\prime}\frac{1}{\lambda}\bm\Sigma \bm{w}=\frac{1}{\lambda} \bm{v}^{\prime}\bm\Sigma \bm{w}
=\frac{1}{\lambda}0=0.\]
\end{proof}

Por tanto, en aplicación de este resultado, \(\textnormal{Cov}(\bm{v}^{\prime}\bm{x},z_1)=0 \Leftrightarrow \bm{v}^{\prime} \bm{v_1}=0.\) Y hemos probado que la condición de ortogonalidad entre los vectores
\(\bm{v_1}\) y \(\bm{v_2}\) es equivalente a la incorrelación de las
componentes \(z_1\) y \(z_2\). En definitiva, podríamos definir la segunda
componente principal como la combinación lineal de las variables de
\(\bm{x}\) que tiene mayor varianza de entre las combinaciones lineales
normalizadas e incorrelacionadas con la primera componente principal.

Respecto a la condición de que \(\lambda_1\neq 0\), hemos de recordar que
\(\lambda_1=\textnormal{Var}(z_1)\). Por tanto si
\(\lambda_1=\textnormal{Var}(z_1)=0\), todos los demás autovalores serán
también cero y \(z_1\) y las demás componentes serán variables
degeneradas.

\begin{theorem}
\protect\hypertarget{thm:unnamed-chunk-75}{}\label{thm:unnamed-chunk-75}La segunda componente principal de \(\bm{x}\) adopta la forma

\[z_2=\bm{v_2}^{\prime}\bm{x}\] siendo \(\lambda_2\) el segundo mayor
autovalor de la matriz de covarianzas \(\bm\Sigma\) y \(\bm{v_2}\) un
autovector asociado a \(\lambda_2\) de norma uno
(\(\bm{v_2}^{\prime}\bm{v_2}=1\)) y ortogonal a \(\bm{v_1}\)
(\(\bm{v_1}^{\prime}\bm{v_2}=0\)). Además

\[\textnormal{Var}(z_2)=\lambda_2.\]
\end{theorem}

\begin{proof}
Procederemos de modo muy similar al caso anterior. Consideramos
\(z=\bm{v}^{\prime}\bm{x}\), cuya varianza es
\(\textnormal{Var}(z)=\bm{v}^{\prime}\bm\Sigma \bm{v}\). Nuestro problema
consiste en

\[\begin{array}{rl}
\max & \bm{v}^{\prime}\bm\Sigma \bm{v} \\
\mbox{ sujeto a} & \bm{v}^{\prime}\bm{v}=1, \\
 & \bm{v_1}^{\prime} \bm{v}=0.
\end{array}\] Resolvemos este problema por el método de los
multiplicadores de Lagrange.

\[L=\bm{v}^{\prime}\bm\Sigma \bm{v}-\lambda_2(\bm{v}^{\prime}\bm{v}-1)-\gamma_2 \bm{v_1}^{\prime} \bm{v}.\]
Derivamos respecto de \(\bm{v}\) e igualamos a cero (momento en el cual
podemos sustituir el valor \(\bm{v_2}\) para mejor comprensión):

\begin{equation}
\left(\frac{\partial L}{\partial \bm{v}}/_{\bm{v}=\bm{v_2}}\right)^{\prime}
=2\bm\Sigma \bm{v_2}-2\lambda_2 \bm{v_2}-\gamma_2\bm{v_1}=0.
\label{eq:lagrange}
\end{equation}

Multiplicando por la izquierda por \(\bm{v_1}^{\prime}\) resulta

\[2\bm{v_1}^{\prime}\bm\Sigma \bm{v_2} - 2\lambda_2 \bm{v_1}^{\prime} \bm{v_2} - \gamma_2 \bm{v_1}^{\prime}\bm{v_1} =0.\]
Las restricciones imponen que \(\bm{v_1}^{\prime}\bm{v_2}=0\). Por la
nota, de ello deducimos que \(\bm{v_1}^{\prime}\bm\Sigma \bm{v_2}=0\). Y
\(\bm{v_1}\) ya verificaba \(\bm{v_1}^{\prime}\bm{v_1}=1\). Entonces la
ecuación anterior se reduce a \(\gamma_2=0\). Y al sustituir ese valor en
la ecuación \eqref{eq:lagrange} se obtiene

\[\bm\Sigma \bm{v_2}=\lambda_2 \bm{v_2}.\] Razonando igual que antes, el
vector \(\bm{v_2}\) que maximice la función sujeto a las restricciones ha
de ser un autovector de \(\bm\Sigma\), y su autovalor asociado es
\(\lambda_2\). La restricción de ortogonalidad respecto de \(\bm{v_1}\) nos
obliga a tomar el segundo mayor autovalor de \(\bm\Sigma\). Como antes,
\(\textnormal{Var}(z_2)=\bm{v_2}^{\prime}\bm\Sigma \bm{v_2}=\lambda_2\).
\end{proof}

Podríamos continuar este proceso extrayendo las componentes principales
de \(\bm{x}\) mediante los autovalores de la matriz de covarianzas
\(\bm\Sigma\) y la base ortonormal de autovectores asociados. De este
modo, obtenemos p componentes principales.

\begin{definition}
\protect\hypertarget{def:unnamed-chunk-77}{}\label{def:unnamed-chunk-77}Se definen las \(d\) componentes principales de \(\bm{x}\) como las
variables aleatorias \((z_1,\ldots,z_d)\) tales que

\[z_1=\bm{v_1}^{\prime}\bm{x},\ldots,z_d=\bm{v_d}^{\prime}\bm{x},\ \ \ \ \ \ \ \bm{v_1},\ldots,\bm{v_d}\in {\Bbb R}^d.\]

\[\begin{aligned}
\textnormal{Var}(z_1)&=\max\{\textnormal{Var}(\bm{v}^{\prime}\bm{x}):\bm{v}\in {\Bbb R}^d,\ \bm{v}^{\prime}\bm{v}=1\}, \\
\textnormal{Var}(z_2)&=\max\{\textnormal{Var}(\bm{v}^{\prime}\bm{x}):\bm{v}\in {\Bbb R}^d,\ \bm{v}^{\prime}\bm{v}=1,\ \bm{v_1}^{\prime}\bm{v}=0\}, \\
&\vdots \\
\textnormal{Var}(z_j)&=\max\{\textnormal{Var}(\bm{v}^{\prime}\bm{x}):\bm{v}\in {\Bbb R}^d,\ \bm{v}^{\prime}\bm{v}=1,
\ \bm{v_1}^{\prime}\bm{v}=0,\ldots,\ \bm{v_{j-1}}^{\prime}\bm{v}=0\}, \\
 &\vdots \\
\textnormal{Var}(z_d)&=\max\{\textnormal{Var}(\bm{v}^{\prime}\bm{x}):\bm{v}\in {\Bbb R}^d,\ \bm{v}^{\prime}\bm{v}=1,
\ \bm{v_1}^{\prime}\bm{v}=0,\ldots,\ \bm{v_{d-1}}^{\prime}\bm{v}=0\}.\end{aligned}\]
\end{definition}

\begin{theorem}
\protect\hypertarget{thm:unnamed-chunk-78}{}\label{thm:unnamed-chunk-78}Las \(d\) componentes principales de \(\bm{x}\) adoptan la forma

\[z_j=\bm{v_j}^{\prime}\bm{x}\ \ \ j\in\{1,\ldots,d\}\] siendo
\(\lambda_1\geq\ldots\geq\lambda_d\geq 0\) los \(d\) autovalores ordenados
de la matriz de covarianzas \(\bm\Sigma\) y \(\bm{v_1},\ldots,\bm{v_d}\) sus
autovectores asociados normalizados, esto es,
\(\{\bm{v_1},\ldots,\bm{v_d}\}\) es una base ortonormal de autovectores.
Además \(\textnormal{Cov}(z_j,z_k)=0\ \ \mbox{si}\ \ j\neq k\) y

\[\textnormal{Var}(z_j)=\lambda_j\ \ \ j\in\{1,\ldots,d\}.\]
\end{theorem}

Este teorema nos permite expresar \[\bm{z}=\bm{V}^{\prime}\bm{x}\]
siendo \(\bm{z}=(z_1,\ldots,z_d)^\prime{}\) y
\(\bm{V}=(\bm{v_1},\ldots,\bm{v_d})\) la matriz cuyas columnas son los
autovectores de \(\bm\Sigma\). Entonces

\[\textnormal{Cov}(\bm{z},\bm{z})=\textnormal{Cov}(\bm{V}^{\prime}\bm{x},\bm{V}^{\prime}\bm{x})=\bm{V}^{\prime}\textnormal{Cov}(\bm{x},\bm{x})\bm{V}=\bm{V}^{\prime}\bm\Sigma \bm{V}.\]
Pero la matriz de covarianzas de las componentes principales es diagonal
y contiene los autovalores de \(\bm\Sigma\) en la diagonal. Luego, el
proceso de extracción de las componentes principales se reduce a la
diagonalización de la matriz de covarianzas del vector aleatorio
\(\bm{x}\).

\textbf{Ejercicio.} Sea \(\bm{x}=(x_1,x_2)^{\prime}\) un vector aleatorio con
vector de medias \(\bm\mu=E(\bm{x})=(4,2)^{\prime}\) y matriz de
covarianzas

\[\bm\Sigma=\left(\begin{array}{cc} 2 & 1 \\ 1 & 2 \end{array}\right)\]
Obtener la descomposición del vector \(\bm{x}\) en sus componentes
principales.

\textbf{\includegraphics[width=1em,height=1em]{bookdown_AM_files/figure-latex/fa-icon-00ef1bb1d064b635425d49f06b331298.pdf} ~\textcolor{excolor}{Ejemplo.}} Se ha examinado a 25 alumnos, aspirantes a ingresar en la
Facultad de Matemáticas, de 5 materias diferentes: Geometría Diferencial
(cuyo resultado se almacena en la variable \emph{geodif}), Análisis Complejo
(\emph{ancompl}), Álgebra (\emph{alg}), Análisis Real (\emph{anreal}) y Estadística
(\emph{estad}). Las puntuaciones obtenidas se
encuentran en el fichero \texttt{aspirantes.txt}.

\footnotesize

\begin{Shaded}
\begin{Highlighting}[]
\SpecialCharTok{\textgreater{}}\NormalTok{ dat }\OtherTok{\textless{}{-}} \FunctionTok{read.table}\NormalTok{(}\StringTok{"data/aspirantes.txt"}\NormalTok{, }\AttributeTok{header =} \ConstantTok{TRUE}\NormalTok{)}
\SpecialCharTok{\textgreater{}}\NormalTok{ dat}
\NormalTok{   geodif ancompl alg anreal estad}
\DecValTok{1}      \DecValTok{36}      \DecValTok{58}  \DecValTok{43}     \DecValTok{36}    \DecValTok{37}
\DecValTok{2}      \DecValTok{62}      \DecValTok{54}  \DecValTok{50}     \DecValTok{46}    \DecValTok{52}
\DecValTok{3}      \DecValTok{31}      \DecValTok{42}  \DecValTok{41}     \DecValTok{40}    \DecValTok{29}
\DecValTok{4}      \DecValTok{76}      \DecValTok{78}  \DecValTok{69}     \DecValTok{66}    \DecValTok{81}
\DecValTok{5}      \DecValTok{46}      \DecValTok{56}  \DecValTok{52}     \DecValTok{56}    \DecValTok{40}
\DecValTok{6}      \DecValTok{12}      \DecValTok{42}  \DecValTok{38}     \DecValTok{38}    \DecValTok{28}
\DecValTok{7}      \DecValTok{39}      \DecValTok{46}  \DecValTok{51}     \DecValTok{54}    \DecValTok{41}
\DecValTok{8}      \DecValTok{30}      \DecValTok{51}  \DecValTok{54}     \DecValTok{52}    \DecValTok{32}
\DecValTok{9}      \DecValTok{22}      \DecValTok{32}  \DecValTok{43}     \DecValTok{28}    \DecValTok{22}
\NormalTok{...}
\end{Highlighting}
\end{Shaded}

\normalsize

El objetivo de este estudio es obtener un ranking global de alumnos para
la entrada en la Facultad de Matemáticas, a través de una puntuación
global, extraída como cierta combinación lineal de las calificaciones en
las cinco materias examinadas. El comando básico para realizar un análisis de componentes principales
en R es \texttt{princomp}. También se pueden obtener resultados similares con
el comando \texttt{prcomp}. La diferencia principal entre uno y otro es que
\texttt{princomp} diagonaliza la matriz \(\bm{S}\), mientras que \texttt{prcomp}
diagonaliza la matriz \(\bm{S_c}\). Esto modifica los autovalores en la
proporción \(n/(n-1)\), pero no supone ningún cambio en los autovectores.
Hemos optado por el comando \texttt{princomp}. A continuación se muestran las
instrucciones que permiten realizar un análisis de componentes
principales para el ejemplo, utilizando R. Mostramos en primer lugar la salida básica del objeto, que muestra las
desviaciones típicas de las componentes que son las raíces cuadradas de
los autovalores.

\footnotesize

\begin{Shaded}
\begin{Highlighting}[]
\SpecialCharTok{\textgreater{}}\NormalTok{ test.pca }\OtherTok{\textless{}{-}} \FunctionTok{princomp}\NormalTok{(dat)  }\CommentTok{\# Análisis de componentes principales}
\SpecialCharTok{\textgreater{}}\NormalTok{ test.pca}
\NormalTok{Call}\SpecialCharTok{:}
\FunctionTok{princomp}\NormalTok{(}\AttributeTok{x =}\NormalTok{ dat)}

\NormalTok{Standard deviations}\SpecialCharTok{:}
\NormalTok{   Comp}\FloatTok{.1}\NormalTok{    Comp}\FloatTok{.2}\NormalTok{    Comp}\FloatTok{.3}\NormalTok{    Comp}\FloatTok{.4}\NormalTok{    Comp}\FloatTok{.5} 
\FloatTok{28.489680}  \FloatTok{9.035471}  \FloatTok{6.600955}  \FloatTok{6.133582}  \FloatTok{3.723358} 

 \DecValTok{5}\NormalTok{  variables and  }\DecValTok{25}\NormalTok{ observations.}
\end{Highlighting}
\end{Shaded}

\normalsize

Con la función \texttt{summary} obtenemos, junto a las desviaciones típicas de
las componentes, la proporción de varianza explicada y sus valores
acumulados.

\footnotesize

\begin{Shaded}
\begin{Highlighting}[]
\SpecialCharTok{\textgreater{}} \FunctionTok{summary}\NormalTok{(test.pca)}
\NormalTok{Importance of components}\SpecialCharTok{:}
\NormalTok{                           Comp}\FloatTok{.1}\NormalTok{     Comp}\FloatTok{.2}\NormalTok{     Comp}\FloatTok{.3}\NormalTok{     Comp}\FloatTok{.4}\NormalTok{     Comp}\FloatTok{.5}
\NormalTok{Standard deviation     }\FloatTok{28.4896795} \FloatTok{9.03547104} \FloatTok{6.60095491} \FloatTok{6.13358179} \FloatTok{3.72335754}
\NormalTok{Proportion of Variance  }\FloatTok{0.8212222} \FloatTok{0.08260135} \FloatTok{0.04408584} \FloatTok{0.03806395} \FloatTok{0.01402668}
\NormalTok{Cumulative Proportion   }\FloatTok{0.8212222} \FloatTok{0.90382353} \FloatTok{0.94790936} \FloatTok{0.98597332} \FloatTok{1.00000000}
\SpecialCharTok{\textgreater{}}\NormalTok{ test.pca}\SpecialCharTok{$}\NormalTok{sdev}\SpecialCharTok{\^{}}\DecValTok{2}
\NormalTok{   Comp}\FloatTok{.1}\NormalTok{    Comp}\FloatTok{.2}\NormalTok{    Comp}\FloatTok{.3}\NormalTok{    Comp}\FloatTok{.4}\NormalTok{    Comp}\FloatTok{.5} 
\FloatTok{811.66184}  \FloatTok{81.63974}  \FloatTok{43.57261}  \FloatTok{37.62083}  \FloatTok{13.86339} 
\end{Highlighting}
\end{Shaded}

\normalsize

Representamos a continuación los autovalores (gráfico de sedimentación)
y el mismo gráfico con líneas, ver Figura
\ref{fig:sed}.

\footnotesize

\begin{Shaded}
\begin{Highlighting}[]
\SpecialCharTok{\textgreater{}} \FunctionTok{screeplot}\NormalTok{(test.pca)}
\SpecialCharTok{\textgreater{}} \FunctionTok{screeplot}\NormalTok{(test.pca, }\AttributeTok{type =} \StringTok{"lines"}\NormalTok{)}
\end{Highlighting}
\end{Shaded}

\begin{figure}

{\centering \resizebox{.4\linewidth}{!}{\includegraphics[width=0.5\linewidth]{bookdown_AM_files/figure-latex/sed-1} } \resizebox{.4\linewidth}{!}{\includegraphics[width=0.5\linewidth]{bookdown_AM_files/figure-latex/sed-2} } 

}

\caption{Gráficos de sedimentación.}\label{fig:sed}
\end{figure}
\normalsize

La función \texttt{loadings} proporciona los coeficientes de las componentes (autovectores), ver Figura
\ref{fig:scores}

\footnotesize

\begin{Shaded}
\begin{Highlighting}[]
\SpecialCharTok{\textgreater{}} \FunctionTok{loadings}\NormalTok{(test.pca)}

\NormalTok{Loadings}\SpecialCharTok{:}
\NormalTok{        Comp}\FloatTok{.1}\NormalTok{ Comp}\FloatTok{.2}\NormalTok{ Comp}\FloatTok{.3}\NormalTok{ Comp}\FloatTok{.4}\NormalTok{ Comp}\FloatTok{.5}
\NormalTok{geodif   }\FloatTok{0.598}  \FloatTok{0.675}  \FloatTok{0.185}  \FloatTok{0.386}       
\NormalTok{ancompl  }\FloatTok{0.361}  \FloatTok{0.245} \SpecialCharTok{{-}}\FloatTok{0.249} \SpecialCharTok{{-}}\FloatTok{0.829} \SpecialCharTok{{-}}\FloatTok{0.247}
\NormalTok{alg      }\FloatTok{0.302} \SpecialCharTok{{-}}\FloatTok{0.214} \SpecialCharTok{{-}}\FloatTok{0.211} \SpecialCharTok{{-}}\FloatTok{0.135}  \FloatTok{0.894}
\NormalTok{anreal   }\FloatTok{0.389} \SpecialCharTok{{-}}\FloatTok{0.338} \SpecialCharTok{{-}}\FloatTok{0.700}  \FloatTok{0.375} \SpecialCharTok{{-}}\FloatTok{0.321}
\NormalTok{estad    }\FloatTok{0.519} \SpecialCharTok{{-}}\FloatTok{0.570}  \FloatTok{0.607}        \SpecialCharTok{{-}}\FloatTok{0.179}

\NormalTok{               Comp}\FloatTok{.1}\NormalTok{ Comp}\FloatTok{.2}\NormalTok{ Comp}\FloatTok{.3}\NormalTok{ Comp}\FloatTok{.4}\NormalTok{ Comp}\FloatTok{.5}
\NormalTok{SS loadings       }\FloatTok{1.0}    \FloatTok{1.0}    \FloatTok{1.0}    \FloatTok{1.0}    \FloatTok{1.0}
\NormalTok{Proportion Var    }\FloatTok{0.2}    \FloatTok{0.2}    \FloatTok{0.2}    \FloatTok{0.2}    \FloatTok{0.2}
\NormalTok{Cumulative Var    }\FloatTok{0.2}    \FloatTok{0.4}    \FloatTok{0.6}    \FloatTok{0.8}    \FloatTok{1.0}
\end{Highlighting}
\end{Shaded}

\normalsize

\footnotesize

\begin{Shaded}
\begin{Highlighting}[]
\SpecialCharTok{\textgreater{}} \FunctionTok{barplot}\NormalTok{(}\FunctionTok{loadings}\NormalTok{(test.pca), }\AttributeTok{beside =} \ConstantTok{TRUE}\NormalTok{)  }\CommentTok{\# Representación coeficientes}
\SpecialCharTok{\textgreater{}}\NormalTok{ test.pca}\SpecialCharTok{$}\NormalTok{scores  }\CommentTok{\# Puntuaciones de los individuos en las componentes}
\NormalTok{           Comp}\FloatTok{.1}\NormalTok{      Comp}\FloatTok{.2}\NormalTok{      Comp}\FloatTok{.3}\NormalTok{      Comp}\FloatTok{.4}\NormalTok{     Comp}\FloatTok{.5}
\NormalTok{ [}\DecValTok{1}\NormalTok{,]  }\SpecialCharTok{{-}}\FloatTok{7.5403215}  \FloatTok{10.2167650}   \FloatTok{2.5374713}  \SpecialCharTok{{-}}\FloatTok{8.6708997} \SpecialCharTok{{-}}\FloatTok{4.3011164}
\NormalTok{ [}\DecValTok{2}\NormalTok{,]  }\FloatTok{20.3610372}  \FloatTok{13.3460340}   \FloatTok{8.9820585}   \FloatTok{6.4124949} \SpecialCharTok{{-}}\FloatTok{1.3548711}
\NormalTok{ [}\DecValTok{3}\NormalTok{,] }\SpecialCharTok{{-}}\FloatTok{19.5031539}   \FloatTok{6.5552439}  \SpecialCharTok{{-}}\FloatTok{1.6414327}   \FloatTok{5.0042015} \SpecialCharTok{{-}}\FloatTok{2.2980498}
\NormalTok{ [}\DecValTok{4}\NormalTok{,]  }\FloatTok{65.9652730}   \FloatTok{1.3136646}   \FloatTok{5.1988159}  \SpecialCharTok{{-}}\FloatTok{5.2093505} \SpecialCharTok{{-}}\FloatTok{1.0457668}
\NormalTok{ [}\DecValTok{5}\NormalTok{,]   }\FloatTok{9.7780565}   \FloatTok{6.0680143}  \SpecialCharTok{{-}}\FloatTok{9.1921582}   \FloatTok{2.9249303} \SpecialCharTok{{-}}\FloatTok{2.1069944}
\NormalTok{ [}\DecValTok{6}\NormalTok{,] }\SpecialCharTok{{-}}\FloatTok{33.0739529}  \SpecialCharTok{{-}}\FloatTok{4.3722312}  \SpecialCharTok{{-}}\FloatTok{3.7345190}  \SpecialCharTok{{-}}\FloatTok{2.6038323} \SpecialCharTok{{-}}\FloatTok{5.3246839}
\NormalTok{ [}\DecValTok{7}\NormalTok{,]   }\FloatTok{1.4212177}  \SpecialCharTok{{-}}\FloatTok{0.7833317}  \SpecialCharTok{{-}}\FloatTok{5.7800406}   \FloatTok{7.8225646} \SpecialCharTok{{-}}\FloatTok{0.4967347}
\NormalTok{ [}\DecValTok{8}\NormalTok{,]  }\SpecialCharTok{{-}}\FloatTok{6.7011638}  \SpecialCharTok{{-}}\FloatTok{0.4667798} \SpecialCharTok{{-}}\FloatTok{13.3936497}  \SpecialCharTok{{-}}\FloatTok{0.3040531}  \FloatTok{2.6525120}
\NormalTok{ [}\DecValTok{9}\NormalTok{,] }\SpecialCharTok{{-}}\FloatTok{36.1916160}   \FloatTok{5.6543609}   \FloatTok{2.9062814}   \FloatTok{5.5458471}  \FloatTok{6.5171961}
\NormalTok{...}
\end{Highlighting}
\end{Shaded}

\begin{figure}

{\centering \resizebox{.4\linewidth}{!}{\includegraphics[width=0.5\linewidth]{bookdown_AM_files/figure-latex/scores-1} } 

}

\caption{Representación de los coeficientes de las componentes principales.}\label{fig:scores}
\end{figure}
\normalsize

El \emph{biplot}, introducido por \citet{gabriel71}, es una representación gráfica simultánea de los individuos
(mediante puntos) y las variables (mediante flechas), en un mismo
sistema de coordenadas bidimensional construido en base a las dos
primeras componentes principales. Permite interpretar el significado de
las componentes (la primera en el eje horizontal y la segunda en el eje
vertical) en base a las direcciones de las flechas.

A su vez, se valoran como parecidos los individuos cuyos puntos están
próximos en el biplot. De igual modo tendrán correlación positiva las
variables con flechas semejantes. Asimismo, los individuos que se
encuentran en la dirección de cierta flecha tendrán observaciones altas
en la variable representada por la flecha.

Representamos el biplot correspondiente a los datos del ejemplo, ver Figura \ref{fig:bip}

\footnotesize

\begin{Shaded}
\begin{Highlighting}[]
\SpecialCharTok{\textgreater{}} \FunctionTok{biplot}\NormalTok{(test.pca)}
\end{Highlighting}
\end{Shaded}

\begin{figure}

{\centering \resizebox{.4\linewidth}{!}{\includegraphics[width=0.5\linewidth]{bookdown_AM_files/figure-latex/bip-1} } 

}

\caption{Biplot}\label{fig:bip}
\end{figure}
\normalsize

Obtenemos los mismos resultados haciendo los cálculos de forma directa:

\footnotesize

\begin{Shaded}
\begin{Highlighting}[]
\SpecialCharTok{\textgreater{}}\NormalTok{ n }\OtherTok{\textless{}{-}} \FunctionTok{nrow}\NormalTok{(dat)}
\SpecialCharTok{\textgreater{}}\NormalTok{ s }\OtherTok{\textless{}{-}} \FunctionTok{cov}\NormalTok{(dat) }\SpecialCharTok{*}\NormalTok{ (n }\SpecialCharTok{{-}} \DecValTok{1}\NormalTok{)}\SpecialCharTok{/}\NormalTok{n}
\SpecialCharTok{\textgreater{}}\NormalTok{ auto }\OtherTok{\textless{}{-}} \FunctionTok{eigen}\NormalTok{(s)}
\SpecialCharTok{\textgreater{}}\NormalTok{ lambda }\OtherTok{\textless{}{-}}\NormalTok{ auto}\SpecialCharTok{$}\NormalTok{values}
\SpecialCharTok{\textgreater{}}\NormalTok{ v }\OtherTok{\textless{}{-}}\NormalTok{ auto}\SpecialCharTok{$}\NormalTok{vectors}
\end{Highlighting}
\end{Shaded}

\normalsize

\section{Propiedades geométricas de las componentes principales}\label{propiedades-geometricas-de-las-componentes-principales}

Esta sección contiene dos propiedades geométricas interesantes, que
cumplen las componentes principales, cuando se comparan con el modelo de
regresión. En primer lugar veremos que las componentes principales
también son aproximaciones mínimo-cuadráticas del vector aleatorio, y en
segundo lugar veremos que en el caso de un vector bivariante
\(\bm{x}=(x_1,x_2)^\prime{}\), la primera componente principal describe
una recta que pasa entre las dos rectas de regresión lineal simple, una
de \(x_2\) sobre \(x_1\) y la otra de \(x_1\) sobre \(x_2\).

\subsection{Las componentes principales como mejores aproximaciones mínimo-cuadráticas}\label{las-componentes-principales-como-mejores-aproximaciones-minimo-cuadraticas}

Recordemos que la componente principal es la variable aleatoria
unidimensional \(z_1=\bm{v_1}^{\prime} \bm{x}\). Al proyectar el vector
\(\bm{x}\) sobre la recta que pasa por el origen con la dirección
\(\bm{v_1}\), se obtiene el vector aleatorio

\[\bm{z_1^m}=\bm{v_1} z_1=\bm{v_1} \bm{v_1}^{\prime} \bm{x}\] que
coincide con la representación de la primera componente \(z_1\) en
\({\Bbb R}^d\) a lo largo de la dirección \(\bm{v_1}\).

Observamos que este vector aleatorio está contenido en un espacio de
dimensión uno. De hecho, su matriz de covarianzas es

\[\textnormal{Cov}(\bm{z_1^m},\bm{z_1^m})=\textnormal{Cov}(\bm{v_1} z_1, \bm{v_1} z_1)=\bm{v_1} \textnormal{Cov}(z_1,z_1) \bm{v_1}^{\prime}= \textnormal{Var}(z_1) \bm{v_1} \bm{v_1}^{\prime}
=\lambda_1 \bm{v_1} \bm{v_1}^{\prime}\] la cual claramente es una matriz
de rango uno y no es más que el primer sumando de la descomposición
espectral de \(\bm\Sigma\).

Centramos el vector \(\bm{z_1^m}\) en el vector de medias de \(\bm{x}\),
\(\bm\mu=E(\bm{x})\).

\[\bm{z_1^{mc}}=\bm{z_1^m}-E(\bm{z_1^m})+E(\bm{x})=\bm\mu + \bm{v_1} z_1^c\]
siendo \(z_1^c=\bm{v_1}^{\prime} (\bm{x}-\bm\mu)\) la primera componente
centrada.

Este vector aleatorio \(\bm{z_1^{mc}}\) es el resultado de proyectar el
vector \(\bm{x}\) sobre la recta que pasa por el vector de medias \(\bm\mu\)
y tiene dirección \(\bm{v_1}\).

\begin{theorem}
\protect\hypertarget{thm:unnamed-chunk-84}{}\label{thm:unnamed-chunk-84}La recta que pasa por el vector de medias \(\bm\mu=E(\bm{x})\) y tiene
dirección \(\bm{v_1}\) es la que hace mínima la distancia cuadrática del
vector aleatorio \(\bm{x}\) a la recta, definiendo esta distancia
cuadrática como la media del cuadrado de la distancia euclídea usual
(distancia ortogonal o menor distancia) del vector aleatorio \(\bm{x}\) a
la recta.
\end{theorem}

\textbf{Observación}. Este resultado debe relacionarse con la recta de
regresión, que hace mínima la media (suma en el caso de una muestra) de
la distancia vertical a la recta. En este caso consideramos la distancia
ortogonal (la menor distancia), pues no hay variable respuesta ni
variable explicativa, sino que todas las variables son tratadas de igual
modo.

\subsection{Posición en el plano de la 1ª componente principal respecto a las rectas de regresión}\label{posicion-en-el-plano-de-la-1a-componente-principal-respecto-a-las-rectas-de-regresion}

Sea \((x,y)\) un vector aleatorio en el plano. Se puede demostrar que al
representar en el plano \(XY\) la primera componente principal, junto con
las dos rectas de regresión de \(y\) sobre \(x\) y de \(x\) sobre \(y\),
obtenemos tres rectas que pasan por el vector de medias, y la primera
componente principal forma una recta ``menos horizontal'' que la recta de
regresión de \(y\) sobre \(x\), pero "más horizontal" que la recta de
regresión de \(x\) sobre \(y\).

\textbf{\includegraphics[width=1em,height=1em]{bookdown_AM_files/figure-latex/fa-icon-00ef1bb1d064b635425d49f06b331298.pdf} ~\textcolor{excolor}{Ejemplo.}} Se ha obtenido la muestra siguiente de un vector aleatorio
\((x,y)\): (-3,7), (1,4), (0,6), (3,5), (1,9), (7,7), (6,9), (8,10),
(7,12), (10,11). Representamos en la Figura \ref{fig:reg} el diagrama de dispersión, junto con la primera
componente principal y las dos rectas de regresión.

\begin{figure}

{\centering \resizebox{.4\linewidth}{!}{\includegraphics[width=0.5\linewidth]{bookdown_AM_files/figure-latex/reg-1} } 

}

\caption{Diagrama de dispersión, junto con la primera componente principal (verde), la recta de regresión de $y$ sobre $x$ (azul) y la recta de regresión $x$ de sobre $y$ (rojo)}\label{fig:reg}
\end{figure}

\section{Descomposición de la variabilidad, proporción de varianza explicada y criterios para la reducción de dimensión}\label{descomposicion-de-la-variabilidad-proporcion-de-varianza-explicada-y-criterios-para-la-reduccion-de-dimension}

Según acabamos de ver, desde un punto de vista matricial podemos
contemplar el proceso de extracción de las componentes principales como
la diagonalización de la matriz de covarianzas

\[\bm\Lambda=\bm{V}^{\prime} \bm\Sigma \bm{V}\] siendo
\(\bm\Sigma=\textnormal{Cov}(\bm{x},\bm{x})\) la matriz de covarianzas del
vector aleatorio original, \(\bm\Lambda=\textnormal{Cov}(\bm{z},\bm{z})\)
la matriz de covarianzas de las componentes (que es diagonal) y
\(\bm{V}=(\bm{v_1}\ldots,\bm{v_d})\) la matriz ortogonal cuyas columnas
son los autovectores, que a su vez constituyen los coeficientes que
definen las componentes.

Multiplicando por \(\bm{V}\) por la izquierda y por \(\bm{V}^{\prime}\) por
la derecha, podemos expresar

\[\bm\Sigma=\bm{V}\bm\Lambda \bm{V}^{\prime}=\sum_{j=1}^d \lambda_j \bm{v_j}\bm{v_j}^{\prime}\]
que se conoce como la descomposición espectral de la matriz \(\bm\Sigma\). Sabemos que

\[\textnormal{rango}(\bm\Sigma)=\textnormal{rango}(\bm\Lambda)=\mbox{``número de autovalores no nulos''}.\]
A su vez también recordamos que el rango de la matriz de covarianzas
coincide con la dimensión del menor espacio lineal que contiene al
vector aleatorio (con probabilidad uno).

Así si el rango es \(k\), el vector aleatorio estaría contenido en un
espacio lineal de dimensión \(k\), los \(k\) primeros autovalores son no
nulos y los \((d-k)\) últimos son cero, y las \((d-k)\) últimas componentes
tienen varianza cero. Se puede por tanto reducir la dimensión hasta \(k\)
sin perder variabilidad. Sólo estaríamos cambiando el sistema de
coordenadas.

Consideremos como medida global de variabilidad del vector aleatorio, la
traza de su matriz de covarianzas. Nótese que no es más que la suma de
las varianzas de las variables del vector, pero además tiene esta
justificación como medida global de variabilidad

\[\textnormal{traza}(\bm\Sigma)=E(\| \bm{x}-E(\bm{x})\|^2).\] Es la
media de las desviaciones cuadráticas respecto a la media y por tanto es
una extensión natural de la varianza de una variable aleatoria.

De la descomposición espectral se deduce

\[\textnormal{traza}(\bm\Sigma)=\sum_{j=1}^d \lambda_j \textnormal{traza}(\bm{v_j}\bm{v_j}^{\prime})=\sum_{j=1}^d \lambda_j.\]
La última igualdad es consecuencia de que los vectores \(\bm{v_j}\) son de
norma uno.

Si extraemos una componente principal, obtenemos la variable aleatoria
unidimensional

\[z_1=\bm{v_1}^{\prime} \bm{x}.\] Si en lugar de recoger
todo el vector aleatorio, sólo aportamos \(z_1\) reduciendo la información
a una variable unidimensional; junto con la simplificación, se produce
una reducción de variabilidad, que pasa a ser

\[\textnormal{Var}(z_1)=\lambda_1.\] Decimos que el cociente

\[\frac{\lambda_1}{\lambda_1+\cdots+\lambda_d}\] es la proporción de
variabilidad explicada por la primera componente principal.

Si en lugar de una única componente principal extraemos \(r\) componentes
resulta

\[\frac{\lambda_1+\cdots+\lambda_r}{\lambda_1+\cdots+\lambda_r+\lambda_{r+1}
+\cdots+\lambda_d}\] es la proporción de variabilidad explicada por las
\(r\) primeras componentes principales.

Debemos decidir entre la simplificación que supone la reducción de la
dimensión y la pérdida de información resultante de la variabilidad no
explicada. Como criterios para tomar esta decisión, se suelen emplear
los siguientes:

\begin{itemize}
\item
  Criterio de la varianza explicada. Consiste en retener el número de
  componentes que conjuntamente expliquen una proporción de varianza
  establecida, habitualmente un 90\% o 95\% del total.
\item
  Criterio del valor propio (Kaiser). Retener sólo las componentes
  cuyos autovalores sean mayores que la media, esto es, mayores que
  \(\textnormal{traza}(\bm\Sigma)/d\). Estas componentes tendrían más
  varianza que la que corresponde en media a una variable, mientras
  que las demás componentes tienen menos.
\item
  Gráfico de sedimentación (\emph{screeplot}). Representar en un gráfico
  los valores propios \(\lambda_1\geq \lambda_2\geq\cdots\geq\lambda_d\)
  en orden decreciente, y buscar un ``codo'' en el gráfico, entendiendo
  por codo un punto a partir del cual los valores propios son
  claramente más pequeños que los anteriores, y muy similares entre
  sí.
\item
  Retener un número preestablecido de componentes principales. Por
  ejemplo, es costumbre retener dos componentes, pues se pueden
  representar fácilmente en el plano.
\end{itemize}

\section{Las componentes principales y los cambios de escala}\label{las-componentes-principales-y-los-cambios-de-escala}

Un problema importante del análisis de componentes principales es que
sus resultados dependen de la escala de medida de las variables
originales, \(\bm{x}\). Así, si se cambia de escala una variable (dejando
las demás fijas) las componentes principales se modifican, cambiando
incluso la posible interpretación de las mismas. En concreto, si se
aumenta la escala de una variable, ésta verá incrementada su varianza y
su aportación proporcional a la varianza total, atrayendo de este modo a
la primera componente principal, que se asemejará a esta variable.

Claro está que si se realiza el mismo cambio de escala en todas las
variables, entonces los resultados del análisis de componentes
principales se mantendrán idénticos.

Este problema se puede solventar de dos maneras:

\begin{itemize}
\item
  Midiendo todas las variables en la misma escala (siempre que sean de
  la misma naturaleza).
\item
  Aplicando el análisis de componentes principales a las variables
  estandarizadas. Esto último equivale a trabajar con la matriz de
  correlaciones, en lugar de la matriz de covarianzas.
\end{itemize}

\section{Interpretación de las componentes principales}\label{interpretacion-de-las-componentes-principales}

Para interpretar las componentes principales,
\(\bm{z}=(z_1,\ldots,z_d)^\prime{}\), es preciso relacionarlas con las
variables originales, \(\bm{x}\). El instrumento natural para establecer
esta relación consiste en calcular la covarianza o la correlación entre
las componentes y las variables originales \(\bm{x}\). Lo obtenemos a
continuación.

\[\textnormal{Cov}(\bm{x},z_k)=\textnormal{Cov}(\bm{x}, \bm{v_k}^{\prime}\bm{x})=\textnormal{Cov}(\bm{x},\bm{x})\bm{v_k}=\bm\Sigma \bm{v_k}=\lambda_k \bm{v_k}.\]

La notación matricial nos permite calcularlas para todas a la vez

\[\textnormal{Cov}(\bm{x},\bm{z})=\textnormal{Cov}(\bm{x}, \bm{V}^{\prime}\bm{x})=\textnormal{Cov}(\bm{x},\bm{x})\bm{V}=\bm\Sigma \bm{V}=\bm{V}\bm\Lambda.\]

En particular para una variable \(x_j\) y una componente \(z_k\):

\[\textnormal{Cov}(x_j,z_k)=\lambda_k v_{jk}.\]

Se trata por tanto de la varianza de la componente \(z_k\) multiplicada
por el coeficiente \(v_{jk}\), que acompaña a \(x_j\) en la construcción de
\(z_k\). El coeficiente de correlación será

\[\textnormal{Corr}(x_j,z_k)=\frac{\textnormal{Cov}(x_j,z_k)}{\sqrt{\textnormal{Var}(x_j) \textnormal{Var}(z_k)}}
=\frac{v_{jk}}{\sigma_j}\sqrt{\lambda_k}.\]

En notación matricial,

\[\begin{aligned}
\textnormal{Corr}(\bm{x},\bm{z})&=\mbox{diag}\left(1/\sigma_1,\ldots,1/\sigma_d\right) \bm{V} \bm\Lambda
\mbox{diag}\left(1/\sqrt{\lambda_1},\ldots, 1/\sqrt{\lambda_d}\right)\\
&=
\mbox{diag}(1/\sigma_1,\ldots,1/\sigma_d) \bm{V} \bm\Lambda^{1/2}.
\end{aligned}\]

Se interpretará cada componente principal en función de las variables
originales que estén correlacionadas con dicha componente. Asimismo, el
signo de cada correlación permite valorar qué variables van en el
sentido de la componente y cuáles en sentido opuesto.

\chapter{Análisis de correspondencias}\label{analisis-de-correspondencias}

El Análisis de Correspondencias estudia la relación entre dos variables
discretas a través de su distribución conjunta de frecuencias. Nos
apoyaremos en el siguiente ejemplo.

\textbf{\includegraphics[width=1em,height=1em]{bookdown_AM_files/figure-latex/fa-icon-00ef1bb1d064b635425d49f06b331298.pdf} ~\textcolor{excolor}{Ejemplo.}} \emph{Encuesta de Población Activa (EPA, 2010). Datos de
actividad y ocupación}. La Encuesta de Población Activa, comúnmente
conocida como EPA, es una encuesta que realiza el Instituto Nacional de
Estadística (INE) y que recoge información sobre el empleo en España.
Así, de esta encuesta se extraen los datos de actividad, ocupación y
tasa de paro, que además se pueden relacionar con características de los
individuos, como su sexo, edad o comunidad autónoma de residencia.

Vamos a estudiar la relación que puede haber entre la comunidad autónoma
y una variable que indicará si el individuo es inactivo, está ocupado o
está parado. Se define como inactivo a un individuo que no realiza
trabajo remunerado, ni está en condiciones de realizarlo, bien por falta
de capacidad (niños, ancianos o discapacitados) o por decisión propia.
Ocupado es aquél que realiza trabajo remunerado, y parado es aquél que
no lo realiza, aunque se encuentra en condiciones de trabajar y busca
empleo. Los datos del fichero \texttt{epa2010.csv} recogen las frecuencias (en miles de
personas) en cada comunidad autónomica y según su actividad (ocupado,
parado o inactivo).

\footnotesize

\begin{Shaded}
\begin{Highlighting}[]
\SpecialCharTok{\textgreater{}}\NormalTok{ dat }\OtherTok{\textless{}{-}} \FunctionTok{read.csv}\NormalTok{(}\StringTok{"data/epa2010.csv"}\NormalTok{, }\AttributeTok{sep =} \StringTok{";"}\NormalTok{, }\AttributeTok{dec =} \StringTok{","}\NormalTok{, }\AttributeTok{fileEncoding =} \StringTok{"latin1"}\NormalTok{)}
\SpecialCharTok{\textgreater{}}\NormalTok{ dat}
\NormalTok{            Comunidad.Autónoma Ocupados Parados Inactivos}
\DecValTok{1}\NormalTok{                    Andalucía   }\FloatTok{2849.1}  \FloatTok{1127.4}    \FloatTok{2794.2}
\DecValTok{2}\NormalTok{                       Aragón    }\FloatTok{540.3}   \FloatTok{103.3}     \FloatTok{467.5}
\DecValTok{3}\NormalTok{      Asturias, Principado de    }\FloatTok{398.0}    \FloatTok{79.6}     \FloatTok{451.8}
\DecValTok{4}\NormalTok{               Balears, Illes    }\FloatTok{450.0}   \FloatTok{128.6}     \FloatTok{320.3}
\DecValTok{5}\NormalTok{                     Canarias    }\FloatTok{771.2}   \FloatTok{314.4}     \FloatTok{663.8}
\DecValTok{6}\NormalTok{                    Cantabria    }\FloatTok{233.8}    \FloatTok{41.0}     \FloatTok{219.6}
\DecValTok{7}\NormalTok{              Castilla y León    }\FloatTok{996.8}   \FloatTok{186.8}     \FloatTok{962.4}
\DecValTok{8}\NormalTok{         Castilla }\SpecialCharTok{{-}}\NormalTok{ La Mancha    }\FloatTok{779.2}   \FloatTok{211.2}     \FloatTok{712.9}
\DecValTok{9}\NormalTok{                     Cataluña   }\FloatTok{3133.5}   \FloatTok{686.8}    \FloatTok{2256.8}
\NormalTok{...}
\end{Highlighting}
\end{Shaded}

\normalsize

Mediante un análisis de correspondencias seremos capaces de describir
las relaciones que pueda haber entre la comunidad autónoma y la
actividad, por ejemplo, indicando que en cierta comunidad autónoma hay
un mayor porcentaje de parados, que en otra abundan especialmente los
inactivos, o que en una tercera predominan los ocupados.

Las tablas de frecuencias, como la vista en el ejemplo, se conocen como
\textbf{tablas de contingencia}. Reflejan la distribución conjunta del par
formado por las dos variables discretas, que denotaremos \((x,y)\).
Introduciendo cierta notación, podríamos expresar la tabla de
contingencia así:

\[\begin{array}{c|ccccc|c}
x\backslash y & y_1 & \ldots & y_j & \ldots & y_s & \\
\hline x_1 & n_{11} & \cdots & n_{1j} & \cdots & n_{1s} & n_{1\bullet} \\
\vdots & \vdots & & \vdots & & \vdots &\vdots \\
x_i & n_{i1} & \cdots & n_{ij} & \cdots & n_{is} & n_{i\bullet} \\
\vdots & \vdots & & \vdots & & \vdots &\vdots \\
x_r & n_{r1} & \cdots & n_{rj} & \cdots & n_{rs} & n_{r\bullet} \\
\hline & n_{\bullet 1} & \cdots & n_{\bullet j} & \cdots & n_{\bullet s} & n \\
\end{array}\]

siendo \(n_{ij}\) la frecuencia absoluta del par \((x_i,y_j)\),
\((n_{1\bullet},\ldots,n_{i\bullet},\ldots,n_{r\bullet})\) la distribución
marginal de la variable \(x\) y
\((n_{\bullet 1},\ldots,n_{\bullet j},\ldots,n_{\bullet s})\) la
distribución marginal de la variable \(y\).

Si dividimos entre el tamaño muestral, \(n\), obtenemos una tabla similar,
pero ahora en términos relativos.

\[\begin{array}{c|ccccc|c}
x\backslash y & y_1 & \ldots & y_j & \ldots & y_s & \\
\hline x_1 & f_{11} & \cdots & f_{1j} & \cdots & f_{1s} & f_{1\bullet} \\
\vdots & \vdots & & \vdots & & \vdots &\vdots \\
x_i & f_{i1} & \cdots & f_{ij} & \cdots & f_{is} & f_{i\bullet} \\
\vdots & \vdots & & \vdots & & \vdots &\vdots \\
x_r & f_{r1} & \cdots & f_{rj} & \cdots & f_{rs} & f_{r\bullet} \\
\hline & f_{\bullet 1} & \cdots & f_{\bullet j} & \cdots & f_{\bullet s} & 1 \\
\end{array}\]

siendo \(f_{ij}\) la frecuencia relativa del par \((x_i,y_j)\),
\((f_{1\bullet},\ldots,f_{i\bullet},\ldots,f_{r\bullet})\) la distribución
marginal de la variable \(x\) y
\((f_{\bullet 1},\ldots,f_{\bullet j},\ldots,f_{\bullet s})\) la
distribución marginal de la variable \(y\).

En primer lugar recordaremos los conceptos de independencia y asociación
(lo contrario a la independencia) entre dos variables discretas. Veremos
que el estadístico ji-cuadrado, que se usa habitualmente para contrastar
la independencia, se puede interpretar como una medida de asociación. En
este sentido, veremos que refleja la variabilidad entre distribuciones
condicionadas.

A continuación abordaremos la representación de la asociación
mediante un procedimiento de componentes principales y la representación simultánea de las categorías de las dos
variables discretas, lo cual culmina los objetivos del análisis de
correspondencias, pues en este gráfico simultáneo identificaremos las
asociaciones entre categorías.

\section{Test ji-cuadrado de independencia y distancia ji-cuadrado entre perfiles}\label{test-ji-cuadrado-de-independencia-y-distancia-ji-cuadrado-entre-perfiles}

Si el análisis de correspondencias pretende estudiar la relación entre
las dos variables, la situación extrema de no relación se producirá
cuando las dos variables sean independientes. Esto es equivalente a que
la distribución conjunta resulte del producto de las marginales. En
términos de las frecuencias observadas, ese producto de marginales da
lugar a lo que se suele conocer como frecuencias esperadas, que son:

\[n_{i\bullet}n_{\bullet j}/n.\]

A partir de aquí se construye el estadístico ji-cuadrado para contrastar
la independencia comparando las frecuencias observadas con las
esperadas:

\[\chi^2=\sum_{i=1}^r\sum_{j=1}^s
\frac{\left(n_{ij}-n_{i\bullet}n_{\bullet j}/n\right)^2}{n_{i\bullet}n_{\bullet j}/n}.\]

Bajo la hipótesis nula de independencia, el estadístico ji-cuadrado
tiene una distribución ji-cuadrado con \((r-1)\times (s-1)\) grados de
libertad.

Además del contraste de independencia, podemos pensar que el estadístico
ji-cuadrado es una medida de discrepancia respecto de la independencia,
y por tanto es una medida de asociación entre las categorías de ambas
variables. Con este fin, vamos a expresar el estadístico ji--cuadrado en
términos de frecuencias relativas:

\[\chi^2=\sum_{i=1}^r\sum_{j=1}^s
\frac{\left(f_{ij}n-f_{i\bullet}f_{\bullet j}n\right)^2}{f_{i\bullet}f_{\bullet j}n}
=n\sum_{i=1}^r\sum_{j=1}^s
\frac{\left(f_{ij}-f_{i\bullet}f_{\bullet j}\right)^2}{f_{i\bullet}f_{\bullet j}}.\]

Ahora observamos que el estadístico ji-cuadrado, representado en cierta
forma, admite una interpretación interesante:

\[\frac{\chi^2}{n}=\sum_{i=1}^r\sum_{j=1}^s
\frac{\left(f_{i\bullet}\frac{f_{ij}}{f_{i\bullet}}
-f_{i\bullet}f_{\bullet j}\right)^2}{f_{i\bullet}f_{\bullet j}}
=\sum_{i=1}^r f_{i\bullet} \sum_{j=1}^s
\left(\frac{f_{ij}}{f_{i\bullet}}-f_{\bullet j}\right)^2
\frac{1}{f_{\bullet j}}=\sum_{i=1}^r f_{i\bullet} d^2(\bm{r_i},\bm{m_y})\]

siendo
\(\bm{r_i}=\left(\frac{f_{i1}}{f_{i\bullet}},\ldots,\frac{f_{is}}{f_{i\bullet}}\right)^\prime{}\)
la distribución de \(y\) condicionada a \(x=x_i\), también conocida como
perfil de fila \(i\)-ésimo;
\(\bm{m_y}=\left(f_{\bullet 1},\ldots,f_{\bullet s}\right)^\prime{}\) la
distribución marginal de \(y\); y \(d^2(\bm{r_i},\bm{m_y})\) una distancia
entre el perfil \(i\)-ésimo y la distribución marginal. A esta distancia
se la conoce como \textbf{distancia ji-cuadrado} y a \(\chi^2/n\) como
\textbf{inercia}.

La tabla siguiente contiene los perfiles de fila, junto con la
distribución marginal de \(y\):

\[\begin{array}{c|ccccc|c}
x\backslash y & y_1 & \ldots & y_j & \ldots & y_s & \\
\hline x_1 & f_{11}/f_{1\bullet} & \cdots & f_{1j}/f_{1\bullet} & \cdots & f_{1s}/f_{1\bullet} & 1 \\
\vdots & \cdots & \cdots & \cdots & \cdots & \cdots & \vdots \\
x_i & f_{i1}/f_{i\bullet} & \cdots & f_{ij}/f_{i\bullet} & \cdots & f_{is}/f_{i\bullet} & 1 \\
\vdots & \cdots & \cdots & \cdots & \cdots & \cdots & \vdots \\
x_r & f_{r1}/f_{r\bullet} & \cdots & f_{rj}/f_{r\bullet} & \cdots & f_{rs}/f_{r\bullet} & 1 \\
\hline & f_{\bullet 1} & \cdots & f_{\bullet j} & \cdots & f_{\bullet s} & 1 \\
\end{array}\]

De este modo, el estadístico ji-cuadrado (dividido entre n) es una media
ponderada de las distancias ji-cuadrado de los perfiles de fila a la
distribución marginal de \(y\). Si observamos que la distribución marginal
se puede obtener como media de los perfiles de fila, veremos el
estadístico ji-cuadrado como una medida de variabilidad de los perfiles
en torno a su media, \(\bm{m_y}\). Bajo hipótesis de independencia, los
perfiles han de ser muy similares entre sí y, por tanto, muy similares a
la distribución marginal, dando lugar a un valor pequeño de la
variabilidad, medida por el estadístico ji-cuadrado.

De igual modo se puede interpretar el estadístico ji-cuadrado en
términos de distancias ji-cuadrado entre perfiles de columna:

\[\frac{\chi^2}{n}=\sum_{i=1}^r\sum_{j=1}^s
\frac{\left(f_{\bullet j}\frac{f_{ij}}{f_{\bullet j}}
-f_{i\bullet}f_{\bullet j}\right)^2}{f_{i\bullet}f_{\bullet j}}
=\sum_{j=1}^s f_{\bullet j} \sum_{i=1}^r
\left(\frac{f_{ij}}{f_{\bullet j}}-f_{i \bullet}\right)^2
\frac{1}{f_{i \bullet}}=\sum_{j=1}^s f_{\bullet j} d^2(\bm{c_j},\bm{m_x})\]

siendo
\(\bm{c_j}=\left(\frac{f_{1j}}{f_{\bullet j}},\ldots,\frac{f_{rj}}{f_{\bullet j}}\right)^\prime{}\)
la distribución de \(x\) condicionada a \(y=y_j\), también conocida como
perfil de columna \(j\)-ésimo;
\(\bm{m_x}=\left(f_{1 \bullet},\ldots,f_{r \bullet}\right)^\prime{}\) la
distribución marginal de \(x\); y \(d^2(\bm{c_j},\bm{m_x})\) la distancia
ji-cuadrado entre el perfil \(j\)-ésimo y la distribución marginal.

La tabla siguiente contiene los perfiles de columna, junto con la
distribución marginal de \(x\):

\[\begin{array}{c|ccccc|c}
x\backslash y & y_1 & \ldots & y_j & \ldots & y_s & \\
\hline x_1 & f_{11}/f_{\bullet 1} & \vdots & f_{1j}/f_{\bullet j} & \vdots & f_{1s}/f_{\bullet s} & f_{1 \bullet} \\
\vdots & \vdots & \vdots & \vdots & \vdots & \vdots & \vdots \\
x_i & f_{i1}/f_{\bullet 1} & \vdots & f_{ij}/f_{\bullet j} & \vdots & f_{is}/f_{\bullet s} & f_{i \bullet} \\
\vdots & \vdots & \vdots & \vdots & \vdots & \vdots & \vdots \\
x_r & f_{r1}/f_{\bullet 1} & \vdots & f_{rj}/f_{\bullet j} & \vdots & f_{rs}/f_{\bullet s} & f_{r \bullet} \\
\hline & 1 & \cdots & 1 & \cdots & 1 & 1 \\
\end{array}\]

El análisis de correspondencias es un procedimiento de representación e
interpretación de la variabilidad resultante del estadístico
ji-cuadrado. Por tanto, si aceptamos la hipótesis de independencia no
procede llevar a cabo un Análisis de Correspondencias. En caso de
rechazo, habremos concluido que existe asociación, y mediante el
Análisis de Correspondencias se realiza un estudio de dicha asociación.

\section{Representación de los perfiles de fila}\label{representacion-de-los-perfiles-de-fila}

\subsection{Perfiles de fila estandarizados}\label{perfiles-de-fila-estandarizados}

Pensemos en la descomposición de la ji-cuadrado a través de los perfiles
de fila. Los perfiles de fila se pueden concebir como vectores del
espacio \({\Bbb R}^s\) y su variabilidad se encuentra medida a través de
la distancia ji-cuadrado en dicho espacio:

\[d^2(\bm{u},\bm{v})=\sum_{j=1}^s (u_j-v_j)^2 \frac{1}{f_{\bullet j}}\]
siendo \(\bm{u}=(u_1,\ldots,u_s)^\prime{}\) y
\(\bm{v}=(v_1,\ldots,v_s)^\prime{}\) vectores de \({\Bbb R}^s\). Si
dividimos las coordenadas de cada perfil de fila por
\(\left(\frac{1}{\sqrt{f_{\bullet 1}}},\ldots,\frac{1}{\sqrt{f_{\bullet s}}}\right)\)
entonces podemos considerar los vectores

\[\bm{r_i^e}=\left(\frac{f_{i1}}{f_{i\bullet}\sqrt{f_{\bullet 1}}},\ldots,
\frac{f_{is}}{f_{i\bullet}\sqrt{f_{\bullet s}}}\right)^\prime{} \qquad i\in \{1,\ldots,r\}\]
dotados ahora de la distancia euclídea ordinaria, ya que la inercia
total se puede expresar como la media ponderada de las distancias
euclídeas de estos vectores a su media: \[\frac{\chi^2}{n}
=\sum_{i=1}^r f_{i\bullet} \sum_{j=1}^s
\left(\frac{f_{ij}}{f_{i\bullet}\sqrt{f_{\bullet j}}}-\sqrt{f_{\bullet j}}\right)^2.\]

Esto coincide con la traza de la matriz de covarianzas de los perfiles
estandarizados, como medida global de variabilidad, y reduce el estudio
a un análisis de componentes principales sobre los perfiles
estandarizados.

\subsection{Análisis de componentes principales de los perfiles estandarizados}\label{analisis-de-componentes-principales-de-los-perfiles-estandarizados}

Formemos una matriz cuyas filas sean los perfiles de fila
estandarizados:

\[\bm{X}=\left(\frac{f_{ij}}{f_{i\bullet}\sqrt{f_{\bullet j}}}\right)_{i,j}
=\left(\begin{array}{c} {\bm{r_1^e}}^\prime{} \\ \vdots \\ {\bm{r_r^e}}^\prime{} \end{array}\right).\]

Si esta matriz se piensa como una matriz de observaciones donde cada
columna es una variable y cada fila es un individuo, ponderado por la
frecuencia marginal \(f_{i\bullet}\) que le corresponda, entonces el
vector de medias es:
\(\bm{m_y^e}=\left(\sqrt{f_{\bullet 1}},\ldots,\sqrt{f_{\bullet s}}\right)^\prime{}\)
y la matriz de covarianzas:

\[\bm{\Sigma_r}=\bm{X}^{\prime} \bm{D_r} \bm{X} - \bm{m_y^e}{\bm{m_y^e}}^\prime{}\]
siendo \(\bm{D_r}=\textnormal{diag}(f_{1\bullet},\ldots,f_{r\bullet})\).

El análisis de componentes principales de los perfiles de fila
estandarizados pasa por diagonalizar la matriz de covarianzas
\(\bm{\Sigma_r}\). Para ello, el resultado que sigue es muy útil.

\textbf{Resultado}. La matriz \(\bm{X^{\prime} D_r X}\) es una matriz simétrica
y semidefinida positiva, y tiene al vector \(\bm{m_y^e}\) como vector
propio de valor propio uno.

En base a este resultado, \(\bm{m_y^e}(\bm{m_y^e})^{\prime}\) no es más
que el primer término de la descomposición espectral de
\(\bm{X^{\prime} D_r X}\), \(\bm{m_y^e}\) es un autovector de
\(\bm{\Sigma_r}\) de autovalor cero, y el resto de autovectores y
autovalores de \(\bm{\Sigma_r}\) y \(\bm{X^{\prime} D_r X}\) coinciden.

Esta pérdida de una dimensión en el espacio de los perfiles de fila se
debe a que los perfiles de fila verifican:

\[\sum_{j=1}^s \frac{f_{i j}}{f_{i\bullet}}=1 \qquad \forall i\in\{1,\ldots,r\}
\qquad \left[\bm{r_i}^\prime{}\bm{1_s}=1\right]\] y, en consecuencia,
los perfiles de fila estandarizados:

\[\sum_{j=1}^s \sqrt{f_{\bullet j}}
\frac{f_{i j}}{\sqrt{f_{\bullet j}}f_{i\bullet}}=1 \qquad \forall i\in\{1,\ldots,r\}
\qquad \left[{\bm{r_i^e}}^\prime{}\bm{m_y^e}=1\right].\]

De este modo, en
ambos casos se encuentran contenidos en un hiperplano, y el vector
característico de dicho hiperplano (\(\bm{m_y^e}\) en el caso de los
perfiles estandarizados) es un autovector cuyo autovalor asociado es
cero.

Para completar la comprensión de este hecho, obtenemos una nueva
expresión para la inercia:

\[\frac{\chi^2}{n}
=\sum_{i=1}^r\sum_{j=1}^s
\frac{\left(f_{ij}-f_{i\bullet}f_{\bullet j}\right)^2}{f_{i\bullet}f_{\bullet j}}
=\sum_{i=1}^r\sum_{j=1}^s
\frac{f_{ij}^2}{f_{i\bullet}f_{\bullet j}}-1=\textnormal{traza} \left(\bm{X^{\prime} D_r X}\right)-1
=\textnormal{traza} \left(\bm{\Sigma_r}\right).\] La segunda igualdad se
obtiene desarrollando el cuadrado. Haber obtenido la traza de la matriz
de covarianzas \(\bm{\Sigma_r}\) no es más que lo que cabía esperar. Así, pensemos en la diagonalización

\[\bm{X^{\prime} D_r X}=\bm{V}\bm{\Lambda}\bm{V}^{\prime}\] siendo
\(\bm{V}=(\bm{m_y^e}, \bm{v_2}, \ldots, \bm{v_s})\) la matriz cuyas
columnas son los autovectores y \(\bm\Lambda\) la matriz diagonal que
contiene los autovalores correspondientes ordenados de mayor a menor, es
decir,
\(\bm\Lambda=\textnormal{diag} \left(1,\lambda_2,\ldots,\lambda_{s}\right)\).
Entonces las coordenadas de las filas respecto de las componentes
principales vendrán dadas por las columnas segunda, tercera y sucesivas
de la matriz

\[\bm{A}=\bm{XV}
=\left(\begin{array}{c} {\bm{r_1^e}}^\prime{} \\ \vdots \\  {\bm{r_r^e}}^\prime{} \end{array}\right)
\left(\bm{m_y^e}, \bm{v_2}, \ldots, \bm{v_s}\right)
=\left(\begin{array}{ccccc}
1 & a_{12} & a_{13} & \cdots & a_{1s} \\
\vdots & \vdots & \vdots & \ddots & \vdots \\
1 & a_{r2} & a_{r3} & \cdots & a_{rs} \end{array}\right).\]

Estas coordenadas verifican las siguientes propiedades:

\begin{itemize}
\item
  Las coordenadas tienen media ponderada cero.
  \(\left(\ \sum_{i=1}^r f_{i\bullet} a_{ij}=0\ \ \forall j\in\{2,\ldots,s\}\ \right).\)
\item
  Las distancias euclídeas entre los vectores de coordenadas de dos
  filas coincide con la distancia euclídea entre los correspondientes
  perfiles estandarizados, la cual a su vez coincidía con la distancia
  ji cuadrado entre los perfiles de fila originarios.
\end{itemize}

\subsection{Contribuciones}\label{contribuciones}

De la segunda propiedad se deduce que la distancia al centroide, esto
es, al perfil medio, de cada perfil de fila se puede expresar:
\[d^2(\bm{r_i},\bm{m_y})=\sum_{k=2}^s a_{ik}^2\] y esto, multiplicado
por \(f_{i\bullet}\), lo concebimos como la contribución del perfil de
fila \(i\)-ésimo a la inercia total.

El análisis de componentes principales de los perfiles estandarizados da
lugar a una descomposición de la inercia como suma de los autovalores

\[\frac{\chi^2}{n}=\textnormal{traza}(\bm{\Sigma_r})=\lambda_2+\lambda_3+\ldots+\lambda_s\]
siendo

\[\lambda_k = \bm{v_k}^{\prime} \bm{X}^{\prime} \bm{D_r X v_k} = \sum_{i=1}^r f_{i\bullet} a_{ik}^2.\]
Para una fácil comprensión de la segunda igualdad, debe recordarse que
\(\bm{X v_k}\) es el vector columna (columna \(k\)-ésima de la matriz
\(\bm{A}\)) que contiene las cooordenadas de cada perfil de fila respecto
de la componente \(k\)-ésima.

La tabla siguiente es una descomposición de la inercia total en perfiles
de fila y componentes:

\[\begin{array}{c|ccccc|c}
 & \mbox{ Componente 2}  & \cdots & \mbox{ Componente k} & \cdots
  & \mbox{ Componente s} & \\ \hline
\mbox{ Perfil 1}
 & f_{1\bullet} a_{12}^2 & \cdots & f_{1\bullet} a_{1k}^2
& \cdots & f_{1\bullet} a_{1s}^2 & \sum_{k=2}^s f_{1\bullet} a_{1k}^2 \\
\vdots & \vdots & & \vdots & & \vdots & \vdots \\
\mbox{ Perfil i}
 & f_{i\bullet} a_{i2}^2  & \cdots & f_{i\bullet} a_{ik}^2
& \cdots & f_{i\bullet} a_{is}^2 & \sum_{k=2}^s f_{i\bullet} a_{ik}^2 \\
\vdots & \vdots  & & \vdots & & \vdots & \vdots \\
\mbox{ Perfil r}
 & f_{r\bullet} a_{r2}^2  & \cdots & f_{r\bullet} a_{rk}^2
& \cdots & f_{r\bullet} a_{rs}^2 & \sum_{k=2}^s f_{1\bullet} a_{rk}^2 \\
\hline & \sum_{i=1}^r f_{i\bullet} a_{i 2}^2 & \cdots & \sum_{i=1}^r f_{i\bullet} a_{i k}^2
& \cdots & \sum_{i=1}^r f_{i\bullet} a_{i s}^2 & \chi^2/n \end{array}\]

La proporción de variabilidad explicada si tomamos \((k-1)\) componentes
(\(k\leq s\)) será:
\[\frac{\lambda_2+\ldots+\lambda_k}{\lambda_2+\ldots+\lambda_{s}}\] y
sería como quedarnos con las \((k-1)\) primeras columnas (de la \(2\) a la
\(k\)) tanto en la matriz \(\bm{A}\) como en la tabla anterior.

\section{Representación de los perfiles de columna}\label{representacion-de-los-perfiles-de-columna}

\subsection{Análisis de componentes principales de los perfiles estandarizados}\label{analisis-de-componentes-principales-de-los-perfiles-estandarizados-2}

El análisis anterior de los perfiles de fila se puede repetir de la
misma manera para los perfiles de columna. Los pasos se llevan a cabo en
perfecta analogía con el caso anterior. Resumiendo, tendríamos que la
inercia total se puede expresar como media ponderada de las distancias
ji-cuadrado de los perfiles de columna (a su media), o equivalentemente,
como las distancias euclídeas de los perfiles estandarizados (a su media
correspondiente):

\[\frac{\chi^2}{n}
=\sum_{j=1}^s f_{\bullet j} \sum_{i=1}^r
\left(\frac{f_{ij}}{f_{\bullet j}\sqrt{f_{i \bullet}}}-\sqrt{f_{i \bullet}}\right)^2
=\sum_{j=1}^s f_{\bullet j} \left\|\bm{c_i^e}-\bm{m_x^e}\right\|^2.\]

Formamos la matriz \(\bm{Y}\), cuyas columnas son los perfiles de columna
estandarizados:

\[\bm{Y}=\left(\frac{f_{ij}}{f_{\bullet j}\sqrt{f_{i \bullet}}}\right)_{i,j}
=\left( \bm{c_1^e}, \ldots,  \bm{c_s^e} \right)\] la cual se concibe
como una matriz de observaciones transpuesta, donde cada fila es una
variable y cada columna es un individuo, ponderado por la frecuencia
marginal \(f_{\bullet j}\) que le corresponda. Entonces el vector de
medias es:
\(\bm{m_x^e}=\left(\sqrt{f_{1 \bullet}},\ldots,\sqrt{f_{r \bullet}}\right)^{\prime}\)
y la matriz de covarianzas:

\[\bm{\Sigma_s}=\bm{Y D_s Y^{\prime}} - \bm{m_x^e} {\bm{m_x^e}}^{\prime}\]
siendo \(\bm{D_s}=\textnormal{diag}(f_{\bullet 1},\ldots,f_{\bullet s})\).

De nuevo se aplica el Análisis de Componentes Principales, mediante la
diagonalización de la matriz de covarianzas \(\bm{\Sigma_s}\). Surge
entonces el resultado siguiente.

\textbf{Resultado}. La matriz \(\bm{Y D_s Y^{\prime}}\) es una matriz simétrica
y semidefinida positiva, y tiene al vector \(\bm{m_x^e}\) como vector
propio de valor propio uno.

En base a este resultado, \(\bm{m_x^e}{\bm{m_x^e}}^{\prime}\) no es más
que el primer término de la descomposición espectral de
\(\bm{Y D_s Y^{\prime}}\), \(\bm{m_x^e}\) es un autovector de
\(\bm{\Sigma_s}\) de autovalor cero, y el resto de autovectores y
autovalores de \(\bm{\Sigma_s}\) e \(\bm{Y D_s Y^{\prime}}\) coinciden.

De nuevo se pierde una dimensión en el espacio de los perfiles de
columna, puesto que verifican una restricción lineal obvia, consecuencia
de que sus componentes suman uno. Esta restricción se convierte en los
perfiles de columna estandarizados, en que pertenecen al hiperplano cuyo
vector característico es \(\bm{m_x^e}\). Así, \[\frac{\chi^2}{n}
=\mbox{traza} \left(\bm{Y D_s Y^{\prime}}\right)-1
=\mbox{traza} \left(\bm{\Sigma_s}\right).\]

Entonces, la diagonalización conduce a

\[\bm{Y D_s Y^{\prime}}=\bm{W}\bm{\Lambda} \bm{W}^{\prime}\] siendo
\(\bm{W}=\left( \bm{m_x^e}, \bm{w_2}, \ldots, \bm{w_r}\right)\) la matriz
cuyas columnas son los autovectores y \(\bm{\Lambda}\) la la matriz
diagonal que contiene los autovalores correspondientes ordenados de
mayor a menor, es decir
\(\bm{\Lambda}=\mbox{diag} \left(1,\lambda_2,\ldots,\lambda_r \right)\).
Ahora las coordenadas de las columnas respecto de las componentes
principales vendrán dadas por las filas segunda, tercera y sucesivas de
la matriz

\[\bm{B}=\bm{W}^{\prime}\bm{Y}
=\left(\begin{array}{c} {\bm{m_x^e}}^{\prime} \\ \bm{w_2}^{\prime} \\ \vdots \\ \bm{w_r}^{\prime} \end{array} \right)
 \left(\bm{c_1^e}, \ldots, \bm{c_s^e} \right)
=\left(\begin{array}{ccc}
1& \cdots & 1 \\
b_{21} & \cdots & b_{2s} \\
b_{31} & \cdots & b_{3s} \\
\vdots & \ddots & \vdots \\
b_{r1} & \cdots & b_{rs} \end{array}\right).\]

Estas coordenadas verifican las siguientes propiedades:

\begin{itemize}
\item
  Las coordenadas tienen media ponderada cero.
  \(\left(\ \sum_{j=1}^s f_{\bullet j} b_{ij}=0\ \ \forall i\in\{2,\ldots,r\}\ \right).\)
\item
  Las distancias euclídeas entre los vectores de coordenadas de dos
  columnas coincide con la distancia euclídea entre los
  correspondientes perfiles estandarizados, la cual a su vez coincidía
  con la distancia ji cuadrado entre los perfiles de columna
  originarios.
\end{itemize}

\subsection{Contribuciones}\label{contribuciones-2}

De la segunda propiedad se deduce que la distancia al centroide, esto
es, al perfil medio, de cada perfil de columna se puede expresar:

\[d^2(\bm{c_j},\bm{m_x})=\sum_{k=2}^r b_{k j}^2\] y esto, multiplicado
por \(f_{\bullet j}\), lo concebimos como la contribución del perfil de
columna \(j\)-ésimo a la inercia total.

El análisis de componentes principales de los perfiles estandarizados da
lugar a una descomposición de la inercia como suma de los autovalores

\[\frac{\chi^2}{n}=\mbox{traza}(\bm{\Sigma_s})=\lambda_2+\lambda_3+\ldots+\lambda_r\]
siendo

\[\lambda_k = \bm{w_k}^{\prime} \bm{Y D_s Y^{\prime}} \bm{w_k} = \sum_{j=1}^r f_{\bullet j} b_{kj}^2.\]
Para una fácil comprensión de la segunda igualdad, debe recordarse que
\(\bm{w_k}^{\prime}\bm{Y}\) es el vector fila (fila \(k\)-ésima de la matriz
\(\bm{B}\)) que contiene las cooordenadas de cada perfil de columna
respecto de la componente \(k\)-ésima.

La tabla siguiente es una descomposición de la inercia total en perfiles
de columna y componentes:

\[\begin{array}{c|ccccc|c}
 & \mbox{ Perfil 1} &  \cdots & \mbox{ Perfil j} & \cdots
  & \mbox{ Perfil s} & \\ \hline
\mbox{ Componente 2}
 & f_{\bullet 1} b_{21}^2 &  \cdots & f_{\bullet j} b_{2j}^2
& \cdots & f_{\bullet s} b_{2s}^2 & \sum_{j=1}^s f_{\bullet j} b_{2j}^2 \\
\vdots & \vdots &  & \vdots & & \vdots & \vdots \\
\mbox{ Componente k}
 & f_{\bullet 1} b_{k1}^2 &  \cdots & f_{\bullet j} b_{kj}^2
& \cdots & f_{\bullet s} b_{ks}^2 & \sum_{j=1}^s f_{\bullet j} b_{kj}^2 \\
\vdots & \vdots &  & \vdots & & \vdots & \vdots \\
\mbox{ Componente r}
 & f_{\bullet 1} b_{r1}^2 &  \cdots & f_{\bullet j} b_{rj}^2
& \cdots & f_{\bullet s} b_{rs}^2 & \sum_{j=1}^s f_{\bullet j} b_{rj}^2 \\
\hline & \sum_{k=2}^r f_{\bullet 1} b_{k 1}^2 & \cdots
 & \sum_{k=2}^r f_{\bullet j} b_{k j}^2
& \cdots & \sum_{k=2}^r f_{\bullet s} b_{k s}^2 & \chi^2/n \end{array}\]

La proporción de variabilidad explicada si tomamos \((k-1)\) componentes
(\(k\leq r\)) será:

\[\frac{\lambda_2+\ldots+\lambda_k}{\lambda_2+\ldots+\lambda_{r}}\] y
sería como quedarnos con las \((k-1)\) primeras columnas (de la \(2\) a la
\(k\)) tanto en la matriz \(\bm{B}\) como en la tabla anterior.

\section{Representación simultánea de los perfiles de fila y columna}\label{representacion-simultanea-de-los-perfiles-de-fila-y-columna}

Consideremos la matriz

\[\bm{Z}=\left( \frac{f_{ij}}{\sqrt{f_{i\bullet}}\sqrt{f_{\bullet j}}} \right)_{i,j}.\]

Observemos que

\[\bm{X}=\bm{D_r}^{-1}\bm{F} \bm{D_s}^{-1/2},\ \ \ \ \bm{Y}=\bm{D_r}^{-1/2} \bm{F} \bm{D_s}^{-1},\ \ \ \ \bm{Z}=\bm{D_r}^{-1/2} \bm{F}\bm{D_s}^{-1/2}\]
siendo \(\bm{F}=\left(f_{ij}\right)_{i,j}\). Entonces

\begin{itemize}
\item
  \(\bm{X}^{\prime} \bm{D_r} \bm{X}=\bm{Z}^{\prime} \bm{Z}\qquad \bm{Y \bm{D_s} \bm{Y}^{\prime}}=\bm{Z}\bm{Z}^{\prime}.\)
\item
  Si \(\bm{v}\) es autovector de \(\bm{Z}^{\prime} \bm{Z}\) con autovalor
  asociado \(\lambda\), entonces \(\bm{w}=\bm{Z}\bm{v}\) es autovector de
  \(\bm{Z}\bm{Z}^{\prime}\) con autovalor asociado \(\lambda\).
\end{itemize}

Por tanto, ambas matrices tienen los mismos autovalores, mientras que
los autovectores mantienen la relación que se acaba de indicar. Como
consecuencia, el número de autovalores no nulos será
\(k_0\leq \min\{r,s\}\). En lo que sigue supondremos una matriz común de
autovalores no nulos \(\bm{D_{\lambda}}\) de orden \(k_0\), y los
autovectores y coordenadas se referirán únicamente a las \(k_0\)
componentes asociadas a estos autovalores no nulos, tanto cuando se
interpreten en el espacio de filas como en el de columnas.

Recordemos la diagonalización

\[\bm{Z}^{\prime} \bm{Z}= \bm{X}^{\prime} \bm{D_r} \bm{X}= \bm{V}\bm{D_{\lambda}}\bm{V}^{\prime}.\]
Entonces las columnas de \(\bm{ZV}\) son los autovectores de
\(\bm{Z}\bm{Z}^{\prime}\), que una vez normalizadas darán lugar a la
matriz ortogonal para la diagonalización
\(\bm{Z}\bm{Z}^{\prime}= \bm{W}\bm{D_{\lambda}} \bm{W}^{\prime}\):

\[\bm{W}=\bm{Z}\bm{VD_{\lambda}}^{-1/2}.\] Finalmente, las coordenadas de
los perfiles de columna eran las columnas de la matriz

\[\bm{B}=\bm{W}^{\prime}\bm{Y}.\] Vamos a deducir la relación que existe
entre estas coordenadas y las coordenadas de los perfiles de fila,
\(\bm{A}=\bm{XV}\).

\[\begin{aligned}
\bm{B}&= \bm{W}^{\prime}\bm{Y}= \left(\bm{Z}\bm{V}\bm{D_{\lambda}}^{-1/2}\right)^{\prime} \bm{Y}
 =\bm{D_{\lambda}}^{-1/2}\bm{V}^{\prime} \bm{Z}^{\prime} \bm{Y} \\
 &=\bm{D_{\lambda}}^{-1/2} \bm{V}^{\prime} \left(\bm{D_r}^{-1/2}\bm{F}\bm{D_s}^{-1/2}\right)^{\prime} \bm{D_r}^{-1/2}\bm{F}\bm{D_s}^{-1}
 =\bm{D_{\lambda}}^{-1/2} \bm{V}^{\prime} \bm{D_s}^{-1/2} \bm{F}^{\prime} \bm{D_r}^{-1} \bm{F}\bm{D_s}^{-1} \\
 &=\bm{D_{\lambda}}^{-1/2} \bm{A}^{\prime} \bm{F}\bm{D_s}^{-1}. \end{aligned}\]
En la última igualdad se ha utilizado que
\(\bm{A}=\bm{X}\bm{V}=\bm{D_r}^{-1}\bm{F}\bm{D_s}^{-1/2}\bm{V}\). Pero el
resultado obtenido se puede escribir así:

\[b_{kj}=\frac{1}{\sqrt{\lambda_k}} \left( \frac{f_{1j}}{f_{\bullet j}} a_{1k}
+\ldots + \frac{f_{rj}}{f_{\bullet j}} a_{rk} \right)\] de modo que la
coordenada del perfil de columna \(j\)-ésimo en la componente \(k\)-ésima es
una combinación ponderada de las coordenadas de cada perfil de fila en
dicha componente, siendo los pesos de tal ponderación las frecuencias
relativas de cada categoría de la variable \(x\) condicionada a la
categoría j de \(y\). Así, si una categoría i presenta alta frecuencia
condicionada pesará más en esta combinación, aproximando hacia sí la
coordenada de la categoría \(j\) de la variable \(y\). De igual modo, se puede ver que

\[\bm{A}=\bm{D_r}^{-1}\bm{F}\bm{B}^{\prime}\bm{D_{\lambda}}^{-1/2}.\]

\textbf{\includegraphics[width=1em,height=1em]{bookdown_AM_files/figure-latex/fa-icon-00ef1bb1d064b635425d49f06b331298.pdf} ~\textcolor{excolor}{Ejemplo.}} En el ejemplo sobre los datos de la EPA (2010), vamos a
representar en un mismo gráfico bidimensional los perfiles de fila y
columna.

\footnotesize

\begin{Shaded}
\begin{Highlighting}[]
\SpecialCharTok{\textgreater{}} \FunctionTok{library}\NormalTok{(MASS)  }\CommentTok{\# Cargamos el paquete MASS (Venables y Ripley)}
\SpecialCharTok{\textgreater{}}\NormalTok{ m }\OtherTok{\textless{}{-}} \FunctionTok{cbind}\NormalTok{(dat}\SpecialCharTok{$}\NormalTok{Ocupados, dat}\SpecialCharTok{$}\NormalTok{Parados, dat}\SpecialCharTok{$}\NormalTok{Inactivos)  }\CommentTok{\# Frecuencias absolutas}
\SpecialCharTok{\textgreater{}} \FunctionTok{rownames}\NormalTok{(m) }\OtherTok{\textless{}{-}}\NormalTok{ dat}\SpecialCharTok{$}\NormalTok{Comunidad.Autónoma  }\CommentTok{\# Nombres de filas}
\SpecialCharTok{\textgreater{}} \FunctionTok{colnames}\NormalTok{(m) }\OtherTok{\textless{}{-}} \FunctionTok{names}\NormalTok{(dat)[}\DecValTok{2}\SpecialCharTok{:}\DecValTok{4}\NormalTok{]  }\CommentTok{\# Nombres de las columnas}
\SpecialCharTok{\textgreater{}}\NormalTok{ r }\OtherTok{\textless{}{-}} \FunctionTok{nrow}\NormalTok{(m)  }\CommentTok{\# Número de filas}
\SpecialCharTok{\textgreater{}}\NormalTok{ s }\OtherTok{\textless{}{-}} \FunctionTok{ncol}\NormalTok{(m)  }\CommentTok{\# Número de columnas}
\SpecialCharTok{\textgreater{}}\NormalTok{ n }\OtherTok{\textless{}{-}} \FunctionTok{sum}\NormalTok{(m)  }\CommentTok{\# Frecuencia total (tamaño muestral)}
\SpecialCharTok{\textgreater{}}\NormalTok{ f }\OtherTok{\textless{}{-}}\NormalTok{ m}\SpecialCharTok{/}\NormalTok{n  }\CommentTok{\# Frecuencias relativas (conjuntas)}
\SpecialCharTok{\textgreater{}}\NormalTok{ mx }\OtherTok{\textless{}{-}} \FunctionTok{apply}\NormalTok{(f, }\DecValTok{1}\NormalTok{, sum)  }\CommentTok{\# Frecuencias marginales de columna}
\SpecialCharTok{\textgreater{}}\NormalTok{ my }\OtherTok{\textless{}{-}} \FunctionTok{apply}\NormalTok{(f, }\DecValTok{2}\NormalTok{, sum)  }\CommentTok{\# Frecuencias marginales de fila}
\end{Highlighting}
\end{Shaded}

\normalsize

Calculamos los perfiles de fila y columna.

\footnotesize

\begin{Shaded}
\begin{Highlighting}[]
\SpecialCharTok{\textgreater{}}\NormalTok{ ri }\OtherTok{\textless{}{-}} \FunctionTok{diag}\NormalTok{(}\DecValTok{1}\SpecialCharTok{/}\NormalTok{mx) }\SpecialCharTok{\%*\%}\NormalTok{ f  }\CommentTok{\# Perfiles de fila}
\SpecialCharTok{\textgreater{}} \FunctionTok{rownames}\NormalTok{(ri) }\OtherTok{\textless{}{-}}\NormalTok{ dat}\SpecialCharTok{$}\NormalTok{Comunidad.Autónoma}
\SpecialCharTok{\textgreater{}}\NormalTok{ cj }\OtherTok{\textless{}{-}}\NormalTok{ f }\SpecialCharTok{\%*\%} \FunctionTok{diag}\NormalTok{(}\DecValTok{1}\SpecialCharTok{/}\NormalTok{my)  }\CommentTok{\# Perfiles de columna}
\SpecialCharTok{\textgreater{}} \FunctionTok{colnames}\NormalTok{(cj) }\OtherTok{\textless{}{-}} \FunctionTok{names}\NormalTok{(dat)[}\DecValTok{2}\SpecialCharTok{:}\DecValTok{4}\NormalTok{]}
\end{Highlighting}
\end{Shaded}

\normalsize

A continuación, calculamos las componentes principales de los perfiles
de fila estandarizados.

\footnotesize

\begin{Shaded}
\begin{Highlighting}[]
\SpecialCharTok{\textgreater{}}\NormalTok{ x }\OtherTok{\textless{}{-}}\NormalTok{ ri }\SpecialCharTok{\%*\%} \FunctionTok{diag}\NormalTok{(}\DecValTok{1}\SpecialCharTok{/}\FunctionTok{sqrt}\NormalTok{(my))  }\CommentTok{\# Perfiles de fila estandarizados}
\SpecialCharTok{\textgreater{}}\NormalTok{ mye }\OtherTok{\textless{}{-}} \FunctionTok{sqrt}\NormalTok{(my)  }\CommentTok{\# Marginal de fila estandarizada}
\SpecialCharTok{\textgreater{}}\NormalTok{ sr }\OtherTok{\textless{}{-}} \FunctionTok{t}\NormalTok{(x) }\SpecialCharTok{\%*\%} \FunctionTok{diag}\NormalTok{(mx) }\SpecialCharTok{\%*\%}\NormalTok{ x }\SpecialCharTok{{-}}\NormalTok{ mye }\SpecialCharTok{\%*\%} \FunctionTok{t}\NormalTok{(mye)  }\CommentTok{\# Matriz Sigma\_r}
\SpecialCharTok{\textgreater{}}\NormalTok{ auto }\OtherTok{\textless{}{-}} \FunctionTok{eigen}\NormalTok{(}\FunctionTok{t}\NormalTok{(x) }\SpecialCharTok{\%*\%} \FunctionTok{diag}\NormalTok{(mx) }\SpecialCharTok{\%*\%}\NormalTok{ x)  }\CommentTok{\# Diagonalización }
\SpecialCharTok{\textgreater{}}\NormalTok{ lam }\OtherTok{\textless{}{-}}\NormalTok{ auto}\SpecialCharTok{$}\NormalTok{values  }\CommentTok{\# Autovalores}
\SpecialCharTok{\textgreater{}}\NormalTok{ v }\OtherTok{\textless{}{-}}\NormalTok{ auto}\SpecialCharTok{$}\NormalTok{vectors  }\CommentTok{\# Autovectores}
\end{Highlighting}
\end{Shaded}

\normalsize

Calculamos ahora las coordenadas de los perfiles de fila estandarizados
en las componentes y hacemos la representación (Figura
\ref{fig:co}).

\footnotesize

\begin{Shaded}
\begin{Highlighting}[]
\SpecialCharTok{\textgreater{}}\NormalTok{ a }\OtherTok{\textless{}{-}}\NormalTok{ x }\SpecialCharTok{\%*\%}\NormalTok{ v}
\SpecialCharTok{\textgreater{}} \FunctionTok{plot}\NormalTok{(a[, }\DecValTok{2}\NormalTok{], a[, }\DecValTok{3}\NormalTok{])}
\SpecialCharTok{\textgreater{}} \FunctionTok{text}\NormalTok{(a[, }\DecValTok{2}\NormalTok{], a[, }\DecValTok{3}\NormalTok{], }\AttributeTok{labels =}\NormalTok{ dat}\SpecialCharTok{$}\NormalTok{Comunidad.Autónoma)}
\SpecialCharTok{\textgreater{}} \FunctionTok{abline}\NormalTok{(}\AttributeTok{h =} \DecValTok{0}\NormalTok{, }\AttributeTok{lty =} \DecValTok{2}\NormalTok{)}
\SpecialCharTok{\textgreater{}} \FunctionTok{abline}\NormalTok{(}\AttributeTok{v =} \DecValTok{0}\NormalTok{, }\AttributeTok{lty =} \DecValTok{2}\NormalTok{)}
\end{Highlighting}
\end{Shaded}

\begin{figure}

{\centering \resizebox{.4\linewidth}{!}{\includegraphics[width=0.5\linewidth]{bookdown_AM_files/figure-latex/co-1} } 

}

\caption{Gráficos de perfiles de fila.}\label{fig:co}
\end{figure}
\normalsize

A continuación, calculamos las componentes principales de los perfiles
de columna estandarizados.

\footnotesize

\begin{Shaded}
\begin{Highlighting}[]
\SpecialCharTok{\textgreater{}}\NormalTok{ y }\OtherTok{\textless{}{-}} \FunctionTok{diag}\NormalTok{(}\DecValTok{1}\SpecialCharTok{/}\FunctionTok{sqrt}\NormalTok{(mx)) }\SpecialCharTok{\%*\%}\NormalTok{ cj  }\CommentTok{\# Perfiles de columna estandarizados}
\SpecialCharTok{\textgreater{}}\NormalTok{ mxe }\OtherTok{\textless{}{-}} \FunctionTok{sqrt}\NormalTok{(mx)  }\CommentTok{\# Marginal de columna estandarizada}
\SpecialCharTok{\textgreater{}}\NormalTok{ sigs }\OtherTok{\textless{}{-}}\NormalTok{ y }\SpecialCharTok{\%*\%} \FunctionTok{diag}\NormalTok{(my) }\SpecialCharTok{\%*\%} \FunctionTok{t}\NormalTok{(y) }\SpecialCharTok{{-}}\NormalTok{ mxe }\SpecialCharTok{\%*\%} \FunctionTok{t}\NormalTok{(mxe)  }\CommentTok{\# Matriz Sigma\_s}
\SpecialCharTok{\textgreater{}}\NormalTok{ auto }\OtherTok{\textless{}{-}} \FunctionTok{eigen}\NormalTok{(y }\SpecialCharTok{\%*\%} \FunctionTok{diag}\NormalTok{(my) }\SpecialCharTok{\%*\%} \FunctionTok{t}\NormalTok{(y))  }\CommentTok{\# Diagonalización }
\SpecialCharTok{\textgreater{}}\NormalTok{ lam }\OtherTok{\textless{}{-}}\NormalTok{ auto}\SpecialCharTok{$}\NormalTok{values  }\CommentTok{\# Autovalores}
\SpecialCharTok{\textgreater{}}\NormalTok{ w }\OtherTok{\textless{}{-}}\NormalTok{ auto}\SpecialCharTok{$}\NormalTok{vectors  }\CommentTok{\# Autovectores}
\end{Highlighting}
\end{Shaded}

\normalsize

Calculamos ahora las coordenadas de los perfiles de columna
estandarizados en las componentes y hacemos la representación (Figura
\ref{fig:cod}).

\footnotesize

\begin{Shaded}
\begin{Highlighting}[]
\SpecialCharTok{\textgreater{}}\NormalTok{ b }\OtherTok{\textless{}{-}} \FunctionTok{t}\NormalTok{(w) }\SpecialCharTok{\%*\%}\NormalTok{ y}
\SpecialCharTok{\textgreater{}} \FunctionTok{plot}\NormalTok{(b[}\DecValTok{2}\NormalTok{, ], b[}\DecValTok{3}\NormalTok{, ])}
\SpecialCharTok{\textgreater{}} \FunctionTok{text}\NormalTok{(b[}\DecValTok{2}\NormalTok{, ], b[}\DecValTok{3}\NormalTok{, ], }\FunctionTok{names}\NormalTok{(dat)[}\DecValTok{2}\SpecialCharTok{:}\DecValTok{4}\NormalTok{])}
\SpecialCharTok{\textgreater{}} \FunctionTok{abline}\NormalTok{(}\AttributeTok{h =} \DecValTok{0}\NormalTok{, }\AttributeTok{lty =} \DecValTok{2}\NormalTok{)}
\SpecialCharTok{\textgreater{}} \FunctionTok{abline}\NormalTok{(}\AttributeTok{v =} \DecValTok{0}\NormalTok{, }\AttributeTok{lty =} \DecValTok{2}\NormalTok{)}
\end{Highlighting}
\end{Shaded}

\begin{figure}

{\centering \resizebox{.4\linewidth}{!}{\includegraphics[width=0.5\linewidth]{bookdown_AM_files/figure-latex/cod-1} } 

}

\caption{Gráficos de perfiles de columna.}\label{fig:cod}
\end{figure}
\normalsize

Hacemos ahora la representación simultánea (Figura \ref{fig:cos}).

\footnotesize

\begin{Shaded}
\begin{Highlighting}[]
\SpecialCharTok{\textgreater{}}\NormalTok{ z }\OtherTok{\textless{}{-}} \FunctionTok{diag}\NormalTok{(}\DecValTok{1}\SpecialCharTok{/}\FunctionTok{sqrt}\NormalTok{(mx)) }\SpecialCharTok{\%*\%}\NormalTok{ f }\SpecialCharTok{\%*\%} \FunctionTok{diag}\NormalTok{(}\DecValTok{1}\SpecialCharTok{/}\FunctionTok{sqrt}\NormalTok{(my))}
\SpecialCharTok{\textgreater{}}\NormalTok{  k0 }\OtherTok{\textless{}{-}} \DecValTok{3}
\SpecialCharTok{\textgreater{}}\NormalTok{  w }\OtherTok{\textless{}{-}}\NormalTok{ z }\SpecialCharTok{\%*\%}\NormalTok{ v[, }\DecValTok{1}\SpecialCharTok{:}\NormalTok{k0] }\SpecialCharTok{\%*\%} \FunctionTok{diag}\NormalTok{(}\DecValTok{1}\SpecialCharTok{/}\FunctionTok{sqrt}\NormalTok{(lam[}\DecValTok{1}\SpecialCharTok{:}\NormalTok{k0]))}
\SpecialCharTok{\textgreater{}}\NormalTok{  b }\OtherTok{\textless{}{-}} \FunctionTok{t}\NormalTok{(w) }\SpecialCharTok{\%*\%}\NormalTok{ y}
\SpecialCharTok{\textgreater{}} \FunctionTok{plot}\NormalTok{(}\FunctionTok{c}\NormalTok{(a[, }\DecValTok{2}\NormalTok{], b[}\DecValTok{2}\NormalTok{, ]), }\FunctionTok{c}\NormalTok{(a[, }\DecValTok{3}\NormalTok{], b[}\DecValTok{3}\NormalTok{, ]))}
\SpecialCharTok{\textgreater{}} \FunctionTok{text}\NormalTok{(}\FunctionTok{c}\NormalTok{(a[, }\DecValTok{2}\NormalTok{], b[}\DecValTok{2}\NormalTok{, ]), }\FunctionTok{c}\NormalTok{(a[, }\DecValTok{3}\NormalTok{], b[}\DecValTok{3}\NormalTok{, ]), }
\SpecialCharTok{+}      \AttributeTok{labels =} \FunctionTok{c}\NormalTok{(}\FunctionTok{as.character}\NormalTok{(dat}\SpecialCharTok{$}\NormalTok{Comunidad.Autónoma), }\FunctionTok{names}\NormalTok{(dat)[}\DecValTok{2}\SpecialCharTok{:}\DecValTok{4}\NormalTok{]), }
\SpecialCharTok{+}      \AttributeTok{col =} \FunctionTok{c}\NormalTok{(}\FunctionTok{rep}\NormalTok{(}\StringTok{"blue"}\NormalTok{, r), }\FunctionTok{rep}\NormalTok{(}\StringTok{"red"}\NormalTok{, s)), }\AttributeTok{cex =} \FloatTok{0.7}\NormalTok{)}
\SpecialCharTok{\textgreater{}} \FunctionTok{abline}\NormalTok{(}\AttributeTok{h =} \DecValTok{0}\NormalTok{, }\AttributeTok{lty =} \DecValTok{2}\NormalTok{)}
\SpecialCharTok{\textgreater{}} \FunctionTok{abline}\NormalTok{(}\AttributeTok{v =} \DecValTok{0}\NormalTok{, }\AttributeTok{lty =} \DecValTok{2}\NormalTok{)}
\end{Highlighting}
\end{Shaded}

\begin{figure}

{\centering \resizebox{.4\linewidth}{!}{\includegraphics[width=0.5\linewidth]{bookdown_AM_files/figure-latex/cos-1} } 

}

\caption{Representación simultánea de los perfiles de fila y columna para los datos de la EPA.}\label{fig:cos}
\end{figure}
\normalsize

También podemos hacer la representación utilizando el comando \texttt{corresp}.

\footnotesize

\begin{Shaded}
\begin{Highlighting}[]
\SpecialCharTok{\textgreater{}}\NormalTok{ m }\OtherTok{\textless{}{-}} \DecValTok{10} \SpecialCharTok{*}\NormalTok{ m}
\SpecialCharTok{\textgreater{}}\NormalTok{ r.out }\OtherTok{\textless{}{-}} \FunctionTok{corresp}\NormalTok{(m, }\AttributeTok{nf =} \DecValTok{2}\NormalTok{)}
\SpecialCharTok{\textgreater{}} \FunctionTok{biplot}\NormalTok{(}\FunctionTok{corresp}\NormalTok{(m, }\AttributeTok{nf =} \DecValTok{2}\NormalTok{), }\AttributeTok{cex =} \FunctionTok{c}\NormalTok{(}\FloatTok{0.7}\NormalTok{, }\DecValTok{1}\NormalTok{))}
\end{Highlighting}
\end{Shaded}

\begin{figure}

{\centering \resizebox{.4\linewidth}{!}{\includegraphics[width=0.5\linewidth]{bookdown_AM_files/figure-latex/co2-1} } 

}

\caption{Representación simultánea de los perfiles de fila y columna para los datos de la EPA (función corresp).}\label{fig:co2}
\end{figure}
\normalsize

El resultado se muestra en la Figura \ref{fig:co2}. El eje horizontal corresponde a la primera
componente y el eje vertical a la segunda. La proximidad entre los
puntos de distintas comunidades autónomas indica que presentan perfiles
similares. Asimismo, si una comunidad autónoma está próxima al punto
correspondiente a los Inactivos, esto indica que en su perfil hay un
porcentaje alto de población inactiva. Igualmente para las categorías de
Ocupados y Parados.

\chapter{Análisis discriminante}\label{analisis-discriminante}

Empezamos definiendo el problema discriminante. Resumidamente se puede
decir que el problema discriminante consiste en establecer reglas de
diagnóstico que permitan clasificar a un individuo en alguno de dos o
más grupos bien definidos. El ejemplo paradigmático sería la diagnosis
entre enfermos y sanos en Medicina. Para fundamentar la decisión de
diagnóstico se emplean ciertas observaciones del individuo: presencia de
síntomas, análisis de sangre y otras pruebas diagnósticas. Estas pruebas
son útiles para el diagnóstico en la medida en que aportan resultados
diferentes para los individuos enfermos respecto de los sanos.

Esta situación da lugar a varias tareas, según el objetivo que se
persiga:

\begin{enumerate}
\def\labelenumi{\arabic{enumi}.}
\item
  Determinar la presencia de diferencias entre los grupos.
\item
  Determinar qué variables permiten discriminar mejor entre los
  grupos, en base a las diferencias que presentan entre ellos.
\item
  Construir la regla discriminante en base a una muestra de
  entrenamiento.
\item
  Asignar un individuo a cierto grupo en base a la regla
  discriminante, bien la regla óptima o bien estimada en base a una
  muestra de entrenamiento.
\end{enumerate}

La primera tarea se resuelve mediante un Análisis (multivariante) de la
varianza, mientras que las otras tres son específicas del Análisis
discriminante.

Realmente la construcción de una regla discriminante y su aplicación
constituye una cuestión de gran calado dentro de la teoría de la
decisión y la inferencia estadística. Por ello, aunque el Análisis
discriminante se puede considerar como una técnica de Análisis
multivariante, su importancia supera ampliamente a la de otras técnicas,
como las estudiadas en temas anteriores. El Análisis discriminante se
puede considerar en sí mismo como una rama de la Estadística.

\section{Dos grupos con distribuciones conocidas}\label{dos-grupos-con-distribuciones-conocidas}

Suponemos la existencia de dos grupos \(G_1\) y \(G_2\), definidos de manera
inequívoca. Consideramos un vector aleatorio \(\bm{x}\) que presenta
funciones de densidad o de probabilidad \(f_1\) si el individuo procede
del grupo \(G_1\) y \(f_2\) si procede del grupo \(G_2\). Utilizaremos el
valor observado del vector aleatorio \(\bm{x}\) para decidir si un
individuo pertenece a uno u otro grupo. Esto supone dividir el soporte
del vector \(\bm{x}\) mediante una partición en dos regiones \(R_1\) y
\(R_2\): \(R=R_1\cup R_2\) y \(R_1\cap R_2=\emptyset\), de modo que
clasificaremos al individuo en el grupo \(G_1\) si \(\bm{x}\in R_1\) y en el
grupo \(G_2\) si \(\bm{x}\in R_2\). En realidad, la regla discriminante
anterior sería una regla determinista, de modo que una definición
general de regla discriminante, que incluya la posibilidad de
aleatorización, sería la siguiente.

\begin{definition}
\protect\hypertarget{def:unnamed-chunk-90}{}\label{def:unnamed-chunk-90}Una regla discriminante (aleatorizada) es una aplicación

\[\begin{array}{rccl}
\varphi : & R & \longrightarrow & [0,1] \\
          & \bm{x} & \longrightarrow & \varphi(\bm{x})=P(\mbox{Clasificar en }G_1/\bm{x}=\bm{x}).
\end{array}\]

Una regla discriminante no aleatorizada es de la forma

\[\varphi(\bm{x})=\left\{\begin{array}{ll}
1 & \mbox{si } \bm{x}\in R_1, \\
0 & \mbox{si } \bm{x}\in R_2. \end{array}\right.\]
\end{definition}

En el procedimiento de clasificación podemos cometer dos tipos de
errores: clasificar a un individuo que procede del grupo \(G_1\) en el
grupo \(G_2\), y clasificar a un individuo del grupo \(G_2\) en el grupo
\(G_1\). Las probabilidades condicionales de cada uno de estos errores
son

\[P(2/1)=\int \left(1-\varphi(\bm{x})\right)f_1(\bm{x})\,d\bm{x}\ \ \ \ \ \ \mbox{y}
\ \ \ \ \ \ P(1/2)=\int \varphi(\bm{x}) f_2(\bm{x})\,d\bm{x}\] donde las
integrales son sumas en el caso discreto. Si la regla discriminante es
no aleatorizada, entonces las probabilidades anteriores se pueden
expresar así:

\[P(2/1)=\int_{R_2} f_1(\bm{x})\,d\bm{x}\ \ \ \ \ \ \mbox{y}
\ \ \ \ \ \ P(1/2)=\int_{R_1} f_2(\bm{x})\,d\bm{x}.\]

Nótese cómo el problema discriminante se puede encuadrar dentro de la
teoría de la decisión. Un paso más en esta dirección nos conduce a la
comparación de reglas discriminantes.

\begin{definition}
\protect\hypertarget{def:unnamed-chunk-91}{}\label{def:unnamed-chunk-91}Una regla discriminante \(\varphi\) es preferible a otra \(\varphi'\) si

\[P_{\varphi}(2/1)\leq P_{\varphi'}(2/1)\] y

\[P_{\varphi}(1/2)\leq P_{\varphi'}(1/2).\]
\end{definition}

\textbf{Nota}. La relación ``ser preferible'' es una relación de orden parcial.

\begin{definition}
\protect\hypertarget{def:unnamed-chunk-92}{}\label{def:unnamed-chunk-92}Una regla discriminante se dice admisible si no hay ninguna otra regla
discriminante que sea preferible (estrictamente) a ella. Por ``preferible
estrictamente'' entendemos que, además de ser preferible, en alguno de
los dos errores de clasificación sea estrictamente mejor.
\end{definition}

\begin{theorem}
\protect\hypertarget{thm:unnamed-chunk-93}{}\label{thm:unnamed-chunk-93}Sea \(f_1\) la función de densidad o de probabilidad, según el caso, del
vector aleatorio \(\bm{x}\) condicionada al grupo \(G_1\) y \(f_2\) la
correspondiente función de densidad o probabilidad condicionada al grupo
\(G_2\).

Las reglas discriminantes de la forma

\[\varphi(\bm{x})=\left\{
{\displaystyle \begin{array}{ll}
1 & \mbox{si }\ {\displaystyle \frac{f_1(\bm{x})}{f_2(\bm{x})} > c, } \\
\gamma(\bm{x}) & \mbox{si }\ {\displaystyle \frac{f_1(\bm{x})}{f_2(\bm{x})} = c,} \\
0 & \mbox{si }\ {\displaystyle \frac{f_1(\bm{x})}{f_2(\bm{x})} < c.}
\end{array}}\right.\] son admisibles. Además, son las únicas reglas admisibles.
\end{theorem}

\begin{proof}
Es la misma que para el Lema de Neyman-Pearson, en el caso del
contraste de una hipótesis nula simple frente a una alternativa
simple.
\end{proof}

En lo que sigue veremos cómo distintos criterios permiten la elección de
una regla discriminante.

\section{Criterios para elegir una regla discriminante}\label{criterios-para-elegir-una-regla-discriminante}

\subsection{Minimizar la probabilidad total de clasificación incorrecta}\label{minimizar-la-probabilidad-total-de-clasificacion-incorrecta}

Pensemos en una población que se encuentra dividida en individuos
pertenecientes al grupo \(G_1\) e individuos pertenecientes al grupo
\(G_2\), con proporciones respectivas \(\pi_1\) y \(\pi_2=1-\pi_1\). Así,
\(\pi_1\) y \(\pi_2\) se convierten en probabilidades originales de
pertenencia a cada grupo, antes de observar el vector aleatorio
\(\bm{x}\). En consecuencia, la probabilidad total de clasificación
incorrecta es

\[P(R,f)=\pi_1 P(2/1) + \pi_2 P(1/2)\] denotando
\({R}=(R_1,R_2)^{\prime}\) al vector de regiones de clasificación y
\({f}=(f_1,f_2)^{\prime}\) al vector de densidades. Nuestro objetivo será
adoptar una regla de clasificación que presente la menor probabilidad
total de clasificación incorrecta. Recurrimos a un lema que nos ayudará
en este propósito.

\begin{lemma}
\protect\hypertarget{lem:unnamed-chunk-95}{}\label{lem:unnamed-chunk-95}La integral \(\int_{{R}} g(\bm{x})\,d\bm{x}\) alcanza su mínimo como
función de \({R}\) en

\[R^-=\{\bm{x}: g(\bm{x})<0\}.\]
\end{lemma}

\begin{proof}
Denotemos \(R^+=(R^-)^c=\{\bm{x}:g(\bm{x})\geq 0\}\). Para
cualquier conjunto \(R\), se tiene \[\begin{aligned}
\int_R g(\bm{x})\,d\bm{x} & = & \int_{R\cap R^-} g(\bm{x})\,dx + \int_{R\cap R^+} g(\bm{x})\,d\bm{x}
\geq  \int_{R\cap R^-} g(\bm{x})\,d\bm{x} \\
& = & \int_{R^-} g(\bm{x})\,d\bm{x} - \int_{R^c\cap R^-} g(\bm{x})\,d\bm{x}
\geq \int_{R^-} g(\bm{x})\,d\bm{x},\end{aligned}\] lo cual concluye la
demostración del presente lema.
\end{proof}

Desarrollando la probabilidad total de clasificación incorrecta,
obtenemos

\[\begin{aligned}
P({R},{f})&=\pi_1P(2/1)+\pi_2P(1/2)\\
&=\pi_1 (1-\int_{R_1} f_1(\bm{x})\,d\bm{x})
+\pi_2\int_{R_1} f_2(\bm{x})\,d\bm{x} \\
&= \pi_1+\int_{R_1} [\pi_2 f_2(\bm{x})-\pi_1 f_1(\bm{x})]\,d\bm{x}\end{aligned}\]

Por el lema anterior, como función de \(R_1\) alcanza su mínimo en
\(R_1=\{\bm{x}:\pi_2 f_2(\bm{x})-\pi_1 f_1(\bm{x})<0\}\), que se
corresponde con la regla de clasificación que asigna al grupo \(G_1\) si

\[\frac{f_1(\bm{x})}{f_2(\bm{x})}>\frac{\pi_2}{\pi_1}\] y al grupo \(G_2\)
en otro caso. La asignación en la frontera de \(R_1\) no debería ser
relevante, ya que no afecta a la probabilidad total. Sin embargo, esto
puede dar lugar a varias reglas con la misma probabilidad total. Pero si
el vector \(\bm{x}\) es continuo y se cumple alguna otra condición más,
\(P(f_1(\bm{x})/f_2(\bm{x})=\pi_2/\pi_1)=0\), por lo que serían reglas
equivalentes.

\textbf{\includegraphics[width=1em,height=1em]{bookdown_AM_files/figure-latex/fa-icon-00ef1bb1d064b635425d49f06b331298.pdf} ~\textcolor{excolor}{Ejemplo.}} Supongamos que \(f_i\) es la densidad de
\(N_d(\bm{\mu_i}, \bm{\Sigma_i})\), con
\(\bm{\Sigma_1}=\bm{\Sigma_2}=\bm{\Sigma}\). Vamos a calcular la regla de
clasificación óptima con el criterio anterior y las correspondientes
probabilidades de clasificación incorrecta. Las densidades adoptan la
forma

\[f_i(\bm{x})=(2\pi)^{-p/2}|\bm\Sigma|^{-1/2}
\exp\left[-\frac{1}{2} (\bm{x}-\bm{\mu_i})^{\prime}\bm{\Sigma}^{-1}(\bm{x}-\bm{\mu_i})\right]\]
y su cociente será

\[\begin{aligned}
\frac{f_1(\bm{x})}{f_2(\bm{x})}&=\exp\left[-\frac{1}{2} (\bm{x}-\bm{\mu_1})^{\prime}\bm{\Sigma}^{-1}(\bm{x}-\bm{\mu_1})
+\frac{1}{2} (\bm{x}-\bm{\mu_2})^{\prime}\bm{\Sigma}^{-1}(\bm{x}-\bm{\mu_2})\right] \\
& =  \exp\left[(\bm{\mu_1}-\bm{\mu_2})^{\prime}\bm{\Sigma}^{-1}\bm{x}
- \frac{1}{2} (\bm{\mu_1}-\bm{\mu_2})^{\prime}\bm{\Sigma}^{-1}(\bm{\mu_1}+\bm{\mu_2})\right].\end{aligned}\]
Tomando logaritmos, la regla de clasificación óptima asigna al grupo
\(G_1\) si

\[D(\bm{x})=\bm{\lambda}^{\prime}\left[\bm{x}-\frac{1}{2}(\bm{\mu_1}+\bm{\mu_2})\right]>\log(\pi_2/\pi_1)\]
siendo \(\bm{\lambda}=\bm\Sigma^{-1} (\bm{\mu_1}-\bm{\mu_2})\). La
ecuación \(D(\bm{x})=\log(\pi_2/\pi_1)\) define un hiperplano que separa
los dos grupos.

Para calcular las probabilidades de clasificación incorrecta, observamos
que, condicionado a cada grupo, la variable aleatoria \(D(\bm{x})\) tiene
distribución normal. Su media y varianza condicionadas son

\[\begin{aligned}
E[D(\bm{x})/G_1]&=\bm\lambda^{\prime}\left[\bm{\mu_1}-\frac{1}{2}(\bm{\mu_1}+\bm{\mu_2})\right]
=\frac{1}{2} \Delta^2, \\
\textnormal{Var}[D(\bm{x})/G_1]&=\textnormal{Var}[\bm\lambda^{\prime}\bm{x}/G_1]=\bm\lambda^{\prime}\bm\Sigma \bm\lambda=\Delta^2\end{aligned}\]
denotando
\(\Delta^2=(\bm{\mu_1}-\bm{\mu_2})^{\prime}\bm\Sigma^{-1}(\bm{\mu_1}-\bm{\mu_2})=\bm\lambda^{\prime}(\bm{\mu_1}-\bm{\mu_2})\)
al cuadrado de la distancia de Mahalanobis entre las medias \(\bm{\mu_1}\)
y \(\bm{\mu_2}\). Entonces

\[\begin{aligned}
P(2/1)&= P[D(\bm{x})\leq \log(\pi_2/\pi_1)/G_1]\\
&=P\left[\frac{D(\bm{x})-\Delta^2/2}{\Delta}\leq \frac{\log(\pi_2/\pi_1)-\Delta^2/2}{\Delta}
/G_1\right] \\
&=\Phi\left(\frac{\log(\pi_2/\pi_1)-\Delta^2/2}{\Delta}\right).\end{aligned}\]
De igual modo, \(E[D(\bm{x})/G_2]=-\frac{1}{2} \Delta^2\) y
\(\textnormal{Var}[D(\bm{x})/G_2]=\Delta^2\), y así:

\[\begin{aligned}
P(1/2)&= P[D(\bm{x})> \log(\pi_2/\pi_1)/G_2]\\
&=P\left[\frac{D(\bm{x})+\Delta^2/2}{\Delta}> \frac{\log(\pi_2/\pi_1)+\Delta^2/2}{\Delta}
/G_2\right] \\
&=1-\Phi\left(\frac{\log(\pi_2/\pi_1)+\Delta^2/2}{\Delta}\right).\end{aligned}\]
La probabilidad total de clasificación incorrecta queda:

\[\begin{aligned}
P(\mbox{Clasificación incorrecta})&=\pi_1 P(2/1)+\pi_2 P(1/2) \\
&=\pi_1 \Phi\left(\frac{\log(\pi_2/\pi_1)-\Delta^2/2}{\Delta}\right)\\
&\qquad +\pi_2 \left[1-\Phi\left(\frac{\log(\pi_2/\pi_1)+\Delta^2/2}{\Delta}\right)\right].\end{aligned}\]
Si \(\pi_1=\pi_2=\frac{1}{2}\), entonces la regla discriminante consiste
en clasificar en \(G_1\) si
\(D(\bm{x})=\bm\lambda^{\prime}[\bm{x}-\frac{1}{2}(\bm{\mu_1}+\bm{\mu_2})]>\log(\pi_2/\pi_1)=0\),
o equivalentemente, si

\[\bm\lambda^{\prime} \bm{x} > \frac{1}{2}\left(\bm\lambda^{\prime}\bm{\mu_1}+\bm\lambda^{\prime}\bm{\mu_2}\right)\]
esto es, si \(\bm\lambda^{\prime} \bm{x}\) está más próximo a
\(\bm\lambda^{\prime}\bm{\mu_1}\) que a \(\bm\lambda^{\prime}\bm{\mu_2}\). A
esto se le conoce como regla discriminante lineal de Fisher (1936).

\textbf{\includegraphics[width=1em,height=1em]{bookdown_AM_files/figure-latex/fa-icon-00ef1bb1d064b635425d49f06b331298.pdf} ~\textcolor{excolor}{Ejemplo.}} En el ejemplo anterior, supongamos que
\(\bm{\Sigma_1}\neq \bm{\Sigma_2}\). Entonces

\[\begin{aligned}
Q(\bm{x})&=\log(f_1(\bm{x})/f_2(\bm{x})) \\
&= \frac{1}{2}\log(|\bm{\Sigma_2}|/|\bm{\Sigma_1}|)
-\frac{1}{2}(\bm{x}-\bm{\mu_1})^{\prime}\bm{\Sigma_1}^{-1} (\bm{x}-\bm{\mu_1})
+\frac{1}{2}(\bm{x}-\bm{\mu_2})^{\prime}\bm{\Sigma_2}^{-1} (\bm{x}-\bm{\mu_2}) \\
&= \frac{1}{2}\log(|\bm{\Sigma_2}|/|\bm{\Sigma_1}|)
-\frac{1}{2} \left[ \bm{x}^{\prime}(\bm{\Sigma_1}^{-1}-\bm{\Sigma_2}^{-1}) \bm{x}
-2\bm{x}^{\prime}(\bm{\Sigma_1}^{-1} \bm{\mu_1} -\bm{\Sigma_2}^{-1} \bm{\mu_2}) \right]\end{aligned}\]
y la regla óptima clasifica en \(G_1\) si \(Q(\bm{x})>\log(\pi_2/\pi_1)\).
Pero, mientras \(D(\bm{x})\) es una función lineal, ahora \(Q(\bm{x})\) es
una función cuadrática, de modo que la frontera entre las dos regiones
ya no es un hiperplano.

\subsection{Razón de verosimilitudes}\label{razon-de-verosimilitudes}

Nos situamos en el contexto en que las probabilidades a priori, \(\pi_1\)
y \(\pi_2\), son desconocidas o sencillamente no se tienen en cuenta. Una
regla de clasificación debería asignar un individuo al grupo más
verosímil, que es aquel que otorga mayor probabilidad a su resultado
muestral. Esto equivale a clasificar en el grupo \(G_1\) si
\(f_1(\bm{x})/f_2(\bm{x})>1\), lo cual coincide con el criterio anterior
para \(\pi_1=\pi_2=1/2\).

\subsection{Minimizar el coste total de clasificación incorrecta}\label{minimizar-el-coste-total-de-clasificacion-incorrecta}

Si asociamos unos costes \(C(2/1)\) y \(C(1/2)\) a cada uno de los errores
de clasificación, entonces el coste total esperado es

\[C=\pi_1P(2/1)C(2/1)+\pi_2P(1/2)C(1/2).\] Podemos repetir los argumentos
que obtuvieron la regla que minimiza la probabilidad total de
clasificación incorrecta sustituyendo \(\pi_1\) y \(\pi_2\) por
\(\pi_1C(2/1)\) y \(\pi_2C(1/2)\), ya que en aquellos argumentos no se usó
que \(\pi_1+\pi_2=1\). Así, obtendremos que la regla que hace mínimo el
coste total esperado consiste en clasificar en el grupo \(G_1\) si

\[\frac{f_1(\bm{x})}{f_2(\bm{x})}>\frac{\pi_2C(1/2)}{\pi_1C(2/1)}.\] Por
supuesto, si \(C(2/1)=C(1/2)\), coincide con el criterio que minimiza la
probabilidad total de clasificación incorrecta.

\subsection{Maximizar la probabilidad a posteriori}\label{maximizar-la-probabilidad-a-posteriori}

Consideramos la probabilidad de que el individuo proceda del grupo
\(G_1\), conocido el valor del vector aleatorio \(\bm{x}\). Dicha
probabilidad se conoce como probabilidad a posteriori y se puede
calcular mediante el teorema de Bayes:

\[P(G_1/\bm{x}=\bm{x_0})=\frac{\pi_1 f_1(\bm{x_0})}{\pi_1 f_1(\bm{x_0})+\pi_2 f_2(\bm{x_0})}.\]
La probabilidad a posteriori del grupo \(G_2\) se define y obtiene de la
misma manera:

\[P(G_2/\bm{x}=\bm{x_0})=\frac{\pi_2 f_2(\bm{x_0})}{\pi_1 f_1(\bm{x_0})+\pi_2 f_2(\bm{x_0})}.\]
Una regla de clasificación razonable asignaría al grupo con mayor
probabilidad a posteriori, esto es, clasificaría en el grupo \(G_1\) si

\[P(G_1/\bm{x})>P(G_2/\bm{x})\] lo cual coincide con
\(f_1(\bm{x})/f_2(\bm{x})>\pi_2/\pi_1\).

\subsection{Criterio minimax}\label{criterio-minimax}

Un procedimiento que minimice la probabilidad total de clasificación
incorrecta puede no prestar suficiente atención a uno de los grupos, en
especial si su probabilidad a priori es pequeña. Por ejemplo, esto
ocurre en el problema de diagnosis de una enfermedad poco frecuente,
donde la probabilidad de un falso negativo puede ser alta sin afectar
mucho a la probabilidad total de clasificación incorrecta.

Esto constituye una justificación para el uso de un criterio más
conservador, como es el minimax. En este caso, escogemos la regla tal
que el máximo de sus dos probabilidades de clasificación incorrecta sea
mínima. Consideremos \(P\) una regla de clasificación cualquiera. Entonces, para todo \(\alpha\in[0,1]\),

\[\begin{aligned}\max\{P(2/1),P(1/2)\}&\geq (1-\alpha)P(2/1)+\alpha P(1/2)\\
&\geq (1-\alpha)P_{\alpha}(2/1)+\alpha P_{\alpha}(1/2),
\end{aligned}\]
siendo \(P_{\alpha}\) la regla de clasificación que asigna a \(G_1\) si
\(f_1(X)/f_2(X)>\alpha/(1-\alpha)\). La primera desigualdad anterior es
inmediata y la segunda se debe a que \(P_{\alpha}\) es la regla que hace
mínima la probabilidad total de clasificación incorrecta bajo la
distribución a priori: \(\pi_2=\alpha\), \(\pi_1=1-\alpha\).

Tomemos \(c=\frac{\alpha_0}{1-\alpha_0}\) (o equivalentemente, \(\alpha_0\))
de modo que \(P_{\alpha_0}(2/1)=P_{\alpha_0}(1/2)\). Entonces

\[\begin{aligned}
\max\{P(2/1),P(1/2)\} & \geq  (1-\alpha_0)P_{\alpha_0}(2/1)+\alpha_0 P_{\alpha_0}(1/2)
=(1-\alpha_0+\alpha_0) P_{\alpha_0}(2/1) \\
&=P_{\alpha_0}(2/1)=\max\{P_{\alpha_0}(2/1),P_{\alpha_0}(1/2)\}\end{aligned}\]
lo cual demuestra que \(P_{\alpha_0}\) es la regla minimax.

En consecuencia, la regla minimax será la que clasifica en el grupo
\(G_1\) si \(f_1(\bm{x})/f_2(\bm{x})>c\), siendo \(c\) tal que
\(P_c(2/1)=P_c(1/2)\).

\textbf{\includegraphics[width=1em,height=1em]{bookdown_AM_files/figure-latex/fa-icon-00ef1bb1d064b635425d49f06b331298.pdf} ~\textcolor{excolor}{Ejemplo.}} Si las dos distribuciones condicionadas a cada grupo son
normales con la misma matriz de covarianzas, la regla de clasificación
minimax asigna al grupo \(G_1\) si

\[D(\bm{x})>\log c\] siendo \(c\) tal que

\[\Phi\left(\frac{\log c -\Delta^2/2}{\Delta}\right)
=1-\Phi\left(\frac{\log c +\Delta^2/2}{\Delta}\right)\] Esta ecuación
tiene solución \(\log c=0\), o lo que es lo mismo, \(c=1\), lo cual coincide
con el criterio de la razón de verosimilitudes.

\textbf{\includegraphics[width=1em,height=1em]{bookdown_AM_files/figure-latex/fa-icon-00ef1bb1d064b635425d49f06b331298.pdf} ~\textcolor{excolor}{Ejemplo.}} Una población se encuentra dividida en dos grupos \(G_1\) y
\(G_2\). Se realiza una observación de una variable que tiene distribución
exponencial, cuyo parámetro depende del grupo. Si procede del grupo
\(G_1\) el parámetro es 1, y si procede del grupo \(G_2\) el parámetro es 2.

a. Supongamos una distribución a priori equiprobable sobre los dos
grupos. Calcula la regla que minimiza la probabilidad de clasificación
incorrecta. Cuál es ese valor mínimo de la probabilidad de clasificación
incorrecta?

b. Calcula la regla de clasificación óptima con el criterio minimax.

\textbf{\includegraphics[width=1em,height=1em]{bookdown_AM_files/figure-latex/fa-icon-00ef1bb1d064b635425d49f06b331298.pdf} ~\textcolor{excolor}{Ejemplo.}} Consideremos una población dividida en dos grupos. A cada
individuo de la población le podemos observar una variable aleatoria
discreta que toma valores 1,2,3,4 y 5, con probabilidades respectivas
\(0.05\), \(0.15\), \(0.5\), \(0.25\), \(0.05\) para el primer grupo y \(0.3\), \(0.1\), \(0.1\), \(0.1\),
\(0.4\) para el segundo grupo.

a. Suponiendo una distribución a priori equiprobable entre los dos
grupos, determina la regla discriminante que hace mínima la probabilidad
total de clasificación incorrecta.

b. Determina la regla discriminante óptima con el criterio minimax.

\section{Estimación de la regla discriminante}\label{estimacion-de-la-regla-discriminante}

La situación más común en la práctica consiste en que desconocemos las
distribuciones de \(\bm{x}\) condicionales a cada uno de los grupos. En su
lugar disponemos de muestras procedentes de cada grupo,

\[\bm{x_{11}},\ldots,\bm{x_{1n_1}}\ \ \mbox{del grupo}\ G_1.\]

\[\bm{x_{21}},\ldots,\bm{x_{2n_2}}\ \ \mbox{del grupo}\ G_2.\] conocidas
como \emph{muestras de entrenamiento}, que nos sirven para estimar la regla
discriminante, por lo general, a través de estimaciones de las
distribuciones respectivas:

\[\hat{R}_1=\left\{\bm{x}:\frac{\hat{f}_1(\bm{x})}{\hat{f}_2(\bm{x})}>c\right\}.\]

Si se supone algún modelo paramétrico para las distribuciones, que
podemos representar en la forma \(f_1(\bm{x})=f_1(\bm{x}/\theta_1)\) y
\(f_2(\bm{x})=f_2(\bm{x}/\theta_2)\), los estimadores serían
\(\hat{f}_1(\bm{x})=f_1(\bm{x}/\hat{\theta}_1)\) y
\(\hat{f}_2(\bm{x})=f_2(\bm{x}/\hat{\theta}_2)\), donde \(\hat{\theta}_1\) y
\(\hat{\theta}_2\) representan estimadores de los parámetros. En ausencia
de modelo paramétrico, \(\hat{f}_1(\bm{x})\) y \(\hat{f}_2(\bm{x})\) serían
algún tipo de estimador no paramétrico de la densidad.

La regla discriminante estimada, \(\hat{R}=(\hat{R}_1,\hat{R}_2)\), se
aplicaría más adelante para clasificar a un nuevo individuo en uno u
otro grupo, en base a su resultado en el vector \(\bm{x}\).

En la sección anterior vimos que la regla discriminante óptima era
\(R=(R_1,R_2)\), donde las regiones de clasificación \(R_1\) y \(R_2\)
dependían de las verdaderas funciones \(f_1\) y \(f_2\). Como \(f_1\) y \(f_2\)
son desconocidas, hemos tenido que sustituirlas por estimaciones
\(\hat{f}_1\) y \(\hat{f}_2\) obteniendo una regla que ya no será óptima.
Nuestra esperanza reside en que sea parecida a la regla óptima.

\subsection{Tasas de error en reglas estimadas}\label{tasas-de-error-en-reglas-estimadas}

Nos interesa conocer cómo afecta la estimación de la regla discriminante
a los errores de clasificación. Debemos distinguir entre:

\begin{description}
\item[Tasas de error óptimas]
Son las probabilidades de error de la regla discriminante óptima.
Las denotamos

\[e_{i,opt}=\int_{R_j}f_i(\bm{x})\,d\bm{x}\qquad i=1,2\ \ j\neq i.\]

\[e_{opt}=\pi_1 e_{1,opt}+\pi_2 e_{2,opt}.\]
\item[Tasas de error efectivas]
Son las probabilidades de error de la regla discriminante estimada.

\[e_{i,efe}=\int_{\hat{R}_j}f_i(\bm{x})\,d\bm{x}\qquad i=1,2\ \ j\neq i.\]

\[e_{efe}=\pi_1 e_{1,efe}+\pi_2 e_{2,efe}.\] Nótese que las regiones
\((\hat{R}_1,\hat{R}_2)\) se obtienen de las muestras de
entrenamiento. En la medida en que estas muestras son aleatorias,
las regiones de clasificación serán aleatorias, y también lo serán
las tasas de error efectivas. Por tanto, tiene sentido calcular su
esperanza.
\item[Esperanza de las tasas de error efectivas]
No son más que la esperanza de las tasas de error anteriores,
calculadas respecto de la distribución de probabilidad de las
muestras de entrenamiento: \(E(e_{i,efe})\) y

\[E(e_{efe})=\pi_1 E(e_{1,efe})+\pi_2 E(e_{2,efe}).\]
\end{description}

Es inmediato que

\[e_{i,opt}\leq e_{i,efe}\qquad\mbox{y}\qquad e_{opt}\leq e_{efe}.\]

Las tasas de error óptimas se refieren a la situación ideal en la que
las distribuciones de cada grupo son conocidas. Sin embargo, en el caso
que nos ocupa, las que tienen mayor interés son las tasas de error
efectivas, pues representan las probabilidades de error en las que se va
a incurrir cuando se clasifique a un nuevo individuo mediante la regla
estimada.

Por desgracia, las tasas de error efectivas son desconocidas ya que
dependen de \(f_1\) y \(f_2\), que representan la distribución de la nueva
observación. A continuación planteamos diversos estimadores de las tasas
de error efectivas.

\subsection{Estimación de las tasas de error efectivas}\label{estimacion-de-las-tasas-de-error-efectivas}

\begin{description}
\item[Estimadores plug-in]
Se obtienen sustituyendo \(f_1\) y \(f_2\) por \(\hat{f}_1\) y
\(\hat{f}_2\).

\[e_{i,pl}=\int_{\hat{R}_j}\hat{f}_i(\bm{x})\,d\bm{x}\qquad i=1,2\ \ j\neq i.\]
\item[Tasas de error aparentes]
Si a cada individuo de la muestra de entrenamiento le aplicamos la
regla discriminante estimada tendremos algunos individuos bien
clasificados y otros mal clasificados. Pues bien, las tasas de error
aparentes son las frecuencias relativas de observaciones mal
clasificadas.

\[e_{i,apa}=\frac{m_i}{n_i}\qquad i=1,2\] siendo \(m_i\) el número de individuos de la muestra de entrenamiento procedente
del grupo \(G_i\) que son mal clasificados por la regla estimada. Las
tasas de error aparentes tienden a infraestimar las tasas de error
efectivas ya que estamos aplicando la regla estimada a los mismos
datos con los que ha sido construida. El mismo fenómeno ocurre con
los estimadores plug-in, ya que entre la regla estimada y
\(\hat{f}_i\) hay una relación de optimalidad que no se cumple con las
tasas de error efectivas. Por este motivo, se proponen correcciones
de las tasas de error aparentes.
\item[Corrección por validación cruzada]
Consiste en construir la regla discriminante con todas las
observaciones menos una, y aplicar después esa regla para clasificar
al individuo omitido. Las tasas de error se calcularían después como
las frecuencias relativas de individuos mal clasificados.
\item[Corrección mediante bootstrap]
Se trata de una corrección del sesgo de las tasas de error
aparentes. Siguiendo las ideas del bootstrap, extraemos muestras
artificiales de la forma:

\[\bm{x_{11}}^*,\ldots,\bm{x_{1n_1}}^*\ \ \mbox{i.i.d. de la distribución equiprobable sobre}
\ \ \bm{x_{11}},\ldots,\bm{x_{1n_1}}.\]

\[\bm{x_{21}}^*,\ldots,\bm{x_{2n_2}}^*\ \ \mbox{i.i.d. de la distribución equiprobable sobre}
\ \ \bm{x_{21}},\ldots,\bm{x_{2n_2}}.\]

Sobre las muestras
artificiales aplicamos los mismos procedimientos de estimación de la
regla discriminante y cálculo de las tasas de error aparentes,
\(m_1^*/n_1\) y \(m_2^*/n_2\), siendo \(m_i^*\) el número de individuos de
la muestra bootstrap, del grupo \(G_i\), que han sido mal clasificados
por la regla estimada bootstrap. Denotemos \(m_i^{orig}\) al número de
individuos de la muestra original que son mal clasificados por la
regla bootstrap. La estimación bootstrap del sesgo sería:
\[E^*\left(\frac{m_i^*}{n_i}-\frac{m_i^{orig}}{n_i}\right)\] donde
\(E^*\) denota la esperanza respecto del muestreo bootstrap, que en la
práctica se puede aproximar mediante la media obtenida tras extraer
muchas muestras bootstrap.

La corrección bootstrap del sesgo daría lugar a las tasas de error
siguientes:
\[\frac{m_i}{n_i}-E^*\left(\frac{m_i^*}{n_i}-\frac{m_i^{orig}}{n_i}\right)\]
\end{description}

A partir de cualquiera de las tasas de error anteriores, también se
puede estimar la probabilidad total de clasificación incorrecta. Para
ello, podemos suponer que las probabilidades a priori \(\pi_1\) y \(\pi_2\)
son conocidas, o en otro caso se pueden estimar mediante las
proporciones \(n_1/(n_1+n_2)\) y \(n_2/(n_1+n_2)\), suponiendo que las dos
muestras han sido extraídas al azar de la población formada al juntar
los dos grupos.

\textbf{\includegraphics[width=1em,height=1em]{bookdown_AM_files/figure-latex/fa-icon-00ef1bb1d064b635425d49f06b331298.pdf} ~\textcolor{excolor}{Ejemplo.}} Supongamos dos muestras de entrenamiento

\[\bm{x_{11}},\ldots,\bm{x_{1n_1}}\in N_d(\bm{\mu_1},\bm\Sigma) \ \mbox{del grupo}\ G_1.\]

\[\bm{x_{21}},\ldots,\bm{x_{2n_2}}\in N_d(\bm{\mu_2},\bm\Sigma) \ \mbox{del grupo}\ G_2.\]
Estamos ante dos distribuciones normales con la misma matriz de
covarianzas, por lo que, según hemos visto en la sección anterior, la
regla de clasificación óptima asignaría a un nuevo individuo al grupo
\(G_1\) si

\[D(\bm{x})=\bm\lambda^{\prime}\left[\bm{x}-\frac{1}{2}(\bm{\mu_1}+\bm{\mu_2})\right]>\log c\]
siendo \(\bm\lambda=\bm\Sigma^{-1} (\bm{\mu_1}-\bm{\mu_2})\).

Como \(\bm{\mu_1}\), \(\bm{\mu_2}\) y \(\bm\Sigma\) son desconocidos, los
sustituimos por sus estimaciones en base a las muestras de entrenamiento
\(\overline{\bm{x_1}}=n_1^{-1}\sum_{i=1}^{n_1} \bm{x_{1i}}\),
\(\overline{\bm{x_2}}=n_2^{-1}\sum_{i=1}^{n_2} \bm{x_{2i}}\) y

\[\bm{S_c}=\frac{(n_1-1)\bm{S_{c1}}+(n_2-1)\bm{S_{c2}}}{n_1+n_2-2}\]
siendo
\(\bm{S_{c1}}=(n_1-1)^{-1}\sum_{i=1}^{n_1} (\bm{x_{1i}}-\overline{\bm{x_1}})(\bm{x_{1i}}-\overline{\bm{x_1}})^\prime{}\),
y análogamente \(\bm{S_{c2}}\).

La regla estimada resulta

\[\hat{D}(\bm{x})=\hat{\bm\lambda}^{\prime}\left[\bm{x}-\frac{1}{2}(\overline{\bm{x_1}}+\overline{\bm{x_2}})\right]>\log c\]
siendo \(\hat{\bm\lambda}=\bm{S_c}^{-1} (\overline{\bm{x_1}}-\overline{\bm{x_2}})\).

Las tasas de error efectivas serían

\[e_{1,efe}=\Phi\left(\frac{\log c - \hat{D}(\bm{\mu_1})}{\sqrt{\hat{\bm\lambda}^\prime{}\bm\Sigma\hat{\bm\lambda}}}
\right)\qquad\mbox{y}\qquad e_{2,efe}=
\Phi\left(\frac{-\log c + \hat{D}(\bm{\mu_2})}{\sqrt{\hat{\bm\lambda}^\prime{}\bm\Sigma\hat{\bm\lambda}}}
\right).\]

Sustituyendo en las expresiones anteriores \(\bm{\mu_1}\), \(\bm{\mu_2}\) y
\(\bm\Sigma\) por sus estimaciones, obtenemos las estimaciones plug-in de
las tasas de error

\[e_{1,pl}=\Phi\left( \frac{\log c - \hat{\Delta}^2/2}{\hat{\Delta}} \right)
\qquad\mbox{y}\qquad
e_{2,pl}=\Phi\left( \frac{-\log c - \hat{\Delta}^2/2}{\hat{\Delta}} \right)\]
donde
\(\hat{\Delta}^2=\left(\overline{\bm{x_1}}-\overline{\bm{x_2}}\right)^\prime{} \bm{S_c}^{-1} \left(\overline{\bm{x_1}}-\overline{\bm{x_2}}\right)\)
es el cuadrado de la distancia de Mahalanobis en términos muestrales.
Como ocurre en general, los estimadores \(e_{i,pl}\) no son buenos
estimadores de \(e_{i,efe}\). Para corregirlos se ha propuesto utilizar
\(\tilde{\Delta}\) en lugar de \(\hat{\Delta}\) en la expresión de
\(e_{i,pl}\), siendo
\[\tilde{\Delta}^2=\frac{n_1+n_2-d-3}{n_1+n_2-2} \hat{\Delta}^2.\]

Por supuesto, también se pueden llevar a cabo las estimaciones mediante
las tasas de error aparentes, así como sus correcciones por validación
cruzada o por bootstrap.

\section{Varios grupos con distribuciones conocidas}\label{varios-grupos-con-distribuciones-conocidas}

Supongamos que están definidos \(I\) grupos, y queremos construir una
regla discriminante que permita clasificar en alguno de los grupos en
función del resultado de un vector aleatorio \(\bm{x}\). Consideramos
\(\pi_1,\ldots,\pi_I\) la distribucion a priori sobre los \(I\) grupos, y
\(f_1,\ldots,f_I\) las funciones de densidad o probabilidad de \(\bm{x}\)
condicionadas a cada uno de los grupos. La regla discriminante óptima
consistirá en

\[\mbox{Clasificar en el grupo}\ \  g \ \ \mbox{si}\qquad \pi_g f_g(\bm{x})=\max_{i\in\{1,\ldots,I\}} \pi_i f_i(\bm{x}).\] Esta regla es la que clasifica en el grupo con
mayor probabilidad a posteriori, y es la regla de máxima verosimilitud
si se omiten las probabilidades a priori.

\subsection{Clasificación óptima de varias poblaciones normales}\label{clasificacion-optima-de-varias-poblaciones-normales}

Pensemos que los grupos tienen las mismas probabilidades a priori, y que
las distribuciones \(f_i\) son normales de medias
\(\bm{\mu_1},\ldots,\bm{\mu_I}\), y con la misma matriz de covarianzas
\(\bm\Sigma\). Entonces la regla anterior se reduce a

\begin{equation}
\mbox{Clasificar en el grupo}\ \  g \ \ \mbox{si}\ \left(\bm{x}-\bm{\mu_g}\right)^\prime{}\bm{\Sigma}^{-1}\left(\bm{x}-\bm{\mu_g}\right)
=\min_{i\in\{1,\ldots,I\}} \left(\bm{x}-\bm{\mu_i}\right)^\prime{}\bm{\Sigma}^{-1}\left(\bm{x}-\bm{\mu_i}\right)
\label{eq:reglaoptima}
\end{equation}
esto es, se clasifica la observación \(\bm{x}\) en el grupo cuya media
está más próxima, bajo la distancia de Mahalanobis.

Cuando sólo hay dos grupos, hemos visto que esta regla es equivalente a
efectuar la transformación lineal \(\bm{\lambda}^\prime{}\bm{x}\) (siendo
\(\bm{\lambda}=\bm{\Sigma}^{-1}(\bm{\mu_1}-\bm{\mu_2})\)) y clasificar en
el grupo 1 si \(\bm{\lambda}^\prime{}\bm{x}\) está más próximo a
\(\bm{\lambda}^\prime{}\bm{\mu_1}\) que a
\(\bm{\lambda}^\prime{}\bm{\mu_2}\). Cuando hay varios grupos, la regla
discriminante también se convierte en un criterio lineal, aunque
necesitaremos más de una transformación lineal.

Pensemos ahora que se dispone de una muestra procedente de cada grupo

\[\begin{array}{ccccc}
\bm{x_{11}} &  \cdots & \bm{x_{1,n_1}} & \mbox{de una población (o grupo)} & N_d(\bm{\mu_1},\bm{\Sigma_1}) \\
\vdots &  \vdots & \vdots & \vdots & \vdots \\
\bm{x_{I1}}  &  \cdots & \bm{x_{I,n_I}} &\mbox{de una población (o grupo)} & N_d(\bm{\mu_I},\bm{\Sigma_I})
\end{array}\]

Este tipo de observaciones son las que se obtienen en un modelo de
análisis multivariante de la varianza con matrices de covarianzas
desiguales.

En esta situación, suponiendo unas probabilidades a priori
\(\pi_1,\ldots,\pi_I\), la regla óptima dada por la expresión
\eqref{eq:reglaoptima}, se puede reducir a

\[\mbox{Clasificar en el grupo}\ \  g \ \ \mbox{si}\qquad d_g^Q (\bm{x})=\max_{i\in\{1,\ldots,I\}} d_i^Q(\bm{x})\]
siendo

\[d_i^Q(\bm{x})=-\frac{1}{2}\log \left|\bm{\Sigma_i}\right| -\frac{1}{2}\left(\bm{x}-\bm{\mu_i}\right)^\prime{}\bm{\Sigma_i}^{-1}\left(\bm{x}-\bm{\mu_i}\right)
+\log \pi_i\qquad\forall i\in\{1,\ldots,I\}.\]

Las funciones \(d_i^Q(\bm{x})\) se pueden denominar \textbf{funciones de
clasificación}, pues para clasificar a un individuo con observación
\(\bm{x}\), se calculan las funciones de clasificación que le corresponden
en cada grupo, \(d_i^Q(\bm{x})\), y se asigna al grupo cuya función sea
más grande. El superíndice \(Q\) en \(d_i^Q(\bm{x})\) se debe a que estas
funciones son cuadráticas, pues como ya vimos en el tema anterior, si
las matrices de covarianzas son diferentes dentro de cada grupo,
entonces la regla óptima es cuadrática. La muestra anterior servirá para
estimar la regla discriminante, sustituyendo en las funciones de
clasificación los vectores de medias y las matrices de covarianzas por
sus análogos empíricos,

\[\hat{d}_i^Q(\bm{x})=-\frac{1}{2}\log \left|\bm{S_i}\right|
 -\frac{1}{2}\left(\bm{x}-\overline{\bm{x_{i}}}\right)^\prime{}\bm{S_i}^{-1}\left(\bm{x}-\overline{\bm{x_{i}}}\right)
+\log \pi_i\qquad\forall i\in\{1,\ldots,I\}.\]

Si se puede suponer que las matrices de covarianzas dentro de cada grupo
son iguales, esto es, \(\bm{\Sigma_1}=\cdots=\bm{\Sigma_I}=\bm\Sigma\),
entonces en las funciones de clasificación se puede suprimir la parte
cuadrática, pues es la misma para todos los grupos, y resultan las
siguientes funciones de clasificación lineales:

\[d_i(\bm{x})=\bm{\mu_i}^{\prime}\bm{\Sigma}^{-1}\bm{x}-\frac{1}{2}\bm{\mu_i}^{\prime}\bm{\Sigma}^{-1}\bm{\mu_i}+\log \pi_i
\qquad\forall i\in\{1,\ldots,I\}\] que se estiman mediante

\[\hat{d}_i(\bm{x})=\overline{\bm{x_{i}}}^{\prime}\bm{S}^{-1}\bm{x}-\frac{1}{2}\overline{\bm{x_{i}}}^{\prime}\bm{S}^{-1}\overline{\bm{x_{i}}}+\log \pi_i
\qquad\forall i\in\{1,\ldots,I\}\] donde
\(\bm{S}=(n-I)^{-1}\sum_{i=1}^I (n_i-1)\bm{S_i}\) es la estimación de la
matriz de covarianzas común.

En estas funciones lineales los coeficientes vienen dados por los
vectores \(\overline{\bm{x_{i}}}^{\prime}\bm{S}^{-1}\) y las constantes
se obtienen de
\(\log \pi_i-\frac{1}{2}\overline{\bm{x_{i}}}^{\prime}\bm{S}^{-1}\overline{\bm{x_{i}}}\).

\chapter{Técnicas de formación de grupos}\label{tecnicas-de-formacion-de-grupos}

En este capítulo haremos una breve introducción a los procedimientos de
formación de grupos. Lo primero que destacamos es la amplitud y
diversidad de estos métodos, tanto en sus objetivos como en los
procedimientos que emplean. El aspecto común a todos ellos es el
propósito de formar grupos que todavía no están definidos. En esto se
diferencia del análisis discriminante, en el cual los grupos están
perfectamente definidos y el propósito es clasificar a un individuo en
alguno de los grupos. Como decimos, las técnicas de formación de grupos
pretenden formar grupos, reconocer patrones o estructuras dentro de la
población general, en base a cierto tipo de observaciones.

Podemos plantear una clasificación de las técnicas de formación de
grupos en:

\begin{description}
\item[Métodos jerárquicos]
También se conocen como métodos de \emph{taxonomía numérica}, pues tienen
mucha semejanza con la taxonomía de los seres vivos en Biología.
Partiendo de una matriz de distancias entre individuos, el objetivo
es construir un árbol de clasificación en distintos niveles. En el
nivel más bajo estarían grupos de individuos muy próximos o
semejantes (por ejemplo, las especies). A niveles superiores se
crean grupos como resultado de la agregación de los grupos de nivel
inferior (por ejemplo, familias, tipos o reinos).
\item[Métodos de particionamiento]
Al contrario de los métodos jerárquicos, los métodos de
particionamiento sólo pretenden una división de los individuos en un
conjunto de grupos, todos al mismo nivel, que constituyen una
partición de la población original.
\end{description}

Los métodos que estudiaremos en este tema no establecen suposiciones
sobre el modelo estadístico que generó los datos, y en este sentido se
deben interpretar como métodos descriptivos que ayudan a la comprensión
sobre las propiedades que presenta un conjunto de datos multivariantes.
Además, como cualquier método descriptivo, puede servir como punto de
partida que sugiere propiedades de la población, que merezcan un
análisis más detallado.

Aún así, existen también métodos de formación de grupos, que están
basados en un modelo estadístico, que refleja propiedades previsibles
sobre cómo están constituidos los grupos, por ejemplo, la normalidad de
la distribución. En este contexto, se pueden adoptar métodos de
inferencia para la construcción de los grupos. No trataremos aquí estos
métodos basados en el modelo, y nos restringiremos a los métodos más
conocidos de formación de grupos.

\section{Métodos jerárquicos.}\label{metodos-jerarquicos}

Como ya se expresó en la introducción, los métodos jerárquicos parten de
una matriz de distancia entre individuos, y en base a ella pretenden un
agrupamiento de los individuos a distintos niveles. En el nivel más bajo
cada grupo estaría formado por individuos, mientras que los grupos a
niveles superiores serían el resultado de agregar grupos de niveles
inferiores. Los algoritmos que se emplean para crear esta jerarquía
pueden ser de dos tipos:

\begin{description}
\item[Aglomerativos]
Partiendo de los individuos, se construyen grupos formados por
individuos, para después construir grupos, mediante la agregación de
los grupos ya formados en etapas anteriores.
\item[Divisivos]
Partiendo del grupo total formado por todos los individuos, se
genera una división en subgrupos (generalmente en dos), que más
adelante vuelven a ser subdivididos.
\end{description}

Nos centraremos en los algoritmos aglomerativos, que son los más
comunes. El algoritmo constaría de los siguientes pasos:

\begin{itemize}
\item
  \textbf{Paso 1 (Comienzo)}: Se definen los grupos o \emph{clusters}
  \(C_1,\ldots,C_n\), que está formados cada uno por un individuo.
\item
  \textbf{Paso 2}: Se buscan los dos grupos \(C_i\) y \(C_j\) que estén más
  próximos, se juntan y consecuentemente se reduce el número de
  grupos.
\item
  \textbf{Paso 3}: Se recalculan las distancias de todos los demás grupos al
  nuevo grupo, formado al juntar \(C_i\) y \(C_j\).
\item
  \textbf{Paso 4}: Si el número de grupos es uno, se detiene el algoritmo. En
  otro caso, se vuelve al paso 2.
\end{itemize}

En este algoritmo queda por determinar cómo se efectúa el paso 3. De
hecho, bajo esta misma estructura del algoritmo aglomerativo, los
métodos concretos se diferencian en cómo resuelven este paso. Hay
variadas formas de hacerlo, pero vamos a mencionar las tres que
consideramos principales.

\begin{description}
\item[Método del mínimo:]
Consiste en definir
\(d(C_k,C_i\cup C_j)=\min\{d(C_k,C_i),d(C_k,C_j)\}\), donde
\(d(C_r,C_s)\) denota la distancia del grupo \(C_r\) al grupo \(C_s\).
Esto es equivalente a definir la distancia entre dos grupos así:

\[d(C_r,C_s)=\min_{i\in C_r, j\in C_s} d_{ij}\] siendo \(d_{ij}\) la
distancia entre los individuos \(i\) y \(j\), que en este caso
pertenecen a los grupos \(C_r\) y \(C_s\), respectivamente.
\item[Método del máximo:]
Consiste en definir
\(d(C_k,C_i\cup C_j)=\max\{d(C_k,C_i),d(C_k,C_j)\}\). En este caso es
equivalente a definir la distancia entre dos grupos como la mayor
distancia entre sus individuos:

\[d(C_r,C_s)=\max_{i\in C_r, j\in C_s} d_{ij}.\]
\item[Método del promedio:]
Consiste en definir
\(d(C_k,C_i\cup C_j)=\frac{n_i}{n_i+n_j} d(C_k,C_i)+ \frac{n_j}{n_i+n_j} d(C_k,C_j)\),
siendo \(n_i\) y \(n_j\) el número de individuos en los grupos \(C_i\) y
\(C_j\), respectivamente. Es equivalente a definir la distancia entre
dos grupos como el promedio de las distancias entre sus individuos:

\[d(C_r,C_s)=\frac{1}{n_r n_s}  \sum_{i\in C_r, j\in C_s} d_{ij}.\]
\end{description}

\section{Métodos de particionamiento.}\label{metodos-de-particionamiento}

Los métodos de particionamiento, según ya se dijo en la introducción,
pretenden una partición de los individuos de la muestra en \(k\) grupos,
en base a los valores de las variables observadas en cada individuo. Por
supuesto, los grupos se formarán por proximidad en el espacio
\(d\)-dimensional, siendo \(d\) el número de variables.

Un criterio natural para la formación de los grupos, consistiría en
elegir la partición en grupos que haga mínima la variabilidad dentro de
cada grupo, medida por la suma de cuadrados intra-grupo. El criterio
parece sencillo pero el problema será la imposibilidad de recorrer todas
las particiones posibles de \(n\) individuos en \(k\) grupos. Sólo \(n=15\)
datos y \(k=3\) grupos son suficientes para generar más de dos millones de
posibles particiones. Con valores de \(n\) y \(k\) algo más grandes, el
número de particiones crece de manera desorbitada.

Ante esta situación, es preciso adoptar un algoritmo que partiendo de
una solución inicial razonable, llegue a una solución mejor de acuerdo
con el criterio, mediante pasos que aporten mejoras sucesivas. Nótese
que, a pesar de existir muchas particiones, la inmensa mayoría son
absurdas como solución del problema, pues carecen de sentido todos los
grupos que se formen tomando observaciones "salteadas" en el espacio.

Vamos a considerar únicamente el algoritmo de las \(k\)-medias, del que
hay múltiples versiones, aunque su forma más sencilla podría constar de
los pasos siguientes.

\begin{itemize}
\item
  \textbf{Paso 1}: Crear una partición inicial en \(k\) grupos.
\item
  \textbf{Paso 2}: Asignar cada individuo al grupo cuyo centro, que
  denominaremos \textbf{centroide}, le quede más próximo. Recalcular los
  centroides en base a los grupos modificados.
\item
  \textbf{Paso 3}: Repetir el paso 2 hasta que no haya más reasignaciones.
\end{itemize}

El primer paso se puede obtener proporcionando \(k\) centroides iniciales
escogidos adecuadamente a la vista de la muestra. De hecho, pueden ser
individuos de la muestra. También es posible conformar la partición
inicial, aplicando una técnica de agrupamiento jerárquico, y escogiendo
el nivel de agrupamiento en el cual se generan \(k\) grupos.

El paso 2 admite muchas variantes dependiendo del algoritmo concreto.
Planteado en términos generales, la idea es considerar posibles cambios
de algunos individuos de un grupo a otro. Serían candidatos a cambiar de
grupo los individuos que ocupan posiciones fronterizas. Para cada
posible cambio, se evalúan las consecuencias que tendría sobre el
criterio (por ejemplo, suma de cuadrados intra-grupo), y se efectúa el
cambio que más contribuye a mejorar el criterio. Se detiene el algoritmo
cuando no hay cambios que redunden en mejoras apreciables del criterio.

Los algoritmos de las \(k\)-medias presentan ciertas limitaciones, como
puede ser la dependencia de la partición inicial que se escoja, así como
la tendencia a formar grupos esféricos, ya que la distancia que se
considera es la distancia euclídea usual. Existen muchas propuestas de
otros métodos que permiten superar estas limitaciones y que se adaptan a
propósitos diversos. En cualquier caso, el algoritmo de las \(k\)-medias
es un punto de referencia inicial, que se debe tener presente dentro de
los métodos de particionamiento.

La función \emph{kmeans} de \(R\) permite escoger entre varios métodos de las
\(k\)-medias. Si no se indica el método, por defecto aplica el algoritmo
de \citet{Hartigan1979}.

\appendix


\chapter{Distribuciones notables multidimensionales}\label{anexoDist}

Al igual que en el caso univariante, en el que el análisis del
comportamiento muestral de determinados estadísticos da lugar a la
aparición de distribuciones como la \(t\) de Student, la ji-cuadrado o la
\(F\) de Snédecor, también en el contexto multidimensional surgen nuevas
distribuciones como consecuencia del estudio de estadísticos
importantes. En este tema estudiaremos distribuciones que surgen al
hacer operaciones naturales de cara a la inferencia sobre vectores
normales multivariantes.

\section{Distribución normal multivariante}\label{distribucion-normal-multivariante}

Un vector aleatorio \(d\)-dimensional \(\bm{x}\) tiene distribución normal
multivariante con vector de medias \(\bm{\mu}\) y matriz de covarianzas
\(\bm\Sigma\) si su función de densidad viene dada por

\[f(\bm{x})=\frac{1}{(\sqrt{2\pi})^d\left|\bm\Sigma\right|^{1/2}}e^\frac{{-(\bm{x}-\bm{\mu})^\prime\bm{\Sigma}^{-1}(\bm{x}-\bm{\mu})}}{2}.\]
Decimos en este caso que \(\bm{x}\) tiene distribución
\(N_d(\bm{\mu},\bm\Sigma)\).

\subsection{Propiedades de la distribución normal multivariante}\label{propiedades-de-la-distribucion-normal-multivariante}

Presentamos a continuación algunas propiedades conocidas de la
distribución normal multivariante.

\begin{itemize}
\item
  Sean \(a_1,\ldots,a_d\) constantes y consideremos la combinación
  lineal \(z=a_1x_1+\ldots+a_dx_d=\bm{a}^\prime\bm{x}.\) Si \(\bm{x}\) es
  \(N_d(\bm{\mu},\bm\Sigma)\), entonces \(z\) es
  \(N(\bm{a}^\prime\bm{\mu}, \bm{a}^\prime\bm{\Sigma}\bm{a})\).
\item
  Si \(\bm{A}\) es una matriz \(k\times d\) de rango \(k\) (siendo
  \(k\leq d\)), entonces las \(k\) combinaciones lineales en
  \(\bm{A}\bm{x}\) tienen distribución normal multivariante, es decir,
  si \(\bm{x}\) es \(N_d(\bm{\mu},\bm\Sigma)\), entonces \(\bm{A}\bm{x}\) es
  \(N_k(\bm{A}\bm{\mu}, \bm{A}\bm{\Sigma}\bm{A}^\prime)\).
\item
  Recordamos que si \(\bm{x}\) es \(N_d(\bm{\mu},\bm\Sigma)\), entonces
  \(\bm{z}=\bm{\Sigma^{-1/2}}(\bm{x}-\bm\mu)\) tiene distribución \(N_d(\bm{0},\bm{I_d}).\)
  Además,

  \[\bm{z}^\prime{}\bm{z}=\sum_{i=1}^dz_i^2\] es la suma de
  los cuadrados de \(d\) variables normales estándar e independientes y,
  por lo tanto, tiene distribución \(\chi_d^2\). Es decir

  \[(\bm{x}-\bm\mu)^\prime{}\bm{\Sigma}^{-1}(\bm{x}-\bm\mu) \textnormal{ tiene distribución } \chi_d^2.\]
\item
  Cualquier subconjunto de variables de un vector \(\bm{x}\) con
  distribución normal \(N_d(\bm{\mu},\bm\Sigma)\), tiene a su vez
  distribución normal (con vector de medias y matriz de covarianzas
  correspondiente a la partición de \(\bm{x}\)). Por ejemplo, si
  \(\bm{x_1}=(x_1,\ldots,x_p)^\prime{}\) denota el subvector de las
  primeras \(p\) componentes de \(\bm{x}\) y
  \(\bm{x_2}=(x_{p+1},\ldots,x_d)^\prime{}\) denota el subvector con las
  restantes \(d-p\) componentes, podemos escribir:

  \[\bm{x}=\left(
  \begin{array}{c}
  \bm{x_1}\\ 
  \bm{x_2}
  \end{array}
  \right),\ \ \ \bm{\mu}=\left(\begin{array}{c}\bm{\mu_1}\\ \bm{\mu_2}\end{array}\right),\ \ \ \bm{\Sigma}=\left(\begin{array}{cc}\bm{\Sigma_{11}}&\bm{\Sigma_{12}}\\ \bm{\Sigma_{21}}&\bm{\Sigma_{22}}\end{array}\right).\]
  Además

  \[\bm{x_1} \textnormal{ tiene distribución }  N_p(\bm{\mu_1},\bm{\Sigma_{11}}),\]

  \[\bm{x_2} \textnormal{ tiene distribución }  N_{d-p}(\bm{\mu_2},\bm{\Sigma_{22}}).\]
  En particular cada componente \(x_i\) del vector \(\bm{x}\)
  (\(i=1,\ldots,d\)) tiene distribución \(N(\mu_i,\sigma_i^2)\). Es
  importante notar que el recíproco no es cierto, es decir, si la
  densidad de cada \(x_i\) es normal no se sigue necesariamente que
  \(\bm{x}\) sea normal multivariante.
\item
  Consideremos ahora un vector particionado en dos subvectores
  \(\bm{x}\) de dimensión \(p\times 1\) e \(\bm{y}\) de dimensión
  \(q\times 1\), de forma que:

  \[\left(\begin{array}{c}
  \bm{x}\\ 
  \bm{y}
  \end{array}
  \right) \textnormal{ tiene distribución }  N_{p+q}\left(\left(\begin{array}{c}
  \bm{\mu_x}\\ 
  \bm{\mu_y}
  \end{array}\right),\left(\begin{array}{cc}\bm{\Sigma_{xx}}&\bm{\Sigma_{xy}}\\ \bm{\Sigma_{yx}}&\bm{\Sigma_{yy}}\end{array}\right)\right).\]

  \begin{itemize}
  \item
    Los subvectores \(\bm{x}\) e \(\bm{y}\) son independientes si y sólo
    si \(\bm{\Sigma_{xy}}=\bm{0}\).
  \item
    Si \(\bm{x}\) e \(\bm{y}\) no son independientes, entonces
    \(\bm{\Sigma_{xy}}\neq\bm{0}\) y la distribución condicional del
    vector \(\bm{x}\) dados los valores del vector \(\bm{y}\) también es
    normal (\(p\)-dimensional) con vector de medias

    \[E(\bm{x}/\bm{y})=\bm{\mu_x}+\bm{\Sigma_{xy}}\bm{\Sigma_{yy}}^{-1}(\bm{y}-\bm{\mu_y})\]
    y con matriz de covarianzas

    \[\bm{\Sigma_{x/y}}=\bm{\Sigma_{xx}}-\bm{\Sigma_{xy}}\bm{\Sigma_{yy}}^{-1}\bm{\Sigma_{yx}}.\]
  \end{itemize}
\item
  Si \(\bm{x}\) e \(\bm{y}\) son vectores con distribución normal, tienen
  la misma dimensión \(p\times 1\) y son independientes, entonces

  \[\bm{x}+\bm{y}\textnormal{ tiene distribución }  N_p(\bm{\mu_x}+\bm{\mu_y},\bm{\Sigma_{xx}}+\bm{\Sigma_{yy}}).\]

  \[\bm{x}-\bm{y}\textnormal{ tiene distribución }  N_p(\bm{\mu_x}-\bm{\mu_y},\bm{\Sigma_{xx}}+\bm{\Sigma_{yy}}).\]
\end{itemize}

\section{Distribución de Wishart}\label{distribucion-de-wishart}

\subsection{Definición y propiedades}\label{definicion-y-propiedades}

La distribución de Wishart surge como una extensión al caso
multivariante de la distribución ji-cuadrado. Recuerda que si
\(x_1,\ldots, x_m\) son variables aleatorias independientes con
distribución \(N(0,1)\), entonces la variable \(\sum_{i=1}^m{{x_i}^2}\)
tiene distribución ji-cuadrado con \(m\) grados de libertad.

Cuando \(\bm{x_1},\ldots, \bm{x_m}\) son vectores aleatorios en lugar de
variables aleatorias, por ejemplo \(\bm{x_i}\in\mathbb{R}^d\) con
\(\bm{x_i}\in N_d(\bm{0},\bm{I_d})\), una posible forma de generalizar la
suma de cuadrados anterior sería construir la matriz \(d\times d\)
semidefinida positiva \(\bm{M}=\sum_{i=1}^m\bm{x_i}\bm{x_i}^{\prime}\).
Mientras la distribución ji-cuadrado se atribuye a estimadores de la
varianza, como la varianza muestral, la distribución de Wishart
corresponde a matrices de covarianza muestrales.

La representación de \(\bm{M}\) da lugar, en esencia, a la definición de
la distribución de Wishart. Como se trata de una matriz aleatoria, no
nos vamos a preocupar por una supuesta función de densidad para la
distribución de Wishart, ni vamos a considerar tablas de cuantiles.
Simplemente vamos a estudiar las propiedades que presenta una matriz de
covarianzas aleatoria, que serán de utilidad más adelante para las
regiones de confianza multivariantes.

\begin{definition}
\protect\hypertarget{def:unnamed-chunk-97}{}\label{def:unnamed-chunk-97}Sean \(\bm{x_1},\ldots,\bm{x_m}\in N_d(\bm{0},\bm\Sigma)\) independientes.
Diremos que la matriz aleatoria

\[\bm{M}=\sum_{i=1}^m \bm{x_i} \bm{x_i}^{\prime}\] tiene una
distribución Wishart con matriz \(\bm\Sigma\) y \(m\) grados de libertad.
Denotamos \(\bm{M}\in W_d(\bm\Sigma,m)\).
\end{definition}

\textbf{Observaciones}

\begin{enumerate}
\def\labelenumi{\arabic{enumi}.}
\item
  Si construimos la matriz de observaciones

  \[\bm{X}=\left(\begin{array}{c} \bm{x_1}^{\prime} \\ \vdots \\ \bm{x_m}^{\prime} \end{array}\right)
  =\left(\begin{array}{cccc} x_{11} & x_{12} & \cdots & x_{1d}\\
  x_{21} & x_{22} & \cdots & x_{2d}\\
  \vdots & \vdots & \ddots & \vdots \\
  x_{m1} & x_{m2} & \cdots & x_{md}\end{array}\right)\] de orden
  \(m\times d\), entonces podemos escribir

  \[\bm{M}=\sum_{i=1}^m \bm{x_i} \bm{x_i}^{\prime}=\left(\bm{x_1},\ldots,\bm{x_m}\right)
  \left(\begin{array}{c} \bm{x_1}^{\prime} \\ \vdots \\ \bm{x_m}^{\prime} \end{array}\right)
  =\bm{X}^{\prime}\bm{X}=\left(\sum_{i=1}^m x_{ij} x_{ik}\right)_{j,k\in\{1,\ldots,d\}}.\]
\item
  Si \(\bm\Sigma=\bm{I_d}\) diremos que la distribución tiene forma
  estándar.
\item
  Si \(d=1\), entonces \(\bm\Sigma=\sigma^2\) y la distribución de Wishart
  es una ji-cuadrado con \(m\) grados de libertad y parámetro de escala
  \(\sigma^2\). Es decir, \(x_1,\ldots,x_m\in N(0,\sigma^2)\)
  independientes y

  \[\bm{M}=\sum_{i=1}^m x_i x_i^{\prime}=\sum_{i=1}^m x_i^2
  =\sigma^2 \sum_{i=1}^m \left(\frac{x_i}{\sigma}\right)^2 \in \sigma^2 \chi_m^2.\]
\end{enumerate}

\begin{theorem}
\protect\hypertarget{thm:transW}{}\label{thm:transW}Si \(\bm{M}\in W_d(\bm\Sigma,m)\) y
\(\bm{A}\) es una matriz \(q\times d\) de rango \(q\) (\(q\leq d\)), entonces

\[\bm{A}\bm{M}\bm{A}^{\prime}\in W_q\left(\bm{A}\bm\Sigma \bm{A}^{\prime},m\right).\]
\end{theorem}

\begin{proof}
Si \(\bm{M}\in W_d(\bm\Sigma,m)\), entonces lo podemos expresar así

\[\bm{M}=\sum_{i=1}^m\bm{x_i}\bm{x_i}^{\prime}\] siendo \(\bm{x_1},\ldots,\bm{x_m}\in N_d(\bm{0},\bm\Sigma)\) independientes. En tal caso, resulta que

\[\bm{A}\bm{M}\bm{A}^{\prime}=\sum_{i=1}^m \bm{A}\bm{x_i} \bm{x_i}^{\prime}\bm{A}^{\prime}
=\sum_{i=1}^m \left(\bm{A}\bm{x_i}\right) \left(\bm{A}\bm{x_i}\right)^{\prime}\]
y además \(\bm{A}\bm{x_1},\ldots,\bm{A}\bm{x_m}\in N_q(\bm{0},\bm{A}\bm\Sigma \bm{A}^{\prime})\)
independientes. Entonces, atendiendo a la definición de la distribución
de Wishart, se tiene que
\(\bm{A}\bm{M}\bm{A}^{\prime}\in W_q\left(\bm{A}\bm\Sigma \bm{A}^{\prime},m\right)\).
\end{proof}

\begin{corollary}
\protect\hypertarget{cor:unnamed-chunk-99}{}\label{cor:unnamed-chunk-99}Si \(\bm{M}\in W_d(\bm\Sigma,m)\), entonces las submatrices diagonales de
\(\bm{M}\) también tienen distribuciones de Wishart, con parámetros la
submatriz correspondiente y \(m\) grados de libertad.
\end{corollary}

\begin{corollary}
\protect\hypertarget{cor:colW}{}\label{cor:colW}Si \(\bm{M}\in W_d(\bm\Sigma,m)\) y
\(\bm{a}\) es un vector \(d\)-dimensional tal que
\(\bm{a}^\prime{}\bm\Sigma \bm{a}>0\), entonces

\[\bm{a}^\prime{}\bm{M}\bm{a}\in \left(\bm{a}^\prime{}\bm\Sigma \bm{a}\right)\chi_m^2.\]
\end{corollary}

En particular, por el Corolario \ref{cor:colW} se tiene que dado
\(\bm{e_j}=(0,\ldots,0,1,0,\ldots,0)^\prime{}\),

\[\bm{e_j}^\prime{}\bm{M}\bm{e_j}\equiv m_{jj}\in \sigma_{jj}^2\chi_m^2.\]
Es decir, la distribución marginal de \(m_{jj}\) es ji-cuadrado para
\(j=1,\ldots,d\). Sin embargo, la distribución marginal de \(m_{jk}\)
(\(j\neq k\)) no es ji-cuadrado.

Otra consecuencia del Teorema \ref{thm:transW} es que, fijado \(d\) y \(m\), la familia de
distribuciones \(W_d(\bm\Sigma,m)\) con \(\bm\Sigma\) semidefinida positiva
puede ser generada a partir de la distribución \(W_d(\bm{I_d},m)\) y de
matrices \(d\times d\). Observa que si \(\bm{M}\in W_d(\bm{I_d},m)\) y
\(\bm\Sigma=\bm{A}\bm{A}^\prime{}\), entonces

\[\bm{A}\bm{M}\bm{A}^{\prime}\in W_d\left(\bm{A}\bm{I_d}\bm{A}^{\prime},m\right)=W_d\left(\bm\Sigma,m\right).\]

Veremos a continuación algunas propiedades de la distribución Wishart.

\textbf{Propiedades}

\begin{enumerate}
\def\labelenumi{\arabic{enumi}.}
\item
  Si \(\bm{M}\in W_d(\bm\Sigma,m)\), entonces
  \(E(\bm{M})=m\cdot\bm\Sigma\).
\item
  \(\bm{M}=\bm{x_1}\bm{x_1}^\prime{}+\cdots+\bm{x_m}\bm{x_m}^\prime{}\)
  es una matriz simétrica, semidefinida positiva y de rango la
  dimensión del espacio lineal generado por
  \(\bm{x_1},\ldots,\bm{x_m}\). Si \(\bm{x_1},\ldots,\bm{x_m}\) son
  vectores aleatorios independientes, con la misma distribución y su
  matriz de covarianzas tiene rango \(k\), entonces

  \[P(\mbox{rango}(\bm{M})=\min\{m,k\})=1.\]
\item
  Como consecuencia de la propiedad anterior, si
  \(\bm{M}\in W_d(\bm\Sigma,m)\), \(\bm\Sigma\) es definida positiva y
  \(m\geq d\), entonces \[P(\bm{M}\ \mbox{sea definida positiva})=1.\]
\item
  La familia de distribuciones \(W_d(\bm\Sigma,m)\) es reproductiva
  respecto del parámetro \(m\), esto es, si
  \(\bm{M_1}\in W_d(\bm\Sigma,m_1)\), \(\bm{M_2}\in W_d(\bm\Sigma,m_2)\) y
  además son independientes, entonces

  \[\bm{M_1}+\bm{M_2}\in W_d(\bm\Sigma,m_1+m_2).\]
\end{enumerate}

\subsection{Formas cuadráticas en una muestra de observaciones normales multivariantes}\label{formas-cuadraticas-en-una-muestra-de-observaciones-normales-multivariantes}

En esta sección veremos las extensiones al caso multivariante de los
resultados sobre la distribución de formas cuadráticas de variables
normales. Estos resultados, al igual que en el caso unidimensional,
serán útiles para las tareas de inferencia en base a una muestra de una
distribución normal o en problemas de regresión.

\begin{theorem}
\protect\hypertarget{thm:tc}{}\label{thm:tc}

Sean \(\bm{x_1},\ldots,\bm{x_m}\in N_d(\bm{0},\bm\Sigma)\) independientes y la
matriz de observaciones

\[\bm{X}=\left(\begin{array}{c} \bm{x_1}^\prime{} \\ \vdots \\ \bm{x_m}^\prime{}\end{array}\right)
=\left(\begin{array}{ccc} x_{11} & \cdots & x_{1d} \\
\cdots & \cdots & \cdots \\
x_{m1} & \cdots & x_{md} \end{array}\right).\]

Entonces:

\begin{itemize}
\item
  Si \(\bm{a}\) es un vector de dimensión \(m\) no nulo, entonces
  \[\bm{X}^\prime{}\bm{a}\in N_d\left(\bm{0},\|\bm{a}\|^2\bm\Sigma\right).\]
\item
  Si \(\bm{a}\) y \(\bm{b}\) son dos vectores no nulos y ortogonales,
  \(\bm{a}^\prime{}\bm{b}=0\), entonces \(\bm{X}^\prime{}\bm{a}\) y
  \(\bm{X}^\prime{}\bm{b}\) son independientes.
\end{itemize}

\end{theorem}

\begin{proof}
Probamos a continuación cada uno de los puntos del enunciado.

\(i.\) Observamos que

\[\bm{X}^\prime{}\bm{a}=\left(\bm{x_1},\ldots,\bm{x_m}\right)\bm{a}=\sum_{i=1}^m \bm{x_i} a_i.\]

Entonces, como \(\bm{x_1},\ldots,\bm{x_m}\in N_d(\bm{0},\bm\Sigma)\)
independientes y

\[\mbox{Cov}\left(\bm{X}^\prime{}\bm{a},\bm{X}^\prime{}\bm{a}\right)=\mbox{Cov}\left(\sum_{i=1}^m \bm{x_i} a_i,\sum_{i=1}^m \bm{x_i} a_i\right)=\sum_{i=1}^m\sum_{j=1}^m
a_i a_j \mbox{Cov}\left(\bm{x_i},\bm{x_j}\right)
\stackrel{(a)}{=}\sum_{i=1}^m a_i^2 \bm\Sigma =\|\bm{a}\|^2\Sigma,\] se
concluye que \(\bm{X}^\prime{}\bm{a}\) tiene distribución normal de media
cero y con la matriz de covarianzas que se acaba de obtener. La igualdad
(a) es consecuencia de que los vectores \(\bm{x_1},\ldots,\bm{x_m}\) son
independientes y la matriz de covarianzas de cada uno de ellos es
\(\bm\Sigma\).

\(ii.\) De manera muy similar al apartado anterior, se tiene que

\[\begin{aligned}
\mbox{Cov}\left(\bm{X}^\prime{}\bm{a},\bm{X}^\prime{}\bm{b}\right)&=\mbox{Cov}\left(\sum_{i=1}^m \bm{x_i} a_i,\sum_{i=1}^m \bm{x_i} b_i\right)\\
&=\sum_{i=1}^m\sum_{j=1}^m a_i b_j \mbox{Cov}\left(\bm{x_i},\bm{x_j}\right)\\
&=\sum_{i=1}^m a_ib_i \bm\Sigma\\
&=\left(\bm{a}^\prime{}\bm{b}\right)\bm\Sigma\\
&=\bm{0}.
\end{aligned}\]

De modo que \(\bm{X}^\prime{}\bm{a}\) y \(\bm{X}^\prime{}\bm{b}\) están
incorrelacionados. Como por el apartado anterior, tienen distribución
normal, entonces son independientes.
\end{proof}

\begin{theorem}
\protect\hypertarget{thm:tc2}{}\label{thm:tc2}

Sean
\(\bm{x_1},\ldots,\bm{x_m}\in N_d(\bm{0},\bm\Sigma)\) independientes y
\(\bm{X}\) la correspondiente matriz de observaciones.

\begin{itemize}
\item
  Si \(\bm{A}\) es una matriz simétrica de orden \(m\times m\),
  idempotente (\(\bm{A}^2=\bm{A}\)) y de rango \(r\leq m\), entonces

  \[\bm{X}^\prime{}\bm{A}\bm{X}\in W_d(\bm\Sigma,r).\]
\item
  Si \(\bm{A}\) es una matriz en las condiciones anteriores y \(\bm{b}\)
  es un vector de dimensión \(m\), tal que \(\bm{A}\bm{b}=\bm{0}\),
  entonces

  \[\bm{X}^\prime{}\bm{A}\bm{X}\ \ \mbox{y}\ \  \bm{X}^\prime{}\bm{b}\ \ \mbox{son independientes}.\]
\item
  Si \(\bm{A}\) y \(\bm{B}\) son dos matrices en las condiciones
  anteriores y \(\bm{A}\bm{B}=\bm{0}\), entonces

  \[\bm{X}^\prime{}\bm{A}\bm{X}\ \ \mbox{y}\ \  \bm{X}^\prime{}\bm{B}\bm{X}\ \  \mbox{son independientes}.\]
\end{itemize}

\end{theorem}

\begin{proof}
\(i.\) Por ser \(\bm{A}\) simétrica, idempotente y de rango
\(r\leq m\), tendrá \(r\) autovalores iguales a uno y \((m-r)\) autovalores
iguales a cero. En tal caso, se podrá expresar

\[\bm{A}=\sum_{i=1}^r \bm{v_i}\bm{v_i}^\prime{},\] siendo
\(\bm{v_1},\ldots,\bm{v_r}\) los autovectores asociados a los autovalores
unitarios. Entonces

\begin{equation}
\bm{X}^\prime{}\bm{A}\bm{X}=\sum_{i=1}^r \bm{X}^\prime{}\bm{v_i}\bm{v_i}^\prime{}\bm{X}=\sum_{i=1}^r \left(\bm{X}^\prime{}\bm{v_i}\right)\left(\bm{X}^\prime{}\bm{v_i}\right)^{\prime}.
\label{eq:xpax}
\end{equation}

El apartado (i) del Teorema \ref{thm:tc} nos lleva a que
\(\bm{X}^\prime{}\bm{v_1},\ldots,\bm{X}^\prime{}\bm{v_r}\in N_d(\bm{0},\bm\Sigma)\)
y el apartado (ii) del mismo teorema a que son independientes. En
consecuencia, de la expresión\eqref{eq:xpax} y la
definición de la distribución de Wishart, se deduce que

\[\bm{X}^\prime{}\bm{A}\bm{X}\in W_d(\bm\Sigma,r).\]

\(ii.\) La hipótesis \(\bm{A}\bm{b}=\bm{0}\) implica que
\(\bm{v_j}^\prime{}\bm{b}=0\) para todo \(j\in\{1,\ldots,r\}\). El apartado
(ii) del Teorema \ref{thm:tc}
conduce a que \(\bm{X}^\prime\bm{v_1},\ldots,\bm{X}^\prime\bm{v_r}\) sean
independientes de \(\bm{X}^\prime\bm{b}\). Recordando la expresión
\eqref{eq:xpax}, se
llega a que \(\bm{X}^\prime{}\bm{A}\bm{X}\) y \(\bm{X}^\prime{}\bm{b}\) son
independientes.

\(iii.\) Siendo \(\bm{A}=\sum_{i=1}^r \bm{v_i}\bm{v_i}^\prime{}\) y
\(\bm{B}=\sum_{j=1}^s \bm{\omega_i}\bm{\omega_i}^\prime{},\) con
\(\bm{A}\bm{B}=\bm{0}\), se tiene que \(\bm{v_k}^\prime{}\bm{\omega_l}=0\)
para todo \(k\in\{1,\ldots,r\}\) y todo \(l\in\{1,\ldots,s\}\).

Ahora, como \(\bm{v_k}^\prime{}\bm{\omega_l}=0\), aplicando el apartado
(ii) del Teorema \ref{thm:tc},
tenemos que \(\bm{X}^\prime{}\bm{v_1},\ldots,\bm{X}^\prime{}\bm{v_r}\) son
independientes de
\(\bm{X}^\prime{}\bm{\omega_1},\ldots,\bm{X}^\prime{}\bm{\omega_s}\), y en
base a la expresión \eqref{eq:xpax} y su análogo para \(\bm{X}^\prime{}\bm{B}\bm{X}\),
se concluye que \(\bm{X}^\prime{}\bm{A}\bm{X}\) y
\(\bm{X}^\prime{}\bm{B}\bm{X}\) son independientes.
\end{proof}

El teorema siguiente es la versión multivariante del Teorema de Fisher.

\begin{theorem}
\protect\hypertarget{thm:unnamed-chunk-102}{}\label{thm:unnamed-chunk-102}

Sean \(\bm{x_1},\ldots,\bm{x_n}\in N_d(\mu,\bm\Sigma)\) independientes.
Entonces el vector de medias y la matriz de covarianzas muestrales
verifican

\begin{enumerate}
\def\labelenumi{\arabic{enumi}.}
\item
  \({\overline{\bm{x}}=\frac{1}{n}\sum_{i=1}^n \bm{x_i} \in N_d\left(\bm{\mu},\frac{1}{n}\bm\Sigma\right)}.\)
\item
  \({n\bm{S}=\sum_{i=1}^n \left(\bm{x_i}-\overline{\bm{x}}\right)\left(\bm{x_i}-\overline{\bm{x}}\right)^\prime{}
  \in W_d\left(\bm\Sigma,n-1\right)}.\)
\item
  \(\overline{\bm{x}}\) y \(\bm{S}\) son independientes.
\end{enumerate}

\end{theorem}

\begin{proof}
La demostración es idéntica a la correspondiente al Teorema de
Fisher en el caso unidimensional. Así, se empieza observando que basta
hacer la demostración suponiendo que \(\bm{\mu}=\bm{\bm{0}}\). En el caso
general, bastaría considerar \(\bm{y}=\bm{x}-\bm{\mu}\), que serían
observaciones con media cero, y tener en cuenta que
\(\overline{\bm{y}}=\overline{\bm{x}}-\bm{\mu}\) y que la matriz de
covarianzas muestral no se ve afectada por la resta de \(\bm{\mu}\).

\(i.\) Tomamos el vector \(\bm{b}=(1/n,\ldots,1/n)^\prime{}\) y aplicamos
el Teorema \ref{thm:tc}, apartado
(i), teniendo en cuenta que en este caso \(\|\bm{b}\|^2=1/n\).

\(ii.\) Consideramos la matriz

\[\bm{A}=\bm{M}=\left(\begin{array}{cccc}
1-1/n & -1/n & \cdots & -1/n \\
-1/n & 1-1/n & \ddots & \vdots \\
\vdots & \ddots & \ddots & -1/n \\
-1/n & \cdots & -1/n & 1-1/n
\end{array}\right)=\bm{I_n}-(1/n)\bm{1_n}\bm{1_n}^\prime{}\] que tiene
el valor \((1-1/n)\) en la diagonal y el valor \((-1/n)\) fuera de la
diagonal. La notación \(\bm{1_n}\) significa
\(\bm{1_n}=(1,\ldots,1)^\prime{}\), esto es, un vector compuesto por unos
en todas sus componentes.

La matriz \(\bm{M}\) es muy importante en la Estadística. Su efecto sobre
el conjunto de observaciones es el siguiente:

\[\bm{M}\bm{X}=\left(\begin{array}{c}
\left(\bm{x_1}-\overline{\bm{x}}\right)^{\prime} \\
\vdots \\
\left(\bm{x_n}-\overline{\bm{x}}\right)^{\prime}
\end{array}\right).\] De modo que \(\bm{M}\) es una matriz centralizadora,
pues aplicada a un conjunto de datos, los centra, al restar a cada dato
la media. Además, si se emplea como forma cuadrática, se tiene que

\[\begin{aligned}
\bm{X}^\prime{}\bm{M}\bm{X}&\stackrel{(a)}{=}\bm{X}^\prime{}\bm{M}\bm{M}\bm{X}\\
&\stackrel{(b)}{=}\ (\bm{M}\bm{X})^\prime{}(\bm{M}\bm{X})\\
&=\ \left(\bm{x_1}-\overline{\bm{x}},\ldots, \bm{x_n}-\overline{\bm{x}}\right)\cdot \left(\begin{array}{c}
\left(\bm{x_1}-\overline{\bm{x}}\right)^{\prime} \\
\vdots \\
\left(\bm{x_n}-\overline{\bm{x}}\right)^{\prime}
\end{array}\right)\\
&=\sum_{i=1}^n \left(\bm{x_i}-\overline{\bm{x}}\right)\left(\bm{x_i}-\overline{\bm{x}}\right)^{\prime}.
\end{aligned}\]

En la igualdad (a) se ha empleado que \(\bm{M}\) es una matriz idempotente
y en la igualdad (b) que es simétrica. Es fácil comprobar que \(\bm{M}\)
es una matriz simétrica, idempotente y de rango \((n-1)\). Pero entonces
el Teorema \ref{thm:tc2},
apartado (i), permite concluir la demostración de que la matriz de
covarianzas muestral tiene distribución Wishart con matriz \(\bm\Sigma\) y
\((n-1)\) grados de libertad.

\(iii.\) La independencia se deduce del Teorema \ref{thm:tc2}, apartado (ii), al
observar que \(\bm{M}\bm{b}=\bm{0}\).
\end{proof}

\section{\texorpdfstring{Distribución \(\Gamma^2\) de Hotelling}{Distribución \textbackslash Gamma\^{}2 de Hotelling}}\label{distribucion-gamma-2-de-hotelling}

\subsection{Definición y propiedades}\label{definicion-y-propiedades-2}

La distribución \(\Gamma^2\) de Hotelling juega un papel semejante a la
distribución \(t\) de Student en el caso multivariante. En concreto, la
distribución \(t\) de Student surge cuando se estandariza la media
muestral mediante la varianza muestral. En el caso de poblaciones
normales multivariantes, se calcula la distancia de la media muestral a
la media teórica usando la distancia de Mahalanobis, que la distancia
asociada a la matriz de covarianzas muestral.

Empezamos con la definición de la distribución \(\Gamma^2\) de Hotelling,
que, igual que ocurría con la \(t\) de Student, es una definición formal,
preparada para adaptarse al problema de inferencia.

\begin{definition}
\protect\hypertarget{def:unnamed-chunk-104}{}\label{def:unnamed-chunk-104}Sean \(\bm{Z}\in N_d(\bm{0},\bm{I_d})\) y \(\bm{R}\in W_d(\bm{I_d},m)\)
independientes, con \(m\geq d\). Diremos que la variable aleatoria
unidimensional \[m\bm{Z}^{\prime} \bm{R}^{-1} \bm{Z} \in \Gamma^2(d,m)\]
tiene una distribución \(\Gamma^2\) de Hotelling de parámetros \(d\) y \(m\).
\end{definition}

En el caso particular \(d=1\), resulta
\(\Gamma^2(1,m)\stackrel{d}{=}t_{m}^2\), por lo cual decimos que la
\(\Gamma^2\) de Hotelling generaliza la distribución \(t\) de Student.
Nótese que en la definición anterior, tanto \(\bm{Z}\) como \(\bm{R}\) están
estandarizados. Así, \(\bm{Z}\) tiene media cero y como matriz de
covarianzas la identidad. \(\bm{R}\) tiene forma estándar, pues su matriz
de covarianzas asociada es también la identidad.

En el teorema siguiente vemos cómo se obtiene una distribución
\(\Gamma^2\) de Hotelling a partir de variables no estandarizadas. El
único requisito es que el vector aleatorio y la matriz aleatoria tengan
la misma matriz de covarianzas de referencia.

\begin{theorem}
\protect\hypertarget{thm:unnamed-chunk-105}{}\label{thm:unnamed-chunk-105}Si \(\bm{x}\in N_d(\bm\mu,\bm\Sigma)\) y \(\bm{M}\in W_d(\bm\Sigma,m)\) son
independientes, entonces

\[m \left(\bm{x}-\bm\mu\right)^{\prime} \bm{M}^{-1} \left(\bm{x}-\bm\mu\right) \in \Gamma^2(d,m).\]
\end{theorem}

\begin{proof}
Para demostrar este teorema vamos a estandarizar el vector
aleatorio \(\bm{x}\) y la matriz aleatoria \(\bm{M}\). Así, definiendo como
\(\bm{Z}\) y \(\bm{R}\) a sus estandarizaciones, se tiene que

\[\begin{aligned}
\bm{Z}&=&\bm\Sigma^{-1/2}(\bm{x}-\bm\mu)\in N_d(\bm{0},\bm{I_d}). \\
\bm{R}&=&\bm\Sigma^{-1/2}\bm{M}\bm\Sigma^{-1/2} \in W_d(\bm{I_d},m).\end{aligned}\]

Para la distribución de \(\bm{R}\) se aplicó el Teorema \ref{thm:transW}.
\end{proof}

El siguiente teorema es muy útil, pues permite convertir una variable
con distribución \(\Gamma^2\) de Hotelling en otra con distribución \(F\) de
Snédecor. De este modo se evita la necesidad de considerar nuevas tablas
de cuantiles, o en general cualquier tipo de cálculos relacionados con
la \(\Gamma^2\) de Hotelling. Simplemente se efectuaría la transformación
y se apelaría a las tablas y operaciones conocidas para la \(F\) de
Snédecor.

\begin{theorem}
\protect\hypertarget{thm:gf}{}\label{thm:gf}\[\Gamma^2(d,m)\stackrel{d}{=} \frac{md}{m-d+1} F_{d,m-d+1}\]
\end{theorem}

\begin{proof}
Se omite aquí la demostración de este teorema. Se puede
encontrar en \citet{mkb79}.
\end{proof}

\subsection{\texorpdfstring{El estadístico \(\Gamma^2\) de Hotelling para inferencia sobre el vector de medias}{El estadístico \textbackslash Gamma\^{}2 de Hotelling para inferencia sobre el vector de medias}}\label{el-estadistico-gamma-2-de-hotelling-para-inferencia-sobre-el-vector-de-medias}

El siguiente corolario establece que la distancia de la media muestral a
la media poblacional, medida a partir de la matriz de covarianzas
muestral, tiene distribución de Hotelling. De hecho, este estadístico es
la razón principal por la que se ha definido la distribución de
Hotelling, hasta el punto de que recibe el nombre de estadístico de
Hotelling.

\begin{corollary}
\protect\hypertarget{cor:unnamed-chunk-108}{}\label{cor:unnamed-chunk-108}Sean \(\bm{x_1},\ldots, \bm{x_n}\in N_d(\bm\mu,\bm\Sigma)\)
independientes, y \(\overline{\bm{x}}\) el vector de medias y \(\bm{S}\) la
matriz de covarianzas muestrales obtenidos sobre la muestra anterior.
Entonces

\[(n-1) \left(\overline{\bm{x}}-\bm\mu\right)^{\prime} \bm{S}^{-1} \left(\overline{\bm{x}}-\bm\mu\right)
=n \left(\overline{\bm{x}}-\bm\mu\right)^{\prime} \bm{S_c}^{-1} \left(\overline{\bm{x}}-\bm\mu\right)
\in \Gamma^2(d,n-1)\] siendo \(\bm{S_c}=\frac{n}{n-1}\bm{S}\).
\end{corollary}

\begin{proof}
Por el Teorema de Fisher multivariante, se tiene
que

\[\overline{\bm{x}}=\frac{1}{n}\sum_{i=1}^n \bm{x_i} \in N_d\left(\bm{\mu},\frac{1}{n}\bm\Sigma\right),
\qquad\qquad n\bm{S}=\sum_{i=1}^n \left(\bm{x_i}-\overline{\bm{x}}\right)\left(\bm{x_i}-\overline{\bm{x}}\right)^\prime{}
\in W_d\left(\bm\Sigma,n-1\right)\] y además son independientes. Se
puede deducir que

\[\sqrt{n}\left(\overline{\bm{x}}-\bm{\mu}\right)\in N_d\left(\bm{\mu},\bm\Sigma\right)\]
lo cual sigue siendo independiente de \(n\bm{S}\). Entonces, aplicando el
Teorema 5, se llega a que

\[(n-1)\left(\sqrt{n}\left(\overline{\bm{x}}-\bm{\mu}\right)\right)^\prime{}(n\bm{S})^{-1}
\sqrt{n}\left(\overline{\bm{x}}-\bm{\mu}\right)\in \Gamma^2(d,n-1).\]
Efectuando operaciones se obtiene

\[(n-1)\left(\sqrt{n}\left(\overline{\bm{x}}-\bm{\mu}\right)\right)^\prime{} (n\bm{S})^{-1}
\sqrt{n}\left(\overline{\bm{x}}-\bm{\mu}\right)=(n-1) \left(\overline{\bm{x}}-\bm{\mu}\right)^{\prime} \bm{S}^{-1} \left(\overline{\bm{x}}-\bm{\mu}\right)\]
con lo cual se concluye la demostración.
\end{proof}

Juntando el resultado anterior con el Teorema \ref{thm:gf}, se puede transformar el
estadístico de Hotelling para que tenga distribución \(F\) de Snédecor, lo
cual es más útil a efectos prácticos.

Así, el siguiente corolario contiene el resultado práctico que se suele
emplear para obtener regiones de confianza o realizar contrastes de
hipótesis sobre el vector de medias.

\begin{corollary}
\protect\hypertarget{cor:nose}{}\label{cor:nose}Sean \(\bm{x_1},\ldots, \bm{x_n}\in N_d(\bm{\mu},\bm\Sigma)\)
independientes, y \(\overline{\bm{x}}\) el vector de medias y \(\bm{S}\) la
matriz de covarianzas muestrales obtenidos sobre la muestra anterior.
Entonces

\[\frac{n-d}{d} \left(\overline{\bm{x}}-\bm{\mu}\right)^{\prime} \bm{S}^{-1} \left(\overline{\bm{x}}-\bm{\mu}\right)
\in F_{d,n-d}\]
\end{corollary}

\begin{proof}
Es consecuencia directa del Corolario \ref{cor:nose} y el Teorema \ref{thm:gf}. Así, se
tiene que

\[(n-1) \left(\overline{\bm{x}}-\bm{\mu}\right)^{\prime} \bm{S}^{-1} \left(\overline{\bm{x}}-\bm{\mu}\right)\in \Gamma^2(d,n-1)=\frac{(n-1)d}{n-1-d+1}F_{d,n-1-d+1}.\]
Simplificando se llega a la expresión que figura en este corolario.
\end{proof}

\section{\texorpdfstring{Distribución \(\Lambda\) de Wilks}{Distribución \textbackslash Lambda de Wilks}}\label{distribucion-lambda-de-wilks}

\subsection{Introducción}\label{introduccion}

La distribución \(\Lambda\) de Wilks sirve para extender al caso
multivariante la distribución \(F\) de Snédecor. Así, si la \(F\) de
Snédecor se obtiene mediante el cociente de dos variables de tipo
ji-cuadrado, la \(\Lambda\) de Wilks surge como el cociente de los
determinantes de dos matrices de covarianzas. Iremos viendo cómo se
llega a ese cociente.

El uso más destacable de la distribución \(F\) de Snédecor es para
realizar el test \(F\) de comparación de modelos lineales. En el test \(F\)
se contrastaba una hipótesis nula que constituía una restricción sobre
un modelo más general. El modelo general era la hipótesis alternativa.

Así, si \(RSS\) representaba la suma residual de cuadrados bajo la
hipótesis alternativa y \(RSS_0\) la misma suma residual bajo la hipótesis
nula, entonces necesariamente \(RSS_0>RSS\) y el estadístico de contraste
se calculaba así:

\[\frac{(RSS_0-RSS)/q}{RSS/(n-p)}\in F_{q,n-p}\] donde
\(q\) y \(n-p\) eran los grados de libertad de numerador y denominador (\(q\)
era el número de restricciones, \(n\) el número de individuos y \(p\) el
número de parámetros bajo el modelo general).

En el caso multivariante, esto es, en modelos lineales con respuesta
multivariante, los residuos de la regresión serán vectores, y en lugar
de sumas de cuadrados, generan matrices de covarianzas.

La idea será comparar la matriz de covarianzas residual, que vamos a
denotar por \(\bm{E}\), con la matriz de covarianzas asociada a la
restricción, que denotamos por \(\bm{H}\), y juega el papel de
\((RSS_0-RSS)\). Si \(\bm{H}\) es muy grande, en comparación con \(\bm{E}\),
entonces se rechaza la hipótesis nula. La forma de medir si una matriz
de covarianzas es grande o pequeña será a través de su determinante.

\subsection{Definición y propiedades}\label{definicion-y-propiedades-3}

Vamos a definir la distribución \(\Lambda\) de Wilks como el cociente
entre el determinante de \(\bm{E}\) y el determinante de \(\bm{E}+\bm{H}\).
Recordemos que \(\bm{E}\) estaba asociado a la matriz de covarianzas
residual, y \(\bm{H}\) a la matriz de covarianzas adicional, debida a las
restricciones.

La razón para no comparar directamente \(\bm{H}\) con \(\bm{E}\), sino
\(\bm{E}+\bm{H}\) con \(\bm{E}\) es que \(\bm{H}\) podría tener algún
autovalor igual a cero, y eso haría que su determinante fuera cero. Ese
problema desaparece al sumarlo a \(\bm{E}\), pues \(\bm{E}\) siempre tiene
todos sus autovalores distintos de cero, ya que sus grados de libertad
son \((n-p)\), en los problemas de regresión, donde \(n\) es el número de
individuos y \(p\) los parámetros del modelo.

Lo anterior es la justificación práctica. A continuación se ofrece la
definición formal y genérica de la distribución de Wilks.

\begin{definition}
\protect\hypertarget{def:unnamed-chunk-111}{}\label{def:unnamed-chunk-111}Sean \(\bm{E}\in W_d(\bm{I_d},m_E)\) y \(\bm{H}\in W_d(\bm{I_d},m_H)\)
independientes, y \(m_E\geq d\). Diremos que la variable aleatoria
unidimensional
\[\frac{|\bm{E}|}{|\bm{E}+\bm{H}|}=\left|\bm{I_d}+\bm{E}^{-1}\bm{H}\right|^{-1} \in \Lambda(d,m_H,m_E)\]
tiene una distribución \(\Lambda\) de Wilks de parámetros \(d\), \(m_H\) y
\(m_E\).
\end{definition}

El siguiente resultado, que presentamos sin demostración, permite
simplificar la construcción de tablas para la distribución \(\Lambda\) de
Wilks, pues los dos primeros parámetros se pueden intercambiar en cierta
manera.

\begin{theorem}
\protect\hypertarget{thm:unnamed-chunk-112}{}\label{thm:unnamed-chunk-112}Si denotamos mediante \(\psi(d,m_H,m_E)\) la distribución conjunta de los
autovalores no nulos de \(\bm{E}^{-1}\bm{H}\), con \(m_E\geq d\), \(m_H\geq 1\), \(d\geq 1\), entonces

\[\psi(d,m_H,m_E)=\psi(m_H,d,m_E-d+m_H).\] En consecuencia,

\[\Lambda(d,m_H,m_E)\stackrel{d}{=}\Lambda(m_H,d,m_E-d+m_H).\]
\end{theorem}

A continuación se presenta un compendio de propiedades de las
distribuciones de Hotelling y de Wilks que facilitan la obtención de
cuantiles u otros elementos necesarios para la inferencia.

\section*{Propiedades de la distribución de Hotelling y de Wilks}\label{apuxe9ndice-propiedades-de-la-distribuciuxf3n-de-hotelling-y-de-wilks}
\addcontentsline{toc}{section}{Propiedades de la distribución de Hotelling y de Wilks}

\subsection*{Distribución de Hotelling}\label{distribuciuxf3n-de-hotelling}
\addcontentsline{toc}{subsection}{Distribución de Hotelling}

\[\Gamma^2(d,m)\stackrel{d}{=}\frac{md}{m-d+1}F_{d,m-d+1}.\]

\subsection*{Distribución de Wilks. Casos especiales}\label{distribuciuxf3n-de-wilks.-casos-especiales}
\addcontentsline{toc}{subsection}{Distribución de Wilks. Casos especiales}

\[\Lambda(d,m_H,m_E)\stackrel{d}{=} \Lambda(m_H,d,m_E-d+m_H).\]

Si \(d=1\),

\[\frac{1-\Lambda(1,m_H,m_E)}{\Lambda(1,m_H,m_E)}\stackrel{d}{=} \frac{m_H}{m_E} F_{m_H,m_E}.\]

Si \(m_H=1\),

\[\frac{1-\Lambda(d,1,m_E)}{\Lambda(d,1,m_E)}\stackrel{d}{=} \frac{d}{m_E-d+1} F_{d,m_E-d+1}.\]

Si \(d=2\) (\(m_H\geq 2\)),

\[\frac{1-\sqrt{\Lambda(2,m_H,m_E)}}{\sqrt{\Lambda(2,m_H,m_E)}}\stackrel{d}{=} \frac{m_H}{m_E-1} F_{2m_H,2(m_E-1)}.\]

Si \(m_H=2\) (\(d\geq 2\)),

\[\frac{1-\sqrt{\Lambda(d,2,m_E)}}{\sqrt{\Lambda(d,2,m_E)}}\stackrel{d}{=} \frac{d}{m_E-d+1} F_{2d,2(m_E-d+1)}.\]

Aproximación de Bartlett: Para \(m_E\) grande,

\[-f\ln \Lambda(d,m_H,m_E)\ \mbox{se aproxima a}\ \chi_{dm_H}^2\] donde

\(f=m_E-\frac{1}{2}(d-m_H+1).\)

Aproximación de Rao:

\[\frac{1-\Lambda(d,m_H,m_E)^{1/t}}{\Lambda(d,m_H,m_E)^{1/t}}\frac{ft-g}{dm_H}\textnormal{ es aproximadamente } F_{dm_H,ft-g}\]
siendo

\[f=m_E-\frac{1}{2}(d-m_H+1),\]

\[t=\left(\frac{d^2m_H^2-4}{d^2+m_H^2-5}\right)^{1/2},\]

\[g=\frac{dm_H-2}{2}.\]

\chapter{Conceptos básicos de álgebra matricial}\label{anexoAlg}

Este Anexo introduce conceptos básicos de álgebra matricial y hace una
revisión breve de resultados que serán útiles a lo largo del curso.

\section{Autovalores}\label{autovalores}

Dada una matriz \(\bm{A}\) de tamaño \(n\times n\) con autovalores
\(\lambda_1,\ldots,\lambda_n\), se tiene:

\begin{enumerate}
\def\labelenumi{\arabic{enumi}.}
\item
  \(\textnormal{tr}(\bm{A})=\sum_{i=1}^n\lambda_i.\)
\item
  \(\left|\bm{A}\right|=\prod_{i=1}^n\lambda_i.\)
\item
  Descomposición espectral. Sea \(\bm{A}\) una matriz de tamaño
  \(n\times n\) simétrica. Entonces existe una matriz ortogonal \(\bm{T}\)
  tal que
  \[\bm{T}^\prime{}\bm{A}\bm{T}=\textnormal{diag}(\lambda_1,\ldots,\lambda_n)\]
  donde \(\lambda_1\geq \lambda_2\geq\ldots\geq\lambda_n\) son los
  autovalores de \(\bm{A}\).
\end{enumerate}

\section{Matrices semidefinidas positivas}\label{matrices-semidefinidas-positivas}

Una matriz simétrica \(\bm{A}\) se dice \textbf{semidefinida positiva} si y
solo si \(\bm{x}^\prime{}\bm{A}\bm{x}\geq 0\) para todo \(\bm{x}\).

\begin{enumerate}
\def\labelenumi{\arabic{enumi}.}
\item
  Una matriz simétrica \(\bm{A}\) es semidefinida positiva si y solo si
  todos sus autovalores son no negativos.
\item
  Si \(\bm{A}\) es semidefinida positiva entonces
  \(\textnormal{tr}(\bm{A})\geq 0\).
\item
  Si \(\bm{A}\) es semidefinida positiva entonces
  \(\bm{C}^\prime{}\bm{A}\bm{C}\) es semidefinida positiva.
\item
  Si \(\bm{A}\) es semidefinida positiva entonces existe una matriz
  semidefinida positiva \(\bm{A^{1/2}}\), tal que
  \(\bm{A}=(\bm{A^{1/2}})^2\).
\end{enumerate}

\section{Matrices definidas positivas}\label{matrices-definidas-positivas}

Una matriz simétrica \(\bm{A}\) se dice \textbf{definida positiva} si y solo si
\(\bm{x}^\prime{}\bm{A}\bm{x}> 0\) para todo \(\bm{x}\), con \(\bm{x}\neq 0\).
Observamos que toda matriz definida positiva es también semidefinida
positiva.

\begin{enumerate}
\def\labelenumi{\arabic{enumi}.}
\item
  Una matriz simétrica \(\bm{A}\) es definida positiva si y solo si
  todos sus autovalores son positivos.
\item
  Si \(\bm{A}\) es semidefinida positiva y \(\left|\bm{A}\right|\neq 0\),
  entonces \(\bm{A}\) es definida positiva.
\item
  \(\bm{A}\) es definida positiva si y solo si existe una matriz no
  singular \(\bm{R}\) tal que \(\bm{A}=\bm{R}^\prime{}\bm{R}\).
\item
  Si \(\bm{A}\) es definida positiva entonces existe una matriz definida
  positiva \(\bm{A^{1/2}}\), tal que \(\bm{A}=(\bm{A^{1/2}})^2\).
\item
  Si \(\bm{A}\) es definida positiva también lo es \(\bm{A}^{-1}\).
\item
  Si \(\bm{A}\) es una matriz definida positiva de tamaño \(n\times n\) y
  \(\bm{C}\) es una matriz \(p\times n\) de rango \(p\), entonces
  \(\bm{C}\bm{A}\bm{C}^\prime{}\) es definida positiva.
\item
  Descomposición de Choleski. Si \(\bm{A}\) es definida positiva existe
  una única matriz triangular superior \(\bm{U}\) con elementos
  positivos en la diagonal tal que \(\bm{A}=\bm{U}^\prime{}\bm{U}\).
  También existe una única matriz triangular inferior \(\bm{L}\) con
  elementos positivos en la diagonal tal que
  \(\bm{A}=\bm{L}^\prime{}\bm{L}\).
\end{enumerate}

\nocite{*}
\bibliography{book.bib,packages.bib}

\printindex

\end{document}
